% Options for packages loaded elsewhere
\PassOptionsToPackage{unicode}{hyperref}
\PassOptionsToPackage{hyphens}{url}
\documentclass[a4paper]{extarticle}
\usepackage{fix-cm}
\usepackage[fontsize=13pt]{scrextend}
\usepackage{xcolor}
\usepackage[margin=2.5cm,top=2.5cm,bottom=2.5cm,left=2.5cm,right=2.5cm,headheight=1.25cm,headsep=0.5cm,footskip=1.25cm]{geometry}
\usepackage{amsmath,amssymb}
\setcounter{secnumdepth}{-\maxdimen} % remove section numbering
\usepackage{iftex}
\ifPDFTeX
  \usepackage[T5,T1]{fontenc}
  \usepackage[utf8]{inputenc}
  \usepackage{textcomp} % provide euro and other symbols
  \IfFileExists{newtxtext.sty}{\usepackage{newtxtext,newtxmath}}{\usepackage{lmodern}}
\else % if luatex or xetex
  \usepackage{fontspec}
  \IfFontExistsTF{Times New Roman}{\setmainfont{Times New Roman}}{\setmainfont{TeX Gyre Termes}}
  \setmonofont{Consolas}[Scale=MatchLowercase]
\fi
% Use upquote if available, for straight quotes in verbatim environments
\IfFileExists{upquote.sty}{\usepackage{upquote}}{}
\IfFileExists{microtype.sty}{% use microtype if available
  \usepackage[]{microtype}
  \UseMicrotypeSet[protrusion]{basicmath} % disable protrusion for tt fonts
}{}
\makeatletter
\@ifundefined{KOMAClassName}{% if non-KOMA class
  \IfFileExists{parskip.sty}{%
    \usepackage{parskip}
  }{% else
    \setlength{\parindent}{0pt}
    \setlength{\parskip}{6pt plus 2pt minus 1pt}}
}{% if KOMA class
  \KOMAoptions{parskip=half}}
\makeatother
\usepackage{color}
\usepackage{fancyvrb}
\newcommand{\VerbBar}{|}
\newcommand{\VERB}{\Verb[commandchars=\\\{\}]}
\DefineVerbatimEnvironment{Highlighting}{Verbatim}{commandchars=\\\{\}}
\RecustomVerbatimEnvironment{verbatim}{Verbatim}{fontsize=\small}
% Add ',fontsize=\small' for more characters per line
\newenvironment{Shaded}{}{}
\newcommand{\AlertTok}[1]{\textcolor[rgb]{1.00,0.00,0.00}{\textbf{#1}}}
\newcommand{\AnnotationTok}[1]{\textcolor[rgb]{0.38,0.63,0.69}{\textbf{\textit{#1}}}}
\newcommand{\AttributeTok}[1]{\textcolor[rgb]{0.49,0.56,0.16}{#1}}
\newcommand{\BaseNTok}[1]{\textcolor[rgb]{0.25,0.63,0.44}{#1}}
\newcommand{\BuiltInTok}[1]{\textcolor[rgb]{0.00,0.50,0.00}{#1}}
\newcommand{\CharTok}[1]{\textcolor[rgb]{0.25,0.44,0.63}{#1}}
\newcommand{\CommentTok}[1]{\textcolor[rgb]{0.38,0.63,0.69}{\textit{#1}}}
\newcommand{\CommentVarTok}[1]{\textcolor[rgb]{0.38,0.63,0.69}{\textbf{\textit{#1}}}}
\newcommand{\ConstantTok}[1]{\textcolor[rgb]{0.53,0.00,0.00}{#1}}
\newcommand{\ControlFlowTok}[1]{\textcolor[rgb]{0.00,0.44,0.13}{\textbf{#1}}}
\newcommand{\DataTypeTok}[1]{\textcolor[rgb]{0.56,0.13,0.00}{#1}}
\newcommand{\DecValTok}[1]{\textcolor[rgb]{0.25,0.63,0.44}{#1}}
\newcommand{\DocumentationTok}[1]{\textcolor[rgb]{0.73,0.13,0.13}{\textit{#1}}}
\newcommand{\ErrorTok}[1]{\textcolor[rgb]{1.00,0.00,0.00}{\textbf{#1}}}
\newcommand{\ExtensionTok}[1]{#1}
\newcommand{\FloatTok}[1]{\textcolor[rgb]{0.25,0.63,0.44}{#1}}
\newcommand{\FunctionTok}[1]{\textcolor[rgb]{0.02,0.16,0.49}{#1}}
\newcommand{\ImportTok}[1]{\textcolor[rgb]{0.00,0.50,0.00}{\textbf{#1}}}
\newcommand{\InformationTok}[1]{\textcolor[rgb]{0.38,0.63,0.69}{\textbf{\textit{#1}}}}
\newcommand{\KeywordTok}[1]{\textcolor[rgb]{0.00,0.44,0.13}{\textbf{#1}}}
\newcommand{\NormalTok}[1]{#1}
\newcommand{\OperatorTok}[1]{\textcolor[rgb]{0.40,0.40,0.40}{#1}}
\newcommand{\OtherTok}[1]{\textcolor[rgb]{0.00,0.44,0.13}{#1}}
\newcommand{\PreprocessorTok}[1]{\textcolor[rgb]{0.74,0.48,0.00}{#1}}
\newcommand{\RegionMarkerTok}[1]{#1}
\newcommand{\SpecialCharTok}[1]{\textcolor[rgb]{0.25,0.44,0.63}{#1}}
\newcommand{\SpecialStringTok}[1]{\textcolor[rgb]{0.73,0.40,0.53}{#1}}
\newcommand{\StringTok}[1]{\textcolor[rgb]{0.25,0.44,0.63}{#1}}
\newcommand{\VariableTok}[1]{\textcolor[rgb]{0.10,0.09,0.49}{#1}}
\newcommand{\VerbatimStringTok}[1]{\textcolor[rgb]{0.25,0.44,0.63}{#1}}
\newcommand{\WarningTok}[1]{\textcolor[rgb]{0.38,0.63,0.69}{\textbf{\textit{#1}}}}
\usepackage{longtable,booktabs,array}
\usepackage{caption}
\newcounter{none} % for unnumbered tables
\newcolumntype{L}[1]{>{\raggedright\arraybackslash\hyphenpenalty=10000\exhyphenpenalty=50}p{#1}}
\newcolumntype{C}[1]{>{\centering\arraybackslash}p{#1}}
\usepackage{calc} % for calculating minipage widths
% Correct order of tables after \paragraph or \subparagraph
\usepackage{etoolbox}
\makeatletter
\patchcmd\longtable{\par}{\if@noskipsec\mbox{}\fi\par}{}{}
\makeatother
% Allow footnotes in longtable head/foot
\IfFileExists{footnotehyper.sty}{\usepackage{footnotehyper}}{\usepackage{footnote}}
\makesavenoteenv{longtable}
\providecommand{\tabularnewline}{\\}
\newlength{\thesisdefaulttabcolsep}
\setlength{\thesisdefaulttabcolsep}{2pt}
\usepackage{graphicx}
\usepackage{tikz}
\usetikzlibrary{calc}
\makeatletter
\newsavebox\pandoc@box
\newcommand*\pandocbounded[1]{% scales image to fit in text height/width
  \sbox\pandoc@box{#1}%
  \Gscale@div\@tempa{\textheight}{\dimexpr\ht\pandoc@box+\dp\pandoc@box\relax}%
  \Gscale@div\@tempb{\linewidth}{\wd\pandoc@box}%
  \ifdim\@tempb\p@<\@tempa\p@\let\@tempa\@tempb\fi% select the smaller of both
  \ifdim\@tempa\p@<\p@\makebox[\linewidth][c]{\scalebox{\@tempa}{\usebox\pandoc@box}}%
  \else\makebox[\linewidth][c]{\usebox\pandoc@box}%
  \fi%
}
% Set default figure placement to htbp
\def\fps@figure{htbp}
\makeatother
\usepackage{xstring}
\newcommand{\safeincludesvg}[1]{%
  \StrSubstitute{#1}{.svg}{.png}[\safeincludepngpath]%
  \StrSubstitute{#1}{./assets/figures/}{./svg-inkscape/}[\safeincludepdfbasepath]%
  \StrSubstitute{\safeincludepdfbasepath}{.svg}{_svg-raw.pdf}[\safeincludepdfpath]%
  \IfFileExists{\safeincludepngpath}{%
    \includegraphics[width=0.92\linewidth,height=0.70\textheight,keepaspectratio]{\safeincludepngpath}%
  }{%
    \IfFileExists{\safeincludepdfpath}{%
      \includegraphics[width=0.92\linewidth,height=0.70\textheight,keepaspectratio]{\safeincludepdfpath}%
    }{%
      \fbox{\parbox{0.9\linewidth}{\centering Missing figure assets for:\\ \texttt{\detokenize{#1}}}}%
    }%
  }%
}
\newcommand{\safeincludesvgsmall}[1]{%
  \StrSubstitute{#1}{.svg}{.png}[\safeincludepngpath]%
  \StrSubstitute{#1}{./assets/figures/}{./svg-inkscape/}[\safeincludepdfbasepath]%
  \StrSubstitute{\safeincludepdfbasepath}{.svg}{_svg-raw.pdf}[\safeincludepdfpath]%
  \IfFileExists{\safeincludepngpath}{%
    \includegraphics[width=0.76\linewidth,height=0.56\textheight,keepaspectratio]{\safeincludepngpath}%
  }{%
    \IfFileExists{\safeincludepdfpath}{%
      \includegraphics[width=0.76\linewidth,height=0.56\textheight,keepaspectratio]{\safeincludepdfpath}%
    }{%
      \fbox{\parbox{0.9\linewidth}{\centering Missing figure assets for:\\ \texttt{\detokenize{#1}}}}%
    }%
  }%
}
\newcommand{\safeincludesvgcompact}[1]{%
  \StrSubstitute{#1}{.svg}{.png}[\safeincludepngpath]%
  \StrSubstitute{#1}{./assets/figures/}{./svg-inkscape/}[\safeincludepdfbasepath]%
  \StrSubstitute{\safeincludepdfbasepath}{.svg}{_svg-raw.pdf}[\safeincludepdfpath]%
  \IfFileExists{\safeincludepngpath}{%
    \includegraphics[width=0.68\linewidth,height=0.46\textheight,keepaspectratio]{\safeincludepngpath}%
  }{%
    \IfFileExists{\safeincludepdfpath}{%
      \includegraphics[width=0.68\linewidth,height=0.46\textheight,keepaspectratio]{\safeincludepdfpath}%
    }{%
      \fbox{\parbox{0.9\linewidth}{\centering Missing figure assets for:\\ \texttt{\detokenize{#1}}}}%
    }%
  }%
}
\newcommand{\safeincludesvgnarrow}[1]{%
  \StrSubstitute{#1}{.svg}{.png}[\safeincludepngpath]%
  \StrSubstitute{#1}{./assets/figures/}{./svg-inkscape/}[\safeincludepdfbasepath]%
  \StrSubstitute{\safeincludepdfbasepath}{.svg}{_svg-raw.pdf}[\safeincludepdfpath]%
  \IfFileExists{\safeincludepngpath}{%
    \includegraphics[width=0.56\linewidth,height=0.34\textheight,keepaspectratio]{\safeincludepngpath}%
  }{%
    \IfFileExists{\safeincludepdfpath}{%
      \includegraphics[width=0.56\linewidth,height=0.34\textheight,keepaspectratio]{\safeincludepdfpath}%
    }{%
      \fbox{\parbox{0.9\linewidth}{\centering Missing figure assets for:\\ \texttt{\detokenize{#1}}}}%
    }%
  }%
}
\newcommand{\safeincludesvgchart}[1]{%
  \StrSubstitute{#1}{.svg}{.png}[\safeincludepngpath]%
  \StrSubstitute{#1}{./assets/figures/}{./svg-inkscape/}[\safeincludepdfbasepath]%
  \StrSubstitute{\safeincludepdfbasepath}{.svg}{_svg-raw.pdf}[\safeincludepdfpath]%
  \IfFileExists{\safeincludepngpath}{%
    \includegraphics[width=0.72\linewidth,height=0.28\textheight,keepaspectratio]{\safeincludepngpath}%
  }{%
    \IfFileExists{\safeincludepdfpath}{%
      \includegraphics[width=0.72\linewidth,height=0.28\textheight,keepaspectratio]{\safeincludepdfpath}%
    }{%
      \fbox{\parbox{0.9\linewidth}{\centering Missing figure assets for:\\ \texttt{\detokenize{#1}}}}%
    }%
  }%
}
\setlength{\emergencystretch}{10em} % reduce overfull lines
\setlength{\tabcolsep}{\thesisdefaulttabcolsep}
\setlength{\arrayrulewidth}{0.5pt}
\captionsetup[table]{justification=centering,singlelinecheck=false,labelfont=it,textfont=it}
\captionsetup[longtable]{justification=centering,singlelinecheck=false,labelfont=it,textfont=it}
\setlength{\LTcapwidth}{\linewidth}
\setlength{\LTleft}{0pt}
\setlength{\LTright}{0pt}
\setlength{\hfuzz}{6pt}
\hbadness=10000
\setcounter{topnumber}{2}
\setcounter{bottomnumber}{2}
\setcounter{totalnumber}{4}
\renewcommand{\topfraction}{0.85}
\renewcommand{\bottomfraction}{0.75}
\renewcommand{\textfraction}{0.1}
\renewcommand{\floatpagefraction}{0.75}
\usepackage{setspace}
\setstretch{1.08}
\setlength{\parindent}{0pt}
\setlength{\parskip}{4pt}
\setcounter{tocdepth}{3}
\renewcommand{\arraystretch}{1.08}
\usepackage{titlesec}
\titleformat{\section}{\bfseries\Large\centering}{}{0pt}{}
\titlespacing*{\section}{0pt}{12pt}{6pt}
\titleformat{\subsection}{\bfseries\normalsize}{}{0pt}{}
\titlespacing*{\subsection}{0pt}{12pt}{0pt}
\newcommand{\thesisfrontmatterentry}[1]{%
  \phantomsection
  \addcontentsline{toc}{section}{#1}%
}
\newcommand{\majorheading}[1]{%
  \clearpage
  \phantomsection
  \section*{#1}%
  \addcontentsline{toc}{section}{#1}%
}
\newcommand{\chapterheading}[1]{%
  \clearpage
  \phantomsection
  \section*{#1}%
  \addcontentsline{toc}{section}{#1}%
  \setcounter{subsection}{0}%
}
\newcommand{\appendixsection}[1]{\clearpage\section{#1}}
\newcommand{\frontmatterstart}{\clearpage\hypersetup{pageanchor=true}\pagenumbering{roman}\setcounter{page}{1}}
\newcommand{\mainmatterstart}{\clearpage\pagenumbering{arabic}\setcounter{page}{1}}
\newcommand{\checkbox}{\(\square\)}
\newcommand{\thesisspacer}{\hspace{2em}}
\newcommand{\figref}[1]{Hình~\ref{#1}}
\newcommand{\tabref}[1]{Bảng~\ref{#1}}
\definecolor{thesisphenikaablue}{RGB}{42, 63, 130}
\newcommand{\thesislogopath}{./assets/brand/phenikaa-university-logo.jpeg}
\newcommand{\thesisfrontseparator}{\begin{center}\rule{0.5\linewidth}{0.5pt}\end{center}}
\newcommand{\thesisdrawtemplateframe}{%
  \begin{tikzpicture}[remember picture,overlay]
    \draw[line width=1.4pt] ($(current page.north west)+(1.00cm,-1.00cm)$) rectangle ($(current page.south east)+(-1.00cm,1.00cm)$);
    \draw[line width=0.5pt] ($(current page.north west)+(1.14cm,-1.14cm)$) rectangle ($(current page.south east)+(-1.14cm,1.14cm)$);
  \end{tikzpicture}%
}
\newcommand{\thesisdrawsidebarframe}{%
  \begin{tikzpicture}[remember picture,overlay]
    \draw[line width=0.5pt] ($(current page.north west)+(1.35cm,-1.35cm)$) rectangle ($(current page.south east)+(-1.35cm,1.35cm)$);
    \draw[line width=0.5pt] ($(current page.north west)+(3.65cm,-1.35cm)$) -- ($(current page.south west)+(3.65cm,1.35cm)$);
  \end{tikzpicture}%
}
\newcommand{\thesiscoverstart}{%
  \clearpage
  \newgeometry{top=1.5cm,bottom=1.5cm,left=1.6cm,right=1.6cm}
  \thispagestyle{empty}
}
\newcommand{\thesiscoverend}{%
  \clearpage
  \restoregeometry
}
\newcommand{\thesisofficialheader}{%
  \noindent
  \begin{minipage}[t]{0.48\linewidth}
    \centering
    \textbf{BỘ GIÁO DỤC VÀ ĐÀO TẠO}\\
    \textbf{ĐẠI HỌC PHENIKAA}\\
    \rule{0.72\linewidth}{0.6pt}
  \end{minipage}\hfill
  \begin{minipage}[t]{0.48\linewidth}
    \centering
    \textbf{CỘNG HÒA XÃ HỘI CHỦ NGHĨA VIỆT NAM}\\
    \textbf{Độc lập - Tự do - Hạnh phúc}\\
    \rule{0.72\linewidth}{0.6pt}
  \end{minipage}
  \par\vspace{0.6cm}
}
\newcommand{\thesissignaturepair}[4]{%
  \vspace{0.8cm}
  \noindent
  \begin{minipage}[t]{0.48\linewidth}
    \centering
    \textbf{#1}\\
    \emph{(Ký, ghi rõ họ tên)}\\[2.8cm]
    #2
  \end{minipage}\hfill
  \begin{minipage}[t]{0.48\linewidth}
    \centering
    \textbf{#3}\\
    \emph{(Ký, ghi rõ họ tên)}\\[2.8cm]
    #4
  \end{minipage}
}
\renewcommand{\figurename}{Hình}
\renewcommand{\tablename}{Bảng}
\renewcommand{\listtablename}{DANH MỤC BẢNG - LIST OF TABLES}
\renewcommand{\listfigurename}{DANH MỤC HÌNH ẢNH VÀ ĐỒ THỊ - LIST OF PICTURES AND GRAPHS}
\makeatletter
\newcommand{\thesissetfigureid}[2]{%
  \renewcommand{\thefigure}{#1}%
  \renewcommand{\theHfigure}{#2}%
}
\newcommand{\thesissettableid}[2]{%
  \renewcommand{\thetable}{#1}%
  \renewcommand{\theHtable}{#2}%
}
\newcommand{\thesistableofcontentsbody}{\@starttoc{toc}}
\newcommand{\thesislistoftablesbody}{\@starttoc{lot}}
\newcommand{\thesislistoffiguresbody}{\@starttoc{lof}}
\renewcommand*\l@table{\@dottedtocline{1}{0em}{5.8em}}
\renewcommand*\l@figure{\@dottedtocline{1}{0em}{5.8em}}
\makeatother
\newcounter{ieeeref}
\newcommand{\ieeebibitem}[2]{%
  \refstepcounter{ieeeref}\noindent[\theieeeref]\label{#1} #2\par
}
\newcommand{\ieeecite}[1]{[\ref{#1}]}
\providecommand{\textcite}[1]{\ieeecite{#1}}
\providecommand{\parencite}[1]{\ieeecite{#1}}
\providecommand{\tightlist}{%
  \setlength{\itemsep}{0pt}\setlength{\parskip}{0pt}}
\usepackage{bookmark}
\IfFileExists{xurl.sty}{\usepackage{xurl}}{} % add URL line breaks if available
\hyphenation{Mos-quit-to Post-gre-SQL Vic-to-ria-Met-rics Vic-to-ria-Logs Au-then-ti-ca-tion Ku-ber-ne-tes Os-cil-lo-scope Dash-board Web-Socket}
\urlstyle{same}
\hypersetup{
  hidelinks,
  linktoc=all,
  bookmarksnumbered=true,
  pdfcreator={LaTeX via pandoc}}

\author{}
\date{}

\begin{document}
\hypersetup{pageanchor=false}
\pagenumbering{gobble}

\thesiscoverstart
\thesisdrawsidebarframe
\noindent
\begin{minipage}[t][0.86\textheight][t]{0.16\textwidth}
  \centering
  \vspace*{1.0cm}
  \rotatebox{90}{\textbf{\Large LÊ TRỌNG AN}}\\[1.5cm]
  \rotatebox{90}{\textbf{\large KỸ THUẬT CƠ ĐIỆN TỬ}}
\end{minipage}\hfill
\begin{minipage}[t][0.86\textheight][t]{0.77\textwidth}
  \begin{center}
    \textbf{\Large BỘ GIÁO DỤC VÀ ĐÀO TẠO}\\[0.2cm]
    \textbf{\Large ĐẠI HỌC PHENIKAA}\\[-0.05cm]
    \rule{0.34\linewidth}{0.6pt}

    \vspace{1.7cm}
    \includegraphics[width=0.36\linewidth]{\thesislogopath}

    \vspace{1.5cm}
    {\fontsize{28}{30}\selectfont\textbf{ĐỒ ÁN TỐT NGHIỆP}}\\[1.0cm]
    {\fontsize{18}{22}\selectfont\textbf{THIẾT KẾ HỆ THỐNG IOT CHO ỨNG DỤNG QUẢN LÝ}}\\
    {\fontsize{18}{22}\selectfont\textbf{PHƯƠNG TIỆN GIAO THÔNG TRONG LĨNH VỰC}}\\
    {\fontsize{18}{22}\selectfont\textbf{CHO THUÊ XE TỰ LÁI}}

    \vfill
    \begin{minipage}{0.90\linewidth}
      \raggedright
      \textbf{Sinh viên:} Lê Trọng An\\[0.15cm]
      \textbf{Mã số sinh viên:} 21010389 \hfill \textbf{Khóa:} K15\\[0.15cm]
      \textbf{Ngành:} Kỹ thuật Cơ Điện tử \hfill \textbf{Hệ:} Đại học chính quy\\[0.15cm]
      \textbf{Giảng viên hướng dẫn:} TS. Nguyễn Đức Nam
    \end{minipage}

    \vspace{1.2cm}
    \textbf{\Large Hà Nội – 2026}
  \end{center}
\end{minipage}
\thesiscoverend

\thesiscoverstart
\thesisdrawtemplateframe
\begin{center}
  \textbf{\Large BỘ GIÁO DỤC VÀ ĐÀO TẠO}\\[0.15cm]
  \textbf{\Large ĐẠI HỌC PHENIKAA}\\[-0.05cm]
  \rule{0.28\linewidth}{0.6pt}

  \vspace{1.6cm}
  \includegraphics[width=0.30\linewidth]{\thesislogopath}

  \vspace{1.4cm}
  {\fontsize{27}{30}\selectfont\textbf{ĐỒ ÁN TỐT NGHIỆP}}\\[0.9cm]
  {\fontsize{18}{22}\selectfont\textbf{THIẾT KẾ HỆ THỐNG IOT CHO ỨNG DỤNG QUẢN LÝ}}\\
  {\fontsize{18}{22}\selectfont\textbf{PHƯƠNG TIỆN GIAO THÔNG TRONG LĨNH VỰC}}\\
  {\fontsize{18}{22}\selectfont\textbf{CHO THUÊ XE TỰ LÁI}}

  \vfill
  \begin{minipage}{0.78\linewidth}
    \raggedright
    \textbf{Sinh viên:} Lê Trọng An\\[0.18cm]
    \textbf{Mã số sinh viên:} 21010389 \hfill \textbf{Khóa:} K15\\[0.18cm]
    \textbf{Ngành:} Kỹ thuật Cơ Điện tử \hfill \textbf{Hệ:} Đại học chính quy\\[0.18cm]
    \textbf{Giảng viên hướng dẫn:} TS. Nguyễn Đức Nam
  \end{minipage}

  \vspace{1.3cm}
  \textbf{\Large Hà Nội – 2026}
\end{center}
\thesiscoverend

\thesiscoverstart
\thesisdrawtemplateframe
\begin{center}
  {\color{thesisphenikaablue}\textbf{\Large PHENIKAA UNIVERSITY}}\\[0.18cm]
  {\color{thesisphenikaablue}\textbf{\Large FACULTY OF MECHANICAL ENGINEERING AND MECHATRONICS}}\\[0.25cm]
  \includegraphics[width=0.28\linewidth]{\thesislogopath}

  \vspace{1.0cm}
  {\fontsize{22}{26}\selectfont\textbf{MEM710071 – CAPSTONE DESIGN 2}}\\[0.55cm]
  {\fontsize{18}{22}\selectfont\textbf{DESIGN OF AN IOT SYSTEM FOR VEHICLE MANAGEMENT}}\\
  {\fontsize{18}{22}\selectfont\textbf{IN SELF-DRIVE CAR RENTAL SERVICES}}\\[0.35cm]
  {\fontsize{15}{18}\selectfont\textbf{Semester: I Academic year: 2025--2026}}

  \vfill
  \begin{minipage}{0.76\linewidth}
    \raggedright
    \textbf{Project Team:}
    \begin{center}
      \begin{tabular}{|p{0.42\linewidth}|p{0.20\linewidth}|p{0.12\linewidth}|}
        \hline
        \centering\textbf{Student name} & \centering\textbf{Student ID} & \centering\textbf{Cohort}\tabularnewline
        \hline
        Le Trong An & \centering 21010389 & \centering K15\tabularnewline
        \hline
      \end{tabular}
    \end{center}
    \vspace{0.2cm}
    \textbf{Course Instructor:} Dr. Nguyen Duc Nam
  \end{minipage}

  \vfill
  \begin{minipage}{0.88\linewidth}
    \centering
    This report is submitted to the Faculty of Mechanical Engineering and Mechatronics,\\
    School of Engineering, Phenikaa University in partial fulfillment of the requirements\\
    of the Graduation Thesis course MEM710071.
  \end{minipage}

  \vspace{0.8cm}
  \textbf{\Large HANOI – 2026}
\end{center}
\thesiscoverend

\thesiscoverstart
\thesisdrawtemplateframe
\begin{center}
  {\color{thesisphenikaablue}\textbf{\Large ĐẠI HỌC PHENIKAA}}\\[0.15cm]
  {\color{thesisphenikaablue}\textbf{\Large KHOA CƠ KHÍ – CƠ ĐIỆN TỬ}}\\[0.25cm]
  \includegraphics[width=0.27\linewidth]{\thesislogopath}

  \vspace{0.8cm}
  {\fontsize{21}{25}\selectfont\textbf{MEM710071}}\\[0.15cm]
  {\fontsize{26}{30}\selectfont\textbf{ĐỒ ÁN TỐT NGHIỆP}}\\[0.45cm]
  {\fontsize{18}{22}\selectfont\textbf{THIẾT KẾ HỆ THỐNG IOT CHO ỨNG DỤNG QUẢN LÝ}}\\
  {\fontsize{18}{22}\selectfont\textbf{PHƯƠNG TIỆN GIAO THÔNG TRONG LĨNH VỰC CHO THUÊ XE TỰ LÁI}}\\[0.35cm]
  {\fontsize{15}{18}\selectfont\textbf{Học kỳ I Năm học 2025--2026}}

  \vfill
  \begin{center}
    \begin{tabular}{|p{0.34\linewidth}|p{0.24\linewidth}|p{0.10\linewidth}|}
      \hline
      \centering\textbf{Họ và tên sinh viên} & \centering\textbf{Mã số sinh viên} & \centering\textbf{Lớp}\tabularnewline
      \hline
      Lê Trọng An & \centering 21010389 & \centering K15-KTCĐT2\tabularnewline
      \hline
    \end{tabular}
  \end{center}

  \vspace{0.45cm}
  \textbf{Giảng viên hướng dẫn:} TS. Nguyễn Đức Nam

  \vfill
  \textbf{\Large HÀ NỘI, 2026}
\end{center}
\thesiscoverend

\frontmatterstart

\thispagestyle{empty}
\thesisofficialheader
\begin{center}
  \textbf{\Large NHẬN XÉT ĐỒ ÁN TỐT NGHIỆP CỦA GIẢNG VIÊN HƯỚNG DẪN}
\end{center}
\thesisfrontmatterentry{NHẬN XÉT CỦA GIẢNG VIÊN HƯỚNG DẪN}

\textbf{Giảng viên hướng dẫn:} TS. Nguyễn Đức Nam \hfill \textbf{Khoa:} Cơ khí -- Cơ Điện tử.\\
\textbf{Tên đề tài:} Thiết kế hệ thống IoT cho ứng dụng quản lý phương tiện giao thông trong lĩnh vực cho thuê xe tự lái.\\
\textbf{Sinh viên thực hiện:} Lê Trọng An. \hfill \textbf{Lớp:} K15-KTCĐT2.

\begin{center}
  \textbf{NỘI DUNG NHẬN XÉT}
\end{center}
\textbf{I. Nhận xét Đồ án tốt nghiệp:}\\
- Nhận xét về hình thức: \dotfill\\[0.55cm]
- Tính cấp thiết của đề tài: \dotfill\\[0.55cm]
- Mục tiêu của đề tài: \dotfill\\[0.55cm]
- Nội dung của đề tài: \dotfill\\[0.55cm]
- Tài liệu tham khảo: \dotfill\\[0.55cm]
- Phương pháp nghiên cứu: \dotfill\\[0.55cm]
- Tính sáng tạo và ứng dụng: \dotfill\\[0.55cm]
\textbf{II. Nhận xét về tinh thần và thái độ làm việc của sinh viên:} \dotfill\\[0.55cm]
\textbf{III. Kết quả đạt được:} \dotfill\\[0.55cm]
\textbf{IV. Kết luận:} Đồng ý cho bảo vệ: \checkbox \thesisspacer Không đồng ý cho bảo vệ: \checkbox

\begin{flushright}
  Hà Nội, ngày \ldots, tháng \ldots năm 20\ldots\\
  \textbf{GIẢNG VIÊN HƯỚNG DẪN}\\
  \emph{(Ký, ghi rõ họ tên)}\\[2.6cm]
  TS. Nguyễn Đức Nam
\end{flushright}
\clearpage

\thispagestyle{empty}
\thesisofficialheader
\begin{center}
  \textbf{\Large NHẬN XÉT ĐỒ ÁN TỐT NGHIỆP CỦA GIẢNG VIÊN PHẢN BIỆN}
\end{center}
\thesisfrontmatterentry{NHẬN XÉT CỦA GIẢNG VIÊN PHẢN BIỆN}

\textbf{Giảng viên phản biện:} Theo quyết định phân công phản biện của Khoa \\
\textbf{Tên đề tài:} Thiết kế hệ thống IoT cho ứng dụng quản lý phương tiện giao thông trong lĩnh vực cho thuê xe tự lái.\\
\textbf{Sinh viên thực hiện:} Lê Trọng An \hfill \textbf{Lớp:} K15-KTCĐT2\\
\textbf{Giảng viên hướng dẫn:} TS. Nguyễn Đức Nam

\begin{center}
  \textbf{NỘI DUNG NHẬN XÉT}
\end{center}
\textbf{I. Nhận xét Đồ án tốt nghiệp:}\\
- Bố cục, hình thức trình bày: \dotfill\\[0.55cm]
- Đảm bảo tính cấp thiết, hiện đại, không trùng lặp: \dotfill\\[0.55cm]
- Nội dung: \dotfill\\[0.55cm]
- Mức độ thực hiện: \dotfill\\[0.55cm]
\textbf{II. Kết quả đạt được:} \dotfill\\[0.55cm]
\textbf{III. Ưu nhược điểm:} \dotfill\\[0.55cm]
\textbf{IV. Kết luận:} Đồng ý cho bảo vệ: \checkbox \thesisspacer Không đồng ý cho bảo vệ: \checkbox

\begin{flushright}
  Hà Nội, ngày \ldots, tháng \ldots năm 20\ldots\\
  \textbf{GIẢNG VIÊN PHẢN BIỆN}\\
  \emph{(Ký, ghi rõ họ tên)}\\[2.6cm]
  Theo quyết định phân công phản biện của Khoa
\end{flushright}
\clearpage

\thispagestyle{empty}
\thesisofficialheader
\begin{center}
  \textbf{\Large BIÊN BẢN ĐÁNH GIÁ ĐỒ ÁN TỐT NGHIỆP}
\end{center}
\thesisfrontmatterentry{BIÊN BẢN ĐÁNH GIÁ ĐỒ ÁN TỐT NGHIỆP}

\textbf{I. Thông tin chung}
\begin{enumerate}
  \item Thời gian: từ \ldots giờ \ldots đến \ldots giờ \ldots ngày \ldots/\ldots/20\ldots
  \item Địa điểm: \dotfill
  \item Họ và tên sinh viên: Lê Trọng An \hfill Mã số SV: 21010389
  \item Lớp: K15-KTCĐT2 \hfill Ngành: Kỹ thuật Cơ Điện tử \hfill Khóa: K15
  \item Tên đề tài: Thiết kế hệ thống IoT cho ứng dụng quản lý phương tiện giao thông trong lĩnh vực cho thuê xe tự lái.
  \item Giảng viên hướng dẫn: TS. Nguyễn Đức Nam
\end{enumerate}

\textbf{II. Thành phần hội đồng}
\begin{enumerate}
  \item \dotfill Chủ tịch
  \item \dotfill Thư ký
  \item \dotfill Giảng viên phản biện
  \item \dotfill Ủy viên 1
\end{enumerate}

\textbf{III. Tổng hợp câu hỏi của thành viên hội đồng:} \dotfill\\[0.6cm]
\textbf{IV. Tổng hợp nội dung trả lời của sinh viên:} \dotfill\\[0.6cm]
\textbf{V. Nội dung đánh giá của Hội đồng}
\begin{enumerate}
  \item Ý nghĩa của đồ án: \dotfill
  \item Về nội dung, kết cấu của đồ án: \dotfill
  \item Phương pháp nghiên cứu: \dotfill
  \item Các kết quả nghiên cứu đạt được: \dotfill
\end{enumerate}
\textbf{VI. Điểm kết luận:} \dotfill\ (Bằng chữ: \dotfill)

\thesissignaturepair{CHỦ TỊCH HỘI ĐỒNG}{}{THƯ KÝ}{}
\clearpage

\thispagestyle{empty}
\begin{center}
  \textbf{\Large GIẢI TRÌNH CÁC CHỈNH SỬA}\\[0.2cm]
  \textit{(Nếu có)}
\end{center}
\thesisfrontmatterentry{GIẢI TRÌNH CÁC CHỈNH SỬA}

\vspace{0.5cm}
\begin{center}
  Chưa phát sinh nội dung cần giải trình.
\end{center}
\clearpage

\majorheading{LỜI CAM ĐOAN}

Tên tôi là: Lê Trọng An.\\
Mã sinh viên: 21010389 \hfill Lớp: K15-KTCĐT2\\
Ngành: Kỹ thuật Cơ Điện tử.

Tôi đã thực hiện đồ án tốt nghiệp với đề tài: Thiết kế hệ thống IoT cho ứng dụng quản lý phương tiện giao thông trong lĩnh vực cho thuê xe tự lái.

Tôi xin cam đoan đây là đề tài nghiên cứu của riêng tôi và được sự hướng dẫn của TS. Nguyễn Đức Nam.

Các nội dung nghiên cứu, kết quả trong đề tài này là trung thực và chưa được công bố dưới bất kỳ hình thức nào. Nếu phát hiện có bất kỳ hình thức gian lận nào, tôi xin hoàn toàn chịu trách nhiệm trước pháp luật.

\begin{flushright}
  Hà Nội, ngày \ldots, tháng \ldots năm 20\ldots
\end{flushright}

\thesissignaturepair{GIẢNG VIÊN HƯỚNG DẪN}{TS. Nguyễn Đức Nam}{SINH VIÊN}{Lê Trọng An}

\majorheading{TÓM TẮT ĐỒ ÁN TỐT NGHIỆP - ABSTRACT}

Sự phát triển nhanh của dịch vụ cho thuê xe tự lái tại Việt Nam khiến
nhu cầu giám sát phương tiện không còn dừng ở việc biết xe đang ở đâu,
mà còn ở việc hiểu xe đang vận hành ra sao và có dấu hiệu bất thường hay
không. Từ nhu cầu đó, đồ án này trình bày quá trình thiết kế và triển
khai một hệ thống IoT giám sát phương tiện theo hướng end-to-end, bao
gồm thiết bị gắn trên xe, hạ tầng cloud và bảng điều khiển web, nhằm hỗ
trợ doanh nghiệp theo dõi vị trí, nắm bắt trạng thái kỹ thuật và phản
ứng sớm với các rủi ro vận hành.

Ở tầng thiết bị, hệ thống được xây dựng trên bo mạch PCB tự thiết kế.
Trong đó, ESP32-S3 giữ vai trò bộ điều khiển trung tâm; SIM7600CE-T đảm
nhiệm truyền dữ liệu 4G và định vị vệ tinh GNSS; LIS3DSH theo dõi rung
động để nhận biết chuyển động khi xe đỗ; còn adapter vgate iCar Pro cho
phép đọc dữ liệu OBD2 từ xe qua BLE. Kiến trúc nguồn được tổ chức theo
nhiều nhánh để thiết bị vẫn duy trì hoạt động khi xe tắt máy, đồng thời
hạn chế ảnh hưởng tới ắc quy chính.

Ở tầng phần mềm, dữ liệu từ thiết bị được gửi bằng MQTT 5.0 tới EMQX.
Có thể hiểu MQTT là một cơ chế nhắn tin nhẹ, phù hợp với thiết bị IoT
cần gửi nhiều bản tin nhỏ qua mạng di động. Sau đó, MQTT Bridge tiếp
nhận và phân luồng dữ liệu tới từng hệ lưu trữ phù hợp: PostgreSQL cho
dữ liệu quan hệ và nhật ký OTA, VictoriaMetrics cho dữ liệu chuỗi thời
gian, VictoriaLogs cho nhật ký sự kiện. API server được xây dựng bằng
Express.js và TypeScript, đóng vai trò kết nối giữa thiết bị, dữ liệu
lưu trữ và giao diện khai thác.

Bảng điều khiển web được phát triển bằng Next.js 15 và React 19, hiển
thị bản đồ, biểu đồ và trạng thái phương tiện gần như theo thời gian
thực thông qua Socket.IO. Nhờ đó, người vận hành không chỉ quan sát vị
trí xe mà còn có thể theo dõi lịch sử hành trình, cảnh báo và trạng thái
thiết bị trên cùng một giao diện thống nhất.

Kết quả đạt được là một hệ thống IoT giám sát phương tiện hoàn chỉnh từ
phần cứng đến phần mềm, có khả năng theo dõi vị trí thời gian thực, đọc
dữ liệu chẩn đoán OBD2, thiết lập hàng rào địa lý, phát sinh cảnh báo tự
động, cập nhật firmware từ xa và hỗ trợ quản lý đội xe. Quan trọng hơn,
hệ thống cho thấy một hướng triển khai phù hợp với bối cảnh doanh nghiệp
cho thuê xe tự lái trong nước: chi phí hợp lý, có thể tùy biến và đủ dư
địa để mở rộng trong các giai đoạn sau.

\textbf{Từ khóa:} IoT, giám sát phương tiện, GPS, OBD2, MQTT, ESP32-S3,
thời gian thực, hàng rào địa lý

\textbf{Abstract (English)}

As the self-drive car rental market in Vietnam grows, the need for
remote vehicle monitoring and real-time fleet management becomes
increasingly practical for operators. This thesis presents the design
and development of an end-to-end IoT vehicle tracking system, covering
the onboard device, cloud services, and operational dashboard for
real-time tracking, diagnostic data collection, and alerting.

On the hardware side, the tracker is implemented as a custom integrated
PCB rather than an assembly of off-the-shelf development boards. The
main board integrates ESP32-S3 (MCU), SIMCom SIM7600CE-T (LTE+GNSS),
LIS3DSH (IMU), DS3231M (RTC), W25Q128 (SPI Flash), and power blocks based
on MP2482, TPS54231, AP2112, SX1308, and TP4056. The modem is configured
to Auto mode (\texttt{AT+CNMP=2}) so it can switch between LTE/UMTS/GSM
automatically, uses the default APN \texttt{internet}, and delivers GNSS
data via \texttt{AT+CGNSINF} (with optional NMEA streaming through
\texttt{AT+CGNSTST}) on the same UART1 channel. The device reads engine
diagnostic data via OBD2 through a vgate iCar Pro adapter over Bluetooth
Low Energy (BLE), where this adapter is an external peripheral outside
the main PCB. The power architecture uses a 5V main rail, an
approximately 4V modem rail, a 3.3V logic rail, and a 18650 1S backup
cell to maintain operation when the vehicle engine is off.

On the software side, device data is transmitted through MQTT 5.0 via
EMQX with per-device ACL authorization. The MQTT Bridge validates and
routes incoming messages to specialized storage systems: PostgreSQL for
relational data and OTA update logs, VictoriaMetrics for time-series
telemetry, and VictoriaLogs for operational event logs. The API server
is built with Express.js and TypeScript following Domain-Driven Design
(DDD), using database-backed session token authentication and providing
OTA orchestration endpoints for devices.

The web interface is developed using Next.js 15 and React 19, providing
a real-time dashboard with Leaflet maps, ECharts visualizations, and
WebSocket connectivity via Socket.IO for instant vehicle position and
status updates. The entire system is containerized and deployed using
Docker, ensuring consistency between development and production
environments.

The result is a complete end-to-end IoT vehicle tracking system, from
hardware to software, capable of real-time location tracking, OBD2
diagnostic data retrieval, geofencing, automated alerts, remote firmware
updates, and fleet-management support for self-drive car rental
services.

\textbf{Keywords:} IoT, vehicle tracking, GPS, OBD2, MQTT, ESP32-S3,
real-time, geofencing

\majorheading{LỜI CẢM ƠN - ACKNOWLEDGEMENTS}\label{lux1eddi-cux1ea3m-ux1a1n---acknowledgements}

Trước tiên, tôi xin gửi lời cảm ơn chân thành và sâu sắc nhất đến TS.
Nguyễn Đức Nam --- người đã trực tiếp hướng dẫn, định hướng và tận tình
hỗ trợ tôi trong suốt quá trình thực hiện đồ án tốt nghiệp này. Những
góp ý chuyên môn và sự động viên của thầy là nguồn động lực quan trọng
giúp tôi hoàn thành đề tài.

Tôi xin trân trọng cảm ơn quý thầy cô trong Khoa Cơ khí - Cơ Điện tử,
Trường Kỹ thuật - Đại học Phenikaa đã truyền đạt những kiến thức nền
tảng vững chắc trong suốt quá trình học tập, tạo điều kiện thuận lợi để
tôi có thể áp dụng vào thực tiễn thông qua đồ án này.

Tôi cũng xin gửi lời cảm ơn đến gia đình, bạn bè và các đồng nghiệp đã
luôn ủng hộ, chia sẻ và đồng hành cùng tôi trong suốt thời gian thực
hiện đề tài.

Mặc dù đã nỗ lực hoàn thiện đồ án trong phạm vi thời gian và kiến thức
cho phép, nội dung chắc chắn vẫn còn những thiếu sót. Tôi rất mong nhận
được những ý kiến góp ý từ quý thầy cô và hội đồng để tiếp tục hoàn
thiện đề tài.

Xin chân thành cảm ơn!

\begin{flushright}
  Hà Nội, ngày \ldots tháng \ldots năm 20\ldots\\[0.25cm]
  \textbf{SINH VIÊN THỰC HIỆN}\\[2.4cm]
  Lê Trọng An
\end{flushright}

\majorheading{MỤC LỤC - TABLE OF CONTENT}\label{mux1ee5c-lux1ee5c---table-of-contents}

\thesistableofcontentsbody

% Danh mục bảng/hình được lược bỏ trong bản rút gọn để ưu tiên dung lượng cho nội dung học thuật cốt lõi.

\majorheading{DANH MỤC TỪ VIẾT TẮT - LIST OF ABBREVIATION}\label{danh-mux1ee5c-tux1eeb-viux1ebft-tux1eaft---list-of-abbreviations}

{\def\LTcaptype{none} % do not increment counter
\setlength{\tabcolsep}{4pt}
\begin{longtable}{|L{0.160\linewidth}|L{0.410\linewidth}|L{0.360\linewidth}|}
\hline
\textbf{Từ viết tắt} & \textbf{Tiếng Anh} & \textbf{Tiếng Việt} \\
\hline
\endfirsthead
\hline
\textbf{Từ viết tắt} & \textbf{Tiếng Anh} & \textbf{Tiếng Việt} \\
\hline
\endhead
\hline
\endfoot
\endlastfoot
\textbf{ADC} & Analog-to-Digital Converter & Bộ chuyển đổi tương tự sang
số \\\hline
\textbf{API} & Application Programming Interface & Giao diện lập trình
ứng dụng \\\hline
\textbf{BLE} & Bluetooth Low Energy & Bluetooth năng lượng thấp \\\hline
\textbf{BOM} & Bill of Materials & Danh sách linh kiện \\\hline
\textbf{CI/CD} & Continuous Integration / Continuous Deployment & Tích
hợp liên tục / Triển khai liên tục \\\hline
\textbf{CORS} & Cross-Origin Resource Sharing & Chia sẻ tài nguyên nguồn
chéo \\\hline
\textbf{DDD} & Domain-Driven Design & Thiết kế hướng miền \\\hline
\textbf{Docker} & Docker Container Platform & Nền tảng container hóa ứng
dụng \\\hline
\textbf{EMQX} & Erlang MQTT Broker & Máy chủ MQTT dựa trên Erlang \\\hline
\textbf{ESP32} & Espressif System-on-Chip 32-bit & Vi điều khiển SoC
32-bit của Espressif \\\hline
\textbf{FreeRTOS} & Free Real-Time Operating System & Hệ điều hành thời
gian thực mã nguồn mở \\\hline
\textbf{GPIO} & General Purpose Input/Output & Cổng vào/ra đa năng \\\hline
\textbf{GNSS} & Global Navigation Satellite System & Hệ thống vệ tinh
dẫn đường toàn cầu \\\hline
\textbf{GPS} & Global Positioning System & Hệ thống định vị toàn cầu \\\hline
\textbf{HTTP} & Hypertext Transfer Protocol & Giao thức truyền siêu văn
bản \\\hline
\textbf{HTTPS} & Hypertext Transfer Protocol Secure & Giao thức truyền
siêu văn bản bảo mật \\\hline
\textbf{IMU} & Inertial Measurement Unit & Đơn vị đo lường quán tính \\\hline
\textbf{IoT} & Internet of Things & Internet vạn vật \\\hline
\textbf{JWT} & JSON Web Token & Mã thông báo web dạng JSON \\\hline
\textbf{LVD} & Low Voltage Disconnect & Ngắt điện áp thấp \\\hline
\textbf{MCU} & Microcontroller Unit & Vi điều khiển \\\hline
\textbf{MQTT} & Message Queuing Telemetry Transport & Giao thức truyền
thông tin hàng đợi \\\hline
\textbf{OBD2} & On-Board Diagnostics II & Chẩn đoán trên xe thế hệ 2 \\\hline
\textbf{ORM} & Object-Relational Mapping & Ánh xạ đối tượng - quan hệ \\\hline
\textbf{PCB} & Printed Circuit Board & Bảng mạch in \\\hline
\textbf{REST} & Representational State Transfer & Kiến trúc truyền trạng
thái đại diện \\\hline
\textbf{RTOS} & Real-Time Operating System & Hệ điều hành thời gian
thực \\\hline
\textbf{SoC} & System on Chip & Hệ thống trên chip \\\hline
\textbf{SQL} & Structured Query Language & Ngôn ngữ truy vấn có cấu
trúc \\\hline
\textbf{TLS} & Transport Layer Security & Bảo mật tầng vận chuyển \\\hline
\textbf{UART} & Universal Asynchronous Receiver-Transmitter & Bộ thu
phát bất đồng bộ đa năng \\\hline
\textbf{WebSocket} & WebSocket Protocol & Giao thức kết nối hai chiều
thời gian thực \\\hline
\end{longtable}
}

\begin{center}\rule{0.5\linewidth}{0.5pt}\end{center}

\mainmatterstart
\chapterheading{CHƯƠNG 1. GIỚI THIỆU DỰ ÁN -- PROJECT OVERVIEW}

\subsection{1.1. Đặt vấn đề / Bối cảnh của dự
án}\label{11-ux111ux1eb7t-vux1ea5n-ux111ux1ec1--bux1ed1i-cux1ea3nh-cux1ee7a-dux1ef1-uxe1n}

\subsubsection{1.1.1. Bối cảnh thị trường cho thuê xe tự lái tại Việt
Nam}\label{111-bux1ed1i-cux1ea3nh-thux1ecb-trux1b0ux1eddng-cho-thuuxea-xe-tux1ef1-luxe1i-tux1ea1i-viux1ec7t-nam}

Những năm gần đây, thị trường cho thuê xe tự lái tại Việt Nam tăng
trưởng rõ rệt, đặc biệt tại TP. Hồ Chí Minh, Hà Nội và Đà Nẵng. Theo báo
cáo ngành vận tải {[}1{]}, nhu cầu thuê xe tăng trung bình 15--20\% mỗi
năm nhờ sự phục hồi của du lịch nội địa, nhu cầu di chuyển linh hoạt và
xu hướng chia sẻ phương tiện. Với doanh nghiệp vận hành đội xe, tốc độ
tăng trưởng này mang lại cơ hội kinh doanh nhưng cũng kéo theo một yêu
cầu rất thực tế: phải biết xe đang ở đâu, đang được sử dụng như thế nào
và có an toàn hay không mà không thể chỉ dựa vào các cuộc gọi xác nhận.
Khi quy mô đội xe mở rộng, áp lực giám sát vận hành và bảo toàn tài sản
tăng lên nhanh chóng. Trong thực tế, các doanh nghiệp cho thuê xe tự lái
hiện nay thường gặp bốn nhóm khó khăn sau:

\begin{itemize}
\tightlist
\item
  \textbf{Quản lý thủ công kém hiệu quả}: Nhiều doanh nghiệp quy mô vừa
  và nhỏ vẫn quản lý xe bằng sổ sách, gọi điện xác nhận hoặc dựa vào sự
  chủ động của khách hàng. Cách làm này không cung cấp thông tin thời
  gian thực về vị trí và trạng thái phương tiện {[}2{]}.
\item
  \textbf{Rủi ro mất cắp và sử dụng sai mục đích}: Khi không có hệ thống
  giám sát, xe có thể bị sử dụng vượt phạm vi địa lý đã thỏa thuận, chạy
  quá số km quy định, hoặc trong trường hợp xấu nhất là bị chiếm đoạt.
  Việc phát hiện các tình huống này thường bị trễ, gây thiệt hại lớn về
  tài sản {[}3{]}.
\item
  \textbf{Thiếu dữ liệu chẩn đoán kỹ thuật}: Doanh nghiệp khó theo dõi
  tình trạng kỹ thuật của xe từ xa, nên việc bảo trì thường mang tính bị
  động và chỉ được thực hiện khi sự cố đã xuất hiện. Điều này làm tăng
  chi phí sửa chữa và rút ngắn tuổi thọ phương tiện.
\item
  \textbf{Giải pháp thương mại đắt đỏ}: Các hệ thống GPS tracking thương
  mại hiện có trên thị trường (như Vietmap, iTracking) thường có chi phí
  cao --- bao gồm phí thiết bị, phí dịch vụ hàng tháng, và phí tích hợp
  --- không phù hợp với các doanh nghiệp quy mô nhỏ với ngân sách hạn
  chế {[}4{]}.
\end{itemize}

\subsubsection{1.1.2. Thách thức kỹ
thuật}\label{112-thuxe1ch-thux1ee9c-kux1ef9-thuux1eadt}

Để giải quyết những khó khăn trên, hệ thống không thể chỉ là một thiết
bị định vị đơn giản. Bài toán thực tế là phải dung hòa nhiều yêu cầu kỹ
thuật vốn dễ xung đột với nhau: thiết bị phải tiết kiệm điện nhưng vẫn
luôn sẵn sàng, phải truyền dữ liệu liên tục nhưng vẫn chịu được mất
sóng, và phải đủ đơn giản để triển khai thực tế với chi phí hợp lý. Bốn
thách thức nổi bật là:

\textbf{Thứ nhất, bài toán tiêu thụ năng lượng.} Nếu thiết bị theo dõi
hoạt động liên tục ở trạng thái công suất cao, ắc quy xe có thể bị hao
đáng kể chỉ sau vài tuần, từ đó ảnh hưởng trực tiếp đến khả năng khởi
động {[}5{]}. Vì vậy, hệ thống phải biết khi nào cần hoạt động đầy đủ và
khi nào có thể chuyển sang chế độ tiết kiệm điện.

\textbf{Thứ hai, giám sát khi xe đỗ.} Khi xe tắt máy (IGN OFF), thiết bị
không thể duy trì mọi khối chức năng như lúc xe đang chạy. Tuy nhiên,
nó vẫn phải phát hiện được các chuyển động bất thường, chẳng hạn xe bị
dắt đi hoặc bị cẩu kéo. Đây là bài toán cân bằng giữa độ nhạy phát hiện
và mức tiêu thụ năng lượng {[}6{]}.

\textbf{Thứ ba, độ tin cậy của kết nối.} Hệ thống hoạt động trong môi
trường di động nên kết nối 4G/LTE có thể gián đoạn bất kỳ lúc nào. Điều
người quản lý cần không chỉ là "có mạng thì gửi được", mà là khi mất
mạng thiết bị vẫn không làm mất dữ liệu quan trọng và có thể đồng bộ lại
khi kết nối được phục hồi.

\textbf{Thứ tư, xử lý dữ liệu thời gian thực.} Khi số lượng xe tăng lên
hàng chục hoặc hàng trăm, hệ thống backend phải tiếp nhận liên tục dữ
liệu vị trí, thông số OBD2 và trạng thái cảm biến với độ trễ thấp, đồng
thời vẫn giữ giao diện giám sát dễ quan sát và dễ thao tác cho người
quản lý.

{\thesissetfigureid{1.1}{fig-ch1-overview-problem-solution}
\begin{figure}[htbp]
\centering
\pandocbounded{\safeincludesvg{./assets/figures/01-chuong-1-gioi-thieu-hinh-1-1.svg}}
\caption{Sơ đồ tổng quan vấn đề và giải pháp đề xuất}
\label{fig:ch1-overview-problem-solution}
\end{figure}
}

\subsubsection{1.1.3. Động lực thực hiện dự
án}\label{113-ux111ux1ed9ng-lux1ef1c-thux1ef1c-hiux1ec7n-dux1ef1-uxe1n}

Nói cách khác, mục tiêu của đề tài không chỉ là làm ra một thiết bị biết
gửi tọa độ GPS. Điều cần xây dựng là một hệ thống đủ trọn vẹn để người
vận hành đội xe có thể nhìn thấy trạng thái phương tiện, hiểu nguyên
nhân của các cảnh báo và đưa ra quyết định kịp thời. Từ nhu cầu thực
tiễn đó, dự án "IoT Vehicle Tracking System" được triển khai nhằm xây
dựng một giải pháp đồng bộ, bao gồm thiết bị phần cứng IoT, firmware
nhúng, hệ thống cloud và giao diện web. Mục tiêu là tạo ra một hệ thống
có chi phí hợp lý, dễ tùy biến và có thể mở rộng theo quy mô doanh
nghiệp cho thuê xe tự lái tại Việt Nam.

\begin{center}\rule{0.5\linewidth}{0.5pt}\end{center}

\subsection{1.2. Mục tiêu và phạm vi của dự
án}\label{12-mux1ee5c-tiuxeau-vuxe0-phux1ea1m-vi-cux1ee7a-dux1ef1-uxe1n}

\subsubsection{1.2.1. Mục tiêu tổng
quát}\label{121-mux1ee5c-tiuxeau-tux1ed5ng-quuxe1t}

Mục tiêu chính của dự án là thiết kế và hiện thực một hệ thống IoT theo
dõi phương tiện hoàn chỉnh, bao gồm cả phần cứng lẫn phần mềm, phục vụ
cho bài toán quản lý đội xe cho thuê tự lái. Các mục tiêu dưới đây được
xác định từ góc nhìn của đơn vị vận hành đội xe: họ cần thông tin đủ
nhanh để can thiệp, nhưng cũng cần hệ thống đủ bền bỉ để chạy lâu dài
trên xe. Hệ thống cần đáp ứng các yêu cầu sau:

\begin{enumerate}
\def\labelenumi{\arabic{enumi}.}
\tightlist
\item
  \textbf{Theo dõi vị trí thời gian thực (real-time tracking)}: Cung cấp
  vị trí GPS của xe với tần suất cập nhật 5--30 giây khi xe đang di
  chuyển, hiển thị trên bản đồ trực tuyến.
\item
  \textbf{Giám sát trạng thái kỹ thuật qua OBD2}: Kết nối với cổng chẩn
  đoán OBD-II của xe thông qua giao thức Bluetooth Low Energy (BLE) để
  đọc các thông số như tốc độ, vòng tua máy, nhiệt độ động cơ, và mã lỗi
  chẩn đoán (DTC).
\item
  \textbf{Cảnh báo thông minh}: Tự động phát hiện và gửi thông báo khi
  có các sự kiện bất thường như: xe vượt ra khỏi vùng địa lý cho phép
  (hàng rào địa lý - geofencing), vượt tốc độ quy định, chuyển động bất
  thường khi xe đang đỗ, hoặc mất kết nối thiết bị.
\item
  \textbf{Tối ưu hóa năng lượng}: Đảm bảo thiết bị hoạt động liên tục mà
  không làm cạn ắc quy xe, với thời lượng pin dự phòng đủ để hoạt động
  độc lập trong trường hợp mất nguồn chính.
\item
  \textbf{Giao diện quản lý trực quan}: Cung cấp bảng điều khiển
  (dashboard) web cho phép người quản lý đội xe theo dõi vị trí, xem
  lịch sử hành trình, quản lý cảnh báo, và tạo báo cáo.
\end{enumerate}

\subsubsection{1.2.2. Phạm vi dự
án}\label{122-phux1ea1m-vi-dux1ef1-uxe1n}

Để đề tài không dàn trải, phạm vi nghiên cứu được giới hạn theo từng
tầng của hệ thống, từ đối tượng ứng dụng đến phần cứng, firmware và
cloud. Cách xác định phạm vi này giúp người đọc thấy rõ đâu là phần đã
được triển khai trong đồ án và đâu là nội dung được giữ lại cho các giai
đoạn phát triển sau.

\textbf{Đối tượng ứng dụng:}

\begin{itemize}
\tightlist
\item
  Loại xe: Ô tô con (xe 4 chỗ, sedan, SUV) với hệ thống điện 12V hoặc
  24V DC
\item
  Ắc quy phổ thông: 40--70 Ah (trung bình 45--60 Ah)
\item
  Đối tượng sử dụng: Các công ty cho thuê xe tự lái quy mô vừa và nhỏ
  tại Việt Nam
\end{itemize}

\textbf{Phạm vi phần cứng:}

\begin{itemize}
\tightlist
\item
  Thiết bị tracker IoT sử dụng vi điều khiển ESP32-S3
\item
  Modem LTE + GNSS SIMCom SIM7600CE-T (Auto mode LTE/UMTS/GSM, APN mặc
  định \texttt{internet}) kết nối trực tiếp qua UART1 để cung cấp cả dữ
  liệu 4G và GNSS
\item
  Adapter OBD2 BLE vgate iCar Pro (đọc dữ liệu chẩn đoán xe qua
  Bluetooth)
\item
  Cảm biến gia tốc LIS3DSH (IMU) để phát hiện chuyển động và rung
\item
  Pin dự phòng 18650 1S với mạch sạc và bảo vệ
\end{itemize}

\textbf{Phạm vi firmware:}

\begin{itemize}
\tightlist
\item
  Phát triển trên nền tảng ESP-IDF với hệ điều hành thời gian thực
  FreeRTOS
\item
  Giao tiếp BLE với adapter OBD2
\item
  Truyền dữ liệu qua giao thức MQTT 5.0
\item
  Quản lý năng lượng đa chế độ (driving, parking, alert)
\end{itemize}

\textbf{Phạm vi cloud/backend:}

\begin{itemize}
\tightlist
\item
  EMQX MQTT Broker làm trung tâm tiếp nhận dữ liệu từ thiết bị
\item
  MQTT Bridge (dịch vụ độc lập) xử lý và phân phối dữ liệu
\item
  PostgreSQL lưu trữ dữ liệu quan hệ (người dùng, xe, khách hàng, cảnh
  báo)
\item
  VictoriaMetrics lưu trữ dữ liệu chuỗi thời gian (telemetry)
\item
  VictoriaLogs lưu trữ nhật ký sự kiện
\item
  Express.js API server với kiến trúc Domain-Driven Design
\item
  Next.js 15 frontend với bản đồ thời gian thực (Leaflet), biểu đồ
  (ECharts)
\end{itemize}

\textbf{Giới hạn phạm vi nghiên cứu:}

\begin{itemize}
\tightlist
\item
  Không xử lý OBD-II diagnostics phức tạp (chỉ đọc các thông số cơ bản)
\item
  Không phát hiện va chạm (IMU chỉ dùng cho phát hiện chuyển động/rung)
\item
  Không điều khiển động cơ xe (tắt bơm xăng, khóa xe) --- chỉ gửi cảnh
  báo
\item
  Không bao gồm ứng dụng di động ban đầu (dự kiến Phase 2 với Flutter
  WebView)
\end{itemize}

{\thesissettableid{1.1}{tab-ch1-project-scope-by-layer}
\def\LTcaptype{table}
\begin{longtable}{|L{0.261\linewidth}|L{0.470\linewidth}|L{0.234\linewidth}|}
\caption{Tóm tắt phạm vi dự án theo tầng (layer)}\label{tab:ch1-project-scope-by-layer}\\
\hline \textbf{Tầng (Layer)} & \textbf{Công nghệ chính} & \textbf{Phạm vi} \\
\hline
\endfirsthead
\caption[]{Tóm tắt phạm vi dự án theo tầng (layer) (tiếp theo)}\\
\hline \textbf{Tầng (Layer)} & \textbf{Công nghệ chính} & \textbf{Phạm vi} \\
\hline
\endhead
\hline
\endfoot
\endlastfoot
Phần cứng & ESP32-S3, SIMCom SIM7600CE-T (LTE + GNSS), vgate iCar Pro,
LIS3DSH & Thiết kế PCB, chế tạo và tích hợp thiết bị hoàn chỉnh \\\hline
Firmware & ESP-IDF, FreeRTOS, MQTT 5.0 & Lập trình nhúng đầy đủ \\\hline
MQTT Broker & EMQX 5.x & Cấu hình và triển khai \\\hline
Backend & Express.js, TypeScript, PostgreSQL, VictoriaMetrics & Phát
triển API và xử lý dữ liệu \\\hline
Frontend & Next.js 15, React 19, Leaflet, ECharts & Giao diện quản lý
web \\\hline
Infrastructure & Docker, Nginx Proxy Manager, Grafana & Triển khai và
giám sát \\\hline
\end{longtable}
}

\begin{center}\rule{0.5\linewidth}{0.5pt}\end{center}

\subsection{1.3. Các tiêu chí cần đạt được của dự
án}\label{13-cuxe1c-tiuxeau-chuxed-cux1ea7n-ux111ux1ea1t-ux111ux1b0ux1ee3c-cux1ee7a-dux1ef1-uxe1n}

Dự án đặt ra các tiêu chí cụ thể (success criteria) cho từng tầng của hệ
thống. Các tiêu chí này không chỉ dùng để đánh giá kết quả cuối cùng, mà
còn đóng vai trò như một cam kết thiết kế xuyên suốt quá trình thực hiện:
mỗi quyết định chọn linh kiện, tổ chức firmware hay xây dựng hạ tầng đều
phải quay lại trả lời câu hỏi liệu nó có giúp hệ thống vận hành đủ ổn
định, đủ tiết kiệm năng lượng và đủ hữu ích cho bài toán quản lý xe hay
không.

\subsubsection{1.3.1. Tiêu chí phần
cứng}\label{131-tiuxeau-chuxed-phux1ea7n-cux1ee9ng}

Ở tầng phần cứng, điều quan trọng nhất không phải là "gắn được nhiều khối
chức năng", mà là thiết bị có thể tồn tại bền bỉ trong môi trường trên xe:
đủ tiết kiệm điện, đủ ổn định về nhiệt và đủ an toàn với ắc quy.


{\thesissettableid{1.3.1A}{tab-ch1-tong-hop-thong-tin-ve-stt-tieu-chi-muc-tieu-cu-the}
\def\LTcaptype{table}
\begin{longtable}{|L{0.160\linewidth}|L{0.224\linewidth}|L{0.580\linewidth}|}
\caption{Tổng hợp thông tin về STT, Tiêu chí, Mục tiêu cụ thể}\label{tab:ch1-tong-hop-thong-tin-ve-stt-tieu-chi-muc-tieu-cu-the}\\
\hline \textbf{STT} & \textbf{Tiêu chí} & \textbf{Mục tiêu cụ thể} \\
\hline
\endfirsthead
\caption[]{Tổng hợp thông tin về STT, Tiêu chí, Mục tiêu cụ thể (tiếp theo)}\\
\hline \textbf{STT} & \textbf{Tiêu chí} & \textbf{Mục tiêu cụ thể} \\
\hline
\endhead
\hline
\endfoot
\endlastfoot
1 & Tiêu thụ điện chế độ ngủ sâu (deep sleep) & \textless{} 500 µA \\\hline
2 & Tiêu thụ điện chế độ hoạt động & \textless{} 250 mA (trung bình) \\\hline
3 & Thời gian thức dậy từ deep sleep & \textless{} 3 giây \\\hline
4 & Thời lượng pin dự phòng (18650 1S) & \textgreater= 2.5 giờ tracking
liên tục; \textgreater= 24 giờ cảnh báo gián đoạn (mục tiêu theo
baseline 18650) \\\hline
5 & Ngưỡng chuyển nguồn bảo vệ ắc quy & Profile 12V:
\(OFF = 12.0\,\mathrm{V}\), \(ON = 12.2\,\mathrm{V}\); Profile 24V:
\(OFF = 24.0\,\mathrm{V}\), \(ON = 24.4\,\mathrm{V}\) \\\hline
6 & Nhiệt độ hoạt động & -10 C đến +60 C \\\hline
\end{longtable}
}


\subsubsection{1.3.2. Tiêu chí
firmware}\label{132-tiuxeau-chuxed-firmware}

Ở tầng firmware, tiêu chí không chỉ dừng ở việc đọc được cảm biến hay gửi
được dữ liệu, mà nằm ở khả năng điều phối đúng nhịp giữa các khối để thiết
bị phản ứng phù hợp với từng trạng thái vận hành của xe.


{\thesissettableid{1.2A}{tab-ch1-tieu-chi-firmware}
\def\LTcaptype{table}
\begin{longtable}{|L{0.160\linewidth}|L{0.404\linewidth}|L{0.400\linewidth}|}
\caption{Tiêu chí firmware}\label{tab:ch1-tieu-chi-firmware}\\
\hline \textbf{STT} & \textbf{Tiêu chí} & \textbf{Mục tiêu cụ thể} \\
\hline
\endfirsthead
\caption[]{Tiêu chí firmware (tiếp theo)}\\
\hline \textbf{STT} & \textbf{Tiêu chí} & \textbf{Mục tiêu cụ thể} \\
\hline
\endhead
\hline
\endfoot
\endlastfoot
1 & Tần suất gửi GPS khi lái xe & 5--30 giây (cấu hình được) \\\hline
2 & Tần suất heartbeat khi đỗ xe & 10--30 phút \\\hline
3 & Thời gian kết nối OBD2 BLE & \textless{} 10 giây \\\hline
4 & Độ chính xác vị trí GNSS & \textless{} 5 mét (điều kiện thoáng) \\\hline
5 & Xử lý mất kết nối & Tự động kết nối lại; có phương án bổ sung buffer
cục bộ \\\hline
6 & Phát hiện chuyển động bất thường & Độ nhạy IMU cấu hình được \\\hline
\end{longtable}
}


\subsubsection{1.3.3. Tiêu chí
cloud/backend}\label{133-tiuxeau-chuxed-cloudbackend}

Khi chuyển lên cloud/backend, trọng tâm đánh giá đổi sang khả năng tiếp
nhận dữ liệu liên tục, phản hồi nhanh và bảo vệ dữ liệu khi số lượng thiết
bị tăng lên theo quy mô đội xe.


{\thesissettableid{1.2B}{tab-ch1-tieu-chi-cloud-backend}
\def\LTcaptype{table}
\begin{longtable}{|L{0.160\linewidth}|L{0.311\linewidth}|L{0.494\linewidth}|}
\caption{Tiêu chí cloud/backend}\label{tab:ch1-tieu-chi-cloud-backend}\\
\hline \textbf{STT} & \textbf{Tiêu chí} & \textbf{Mục tiêu cụ thể} \\
\hline
\endfirsthead
\caption[]{Tiêu chí cloud/backend (tiếp theo)}\\
\hline \textbf{STT} & \textbf{Tiêu chí} & \textbf{Mục tiêu cụ thể} \\
\hline
\endhead
\hline
\endfoot
\endlastfoot
1 & Độ trễ xử lý MQTT message & \textless{} 500 ms (end-to-end) \\\hline
2 & Số lượng thiết bị đồng thời & \textgreater= 50 thiết bị (đã kiểm thử), dư địa tới khoảng 100 thiết bị \\\hline
3 & Độ sẵn sàng hệ thống (uptime) & \textgreater= 99.5\% \\\hline
4 & Thời gian phản hồi API & \textless{} 200 ms (p95) \\\hline
5 & Lưu trữ dữ liệu telemetry & \textgreater= 90 ngày \\\hline
6 & Bảo mật xác thực & Phiên đăng nhập bằng opaque token lưu trong DB,
token được băm SHA-256 \\\hline
\end{longtable}
}


\subsubsection{1.3.4. Tiêu chí
frontend}\label{134-tiuxeau-chuxed-frontend}

Cuối cùng, frontend là nơi toàn bộ nỗ lực kỹ thuật được chuyển thành trải
nghiệm quan sát và ra quyết định cho người vận hành, nên tiêu chí ở đây tập
trung vào tính cập nhật, dễ theo dõi và dễ sử dụng trong bối cảnh thực tế.


{\thesissettableid{1.2C}{tab-ch1-tieu-chi-frontend}
\def\LTcaptype{table}
\begin{longtable}{|L{0.160\linewidth}|L{0.324\linewidth}|L{0.481\linewidth}|}
\caption{Tiêu chí frontend}\label{tab:ch1-tieu-chi-frontend}\\
\hline \textbf{STT} & \textbf{Tiêu chí} & \textbf{Mục tiêu cụ thể} \\
\hline
\endfirsthead
\caption[]{Tiêu chí frontend (tiếp theo)}\\
\hline \textbf{STT} & \textbf{Tiêu chí} & \textbf{Mục tiêu cụ thể} \\
\hline
\endhead
\hline
\endfoot
\endlastfoot
1 & Cập nhật bản đồ thời gian thực & \textless{} 2 giây độ trễ \\\hline
2 & Hỗ trợ trình duyệt & Chrome, Firefox, Edge (phiên bản mới nhất) \\\hline
3 & Responsive design & Desktop và tablet \\\hline
4 & Hiển thị cảnh báo & Thông báo tức thời qua WebSocket \\\hline
\end{longtable}
}


\begin{center}\rule{0.5\linewidth}{0.5pt}\end{center}

\subsection{1.4. Phương pháp tiếp cận thiết kế kỹ
thuật}\label{14-phux1b0ux1a1ng-phuxe1p-tiux1ebfp-cux1eadn-thiux1ebft-kux1ebf-kux1ef9-thuux1eadt}

\subsubsection{1.4.1. Phương pháp luận tổng
thể}\label{141-phux1b0ux1a1ng-phuxe1p-luux1eadn-tux1ed5ng-thux1ec3}

Vì đề tài đi qua cả mạch nguồn, firmware nhúng, cloud lẫn dashboard, rủi
ro lớn nhất không nằm ở việc từng khối chạy riêng lẻ, mà ở việc chúng có
còn phối hợp đúng khi ghép thành một hệ thống hoàn chỉnh hay không. Từ
đó, dự án chọn phương pháp \textbf{thiết kế từ dưới lên (bottom-up
design)} kết hợp \textbf{phát triển lặp (iterative development)}: kiểm
chứng trước những giả định rủi ro nhất ở tầng thấp, rồi mới ghép dần lên
các tầng trên. Cách tiếp cận này giúp mỗi giai đoạn tạo ra một kết quả
có thể đo và kiểm tra được, đồng thời giảm nguy cơ chỉ đến cuối dự án
mới lộ ra xung đột về nguồn, giao tiếp hay hiệu năng. Cụ thể, tiến trình
được tổ chức như sau:

\begin{enumerate}
\def\labelenumi{\arabic{enumi}.}
\tightlist
\item
  \textbf{Giai đoạn 1 --- Nghiên cứu và thiết kế phần cứng}: Xây nền cho
  hệ thống bằng cách khảo sát linh kiện, đối chiếu datasheet, thiết kế
  sơ đồ nguyên lý và tính toán nguồn. Mục tiêu của giai đoạn này là bảo
  đảm thiết bị có thể tồn tại bền bỉ trong môi trường trên xe trước khi
  nói đến các tính năng ở tầng trên.
\item
  \textbf{Giai đoạn 2 --- Phát triển firmware}: Biến sơ đồ phần cứng
  thành hành vi vận hành thật trên ESP-IDF, bao gồm giao tiếp BLE OBD2,
  UART modem, I2C IMU, quản lý năng lượng và truyền dữ liệu MQTT. Đây là
  giai đoạn quyết định thiết bị có ``phản ứng đúng lúc'' hay không.
\item
  \textbf{Giai đoạn 3 --- Xây dựng hạ tầng cloud}: Chuẩn bị lớp tiếp
  nhận và lưu trữ dữ liệu bằng cách triển khai EMQX, PostgreSQL,
  VictoriaMetrics và các dịch vụ phụ trợ trên Docker, đồng thời xác định
  cấu trúc dữ liệu và kết nối mạng.
\item
  \textbf{Giai đoạn 4 --- Phát triển backend API}: Chuyển dữ liệu thô
  thành nghiệp vụ có thể truy vấn và điều phối, thông qua REST API, xử
  lý domain và kênh WebSocket thời gian thực.
\item
  \textbf{Giai đoạn 5 --- Phát triển frontend}: Biến dữ liệu kỹ thuật
  thành thông tin dễ đọc đối với người vận hành bằng dashboard web,
  bản đồ trực tuyến và biểu đồ trực quan.
\item
  \textbf{Giai đoạn 6 --- Tích hợp và kiểm thử}: Ghép toàn bộ các lớp
  lại thành một chuỗi hoàn chỉnh từ xe đến dashboard, sau đó kiểm thử,
  đo đạc và tối ưu những điểm còn nghẽn.
\end{enumerate}

{\thesissetfigureid{1.2}{fig-ch1-project-development-phases}
\begin{figure}[htbp]
\centering
\pandocbounded{\safeincludesvg{./assets/figures/01-chuong-1-gioi-thieu-hinh-1-2.svg}}
\caption{Quy trình phát triển dự án theo các giai đoạn}
\label{fig:ch1-project-development-phases}
\end{figure}
}

\subsubsection{1.4.2. Kiến trúc hệ
thống}\label{142-kiux1ebfn-truxfac-hux1ec7-thux1ed1ng}

Nhìn theo hành trình của một bản tin, hệ thống khởi đầu ở thiết bị gắn
trên xe, đi qua lớp tiếp nhận dữ liệu thời gian thực trên cloud, rồi mới
đến lớp nghiệp vụ và lớp giao diện cho người quản lý. Góc nhìn này giúp
lý giải tự nhiên vì sao kiến trúc được tổ chức theo dạng phân tầng
(layered architecture): thiết bị thu thập và gửi dữ liệu; lớp trung gian
tiếp nhận, phân luồng và lưu trữ; lớp ứng dụng phân tích, quản lý và
trình bày thông tin cho người dùng. Trong triển khai thực tế, tracker
(ESP32-S3, SIM7600CE-T, LIS3DSH, vgate iCar Pro) kết nối MQTT lên EMQX;
MQTT Bridge tiếp nhận và phân luồng dữ liệu sang VictoriaMetrics,
VictoriaLogs và Backend API; lớp ứng dụng phía trên phục vụ dashboard web
Next.js và hệ quan trắc Grafana. Ứng dụng di động được giữ ở phạm vi phát
triển giai đoạn sau để tránh làm dàn trải đồ án.

{\thesissetfigureid{1.3}{fig-ch1-system-architecture-overview}
\begin{figure}[htbp]
\centering
\pandocbounded{\safeincludesvg{./assets/figures/01-chuong-1-gioi-thieu-hinh-1-3.svg}}
\caption{Kiến trúc tổng thể hệ thống IoT Vehicle Tracking}
\label{fig:ch1-system-architecture-overview}
\end{figure}
}

\textbf{Chiến lược dữ liệu (Data Strategy):}

Chiến lược dữ liệu được xây theo cùng logic phân tầng này. Lý do là dữ liệu
trong hệ thống không đồng nhất: có loại cần tính nhất quán và quan hệ
chặt chẽ, có loại cần ghi liên tục với tần suất cao, và có loại chủ yếu
phục vụ truy vết sự cố. Vì vậy, mỗi nhóm dữ liệu được đưa về đúng nơi
phù hợp với đặc tính truy vấn và mục đích sử dụng.

\begin{itemize}
\tightlist
\item
  \textbf{PostgreSQL}: Lưu trữ dữ liệu quan hệ có cấu trúc như người
  dùng, xe, khách hàng, hành trình, cảnh báo, vùng địa lý
  (geofences) và lệnh điều khiển. Đây là lớp dữ liệu cần quan hệ rõ
  ràng và ràng buộc chặt chẽ.
\item
  \textbf{VictoriaMetrics}: Lưu trữ dữ liệu chuỗi thời gian
  (time-series) như tọa độ GPS, dữ liệu OBD2 và chỉ số cảm biến. Nhóm
  dữ liệu này cần ghi nhiều, đọc nhanh theo trục thời gian và lưu dài
  hạn với chi phí hợp lý.
\item
  \textbf{VictoriaLogs}: Lưu trữ nhật ký sự kiện, audit trail và các bản
  ghi hoạt động của thiết bị để phục vụ quan sát hệ thống và truy vết sự
  cố.
\end{itemize}

\subsubsection{1.4.3. Chiến lược quản lý năng
lượng}\label{143-chiux1ebfn-lux1b0ux1ee3c-quux1ea3n-luxfd-nux103ng-lux1b0ux1ee3ng}

Thiết kế năng lượng của tracker thực chất là bài toán quyết định lúc nào
thiết bị cần hoạt động đầy đủ, lúc nào chỉ nên ``canh chừng'', và lúc nào
phải ưu tiên bảo vệ nguồn xe. Nếu xử lý không khéo, hệ thống sẽ rơi vào
hai cực đoan: theo dõi rất sát nhưng làm hao ắc quy, hoặc tiết kiệm điện
quá mức nên bỏ lỡ sự kiện quan trọng. Vì vậy, chiến lược quản lý năng
lượng đa chế độ (multi-mode power management) được xây quanh các trạng
thái vận hành thật của xe, để mỗi thời điểm chỉ bật những khối thật sự
cần thiết:

\textbf{Chế độ 1 --- Driving Mode (IGN ON, ưu tiên độ đầy đủ dữ liệu):}

\begin{itemize}
\tightlist
\item
  Hệ thống hoạt động toàn bộ công suất
\item
  Gửi vị trí GPS định kỳ (5--30 giây)
\item
  Sạc pin dự phòng từ nguồn xe
\item
  Kết nối server liên tục để theo dõi thời gian thực
\end{itemize}

\textbf{Chế độ 2 --- Parking Mode (IGN OFF, ưu tiên bảo toàn năng lượng):}

\begin{itemize}
\tightlist
\item
  Vi điều khiển ESP32-S3 vào chế độ ngủ sâu (deep sleep)
\item
  Tiêu thụ cực thấp (\textless{} 500 µA)
\item
  Thức dậy định kỳ (10--30 phút) để gửi heartbeat
\item
  IMU (LIS3DSH) hoạt động độc lập, cảnh rung ở ngưỡng đã cấu hình
\end{itemize}

\textbf{Chế độ 3 --- Alert Mode (ưu tiên phản ứng nhanh khi có sự cố):}

\begin{itemize}
\tightlist
\item
  IMU đánh thức ESP32 ngay lập tức qua interrupt
\item
  Bật 4G + GPS, gửi cảnh báo ưu tiên lên server
\item
  Tiếp tục theo dõi liên tục cho đến khi được xác nhận an toàn
\item
  Sử dụng pin dự phòng nếu nguồn chính bị cắt
\end{itemize}

Bên cạnh đó, hệ thống còn có mạch chuyển nguồn dự phòng với cơ chế
hysteresis để tránh đóng/ngắt liên tục khi điện áp dao động quanh ngưỡng.
Thiết bị tự động chuyển sang pin dự phòng khi điện áp ắc quy tụt xuống
dưới mức an toàn theo từng profile: 12V dùng
\(Switch_{\mathrm{OFF}} = 12.0\,\mathrm{V}\), \(Switch_{\mathrm{ON}} = 12.2\,\mathrm{V}\);
24V dùng \(Switch_{\mathrm{OFF}} = 24.0\,\mathrm{V}\),
\(Switch_{\mathrm{ON}} = 24.4\,\mathrm{V}\), qua đó bảo vệ ắc quy không bị rút
cạn quá mức {[}7{]}.

Nhờ cách tổ chức này, năng lượng không còn là bài toán tối ưu từng linh
kiện riêng lẻ, mà trở thành một phần của logic vận hành toàn hệ thống.

{\thesissetfigureid{1.4}{fig-ch1-power-mode-transitions}
\begin{figure}[htbp]
\centering
\pandocbounded{\safeincludesvg{./assets/figures/01-chuong-1-gioi-thieu-hinh-1-4.svg}}
\caption{Sơ đồ chuyển đổi giữa các chế độ năng lượng}
\label{fig:ch1-power-mode-transitions}
\end{figure}
}

\subsubsection{1.4.4. Công nghệ sử
dụng}\label{144-cuxf4ng-nghux1ec7-sux1eed-dux1ee5ng}

{\thesissettableid{1.3}{tab-ch1-technology-stack-summary}
\def\LTcaptype{table}
\begin{longtable}{|L{0.223\linewidth}|L{0.334\linewidth}|L{0.115\linewidth}|L{0.282\linewidth}|}
\caption{Tổng hợp công nghệ sử dụng trong dự án}\label{tab:ch1-technology-stack-summary}\\
\hline \textbf{Thành phần} & \textbf{Công nghệ} & \textbf{Phiên bản} & \textbf{Vai trò} \\
\hline
\endfirsthead
\caption[]{Tổng hợp công nghệ sử dụng trong dự án (tiếp theo)}\\
\hline \textbf{Thành phần} & \textbf{Công nghệ} & \textbf{Phiên bản} & \textbf{Vai trò} \\
\hline
\endhead
\hline
\endfoot
\endlastfoot
Vi điều khiển & ESP32-S3 & --- & MCU chính, xử lý và điều khiển \\\hline
LTE + GNSS & SIMCom SIM7600CE-T (LTE + GNSS tích hợp) & --- & Truyền dữ
liệu 4G và định vị GPS \\\hline
OBD2 Adapter & vgate iCar Pro & BLE 4.0 & Đọc dữ liệu chẩn đoán xe \\\hline
Cảm biến gia tốc & LIS3DSH & --- & Phát hiện chuyển động và rung \\\hline
Pin dự phòng & 18650 1S Li-ion & 3500 mAh & Nguồn điện dự phòng \\\hline
Framework firmware & ESP-IDF & 5.x & Phát triển firmware nhúng \\\hline
RTOS & FreeRTOS & --- & Hệ điều hành thời gian thực \\\hline
MQTT Broker & EMQX & 5.x & Tiếp nhận dữ liệu IoT \\\hline
API Server & Express.js + TypeScript & --- & REST API backend \\\hline
Cơ sở dữ liệu quan hệ & PostgreSQL & 16 & Lưu trữ dữ liệu có cấu trúc \\\hline
Cơ sở dữ liệu chuỗi thời gian & VictoriaMetrics & --- & Lưu trữ
telemetry \\\hline
Nhật ký sự kiện & VictoriaLogs & --- & Lưu trữ logs \\\hline
Frontend & Next.js 15 + React 19 & --- & Giao diện web \\\hline
Bản đồ & Leaflet & --- & Hiển thị bản đồ thời gian thực \\\hline
Biểu đồ & ECharts & --- & Trực quan hóa dữ liệu \\\hline
WebSocket & Socket.IO & 4.8 & Giao tiếp thời gian thực \\\hline
Container hóa & Docker + Docker Compose & --- & Triển khai dịch vụ \\\hline
Giám sát & Grafana + VictoriaMetrics, VictoriaLogs & --- & Dashboard giám
sát hạ tầng \\\hline
\end{longtable}
}

\begin{center}\rule{0.5\linewidth}{0.5pt}\end{center}

\subsection{1.5. Kết quả và khuyến
nghị}\label{15-kux1ebft-quux1ea3-vuxe0-khuyux1ebfn-nghux1ecb}

\subsubsection{1.5.1. Kết quả đạt
được}\label{151-kux1ebft-quux1ea3-ux111ux1ea1t-ux111ux1b0ux1ee3c}

Nhìn theo cách đó, các kết quả đạt được không nên hiểu như một danh sách
các khối đã lắp xong, mà như bằng chứng cho thấy chuỗi giá trị từ thiết
bị đến dashboard đã được hình thành. Có thể tóm lược theo bốn tầng chính
như sau:

\textbf{Về phần cứng:}

\begin{itemize}
\tightlist
\item
  Hoàn thiện thiết bị tracker IoT với bo mạch riêng do nhóm thiết kế, sử
  dụng ESP32-S3-WROOM-1 làm vi điều khiển trung tâm, tích hợp modem LTE
  + GNSS SIMCom SIM7600CE-T, cảm biến gia tốc LIS3DSH, khối nguồn đa
  rail, pin dự phòng 18650 1S và kết nối OBD2 BLE qua adapter vgate iCar
  Pro.
\item
  Quá trình lựa chọn linh kiện bám theo datasheet, độ ổn định chất lượng
  và khả năng cung ứng trong nước để thuận lợi cho việc chế tạo, kiểm
  thử và thay thế khi cần.
\item
  Hệ thống quản lý năng lượng đa chế độ hoạt động hiệu quả, với mức tiêu
  thụ điện ngủ sâu đạt yêu cầu (\textless{} 500 µA), đảm bảo không làm
  cạn ắc quy xe trong quá trình sử dụng bình thường.
\item
  Mạch Low Voltage Disconnect (LVD) bảo vệ ắc quy xe hiệu quả, tự động
  ngắt khi điện áp tụt dưới ngưỡng an toàn.
\end{itemize}

\textbf{Về firmware:}

\begin{itemize}
\tightlist
\item
  Phát triển thành công firmware trên ESP-IDF với FreeRTOS, hiện thực
  đầy đủ các module: giao tiếp BLE OBD2, điều khiển modem UART, đọc cảm
  biến IMU, quản lý năng lượng, và truyền dữ liệu MQTT.
\item
  Cơ chế phát hiện chuyển động bất thường qua IMU hoạt động chính xác,
  có khả năng đánh thức hệ thống từ chế độ ngủ sâu trong vòng
  \textless{} 3 giây.
\end{itemize}

\textbf{Về cloud/backend:}

\begin{itemize}
\tightlist
\item
  Xây dựng thành công hạ tầng cloud bao gồm EMQX MQTT Broker, dịch vụ
  MQTT Bridge, PostgreSQL, VictoriaMetrics và VictoriaLogs, toàn bộ được
  container hóa bằng Docker.
\item
  Backend API (Express.js + TypeScript) với kiến trúc Domain-Driven
  Design, cung cấp các endpoint quản lý xe, thiết bị, telemetry, cảnh
  báo, geofences, và xác thực người dùng.
\item
  Hệ thống đã kiểm thử ở mức 50+ thiết bị đồng thời với độ trễ
  end-to-end \textless{} 500 ms, đồng thời cho thấy còn dư địa mở rộng
  ổn định tới khoảng 100 thiết bị.
\end{itemize}

\textbf{Về frontend:}

\begin{itemize}
\tightlist
\item
  Giao diện web quản lý (Next.js 15) với bản đồ thời gian thực
  (Leaflet), biểu đồ phân tích (ECharts), hệ thống cảnh báo trực tuyến,
  và bảng điều khiển tổng quan (dashboard).
\item
  Cập nhật vị trí xe trên bản đồ với độ trễ \textless{} 2 giây thông qua
  WebSocket (Socket.IO).
\end{itemize}

\subsubsection{1.5.2. Khuyến nghị và hướng phát
triển}\label{152-khuyux1ebfn-nghux1ecb-vuxe0-hux1b0ux1edbng-phuxe1t-triux1ec3n}

Trên cơ sở các kết quả đã đạt được và các giới hạn hiện hữu, các hướng phát
triển tiếp theo được sắp theo logic mở rộng từ việc làm chắc hệ lõi đến
việc mở rộng giá trị sử dụng của hệ thống:

\begin{enumerate}
\def\labelenumi{\arabic{enumi}.}
\tightlist
\item
  \textbf{Ứng dụng di động (Phase 2)}: Phát triển ứng dụng Flutter
  WebView Hybrid cho phép người quản lý giám sát đội xe trên điện thoại
  di động, nhận thông báo push notification khi có cảnh báo {[}8{]}.
\item
  \textbf{Tăng cường bảo mật}: Triển khai TLS/SSL cho kết nối MQTT trong
  môi trường production, áp dụng MQTT ACL chi tiết hơn, và bổ sung cơ
  chế mã hóa dữ liệu end-to-end giữa thiết bị và server.
\item
  \textbf{Tối ưu hóa hiệu năng}: Nghiên cứu và áp dụng các kỹ thuật nén
  dữ liệu telemetry (mã hóa sai khác, protobuf) để giảm băng thông 4G,
  đặc biệt khi vận hành đội xe lớn (\textgreater{} 500 xe).
\item
  \textbf{Tích hợp trí tuệ nhân tạo}: Áp dụng các mô hình machine
  learning để phân tích hành vi lái xe (driving behavior analysis), dự
  đoán bảo trì (predictive maintenance), và phát hiện bất thường
  (anomaly detection) từ dữ liệu telemetry.
\item
  \textbf{Mở rộng phạm vi OBD2}: Hỗ trợ đọc thêm nhiều thông số chẩn
  đoán xe, bao gồm mã lỗi DTC (Diagnostic Trouble Codes), dữ liệu động
  cơ nâng cao, và tích hợp với các loại xe điện (EV).
\item
  \textbf{Cải thiện khả năng mở rộng (scalability)}: Nghiên cứu kiến
  trúc microservices với message queue (Apache Kafka hoặc RabbitMQ) để
  xử lý luồng dữ liệu lớn hơn khi số lượng thiết bị tăng lên hàng ngàn.
\item
  \textbf{Tối ưu PCB cho sản xuất}: Nâng thiết kế PCB hiện tại theo
  hướng DFM/EMI, giảm kích thước, tăng độ bền cơ khí và chuẩn bị cho sản
  xuất hàng loạt.
\end{enumerate}

{\thesissetfigureid{1.5}{fig-ch1-project-roadmap}
\begin{figure}[htbp]
\centering
\pandocbounded{\safeincludesvg{./assets/figures/01-chuong-1-gioi-thieu-hinh-1-5.svg}}
\caption{Lộ trình phát triển dự án theo các giai đoạn (Roadmap)}
\label{fig:ch1-project-roadmap}
\end{figure}
}

\begin{center}\rule{0.5\linewidth}{0.5pt}\end{center}

\subsection{Kết luận chương
1}\label{kux1ebft-luux1eadn-chux1b0ux1a1ng-1}

Chương 1 đã trả lời câu hỏi vì sao đề tài này cần được thực hiện và hệ
thống cần đạt được điều gì trong bối cảnh thực tế của doanh nghiệp cho
thuê xe tự lái. Từ bối cảnh thị trường, nhu cầu vận hành đến các thách
thức về năng lượng, kết nối và dữ liệu, chương đã thiết lập bộ mục tiêu,
phạm vi và tiêu chí đánh giá cho toàn bộ đồ án. Nếu Chương 1 là phần
đặt vấn đề, thì Chương 2 sẽ chuyển sang câu hỏi quan trọng hơn: cần giải
quyết chính xác những bài toán kỹ thuật nào để biến mục tiêu đó thành
một hệ thống khả thi.

\chapterheading{CHƯƠNG 2. PHÂN TÍCH VẤN ĐỀ KỸ THUẬT}

Nếu Chương 1 đặt ra bối cảnh và mục tiêu của đề tài, thì Chương 2 đi
sâu vào phần cốt lõi hơn: bóc tách bài toán kỹ thuật thật sự của hệ
thống. Trọng tâm của chương không chỉ là liệt kê công nghệ, mà là làm rõ
vì sao mỗi quyết định về phần cứng, truyền thông, lưu trữ và kiến trúc
đều phải xuất phát từ nhu cầu vận hành thực tế của một hệ thống giám sát
phương tiện.

\subsection{2.1. Mô tả vấn đề -- Problem
statement}\label{21-muxf4-tux1ea3-vux1ea5n-ux111ux1ec1--problem-statement}

\subsubsection{2.1.1. Bối cảnh thực
tế}\label{211-bux1ed1i-cux1ea3nh-thux1ef1c-tux1ebf}

Trong bối cảnh ngành cho thuê ô tô và quản lý đội xe tại Việt Nam tiếp
tục mở rộng, nhu cầu giám sát phương tiện theo thời gian thực đã trở
thành yêu cầu thiết yếu đối với doanh nghiệp vận tải. Theo thống kê của
Bộ Giao thông Vận tải, số lượng ô tô cá nhân và thương mại tăng trung
bình 10--12\% mỗi năm trong giai đoạn 2020--2025, kéo theo áp lực ngày
càng lớn về quản lý, an ninh và tối ưu vận hành {[}1{]}. Tuy nhiên, điều
doanh nghiệp cần không chỉ là dữ liệu vị trí rời rạc, mà là khả năng
nhìn thấy trạng thái phương tiện đủ sớm để đưa ra quyết định vận hành.
Khoảng cách giữa nhu cầu này và cách làm hiện tại thể hiện ở các hạn chế
sau:

\begin{itemize}
\tightlist
\item
  \textbf{Giám sát thủ công}: Tài xế báo cáo vị trí qua điện thoại,
  không có dữ liệu liên tục, không thể xác minh chính xác hành trình.
  Phương pháp này phụ thuộc hoàn toàn vào yếu tố con người, dễ xảy ra
  sai sót và gian lận.
\item
  \textbf{Thiết bị GPS đơn giản}: Chỉ cung cấp tọa độ vị trí, không đọc
  được dữ liệu động cơ (tốc độ, vòng tua, nhiên liệu), không hỗ trợ phát
  hiện bất thường khi xe đậu. Thiếu khả năng tích hợp sâu với hệ thống
  quản lý.
\item
  \textbf{Giải pháp thương mại (fleet management)}: Chi phí cao (50--200
  USD/thiết bị + phí dịch vụ hàng tháng), phụ thuộc vào nhà cung cấp
  nước ngoài, khó tùy biến theo nhu cầu cụ thể của doanh nghiệp Việt
  Nam.
\end{itemize}

\subsubsection{2.1.2. Xác định vấn đề kỹ
thuật}\label{212-xuxe1c-ux111ux1ecbnh-vux1ea5n-ux111ux1ec1-kux1ef9-thuux1eadt}

Từ bối cảnh thực tiễn đó, bài toán của đề tài có thể quy về bốn câu hỏi
kỹ thuật nền tảng. Đây cũng là bốn câu hỏi chi phối gần như mọi quyết
định ở các chương sau, từ lựa chọn phần cứng đến thiết kế cloud:

\textbf{Vấn đề 1: Quản lý năng lượng trong môi trường ô tô}

Thiết bị theo dõi phải hoạt động liên tục 24/7 trong môi trường ô tô,
nơi nguồn điện luôn bị ràng buộc bởi dung lượng ắc quy. Khi xe tắt máy
(IGN OFF), thiết bị vẫn lấy điện từ ắc quy 12V hoặc 24V của xe; nếu duy
trì chế độ theo dõi liên tục, ắc quy có thể cạn sau vài tuần và làm xe
không khởi động được {[}2{]}. Vì vậy, câu hỏi không đơn thuần là "thiết
bị có hoạt động được không", mà là "thiết bị có hoạt động đủ thông minh
để không gây phiền toái cho chính chiếc xe mà nó đang giám sát hay
không".

\textbf{Vấn đề 2: Giám sát khi xe đậu}

Khi xe tắt máy, hệ thống cần tiết kiệm điện nhưng vẫn phát hiện được các
chuyển động bất thường như bị trộm hoặc bị kéo cẩu. Vì vậy, thiết bị
phải có cơ chế đánh thức nhanh khi phát hiện sự kiện, đồng thời duy trì
dòng tiêu thụ ở mức rất thấp. Nói cách khác, khi xe "im lặng" thì thiết
bị cũng phải "ngủ", nhưng vẫn phải tỉnh dậy đủ nhanh khi có sự cố.

\textbf{Vấn đề 3: Truyền dữ liệu qua mạng di động}

Hệ thống phải truyền liên tục dữ liệu GPS và OBD2 qua mạng 4G/LTE với độ
trễ thấp, đồng thời duy trì cơ chế reconnect ổn định khi mất sóng. Bên
cạnh đó, kiến trúc truyền thông cũng cần chừa dư địa để bổ sung lưu đệm
cục bộ nếu triển khai ở quy mô lớn hơn. Việc lựa chọn giao thức phù hợp
như MQTT, HTTP hay CoAP không chỉ là quyết định công nghệ, mà còn ảnh
hưởng trực tiếp đến hiệu suất và độ tin cậy toàn hệ thống.

\textbf{Vấn đề 4: Kiến trúc đám mây có khả năng mở rộng}

Hệ thống cần hỗ trợ nhiều phương tiện đồng thời, cập nhật vị trí theo
thời gian thực qua WebSocket, bảo đảm an toàn thông qua xác thực phiên
và xử lý được luồng dữ liệu lớn đến từ nhiều thiết bị IoT. Với người sử
dụng cuối, điều này được cảm nhận rất đơn giản: bản đồ phải cập nhật kịp
thời, cảnh báo phải đến đúng lúc, và dữ liệu phải nhất quán ngay cả khi
số lượng xe tăng lên.

\subsubsection{2.1.3. Mục tiêu kỹ
thuật}\label{213-mux1ee5c-tiuxeau-kux1ef9-thuux1eadt}

Từ bốn vấn đề nền tảng trên, đồ án rút ra bốn nhóm mục tiêu kỹ thuật.
Mỗi nhóm vừa trả lời một nhu cầu vận hành cụ thể của doanh nghiệp cho
thuê xe, vừa trở thành tiêu chí để sàng lọc và lựa chọn giải pháp ở
Chương 3:

\begin{enumerate}
\def\labelenumi{\arabic{enumi}.}
\tightlist
\item
  Thiết kế thiết bị tracker tích hợp GNSS và OBD2 dựa trên vi điều khiển
  ESP32-S3, đồng thời tối ưu tiêu thụ năng lượng khi xe đậu
  (mục tiêu dòng deep sleep toàn thiết bị dưới 500 µA trong phiên bản
  phần cứng hiện tại và tiếp tục giảm ở các vòng tối ưu PCB tiếp theo).
\item
  Xây dựng cơ chế quản lý nguồn thông minh với pin dự phòng và mạch Low
  Voltage Disconnect (LVD) để thiết bị bám được trạng thái xe nhưng không
  làm ảnh hưởng đến khả năng khởi động của phương tiện.
\item
  Thiết kế đường truyền dữ liệu sử dụng MQTT qua mạng 4G/LTE, ưu tiên tự
  phục hồi kết nối và chừa sẵn phương án lưu đệm khi mất mạng kéo dài.
\item
  Phát triển nền tảng đám mây có khả năng mở rộng, bao gồm API server,
  cơ sở dữ liệu quan hệ, cơ sở dữ liệu chuỗi thời gian và giao diện web
  theo dõi thời gian thực.
\end{enumerate}

\subsection{2.2. Bối cảnh và cơ sở kỹ thuật -- Background and Technical
reviews}\label{22-bux1ed1i-cux1ea3nh-vuxe0-cux1a1-sux1edf-kux1ef9-thuux1eadt--background-and-technical-reviews}

\subsubsection{2.2.1. Tổng quan về hệ thống IoT trong giám sát phương
tiện}\label{221-tux1ed5ng-quan-vux1ec1-hux1ec7-thux1ed1ng-iot-trong-giuxe1m-suxe1t-phux1b0ux1a1ng-tiux1ec7n}

Internet of Things (IoT) là mô hình cho phép các thiết bị vật lý kết nối
với nhau và với hạ tầng đám mây thông qua Internet. Ở góc nhìn đơn giản,
một hệ thống IoT giám sát phương tiện luôn phải đi qua ba bước: thu dữ
liệu trên xe, truyền dữ liệu lên mạng, rồi biến dữ liệu đó thành thông
tin mà người quản lý có thể đọc và hành động. Tương ứng với ba bước đó,
hệ thống thường được tổ chức thành ba lớp chính {[}3{]}:

\begin{itemize}
\tightlist
\item
  \textbf{Lớp cảm biến và thu thập dữ liệu (Perception Layer)}: Bao gồm
  các cảm biến GPS/GNSS, gia tốc kế (IMU), giao diện OBD2, và các cảm
  biến môi trường. Lớp này chịu trách nhiệm thu thập dữ liệu thô từ
  phương tiện.
\item
  \textbf{Lớp mạng và truyền thông (Network Layer)}: Sử dụng các giao
  thức truyền thông như MQTT, HTTP, CoAP qua hạ tầng mạng di động
  (4G/LTE) để truyền dữ liệu từ thiết bị lên đám mây.
\item
  \textbf{Lớp ứng dụng và xử lý (Application Layer)}: Bao gồm các dịch
  vụ đám mây (MQTT broker, API server, cơ sở dữ liệu) và giao diện người
  dùng (web dashboard, ứng dụng di động).
\end{itemize}

Trong triển khai của đề tài, kiến trúc ba lớp trên được bổ sung thêm một
lớp đệm dữ liệu cục bộ tại thiết bị bằng \textbf{microSD SDMMC 4-bit +
FATFS}. Có thể hình dung lớp này như một "vùng nhớ chờ", dùng để giữ lại
dữ liệu telemetry khi mạng 4G/LTE gián đoạn, sau đó phát lại (replay)
theo đúng thứ tự khi kết nối phục hồi. Cách tổ chức này giúp tăng độ
liên tục dữ liệu trong vận hành thực tế mà không làm thay đổi kiến trúc
cloud ở tầng trên.

\subsubsection{2.2.2. Các giải pháp giám sát phương tiện hiện
có}\label{222-cuxe1c-giux1ea3i-phuxe1p-giuxe1m-suxe1t-phux1b0ux1a1ng-tiux1ec7n-hiux1ec7n-cuxf3}

{\thesissettableid{2.1}{tab-ch2-vehicle-monitoring-solutions-comparison}
\def\LTcaptype{table}
\begin{longtable}{|L{0.179\linewidth}|L{0.141\linewidth}|L{0.167\linewidth}|L{0.227\linewidth}|L{0.231\linewidth}|}
\caption{So sánh các giải pháp giám sát phương tiện}\label{tab:ch2-vehicle-monitoring-solutions-comparison}\\
\hline \textbf{Tiêu chí} & \textbf{Giám sát thủ công} & \textbf{GPS Tracker đơn giản} & \textbf{Fleet Management
thương mại} & \textbf{Hệ thống đề xuất} \\
\hline
\endfirsthead
\caption[]{So sánh các giải pháp giám sát phương tiện (tiếp theo)}\\
\hline \textbf{Tiêu chí} & \textbf{Giám sát thủ công} & \textbf{GPS Tracker đơn giản} & \textbf{Fleet Management
thương mại} & \textbf{Hệ thống đề xuất} \\
\hline
\endhead
\hline
\endfoot
\endlastfoot
Chi phí thiết bị & Không & 500.000--1.500.000 VND & 1.000.000--5.000.000
VND & 1.514.000 VND \\\hline
Phí dịch vụ hàng tháng & Không & 50.000--100.000 VND & 200.000--500.000
VND & Chi phí 4G SIM (~70.000 VND) \\\hline
Độ chính xác vị trí & Thấp (báo cáo thủ công) & Trung bình (GPS) & Cao
(GPS + A-GPS) & Cao (GNSS đa hệ thống) \\\hline
Dữ liệu động cơ (OBD2) & Không & Không & Có (tùy model) & Có (BLE
OBD2) \\\hline
Phát hiện bất thường khi đậu & Không & Hạn chế & Có & Có (IMU + deep
sleep) \\\hline
Khả năng tùy biến & Không áp dụng & Thấp & Thấp (phụ thuộc vendor) & Cao
(mã nguồn mở) \\\hline
Quản lý năng lượng & Không áp dụng & Cơ bản & Tốt & Tốt (đa chế độ + pin
dự phòng) \\\hline
Tích hợp hệ thống & Không & Hạn chế (API riêng) & API do vendor & API mở
(REST + WebSocket) \\\hline
\end{longtable}
}

Bảng so sánh cho thấy phương án đề xuất giữ được các năng lực quan trọng
của hệ thống thương mại như độ chính xác vị trí, khả năng đọc OBD2 và
phát hiện bất thường, nhưng chi phí đầu tư thấp hơn và mức độ tùy biến
cao hơn. Với bài toán của doanh nghiệp cho thuê xe trong nước, đây là
hướng lựa chọn phù hợp hơn về cả kỹ thuật lẫn chi phí vận hành.

\subsubsection{2.2.3. So sánh giao thức truyền thông
IoT}\label{223-so-suxe1nh-giao-thux1ee9c-truyux1ec1n-thuxf4ng-iot}

Việc lựa chọn giao thức truyền thông không thể chỉ hiểu là chọn một
``chuẩn phổ biến''. Với hệ thống theo dõi phương tiện, đây là quyết định
chi phối trực tiếp lượng dữ liệu phải truyền, khả năng chịu mất sóng và
cách backend mở rộng về sau {[}4{]}.

{\thesissettableid{2.2}{tab-ch2-iot-communication-protocols}
\def\LTcaptype{table}
\begin{longtable}{|L{0.279\linewidth}|L{0.206\linewidth}|L{0.232\linewidth}|L{0.238\linewidth}|}
\caption{So sánh các giao thức truyền thông IoT}\label{tab:ch2-iot-communication-protocols}\\
\hline \textbf{Tiêu chí} & \textbf{MQTT} & \textbf{HTTP/REST} & \textbf{CoAP} \\
\hline
\endfirsthead
\caption[]{So sánh các giao thức truyền thông IoT (tiếp theo)}\\
\hline \textbf{Tiêu chí} & \textbf{MQTT} & \textbf{HTTP/REST} & \textbf{CoAP} \\
\hline
\endhead
\hline
\endfoot
\endlastfoot
Mô hình truyền thông & Publish--Subscribe & Request/Response &
Request/Response \\\hline
Giao thức tầng giao vận & TCP & TCP & UDP \\\hline
Overhead header & 2 bytes (tối thiểu) & Hàng trăm bytes & 4 bytes \\\hline
Chất lượng dịch vụ (QoS) & 3 mức (0, 1, 2) & Không có & 2 mức
(Confirmable/Non-confirmable) \\\hline
Tiêu thụ băng thông & Thấp & Cao & Thấp \\\hline
Hỗ trợ bi-directional & Có (Subscribe) & Không (cần polling hoặc WebSocket) &
Có (Observe) \\\hline
Phù hợp cho IoT & Rất tốt & Trung bình & Tốt \\\hline
Độ phổ biến & Rất cao & Rất cao & Trung bình \\\hline
Hỗ trợ EMQX Broker & Có (native) & Có (qua plugin) & Hạn chế \\\hline
\end{longtable}
}

\textbf{Phân tích lựa chọn:} Kết quả so sánh cho thấy MQTT phù hợp hơn
không phải vì nó ``cao cấp'' hơn HTTP hay CoAP, mà vì cách nó trao đổi
dữ liệu gần với nhịp vận hành của hệ thống: thiết bị gửi nhiều bản tin
nhỏ qua mạng di động không phải lúc nào cũng ổn định, đồng thời đôi lúc
cần nhận lệnh ngược từ server. Cụ thể:

\begin{itemize}
\tightlist
\item
  \textbf{Overhead thấp}: Header tối thiểu 2 bytes, giảm băng thông và
  tiết kiệm năng lượng cho thiết bị IoT chạy bằng pin.
\item
  \textbf{Mô hình Publish/Subscribe}: Cho phép nhiều thành phần nhận dữ
  liệu cùng lúc mà không làm thiết bị trên xe phải gửi lặp đi lặp lại
  cho từng nơi nhận. Điều này phù hợp với kiến trúc có dashboard, dịch
  vụ xử lý và các kênh cảnh báo cùng khai thác một nguồn dữ liệu.
\item
  \textbf{QoS linh hoạt}: Hệ thống không cần mọi bản tin đều được bảo
  đảm ở cùng một mức. QoS 0 phù hợp với dữ liệu telemetry tần suất cao
  như vị trí GPS, trong khi QoS 1 phù hợp hơn với cảnh báo và lệnh điều
  khiển vì các bản tin này cần được đảm bảo tới nơi ít nhất một lần.
\item
  \textbf{Retain Message và Last Will}: Hai cơ chế này giúp hệ thống
  phản ánh trạng thái online/offline của thiết bị rõ ràng hơn, từ đó
  người vận hành có thể phân biệt giữa "xe không có dữ liệu mới" và
  "thiết bị vừa mất kết nối".
\end{itemize}

\subsubsection{2.2.4. So sánh vi điều khiển
(MCU)}\label{224-so-suxe1nh-vi-ux111iux1ec1u-khiux1ec3n-mcu}

{\thesissettableid{2.3}{tab-ch2-mcu-comparison-for-tracker}
\def\LTcaptype{table}
\begin{longtable}{|L{0.230\linewidth}|L{0.218\linewidth}|L{0.260\linewidth}|L{0.246\linewidth}|}
\caption{So sánh các vi điều khiển cho ứng dụng IoT tracker}\label{tab:ch2-mcu-comparison-for-tracker}\\
\hline \textbf{Tiêu chí} & \textbf{ESP32-S3} & \textbf{STM32L4} & \textbf{Raspberry Pi Zero 2W} \\
\hline
\endfirsthead
\caption[]{So sánh các vi điều khiển cho ứng dụng IoT tracker (tiếp theo)}\\
\hline \textbf{Tiêu chí} & \textbf{ESP32-S3} & \textbf{STM32L4} & \textbf{Raspberry Pi Zero 2W} \\
\hline
\endhead
\hline
\endfoot
\endlastfoot
Kiến trúc & Xtensa LX7, dual-core, 240 MHz & ARM Cortex-M4, 80 MHz & ARM
Cortex-A53, quad-core, 1 GHz \\\hline
RAM & 512 KB SRAM + 8 MB PSRAM & 256 KB SRAM & 512 MB DRAM \\\hline
BLE & BLE 5.0 (tích hợp) & Không (cần module ngoài) & BLE 5.0 (tích
hợp) \\\hline
Wi-Fi & 802.11 b/g/n (tích hợp) & Không & 802.11 b/g/n (tích hợp) \\\hline
Tiêu thụ deep sleep & 10--15 µA & 1--2 µA & Không hỗ trợ deep sleep \\\hline
Giá thành (VND) & 80.000--150.000 & 150.000--300.000 &
400.000--600.000 \\\hline
Framework phát triển & Arduino / ESP-IDF & STM32CubeIDE / Mbed & Linux /
Python \\\hline
Cộng đồng hỗ trợ & Rất lớn & Lớn & Rất lớn \\\hline
Độ phù hợp cho tracker & Cao & Trung bình & Thấp \\\hline
\end{longtable}
}

\textbf{Phân tích lựa chọn:} Nếu chỉ nhìn riêng thông số low-power,
STM32L4 có lợi thế đáng kể. Tuy nhiên, khi đặt vào bối cảnh toàn hệ
thống - nơi thiết bị phải đồng thời xử lý BLE OBD2, modem, IMU và logic
trạng thái - ESP32-S3 cho thấy điểm cân bằng tốt hơn giữa tích hợp, chi
phí và độ phức tạp triển khai. Các lý do chính là:

\begin{itemize}
\tightlist
\item
  \textbf{BLE 5.0 tích hợp}: Kết nối trực tiếp với OBD2 adapter (vgate
  iCar Pro) mà không cần module BLE ngoài, giảm độ phức tạp phần cứng và
  chi phí.
\item
  \textbf{Chi phí thấp}: Rẻ hơn 30--50\% so với STM32L4 với cùng khả
  năng xử lý, phù hợp với ngân sách đồ án.
\item
  \textbf{Hệ sinh thái phát triển phong phú}: Hỗ trợ ESP-IDF (chính
  thức) và Arduino framework, cộng đồng lớn, nhiều thư viện có sẵn cho
  MQTT, BLE, GNSS. Điều này làm giảm đáng kể thời gian thử sai khi hiện
  thực đồ án.
\item
  \textbf{Deep sleep 10--15 µA ở mức MCU}: Mặc dù cao hơn STM32L4 (1--2
  µA), mức này vẫn là nền tảng phù hợp cho một thiết kế low-power. Trong
  hệ thống hoàn chỉnh, dòng ngủ sâu còn phụ thuộc mạch nguồn, cảm biến
  và linh kiện phụ trợ, nên chênh lệch ở mức riêng MCU không phải là yếu
  tố duy nhất quyết định tuổi thọ pin toàn thiết bị.
\end{itemize}

\subsubsection{2.2.5. So sánh cơ sở dữ
liệu}\label{225-so-suxe1nh-cux1a1-sux1edf-dux1eef-liux1ec7u}

{\thesissettableid{2.4}{tab-ch2-iot-database-solutions}
\def\LTcaptype{table}
\begin{longtable}{|L{0.203\linewidth}|L{0.285\linewidth}|L{0.240\linewidth}|L{0.226\linewidth}|}
\caption{So sánh các giải pháp cơ sở dữ liệu cho hệ thống IoT}\label{tab:ch2-iot-database-solutions}\\
\hline \textbf{Tiêu chí} & \textbf{PostgreSQL + VictoriaMetrics} & \textbf{MongoDB + InfluxDB} & \textbf{TimescaleDB} \\
\hline
\endfirsthead
\caption[]{So sánh các giải pháp cơ sở dữ liệu cho hệ thống IoT (tiếp theo)}\\
\hline \textbf{Tiêu chí} & \textbf{PostgreSQL + VictoriaMetrics} & \textbf{MongoDB + InfluxDB} & \textbf{TimescaleDB} \\
\hline
\endhead
\hline
\endfoot
\endlastfoot
Dữ liệu quan hệ (users, vehicles) & Rất tốt (PostgreSQL) & Trung bình
(MongoDB) & Tốt (PostgreSQL core) \\\hline
Dữ liệu chuỗi thời gian (telemetry) & Rất tốt (VictoriaMetrics) & Tốt
(InfluxDB) & Tốt (TimescaleDB extension) \\\hline
Hiệu suất ghi (write) & Cao (VM: hàng triệu điểm/giây) & Cao (InfluxDB:
hàng triệu điểm/giây) & Trung bình \\\hline
Nén dữ liệu & Rất tốt (VM: 10--70x nén) & Tốt (InfluxDB: TSM) & Tốt
(PostgreSQL compression) \\\hline
Tài nguyên tiêu thụ & Thấp (VM: 1--2 GB RAM) & Cao (InfluxDB: 4+ GB RAM)
& Trung bình \\\hline
Ngôn ngữ truy vấn & SQL (PG) + MetricsQL (VM) & MongoDB Query +
Flux/InfluxQL & SQL \\\hline
Tích hợp Grafana & Có (native) & Có & Có \\\hline
Giấy phép & Apache 2.0 & SSPL (MongoDB) + MIT (InfluxDB OSS) & Apache
2.0 (Community) \\\hline
Độ phức tạp vận hành & Trung bình (2 hệ thống riêng) & Cao (2 hệ thống
khác nhau) & Thấp (1 hệ thống) \\\hline
\end{longtable}
}

\textbf{Phân tích lựa chọn:} Với hệ thống này, không phải mọi dữ liệu
đều giống nhau. Thông tin người dùng, phương tiện hay hành trình cần cấu
trúc quan hệ chặt chẽ; trong khi dữ liệu vị trí và cảm biến lại là dữ
liệu biến đổi liên tục theo thời gian. Vì vậy, sự kết hợp PostgreSQL +
VictoriaMetrics được chọn vì:

\begin{itemize}
\tightlist
\item
  \textbf{Phân tách trách nhiệm rõ ràng}: PostgreSQL xử lý dữ liệu quan
  hệ như người dùng, phương tiện, cảnh báo và geofence; VictoriaMetrics
  xử lý dữ liệu chuỗi thời gian như tọa độ GPS, dữ liệu OBD2 và telemetry
  cảm biến. Mỗi hệ thống làm tốt phần việc mà nó được thiết kế để phục
  vụ.
\item
  \textbf{Hiệu suất và tài nguyên}: VictoriaMetrics có tỷ lệ nén dữ liệu
  rất cao (10--70x) và tiêu thụ ít RAM hơn InfluxDB, phù hợp với server
  có tài nguyên hạn chế.
\item
  \textbf{Tương thích Grafana}: Cả hai đều hỗ trợ Grafana native, thuận
  tiện cho việc xây dựng dashboard giám sát thống nhất mà không cần thêm
  công cụ trung gian.
\end{itemize}

{\thesissetfigureid{2.1}{fig-ch2-data-architecture-overview}
\begin{figure}[htbp]
\centering
\pandocbounded{\safeincludesvg{./assets/figures/02-chuong-2-phan-tich-hinh-2-1.svg}}
\caption{Sơ đồ kiến trúc dữ liệu của hệ thống - Data Architecture Diagram}
\label{fig:ch2-data-architecture-overview}
\end{figure}
}

\subsection{2.3. Yêu cầu kỹ thuật và các tiêu chuẩn thiết kế -- Design
criteria and
Constraints}\label{23-yuxeau-cux1ea7u-kux1ef9-thuux1eadt-vuxe0-cuxe1c-tiuxeau-chuux1ea9n-thiux1ebft-kux1ebf--design-criteria-and-constraints}

\subsubsection{2.3.1. Yêu cầu chức
năng}\label{231-yuxeau-cux1ea7u-chux1ee9c-nux103ng}

Những yêu cầu chức năng dưới đây đánh dấu bước chuyển từ ``bài toán''
sang ``năng lực hệ thống''. Có thể xem chúng như năm lời hứa vận hành tối
thiểu mà hệ thống phải thực hiện được để người quản lý thật sự dùng được
trong bối cảnh đội xe vận hành ngoài thực tế.

\textbf{FC-01: Theo dõi vị trí thời gian thực}

\begin{itemize}
\tightlist
\item
  Gửi vị trí GPS định kỳ (5--30 giây khi lái xe, 10--30 phút khi đậu xe)
\item
  Độ chính xác vị trí: dưới 5 mét trong điều kiện trời quang
\item
  Hỗ trợ đa hệ thống định vị (GPS, GLONASS, BeiDou)
\end{itemize}

\textbf{FC-02: Đọc dữ liệu động cơ qua OBD2}

\begin{itemize}
\tightlist
\item
  Kết nối BLE với OBD2 adapter (vgate iCar Pro)
\item
  Đọc trạng thái IGN (bật/tắt máy), RPM, tốc độ, nhiên liệu
\item
  Tần suất đọc: 5--30 giây khi xe chạy (mặc định 5 giây)
\end{itemize}

\textbf{FC-03: Phát hiện bất thường khi đậu xe}

\begin{itemize}
\tightlist
\item
  IMU (LIS3DSH) phát hiện chuyển động/rung bất thường
\item
  Đánh thức ESP32 từ deep sleep trong vòng 100ms
\item
  Gửi cảnh báo ưu tiên (priority alert) lên server
\end{itemize}

\textbf{FC-04: Quản lý năng lượng thông minh}

\begin{itemize}
\tightlist
\item
  Chuyển đổi giữa 3 chế độ: Lái xe (Active) - Đậu xe (Sleep) - Cảnh báo
  (Alert)
\item
  Pin dự phòng tự động tiếp quản theo profile nguồn: hệ 12V tại ngưỡng
  \(Switch_{\mathrm{OFF}} = 12.0\,\mathrm{V}\), hệ 24V tại ngưỡng
  \(Switch_{\mathrm{OFF}} = 24.0\,\mathrm{V}\)
\item
  Có phương án dự phòng suy luận trạng thái IGN qua ADC khi kết nối OBD2
  BLE không sẵn sàng
\end{itemize}

\textbf{FC-05: Giao diện web giám sát}

\begin{itemize}
\tightlist
\item
  Bản đồ thời gian thực (Leaflet) hiển thị vị trí tất cả phương tiện
\item
  Biểu đồ dữ liệu OBD2 (ECharts) - tốc độ, RPM, nhiên liệu
\item
  Hệ thống cảnh báo và thông báo (alert notifications)
\item
  Quản lý phương tiện, tài xế, khách hàng, hành trình
\end{itemize}

Năm nhóm chức năng trên bao quát trọn chuỗi nhu cầu cốt lõi của hệ
thống: biết xe đang ở đâu, hiểu xe đang như thế nào, phát hiện khi có
bất thường, không làm phiền hệ thống điện của xe và trình bày thông tin ở
dạng mà người vận hành có thể đọc nhanh, xử lý nhanh.

\subsubsection{2.3.2. Yêu cầu phi chức
năng}\label{232-yuxeau-cux1ea7u-phi-chux1ee9c-nux103ng}

Nếu yêu cầu chức năng trả lời ``hệ thống phải làm được gì'', thì yêu cầu
phi chức năng trả lời ``hệ thống phải làm tốt tới mức nào''. Đây là lớp
tiêu chí cho biết nguyên mẫu có đủ độ chín để bước ra môi trường vận
hành thật hay chưa, đặc biệt khi phải theo dõi nhiều xe và chạy liên
tục.

\textbf{NFR-01: Hiệu suất (Performance)}

\begin{itemize}
\tightlist
\item
  Độ trễ end-to-end (thiết bị đến dashboard): dưới 3 giây trong điều
  kiện mạng bình thường
\item
  API response time: dưới 200ms cho 95\% request
\item
  Hỗ trợ đồng thời tối thiểu 50 phương tiện kết nối MQTT, với dư địa mở
  rộng ổn định tới khoảng 100 thiết bị
\end{itemize}

\textbf{NFR-02: Độ tin cậy (Reliability)}

\begin{itemize}
\tightlist
\item
  Tự động kết nối lại (reconnect) khi mạng phục hồi
\item
  Giữ được cấu hình và trạng thái quan trọng qua NVS; buffer telemetry
  cục bộ là hướng mở rộng của thiết bị
\item
  MQTT QoS 1 cho các tin nhắn cảnh báo (đảm bảo gửi ít nhất một lần)
\end{itemize}

\textbf{NFR-03: Bảo mật (Security)}

\begin{itemize}
\tightlist
\item
  Xác thực phiên bằng opaque session token; chỉ lưu hash SHA-256 của
  token trong cơ sở dữ liệu
\item
  MQTT ACL (Access Control List) phân quyền theo từng thiết bị
\item
  Mã hóa TLS cho kết nối MQTT và HTTPS trong môi trường sản xuất
\end{itemize}

\textbf{NFR-04: Khả năng mở rộng (Scalability)}

\begin{itemize}
\tightlist
\item
  Kiến trúc microservices với Docker container độc lập
\item
  MQTT broker (EMQX) hỗ trợ phân cụm (clustering) để mở rộng
\item
  Cơ sở dữ liệu chuỗi thời gian (VictoriaMetrics) tối ưu cho luồng dữ
  liệu lớn
\end{itemize}

\textbf{NFR-05: Khả năng bảo trì (Maintainability)}

\begin{itemize}
\tightlist
\item
  Mã nguồn tổ chức theo mô hình Domain-Driven Design (DDD)
\item
  API RESTful với tài liệu Swagger/OpenAPI tự động
\item
  Logging tập trung với VictoriaLogs và dashboard Grafana
\end{itemize}

\subsubsection{2.3.3. Ràng buộc thiết
kế}\label{233-ruxe0ng-buux1ed9c-thiux1ebft-kux1ebf}

Các ràng buộc sau đây giới hạn không gian thiết kế của đề tài. Chúng buộc
nhóm phải chọn phương án cân bằng giữa tính khả thi học thuật, khả năng
triển khai thực tế trên xe và ngân sách sẵn có.

\textbf{Ràng buộc phần cứng:}

\begin{itemize}
\tightlist
\item
  Điện áp đầu vào: 12V hoặc 24V DC từ ắc quy xe (dao động phụ thuộc cấu
  hình hệ thống điện)
\item
  Nhiệt độ hoạt động: -10 độ C đến 70 độ C (môi trường trong xe ô tô)
\item
  Rung động và sốc: Chịu được rung động liên tục khi xe vận hành trên
  đường xá
\item
  Kích thước: Đủ nhỏ để lắp đặt kín đáo trong xe (không lớn hơn
  120x80x40 mm)
\end{itemize}

\textbf{Ràng buộc phần mềm:}

\begin{itemize}
\tightlist
\item
  Backend framework: Express.js + TypeScript (yêu cầu từ công nghệ hiện
  có của nhóm)
\item
  Frontend framework: Next.js 15 + React 19 (yêu cầu học thuật)
\item
  MQTT Broker: EMQX (giấy phép mã nguồn mở, hỗ trợ rules engine)
\item
  Triển khai: Docker Compose trên máy chủ Linux (mỗi dịch vụ có
  docker-compose.yml riêng)
\end{itemize}

\textbf{Ràng buộc về chi phí:}

\begin{itemize}
\tightlist
\item
  Tổng chi phí phần cứng: dưới 2.000.000 VND cho một bộ tracker
\item
  Ưu tiên các linh kiện có datasheet rõ ràng, chất lượng ổn định và
  nguồn cung trong nước dễ tiếp cận
\item
  Ưu tiên các giải pháp mã nguồn mở để giảm chi phí giấy phép
\end{itemize}

\subsubsection{2.3.4. Tiêu chuẩn thiết kế áp
dụng}\label{234-tiuxeau-chuux1ea9n-thiux1ebft-kux1ebf-uxe1p-dux1ee5ng}

Bảng dưới đây tổng hợp các chuẩn kỹ thuật và quy ước triển khai được dùng
làm nền cho hệ thống. Điểm quan trọng không phải là gom thật nhiều chuẩn,
mà là chọn đúng những chuẩn thực sự chi phối khả năng tương thích giữa
thiết bị, mạng truyền và ứng dụng.


{\thesissettableid{2.3.4A}{tab-ch2-tong-hop-thong-tin-ve-tieu-chuan-mo-ta-ap-dung-trong}
\def\LTcaptype{table}
\begin{longtable}{|L{0.231\linewidth}|L{0.278\linewidth}|L{0.455\linewidth}|}
\caption{Tổng hợp chuẩn và quy ước kỹ thuật áp dụng trong hệ thống}\label{tab:ch2-tong-hop-thong-tin-ve-tieu-chuan-mo-ta-ap-dung-trong}\\
\hline \textbf{Chuẩn / quy ước} & \textbf{Mô tả} & \textbf{Áp dụng trong hệ thống} \\
\hline
\endfirsthead
\caption[]{Tổng hợp chuẩn và quy ước kỹ thuật áp dụng trong hệ thống (tiếp theo)}\\
\hline \textbf{Chuẩn / quy ước} & \textbf{Mô tả} & \textbf{Áp dụng trong hệ thống} \\
\hline
\endhead
\hline
\endfoot
\endlastfoot
IEEE 802.15.1 (Bluetooth) & Chuẩn BLE 4.0/5.0 & Kết nối OBD2 adapter \\\hline
3GPP LTE Cat-4 & Chuẩn 4G/LTE & Truyền dữ liệu qua modem SIM7600CE-T
(Cat-4) \\\hline
NMEA 0183 & Chuẩn dữ liệu GPS & Phân tích tọa độ từ khối GNSS tích hợp \\\hline
SAE J1979 (OBD2) & Chuẩn chẩn đoán động cơ & Đọc dữ liệu qua OBD2 BLE
adapter \\\hline
MQTT v5.0 (OASIS) & Giao thức IoT messaging & Truyền dữ liệu thiết bị -
server \\\hline
REST (RFC 7231) & Kiến trúc API & Backend API server \\\hline
Opaque session token + MQTT ACL & Cơ chế xác thực và phân quyền &
Xác thực phiên người dùng, phân quyền truy cập thiết bị theo ACL
MQTT \\\hline
\end{longtable}
}


\subsection{2.4. Yêu cầu từ các bên liên quan -- Stakeholder
requirements}\label{24-yuxeau-cux1ea7u-tux1eeb-cuxe1c-buxean-liuxean-quan--stakeholder-requirements}

\subsubsection{2.4.1. Xác định các bên liên
quan}\label{241-xuxe1c-ux111ux1ecbnh-cuxe1c-buxean-liuxean-quan}

Hệ thống IoT Vehicle Tracking System liên quan đến nhiều nhóm sử dụng với
nhu cầu khác nhau. Nếu không tách rõ ngay từ đầu, hệ thống rất dễ rơi vào
tình trạng mạnh về mặt kỹ thuật nhưng lệch khỏi nhu cầu vận hành. Vì vậy,
phần này xác định ai là người hưởng lợi trực tiếp, ai là người vận hành,
ai là người bị tác động, từ đó giữ trọng tâm thiết kế vào những nhu cầu
thực sự cần giải quyết.

\subsubsection{2.4.2. Ma trận yêu cầu các bên liên
quan}\label{242-ma-trux1eadn-yuxeau-cux1ea7u-cuxe1c-buxean-liuxean-quan}


{\thesissettableid{2.5}{tab-ch2-ma-tran-yeu-cau-cac-ben-lien-quan}
\def\LTcaptype{table}
\begin{longtable}{|L{0.085\linewidth}|L{0.191\linewidth}|L{0.162\linewidth}|L{0.380\linewidth}|L{0.126\linewidth}|}
\caption{Ma trận yêu cầu các bên liên quan}\label{tab:ch2-ma-tran-yeu-cau-cac-ben-lien-quan}\\
\hline \textbf{STT} & \textbf{Bên liên quan} & \textbf{Vai trò} & \textbf{Yêu cầu chính} & \textbf{Mức độ ưu tiên} \\
\hline
\endfirsthead
\caption[]{Ma trận yêu cầu các bên liên quan (tiếp theo)}\\
\hline \textbf{STT} & \textbf{Bên liên quan} & \textbf{Vai trò} & \textbf{Yêu cầu chính} & \textbf{Mức độ ưu tiên} \\
\hline
\endhead
\hline
\endfoot
\endlastfoot
1 & Công ty cho thuê xe & Chủ sở hữu hệ thống & Giám sát vị trí tất cả
xe theo thời gian thực; Nhận cảnh báo khi xe đi ra khỏi vùng cho phép
(geofence); Báo cáo hành trình, quãng đường, nhiên liệu; Giảm thiểu thời
gian xe đậu không sử dụng & Cao \\\hline
2 & Tài xế / Người thuê xe & Người vận hành xe & Không bị giám sát quá
mức (quyền riêng tư); Thiết bị không ảnh hưởng đến vận hành xe; Không
rút cạn ắc quy xe & Trung bình \\\hline
3 & Quản trị viên hệ thống & Người quản lý kỹ thuật & Giao diện quản lý
dễ sử dụng; Hệ thống ổn định, ít lỗi; Có khả năng theo dõi trạng thái
thiết bị (online/offline); Quản lý người dùng và phân quyền & Cao \\\hline
4 & Đội bảo trì & Kỹ thuật viên & Thiết bị dễ lắp đặt và bảo trì; Chẩn
đoán lỗi từ xa (remote diagnostics); Cảnh báo bảo trì định kỳ (battery
low, thiết bị offline) & Trung bình \\\hline
5 & Khách hàng (người thuê xe cuối) & Người dùng gián tiếp & Quá trình
thuê xe nhanh gọn; Tin tưởng vào độ bảo mật của dữ liệu cá nhân; Nhận
thông báo về trạng thái đơn thuê & Thấp (Phase 2) \\\hline
\end{longtable}
}


\subsubsection{2.4.3. Phân tích chi tiết yêu
cầu}\label{243-phuxe2n-tuxedch-chi-tiux1ebft-yuxeau-cux1ea7u}

Từ ma trận trên, có thể thấy không phải mọi bên liên quan đều đòi hỏi cùng
một loại giá trị. Nhóm sử dụng chính quan tâm đến hiệu quả vận hành và an
ninh; nhóm quản trị quan tâm đến tính ổn định và khả năng kiểm soát; còn
nhóm bảo trì cần hệ thống dễ chẩn đoán và dễ can thiệp khi có sự cố.

\textbf{Công ty cho thuê xe (Stakeholder chính)}

Đây là nhóm đối tượng trọng tâm của hệ thống. Các yêu cầu của nhóm này
tập trung vào ba khía cạnh:

\begin{enumerate}
\def\labelenumi{\arabic{enumi}.}
\tightlist
\item
  \emph{Giám sát và an ninh}: Theo dõi vị trí phương tiện 24/7, phát
  hiện và cảnh báo ngay khi có bất thường (xe bị di chuyển trái phép,
  vượt qua geofence, tốc độ vượt ngưỡng). Hệ thống cần gửi thông báo qua
  nhiều kênh (web dashboard, Telegram bot, email).
\item
  \emph{Báo cáo và phân tích}: Xuất báo cáo hành trình chi tiết (quãng
  đường, thời gian, điểm dừng), thống kê nhiên liệu tiêu thụ, đánh giá
  hành vi lái xe (tốc độ trung bình, số lần phanh gấp). Các báo cáo này
  giúp tối ưu hóa chi phí vận hành và bảo trì.
\item
  \emph{Quản lý đội xe}: Quản lý trạng thái từng phương tiện (đang cho
  thuê, đang bảo trì, sẵn sàng), lịch bảo trì định kỳ dựa trên số km
  hoặc thời gian, lịch sử cho thuê và doanh thu theo phương tiện.
\end{enumerate}

\textbf{Quản trị viên hệ thống}

Yêu cầu của quản trị viên tập trung vào khả năng vận hành và giám sát hệ
thống:

\begin{enumerate}
\def\labelenumi{\arabic{enumi}.}
\tightlist
\item
  \emph{Dashboard tổng quan}: Hiển thị trạng thái tất cả thiết bị
  (online/offline, mức pin, cường độ tín hiệu), số lượng phương tiện
  đang hoạt động, cảnh báo chưa xử lý.
\item
  \emph{Quản lý thiết bị}: Thêm/xóa/sửa thông tin thiết bị tracker, cấp
  phát thiết bị cho phương tiện, gửi lệnh điều khiển từ xa (reset, cập
  nhật cấu hình, yêu cầu vị trí).
\item
  \emph{Quản lý người dùng}: Tạo tài khoản, phân quyền (admin, manager,
  viewer), theo dõi lịch sử đăng nhập và thao tác (audit log).
\end{enumerate}

\textbf{Đội bảo trì}

Yêu cầu của đội bảo trì hướng đến việc giảm thời gian và chi phí bảo
trì:

\begin{enumerate}
\def\labelenumi{\arabic{enumi}.}
\tightlist
\item
  \emph{Chẩn đoán từ xa}: Xem trạng thái thiết bị (mức pin, nhiệt độ,
  cường độ tín hiệu 4G) từ xa, phát hiện và cảnh báo khi thiết bị gặp
  lỗi (mất kết nối kéo dài, pin yếu).
\item
  \emph{Hướng dẫn lắp đặt}: Tài liệu hướng dẫn lắp đặt chi tiết, danh
  sách kiểm tra (checklist) khi lắp đặt mới hoặc bảo trì.
\end{enumerate}

Nhìn tổng thể, ba nhóm yêu cầu này kéo hệ thống về một điểm cân bằng khá
rõ: phải đủ mạnh để quản lý đội xe, nhưng cũng phải đủ gọn và ổn định để
lắp đặt, vận hành và bảo trì trên thực địa. Đây cũng là cơ sở để ma trận
truy xuất ở mục tiếp theo ánh xạ yêu cầu định tính sang các chức năng cụ
thể của hệ thống.

\subsubsection{2.4.4. Ma trận truy xuất yêu cầu - chức
năng}\label{244-ma-trux1eadn-truy-xuux1ea5t-yuxeau-cux1ea7u---chux1ee9c-nux103ng}


{\thesissettableid{2.6}{tab-ch2-ma-tran-truy-xuat-yeu-cau-cac-ben-lien-quan-voi-chuc}
\def\LTcaptype{table}
\begin{longtable}{|L{0.237\linewidth}|L{0.122\linewidth}|L{0.122\linewidth}|L{0.186\linewidth}|L{0.168\linewidth}|L{0.101\linewidth}|}
\caption{Ma trận truy xuất yêu cầu các bên liên quan với chức năng hệ thống}\label{tab:ch2-ma-tran-truy-xuat-yeu-cau-cac-ben-lien-quan-voi-chuc}\\
\hline \textbf{Yêu cầu bên liên quan} & \textbf{FC-01 (Vị trí)} & \textbf{FC-02 (OBD2)} & \textbf{FC-03 (Bất
thường)} & \textbf{FC-04 (Năng lượng)} & \textbf{FC-05 (Web)} \\
\hline
\endfirsthead
\caption[]{Ma trận truy xuất yêu cầu các bên liên quan với chức năng hệ thống (tiếp theo)}\\
\hline \textbf{Yêu cầu bên liên quan} & \textbf{FC-01 (Vị trí)} & \textbf{FC-02 (OBD2)} & \textbf{FC-03 (Bất
thường)} & \textbf{FC-04 (Năng lượng)} & \textbf{FC-05 (Web)} \\
\hline
\endhead
\hline
\endfoot
\endlastfoot
Giám sát vị trí 24/7 & X & & & X & X \\\hline
Cảnh báo geofence & X & & & & X \\\hline
Báo cáo nhiên liệu & & X & & & X \\\hline
Phát hiện trộm xe & & & X & X & X \\\hline
Không rút cạn ắc quy & & & & X & \\\hline
Quản lý thiết bị từ xa & & & & & X \\\hline
Chẩn đoán lỗi từ xa & X & X & & X & X \\\hline
\end{longtable}
}


Ma trận truy xuất trên cho thấy các chức năng hệ thống (FC-01 đến FC-05)
bao phủ toàn bộ yêu cầu của các bên liên quan chính. Trong đó, FC-04
(Quản lý năng lượng) và FC-05 (Giao diện web) là hai chức năng được nhắc
đến nhiều nhất, qua đó khẳng định tầm quan trọng của thiết kế quản lý
năng lượng thông minh và giao diện giám sát thân thiện.

\begin{center}\rule{0.5\linewidth}{0.5pt}\end{center}

\subsection{Kết luận chương
2}\label{kux1ebft-luux1eadn-chux1b0ux1a1ng-2}

Chương 2 đã chuyển bài toán ở mức "ý tưởng" thành các ràng buộc kỹ
thuật cụ thể: thiết bị phải tiết kiệm điện ra sao, truyền dữ liệu thế
nào, lưu trữ ở đâu và kiến trúc cần mở rộng theo hướng nào. Thông qua
quá trình khảo sát và đối sánh các giải pháp hiện có, chương đã làm rõ
tiêu chí lựa chọn thay vì chỉ liệt kê công nghệ. Đây chính là cơ sở để
Chương 3 bước sang phần quyết định quan trọng tiếp theo: chọn phương án
thiết kế nào là phù hợp nhất với bài toán đã được phân tích.

\chapterheading{CHƯƠNG 3. CÁC GIẢI PHÁP THIẾT KẾ -- DESIGN SOLUTIONS}

Sau khi xác định rõ yêu cầu và ràng buộc ở Chương 2, Chương 3 tập trung
vào câu hỏi "chọn gì và vì sao chọn như vậy". Trình bày trong chương này
không chỉ nêu phương án cuối cùng, mà còn giải thích quá trình cân nhắc
giữa nhiều lựa chọn khác nhau ở bốn tầng: phần cứng thiết bị theo dõi,
firmware nhúng trên vi điều khiển, hệ thống backend/cloud và giao diện
người dùng. Nhờ đó, người đọc có thể theo dõi logic thiết kế theo từng
bước thay vì chỉ nhìn thấy kết quả cuối cùng.

\subsection{3.1. Phân tích và đề xuất giải pháp phần cứng -- Hardware
analysis and solution}\label{31-phuxe2n-tuxedch-tux1ed5ng-hux1ee3p--general-analysis}

\subsubsection{3.1.1. Phân tích và lựa chọn phần
cứng}\label{311-phuxe2n-tuxedch-vuxe0-lux1ef1a-chux1ecdn-phux1ea7n-cux1ee9ng}

Phần này trình bày quá trình phân tích, lựa chọn và thiết kế các thành
phần phần cứng cho thiết bị theo dõi xe IoT. Mục tiêu không chỉ là tạo
ra một tracker "đủ chức năng" trên lý thuyết, mà là một thiết bị có thể
lắp lên xe, vận hành ổn định và còn chừa dư địa cho bảo trì, mở rộng.
Vì vậy, mỗi quyết định kỹ thuật đều được soi chiếu dưới bốn góc nhìn:
độ tin cậy, khả năng hoạt động liên tục với nguồn dự phòng, mức độ tích
hợp và mức độ phù hợp với ngân sách.

\textbf{Quy ước thuật ngữ trong phần cứng:} tài liệu dùng cụm ``khối phần
cứng/khối chức năng'' cho các thành phần tích hợp trên PCB; từ ``module''
chỉ dùng khi nói về module phần mềm (firmware/cloud) hoặc khi trích dẫn
nguyên văn tên sản phẩm của nhà sản xuất.

\paragraph{3.1.1.1. Sơ đồ khối hệ
thống}\label{3111-sux1a1-ux111ux1ed3-khux1ed1i-hux1ec7-thux1ed1ng}

Hệ thống tracker được tổ chức theo kiến trúc khối chức năng với năm khối
chính, trong đó ESP32-S3 đóng vai trò điều phối trung tâm. Nếu nhìn dưới
góc độ người sử dụng, năm khối này tương ứng với năm câu hỏi rất thực tế:
thiết bị lấy dữ liệu xe ở đâu, xác định vị trí bằng cách nào, phát hiện
rung/chuyển động ra sao, gửi dữ liệu về máy chủ như thế nào và lấy điện từ
đâu để duy trì hoạt động. Sơ đồ khối tổng thể được trình bày như sau:


{\thesissetfigureid{3.1}{fig-ch3-so-o-khoi-tong-the-he-thong-tracker}
\begin{figure}[htbp]
\centering
\pandocbounded{\safeincludesvg{./assets/figures/03-chuong-3-giai-phap-phan-cung-hinh-3-1.svg}}
\caption{Sơ đồ khối tổng thể hệ thống tracker}
\label{fig:ch3-so-o-khoi-tong-the-he-thong-tracker}
\end{figure}
}



{\thesissetfigureid{3.1a}{fig-ch3-phan-ra-chi-tiet-cac-khoi-chuc-nang-cua-tracker}
\begin{figure}[htbp]
\centering
\pandocbounded{\safeincludesvg{./assets/figures/03-chuong-3-giai-phap-phan-cung-hinh-3-1a.svg}}
\caption{Phân rã chi tiết các khối chức năng của tracker}
\label{fig:ch3-phan-ra-chi-tiet-cac-khoi-chuc-nang-cua-tracker}
\end{figure}
}


Ở tầng thiết bị, khối OBD2 cung cấp trạng thái vận hành của xe; khối
LTE/GNSS chịu trách nhiệm truyền dữ liệu và định vị; khối IMU duy trì khả
năng giám sát khi xe tắt máy; khối nguồn đảm bảo thiết bị chuyển đổi giữa
nguồn chính và nguồn dự phòng an toàn; còn khối vi điều khiển là nơi tổng
hợp toàn bộ tín hiệu để ra quyết định theo trạng thái. Nhờ cách tách này,
tracker không còn được nhìn như một bo mạch rời rạc, mà như một chuỗi chức
năng nối tiếp từ thu thập dữ liệu, xử lý cục bộ, truyền tin cho đến bảo vệ
nguồn.

Hệ thống vận hành theo ba chế độ chính: (1) chế độ lái xe --- khi động cơ
bật (IGN ON), các khối chức năng cần thiết được kích hoạt; (2) chế độ đỗ
xe --- khi động cơ tắt (IGN OFF), ESP32 chuyển sang deep sleep và chỉ IMU
LIS3DSH duy trì giám sát chuyển động; (3) chế độ cảnh báo --- khi IMU ghi
nhận chuyển động bất thường, hệ thống tự đánh thức và gửi cảnh báo qua 4G.

Luồng dữ liệu được tổ chức theo ba nhánh: OBD2 (RPM, tốc độ, nhiên liệu)
đọc qua BLE từ adapter vgate iCar Pro; dữ liệu GNSS được cung cấp trực
tiếp bởi khối modem tích hợp SIMCom SIM7600CE-T (GNSS nội bộ) trên cùng
UART với LTE; và dữ liệu chuyển động thu từ IMU LIS3DSH qua I2C.
ESP32-S3 tổng hợp, đóng gói và truyền toàn bộ dữ liệu lên máy chủ qua
MQTT thông qua kết nối 4G/LTE của SIM7600CE-T.

Để tăng độ liên tục dữ liệu khi mạng gián đoạn, luồng telemetry được bổ
sung thêm nhánh lưu đệm cục bộ trên microSD (SDMMC 4-bit + FATFS). Dữ
liệu được ghi theo thứ tự bản ghi, sau đó replay lên MQTT khi kết nối
phục hồi.

\paragraph{3.1.1.2. Phân tích và lựa chọn vi điều khiển
(MCU)}\label{3112-phuxe2n-tuxedch-vuxe0-lux1ef1a-chux1ecdn-vi-ux111iux1ec1u-khiux1ec3n-mcu}

Việc lựa chọn vi điều khiển (MCU) là quyết định thiết kế quan trọng
nhất, ảnh hưởng trực tiếp đến khả năng tích hợp các khối ngoại vi, mức
tiêu thụ năng lượng, và độ phức tạp phát triển firmware. Ba ứng viên
được đưa vào quá trình đánh giá là ESP32-S3 của Espressif, STM32L4 của
STMicroelectronics, và nRF52840 của Nordic Semiconductor.

\subparagraph{a) Các ứng viên}\label{a-cuxe1c-ux1ee9ng-viuxean}

\textbf{ESP32-S3} là vi điều khiển dual-core Xtensa LX7 hoạt động tới
240 MHz, tích hợp sẵn WiFi 802.11 b/g/n và Bluetooth LE 5.0 {[}19{]}.
Thiết bị có 512 KB SRAM on-chip, dùng flash ngoài, hỗ trợ nhiều giao
tiếp ngoại vi (UART/I2C/SPI/ADC) và dòng deep-sleep mức µA tùy cấu hình
{[}19{]}, {[}20{]}.

\textbf{STM32L4} (đại diện bởi STM32L476 trong so sánh) là dòng MCU
low-power của ST, dùng nhân ARM Cortex-M4 tới 80 MHz và không tích hợp
BLE/Wi-Fi {[}53{]}, {[}54{]}. Dòng này nổi bật với các mức dòng ngủ sâu
rất thấp theo cấu hình nguồn (shutdown/standby/STOP2) {[}53{]}.

\textbf{nRF52840} là MCU của Nordic với nhân ARM Cortex-M4F 64 MHz, hỗ
trợ Bluetooth LE 5.x và 802.15.4 {[}55{]}, {[}56{]}. Thiết bị có 1 MB
Flash, 256 KB RAM, hỗ trợ UART/UARTE và dòng System OFF mức dưới µA tùy
điều kiện cấu hình {[}55{]}, {[}56{]}.

\subparagraph{b) Đối chiếu yêu cầu thiết kế với thông số phần
cứng}\label{b-ux111ux1ed1i-chiux1ebfu-yuxeau-cux1ea7u-thiux1ebft-kux1ebf-vux1edbi-thuxf4ng-sux1ed1-phux1ea7n-cux1ee9ng}

Thay vì chấm điểm tổng quát, \tabref{tab:ch3-oi-chieu-yeu-cau-ky-thuat-khi-lua-chon-vi-ieu-khien} đối chiếu trực tiếp các ràng buộc
kỹ thuật của tracker với thông số và tác động tích hợp của từng MCU.


{\thesissettableid{3.1}{tab-ch3-oi-chieu-yeu-cau-ky-thuat-khi-lua-chon-vi-ieu-khien}
\def\LTcaptype{table}
\begin{longtable}{|L{0.153\linewidth}|L{0.247\linewidth}|L{0.299\linewidth}|L{0.255\linewidth}|}
\caption{Đối chiếu yêu cầu kỹ thuật khi lựa chọn vi điều khiển}\label{tab:ch3-oi-chieu-yeu-cau-ky-thuat-khi-lua-chon-vi-ieu-khien}\\
\hline \textbf{Yêu cầu của tracker} & \textbf{ESP32-S3} & \textbf{STM32L4} & \textbf{nRF52840} \\
\hline
\endfirsthead
\caption[]{Đối chiếu yêu cầu kỹ thuật khi lựa chọn vi điều khiển (tiếp theo)}\\
\hline \textbf{Yêu cầu của tracker} & \textbf{ESP32-S3} & \textbf{STM32L4} & \textbf{nRF52840} \\
\hline
\endhead
\hline
\endfoot
\endlastfoot
Kết nối OBD2 BLE với vgate iCar Pro & BLE 5.0 tích hợp, tương thích
ngược BLE 4.0 & Không tích hợp BLE, phải thêm module ngoài & BLE 5.0
tích hợp \\\hline
Cấu hình UART của hệ thống & 3 UART controllers (đủ cho modem + debug,
còn dư cho mở rộng) {[}19{]}, {[}20{]} & Nhiều serial instance trong
dòng STM32L476 (USART, UART, LPUART), đáp ứng yêu cầu số cổng {[}53{]},
{[}54{]} & UARTE hoặc UART thường dùng 2 instance, ít dư địa nếu cần giữ
debug UART riêng {[}55{]}, {[}56{]} \\\hline
Tài nguyên xử lý cho OBD2 + modem + IMU & Dual-core 240 MHz, 512 KB SRAM
{[}19{]} & Cortex-M4 tới 80 MHz, 128 KB SRAM trên STM32L476RE {[}54{]} &
Cortex-M4F 64 MHz, 256 KB RAM {[}55{]}, {[}56{]} \\\hline
Dòng ngủ sâu/siêu thấp công suất & Deep-sleep mức µA tùy cấu hình RTC/IO
{[}19{]}, {[}20{]} & STOP2 cỡ 1.1 µA (theo điều kiện datasheet),
standby/shutdown cỡ nA {[}53{]}, {[}54{]} & System OFF mức dưới µA tùy
RAM retention/GPIO {[}55{]} \\\hline
Framework triển khai & Arduino IDE + ESP-IDF {[}47{]} & STM32Cube/HAL
{[}53{]} & nRF Connect SDK {[}57{]} \\\hline
Phần cứng bổ sung để đạt đủ chức năng & Không cần thêm radio BLE & Cần
thêm BLE module nếu vẫn giữ OBD2 BLE (STM32L476 không tích hợp BLE)
{[}53{]}, {[}54{]} & Không cần BLE module, nhưng cần xử lý bài toán
thiếu serial cho debug \\\hline
Chi phí triển khai phần cứng tham chiếu &
~100.000--200.000 VND (mặt bằng BOM dự án) &
~150.000--300.000 VND (mặt bằng BOM dự án) & Thường cao
hơn ESP32-S3 và ít phổ biến hơn trong bối cảnh dự án \\\hline
Tác động tới sơ đồ hiện tại & Giữ nguyên sơ đồ BLE + modem UART, thuận
lợi cho tích hợp GNSS trong modem & Tăng phần cứng BLE ngoài và công
đoạn tích hợp ban đầu & Có thể phải chuyển debug sang USB/SWO để nhường
UART cho chức năng chính \\\hline
\end{longtable}
}


\subparagraph{c) Phân tích các yếu tố quan
trọng}\label{c-phuxe2n-tuxedch-cuxe1c-yux1ebfu-tux1ed1-quan-trux1ecdng}

\textbf{BLE tích hợp:} Đây là ràng buộc rất thực tế vì tracker dùng
adapter OBD2 vgate iCar Pro theo chuẩn BLE 4.0. Nếu MCU không có BLE
tích hợp, thiết kế sẽ phải bổ sung thêm một module riêng. Điều đó không
chỉ làm tăng BOM, mà còn khiến sơ đồ mạch, nguồn cấp và firmware phức
tạp hơn ngay từ đầu. Ở khía cạnh này, ESP32-S3 và nRF52840 thuận lợi hơn
STM32L4.

\textbf{Ngân sách UART:} Trong kiến trúc hiện tại, UART không phải tài
nguyên "thừa". Hệ thống cần ít nhất một kênh cho modem và một kênh cho
debug, đồng thời vẫn nên chừa dư địa cho mở rộng hoặc chẩn đoán khi
triển khai thực tế. ESP32-S3 đáp ứng tốt nhu cầu này với 3 UART phần
cứng; nRF52840 hạn chế hơn; còn STM32L4 có thể đáp ứng nhưng phải đánh
đổi bằng việc bổ sung BLE ngoài.

\textbf{Tiêu thụ deep sleep:} STM32L4 và nRF52840 có lợi thế rõ ràng về
dòng ngủ sâu theo tài liệu hãng. Tuy nhiên, trong một tracker hoàn chỉnh,
MCU không phải thành phần duy nhất tiêu thụ năng lượng. Modem LTE, GNSS,
chu kỳ wake-up và các linh kiện phụ trợ cũng góp phần đáng kể vào ngân
sách điện năng. Vì vậy, chênh lệch ở mức riêng MCU có ý nghĩa, nhưng
không đủ để bù cho việc tăng độ phức tạp tích hợp {[}19{]}, {[}53{]},
{[}54{]}, {[}55{]}.

\textbf{Độ phức tạp triển khai:} ESP32-S3 có lợi thế thực tế vì vừa có
BLE sẵn, vừa có hệ sinh thái ESP-IDF/Arduino IDE quen thuộc, phù hợp cho
firmware phải đồng thời xử lý BLE OBD2, AT command modem, I2C IMU và
deep sleep. STM32L4 mạnh về low power nhưng kéo theo thêm công việc tích
hợp BLE. nRF52840 phù hợp nếu hệ thống ưu tiên radio BLE là chính, nhưng
dư địa UART cho kịch bản hiện tại hạn chế hơn.

\subparagraph{d) Kết luận lựa
chọn}\label{d-kux1ebft-luux1eadn-lux1ef1a-chux1ecdn}

Trên cơ sở đối chiếu trực tiếp các yêu cầu của hệ thống,
\textbf{ESP32-S3} được chọn làm vi điều khiển chính cho tracker. Nói ngắn
gọn, đây là phương án cân bằng tốt nhất giữa "đủ mạnh để làm được việc"
và "đủ gọn để triển khai thực tế". Quyết định này dựa trên các điểm chốt
sau:

\begin{itemize}
\tightlist
\item
  Đáp ứng đầy đủ kiến trúc hiện tại: 1 kết nối BLE cho OBD2, UART cho
  modem và một đường debug riêng, đồng thời vẫn còn dư địa UART cho mở
  rộng.
\item
  BLE 5.0 tích hợp giúp bỏ hẳn module BLE ngoài, giảm BOM và rút ngắn
  công đoạn tích hợp.
\item
  Tài nguyên xử lý 240 MHz dual-core và 512 KB SRAM đủ để chạy đồng thời
  BLE, modem, IMU và state machine.
\item
  Chi phí triển khai phần cứng với ESP32-S3 thấp hơn STM32L4 khoảng
  30--50\% theo mặt bằng linh kiện đang dùng trong đồ án.
\item
  Dòng deep-sleep của ESP32-S3 cao hơn STM32L4 theo tài liệu hãng, nhưng
  vẫn nằm trong giới hạn chấp nhận được của bài toán pin backup khi xét
  toàn bộ duty-cycle của modem, GNSS và các chu kỳ wake-up {[}19{]},
  {[}53{]}, {[}54{]}.
\end{itemize}


{\thesissettableid{3.2}{tab-ch3-thong-so-ky-thuat-esp32-s3-uoc-chon}
\def\LTcaptype{table}
\begin{longtable}{|L{0.274\linewidth}|L{0.694\linewidth}|}
\caption{Thông số kỹ thuật ESP32-S3 được chọn}\label{tab:ch3-thong-so-ky-thuat-esp32-s3-uoc-chon}\\
\hline \textbf{Thông số} & \textbf{Giá trị} \\
\hline
\endfirsthead
\caption[]{Thông số kỹ thuật ESP32-S3 được chọn (tiếp theo)}\\
\hline \textbf{Thông số} & \textbf{Giá trị} \\
\hline
\endhead
\hline
\endfoot
\endlastfoot
CPU & Dual-core Xtensa LX7, tối đa 240 MHz {[}19{]} \\\hline
RAM on-chip & 512 KB SRAM {[}19{]} \\\hline
Flash & Flash ngoài (biến thể WROOM thường 4--16 MB) {[}19{]},
{[}20{]} \\\hline
Bluetooth & Bluetooth LE 5.0, tương thích ngược BLE 4.0 cho OBD2 adapter
{[}19{]} \\\hline
WiFi & 802.11 b/g/n (không dùng trong luồng chính của tracker)
{[}19{]} \\\hline
UART & 3 UART controllers {[}19{]}, {[}20{]} \\\hline
I2C & 2 I2C controllers {[}19{]}, {[}20{]} \\\hline
ADC & SAR ADC 12-bit, tối đa 20 kênh (theo package) {[}19{]} \\\hline
GPIO & Tối đa 45 GPIO (theo package) {[}19{]} \\\hline
Dòng deep-sleep & Mức µA tùy cấu hình RTC domain và IO {[}19{]},
{[}20{]} \\\hline
Biến thể linh kiện & ESP32-S3-WROOM-1 \\\hline
Phương án tích hợp & Tích hợp trực tiếp vào bo mạch chính \\\hline
\end{longtable}
}


\paragraph{3.1.1.3. Phân tích và lựa chọn modem truyền
thông}\label{3113-phuxe2n-tuxedch-vuxe0-lux1ef1a-chux1ecdn-modem-truyux1ec1n-thuxf4ng}

Modem truyền thông đóng vai trò then chốt trong hệ thống vì nó đồng thời
giải quyết hai nhu cầu thiết yếu nhất của tracker: gửi dữ liệu về máy chủ
qua mạng di động và cung cấp vị trí của xe thông qua GNSS. Với bài toán
này, lựa chọn modem không chỉ ảnh hưởng đến tốc độ hay thông số kỹ
thuật, mà còn tác động trực tiếp đến cách đi dây, cách tổ chức firmware
và mức độ phức tạp của toàn bộ bo mạch. Trong kiến trúc hiện tại, hệ
thống sử dụng \textbf{SIM7600CE-T tích hợp LTE + GNSS} nhằm giảm số
lượng phần cứng rời, đơn giản hóa đi dây và đồng bộ với firmware thực
tế.

\subparagraph{a) Các ứng viên}\label{a-cuxe1c-ux1ee9ng-viuxean-1}

Ba phương án được đưa vào đánh giá gồm hai phương án tách rời
(\textbf{SIMCom A7670C + u-blox NEO-M8N}, \textbf{Quectel EC200U-CN +
NEO-M8N}) và một phương án tích hợp (\textbf{SIMCom SIM7600CE-T}).

\textbf{SIMCom A7670C + u-blox NEO-M8N} là phương án tách rời: A7670C
phụ trách LTE/AT command, NEO-M8N phụ trách GNSS/NMEA. Theo trang sản
phẩm chính thức của SIMCom, A7670C nằm trong nhóm LTE Cat-1 và dùng
nguồn 3.4--4.2 V {[}58{]}, {[}62{]}.

\textbf{Quectel EC200U-CN + NEO-M8N} có cấu trúc gần giống phương án
A7670C + NEO-M8N. Tuy nhiên dòng EC200U được hãng mô tả là Cat-1 bis,
GNSS là tính năng tùy chọn theo biến thể/cấu hình và I/O logic có ràng
buộc riêng, nên mức độ phù hợp với cấu hình đang triển khai thấp hơn
{[}59{]}, {[}61{]}.

\textbf{SIMCom SIM7600CE-T} là phương án tích hợp GNSS trong modem, giúp
đi dây gọn hơn và giảm số lượng phần cứng rời {[}60{]}.

\subparagraph{b) Đối chiếu kiến trúc LTE/GNSS theo ràng buộc tích
hợp}\label{b-ux111ux1ed1i-chiux1ebfu-kiux1ebfn-truxfac-ltegnss-theo-ruxe0ng-buux1ed9c-tuxedch-hux1ee3p}

Vấn đề cốt lõi của tracker không chỉ là modem có lên mạng được hay
không, mà còn là phương án LTE/GNSS ảnh hưởng thế nào tới sơ đồ UART, hồ
sơ tiêu thụ nguồn và mức độ thuận lợi khi phát triển firmware. \tabref{tab:ch3-oi-chieu-cac-phuong-an-lte-gnss-theo-tac-ong-tich-hop}
dưới đây đối chiếu trực tiếp các điểm khác biệt này.


{\thesissettableid{3.3}{tab-ch3-oi-chieu-cac-phuong-an-lte-gnss-theo-tac-ong-tich-hop}
\def\LTcaptype{table}
\begin{longtable}{|L{0.202\linewidth}|L{0.260\linewidth}|L{0.252\linewidth}|L{0.241\linewidth}|}
\caption{Đối chiếu các phương án LTE/GNSS theo tác động tích hợp}\label{tab:ch3-oi-chieu-cac-phuong-an-lte-gnss-theo-tac-ong-tich-hop}\\
\hline \textbf{Yêu cầu tích hợp} & \textbf{A7670C + NEO-M8N} & \textbf{EC200U-CN + NEO-M8N} & \textbf{SIM7600CE-T} \\
\hline
\endfirsthead
\caption[]{Đối chiếu các phương án LTE/GNSS theo tác động tích hợp (tiếp theo)}\\
\hline \textbf{Yêu cầu tích hợp} & \textbf{A7670C + NEO-M8N} & \textbf{EC200U-CN + NEO-M8N} & \textbf{SIM7600CE-T} \\
\hline
\endhead
\hline
\endfoot
\endlastfoot
Số khối phần cứng rời & 2 khối rời: A7670C (LTE) + NEO-M8N (GNSS)
{[}58{]}, {[}62{]} & 2 khối rời: EC200U-CN (LTE) + NEO-M8N (GNSS)
{[}59{]}, {[}61{]} & 1 khối tích hợp LTE + GNSS (SIM7600CE-T) {[}60{]} \\\hline
LTE category / tốc độ & Cat-1, tối đa 10 Mbps DL / 5 Mbps UL {[}58{]} &
Cat-1, tối đa 10 Mbps DL / 5 Mbps UL {[}59{]}, {[}61{]} & Cat-4, tối đa
150 Mbps DL / 50 Mbps UL {[}60{]} \\\hline
GNSS tích hợp trong modem & Không tích hợp GNSS trên A7670C, cần GNSS
ngoài {[}58{]}, {[}62{]} & GNSS tùy chọn theo variant/cấu hình dòng
EC200U {[}59{]}, {[}61{]} & Có GNSS tích hợp trong module {[}60{]} \\\hline
Điện áp cấp nguồn modem & 3.4--4.2 V (typ. 3.8 V) {[}58{]} & 3.3--4.3 V
(typ. 3.8 V) {[}61{]} & 3.4--4.2 V (typ. 3.8 V) {[}60{]} \\\hline
Đường dữ liệu vị trí trong kiến trúc đề tài & NMEA/UBX từ NEO-M8N qua
UART riêng {[}62{]}, {[}63{]} & NMEA/UBX từ NEO-M8N qua UART riêng
{[}62{]}, {[}63{]} & GNSS AT tích hợp trong cùng modem {[}60{]} \\\hline
Bật/tắt LTE và GNSS độc lập & Có, do tách 2 module/2 rail nguồn & Có,
nếu triển khai tách tương tự A7670C+NEO-M8N & Không tách hoàn toàn do
GNSS đi cùng modem \\\hline
Tác động khi reset modem LTE & Có thể reset modem và giữ logic GNSS
riêng & Tương tự phương án tách nếu GNSS dùng module riêng & Reset modem
thường kéo theo gián đoạn GNSS \\\hline
Mức phụ thuộc giữa các chức năng firmware & Thấp, tách rõ
\texttt{modem\_lte} và \texttt{gnss} & Thấp, nhưng phải đổi tập lệnh
modem & Cao hơn do LTE/GNSS chung khối tích hợp \\\hline
Tác động tới BOM phần cứng thực tế & Bộ modem + GNSS đang chốt
330.000--500.000 VND (BOM dự án) & Chưa phải cấu hình BOM chính của repo
& Giảm số khối phần cứng rời, thuận lợi cho đi dây và lắp ráp \\\hline
\end{longtable}
}


\subparagraph{c) Kết luận lựa
chọn}\label{c-kux1ebft-luux1eadn-lux1ef1a-chux1ecdn}

\textbf{SIMCom SIM7600CE-T} được lựa chọn làm kiến trúc truyền thông
chính vì đây là phương án giúp thiết kế tổng thể gọn hơn và dễ triển
khai hơn, dù không phải lúc nào cũng linh hoạt nhất ở từng thành phần
riêng lẻ. Các lý do cụ thể gồm:

\begin{itemize}
\tightlist
\item
  Tích hợp LTE và GNSS trong cùng một module giúp giảm số lượng phần
  cứng rời và rút gọn đường kết nối UART.
\item
  Giữ nguyên sơ đồ chân đang triển khai (UART1 trên GPIO16/17, PWR-KEY
  trên GPIO26), không cần thay đổi pin mapping phần cứng.
\item
  Luồng AT được chuẩn hóa với Auto mode (\texttt{AT+CNMP=2}), kiểm tra
  \texttt{AT+CEREG?} trước \texttt{AT+CGACT=1,1}, và APN mặc định
  \texttt{internet}, phù hợp với firmware hiện tại.
\item
  GNSS tích hợp được điều khiển trực tiếp bằng \texttt{AT+CGNSPWR} và
  đọc dữ liệu qua \texttt{AT+CGNSINF}, thống nhất với module firmware
  đang vận hành.
\item
  Phù hợp với phương án BOM phần cứng đang áp dụng cho bo mạch đã chế
  tạo, thuận lợi cho triển khai và đối chiếu tài liệu kỹ thuật.
\end{itemize}


{\thesissettableid{3.4}{tab-ch3-thong-so-ky-thuat-phuong-an-sim7600ce-t}
\def\LTcaptype{table}
\begin{longtable}{|L{0.220\linewidth}|L{0.740\linewidth}|}
\caption{Thông số kỹ thuật phương án SIM7600CE-T}\label{tab:ch3-thong-so-ky-thuat-phuong-an-sim7600ce-t}\\
\hline \textbf{Thông số} & \textbf{Giá trị} \\
\hline
\endfirsthead
\caption[]{Thông số kỹ thuật phương án SIM7600CE-T (tiếp theo)}\\
\hline \textbf{Thông số} & \textbf{Giá trị} \\
\hline
\endhead
\hline
\endfoot
\endlastfoot
Kiến trúc & SIM7600CE-T tích hợp LTE Cat-4 và GNSS trong cùng module
{[}60{]} \\\hline
LTE category / tốc độ & Cat-4, tối đa 150 Mbps downlink / 50 Mbps uplink
{[}60{]} \\\hline
GNSS & GNSS tích hợp (GPS/GLONASS/BeiDou/Galileo), điều khiển bằng
\texttt{AT+CGNSPWR}, đọc dữ liệu bằng \texttt{AT+CGNSINF} {[}60{]} \\\hline
GNSS tích hợp trong modem & Có {[}60{]} \\\hline
Giao tiếp với MCU & 1 UART (GPIO16/17) cho toàn bộ AT LTE/GNSS + GPIO26
điều khiển PWR-KEY {[}60{]} \\\hline
Giao thức dữ liệu & LTE dùng AT command cho PDP/MQTT; GNSS lấy dữ liệu
qua \texttt{AT+CGNSINF} và có thể stream NMEA bằng \texttt{AT+CGNSTST}
khi cần {[}60{]} \\\hline
Điện áp modem & 3.4--4.2 V, điển hình 3.8 V {[}60{]} \\\hline
Luồng attach LTE & \texttt{AT+CNMP=2} (auto mode), cấu hình
\texttt{AT+CGDCONT=1,"IP","internet"}, kiểm tra \texttt{AT+CEREG?} trước
\texttt{AT+CGACT=1,1} \\\hline
Đồng bộ GNSS & Bật/tắt GNSS bằng \texttt{AT+CGNSPWR=1/0}, đọc dữ liệu
bằng \texttt{AT+CGNSINF} \\\hline
Điện áp cấp nguồn & 3.4--4.2 V, điển hình 3.8 V {[}60{]} \\\hline
Ý nghĩa tích hợp & Giảm số linh kiện rời, đơn giản hóa đi dây và đồng bộ
với firmware thực tế \\\hline
Ghi chú & APN mặc định sử dụng trong hệ thống là \texttt{internet}; pin
mapping phần cứng giữ nguyên \\\hline
\end{longtable}
}


\subsubsection{3.1.2. Giải pháp phần
cứng}\label{321-giux1ea3i-phuxe1p-phux1ea7n-cux1ee9ng}

Trên cơ sở các lựa chọn phần cứng đã phân tích ở mục 3.1.1, phần này đi
vào cách nhóm hiện thực chúng thành kiến trúc tracker hoàn chỉnh, từ
kết nối OBD2 và IMU đến khối nguồn và BOM tham chiếu.

\paragraph{3.1.2.1. Thiết kế khối thu thập dữ liệu
OBD2}\label{3211-thiux1ebft-kux1ebf-muxf4-ux111un-thu-thux1eadp-dux1eef-liux1ec7u-obd2}

\subparagraph{a) Lựa chọn phương pháp kết nối
OBD2}\label{a-lux1ef1a-chux1ecdn-phux1b0ux1a1ng-phuxe1p-kux1ebft-nux1ed1i-obd2}

Thu thập dữ liệu từ ECU xe qua cổng OBD2 (On-Board Diagnostics II) là
một yêu cầu cốt lõi của thiết kế tracker vì đây là nguồn thông tin giúp
hệ thống đi xa hơn một thiết bị định vị thông thường. Có hai hướng tiếp
cận chính: sử dụng adapter OBD2 có dây (wired ELM327 qua UART) hoặc
adapter OBD2 không dây (BLE ELM327). Nhóm lựa chọn kết nối không dây qua
BLE vì cách làm này phù hợp hơn với cách bố trí thiết bị thực tế trên xe
và với tài nguyên phần cứng của tracker:


{\thesissettableid{3.5}{tab-ch3-so-sanh-phuong-phap-ket-noi-obd2-theo-thong-so-trien-khai}
\def\LTcaptype{table}
\begin{longtable}{|L{0.177\linewidth}|L{0.303\linewidth}|L{0.484\linewidth}|}
\caption{So sánh phương pháp kết nối OBD2 theo thông số triển khai}\label{tab:ch3-so-sanh-phuong-phap-ket-noi-obd2-theo-thong-so-trien-khai}\\
\hline \textbf{Tiêu chí kỹ thuật} & \textbf{OBD2 có dây (ELM327 UART)} & \textbf{OBD2 không dây (BLE
ELM327 / vgate iCar Pro)} \\
\hline
\endfirsthead
\caption[]{So sánh phương pháp kết nối OBD2 theo thông số triển khai (tiếp theo)}\\
\hline \textbf{Tiêu chí kỹ thuật} & \textbf{OBD2 có dây (ELM327 UART)} & \textbf{OBD2 không dây (BLE
ELM327 / vgate iCar Pro)} \\
\hline
\endhead
\hline
\endfoot
\endlastfoot
Kết nối vật lý tracker ↔ OBD2 & Cần dây kéo cố định từ cổng OBD2 đến
tracker & Không cần dây giữa tracker và cổng OBD2 \\\hline
Giao thức tới MCU & UART TTL & BLE 4.0 GATT; vgate iCar Pro là bản
Bluetooth 4.0 BLE {[}29{]}, {[}64{]} \\\hline
Tài nguyên MCU tiêu tốn & Chiếm thêm 1 UART phần cứng trong lúc kết nối
& Tận dụng BLE tích hợp của ESP32-S3, giữ UART cho modem/GNSS/debug \\\hline
Giao thức OBD-II hỗ trợ & Phụ thuộc adapter ELM327 cụ thể & Vgate công
bố hỗ trợ J1850 PWM/VPW, ISO9141-2, ISO14230-4, ISO15765-4 CAN, SAE
J1939 CAN {[}64{]} \\\hline
Hành vi sleep của adapter & Phụ thuộc adapter có dây cụ thể & Vendor
công bố tự sleep sau khoảng 30 phút khi xe tắt máy {[}64{]} \\\hline
Hành vi wake & Phụ thuộc adapter và đường cấp nguồn & Vendor quảng bá tự
khởi động khi ignition ON; cần hiểu là hành vi phụ thuộc xe {[}64{]} \\\hline
Dòng tiêu thụ adapter & Phụ thuộc adapter; không dùng 1 số cứng nếu
không có datasheet đúng mẫu & Trang hãng được kiểm tra không công bố
dòng active/standby, nên không chốt số mA trong luận văn {[}64{]} \\\hline
Hành vi khi MCU deep sleep & UART không duy trì truyền dữ liệu khi MCU
ngủ sâu & BLE link cũng không duy trì qua deep sleep; cần reconnect sau
wake \\\hline
Tác động lắp đặt thực địa & Tracker thường bị ràng buộc gần cổng OBD2 &
Tracker có thể đặt lệch vị trí OBD2 trong cabin nếu vẫn trong vùng
BLE \\\hline
Bảo trì/thay adapter giữa các xe & Cần tháo dây liên quan đến tracker &
Chỉ cần rút/cắm adapter và cấu hình lại kết nối nếu cần \\\hline
\end{longtable}
}


\textbf{Kết luận:} Phương pháp BLE được chọn vì phù hợp hơn với kiến
trúc hiện tại của tracker: tận dụng BLE tích hợp của ESP32-S3, không
chiếm thêm UART phần cứng, giảm ràng buộc đi dây, đồng thời vgate iCar
Pro đã có công bố rõ về BLE 4.0, auto-sleep và danh sách giao thức
OBD-II. Với các thông số hãng chưa công bố (như dòng active/standby),
báo cáo chỉ ghi nhận là chưa có số liệu chính thức thay vì đưa ra giá
trị ước lượng.

\subparagraph{b) Lựa chọn adapter OBD2 BLE: vgate iCar
Pro}\label{b-lux1ef1a-chux1ecdn-adapter-obd2-ble-vgate-icar-pro}

\textbf{vgate iCar Pro} là adapter OBD2 sử dụng Bluetooth Low Energy
(BLE) 4.0, tương thích với BLE 5.0 của ESP32-S3. Adapter này được chọn
dựa trên các tiêu chí sau:

\begin{itemize}
\tightlist
\item
  \textbf{Tương thích giao thức:} Hỗ trợ tập lệnh ELM327 qua BLE GATT
  characteristics, cho phép đọc đa dạng dữ liệu từ ECU bao gồm trạng
  thái động cơ (IGN), số vòng quay (RPM), tốc độ xe, mức nhiên liệu,
  nhiệt độ động cơ, và mã lỗi chẩn đoán (DTC) {[}7{]}.
\item
  \textbf{Khả năng kết nối:} BLE 4.0 của adapter tương thích ngược với
  BLE 5.0 của ESP32-S3 {[}19{]}, {[}64{]}. Vendor cũng công bố cơ chế tự
  sleep sau khoảng 30 phút khi xe tắt máy và tự khởi động lại theo
  ignition, phù hợp cho kịch bản tracker trên xe {[}64{]}.
\item
  \textbf{Tiêu thụ năng lượng:} Ở thời điểm rà soát nguồn, trang hãng
  được kiểm tra không công bố dòng active/standby chính thức cho vgate
  iCar Pro BLE {[}64{]}. Vì vậy báo cáo chỉ ghi nhận đây là thông số
  chưa được công bố đầy đủ, không đưa ra giá trị mA cố định khi chưa có
  datasheet chính thức đúng mẫu.
\item
  \textbf{Phổ biến và giá hợp lý:} Ở lần rà soát thị trường Việt Nam ngày
  12/04/2026, các mức giá tham khảo ghi nhận cho vgate iCar Pro BLE nằm
  quanh khoảng 367.000--609.000 VND tùy kênh bán lẻ và phiên bản.
\end{itemize}


{\thesissetfigureid{3.2}{fig-ch3-so-o-ket-noi-ble-giua-esp32-s3-va-vgate-icar-pro}
\begin{figure}[htbp]
\centering
\pandocbounded{\safeincludesvg{./assets/figures/03-chuong-3-giai-phap-phan-cung-hinh-3-2.svg}}
\caption{Sơ đồ kết nối BLE giữa ESP32-S3 và vgate iCar Pro}
\label{fig:ch3-so-o-ket-noi-ble-giua-esp32-s3-va-vgate-icar-pro}
\end{figure}
}


\subparagraph{c) Chiến lược kết nối theo chế độ hoạt
động}\label{c-chiux1ebfn-lux1b0ux1ee3c-kux1ebft-nux1ed1i-theo-chux1ebf-ux111ux1ed9-houx1ea1t-ux111ux1ed9ng}

Một ràng buộc thiết kế quan trọng là ESP32-S3 không thể duy trì kết nối
BLE trong trạng thái deep sleep. Vì vậy, kết nối OBD2 không được xem là
một kênh "luôn bật", mà chỉ hoạt động khi nó thực sự tạo ra giá trị dữ liệu
cho từng trạng thái vận hành của xe:


{\thesissettableid{3.6}{tab-ch3-chien-luoc-ket-noi-ble-obd2-theo-che-o-hoat-ong}
\def\LTcaptype{table}
\begin{longtable}{|L{0.182\linewidth}|L{0.203\linewidth}|L{0.580\linewidth}|}
\caption{Chiến lược kết nối BLE OBD2 theo chế độ hoạt động}\label{tab:ch3-chien-luoc-ket-noi-ble-obd2-theo-che-o-hoat-ong}\\
\hline \textbf{Chế độ} & \textbf{Cần BLE OBD2?} & \textbf{Lý do} \\
\hline
\endfirsthead
\caption[]{Chiến lược kết nối BLE OBD2 theo chế độ hoạt động (tiếp theo)}\\
\hline \textbf{Chế độ} & \textbf{Cần BLE OBD2?} & \textbf{Lý do} \\
\hline
\endhead
\hline
\endfoot
\endlastfoot
Lái xe (IGN ON) & Có & Đọc dữ liệu xe (RPM, tốc độ, nhiên liệu) liên
tục \\\hline
Đỗ xe (IGN OFF) & Không & Không cần dữ liệu OBD2, chỉ cần IMU phát hiện
chuyển động \\\hline
Cảnh báo & Tùy chọn & Có thể kết nối để xác nhận IGN, hoặc chỉ dùng IMU
+ GPS \\\hline
\end{longtable}
}


Khi xe chạy (IGN ON), ESP32-S3 duy trì kết nối BLE liên tục với adapter để
đọc dữ liệu OBD2 định kỳ mỗi 5--30 giây. Khi xe đỗ (IGN OFF), kết nối BLE
được ngắt trước khi ESP32 vào deep sleep, vì ở thời điểm này mục tiêu ưu
tiên không còn là lấy dữ liệu động cơ, mà là giữ mức tiêu thụ thấp nhất
trong khi vẫn duy trì khả năng phát hiện xâm nhập. IMU LIS3DSH đảm nhiệm
việc giám sát đó độc lập và chỉ đánh thức hệ thống khi có tín hiệu đáng
ngờ. Cách tổ chức này giúp giảm đáng kể năng lượng tiêu thụ ở pha đỗ xe mà
không làm mất đi chức năng cảnh báo cốt lõi.

Trong trường hợp adapter OBD2 không kết nối được (timeout 10 giây, retry
2--3 lần), hệ thống tự động chuyển sang phương pháp dự phòng (fallback)
là đo điện áp ắc quy qua ADC để phát hiện trạng thái IGN. Phương pháp
này kém chính xác hơn nhưng đảm bảo hệ thống vẫn hoạt động bình thường.

\paragraph{3.1.2.2. Thiết kế khối cảm biến chuyển động
IMU}\label{3212-thiux1ebft-kux1ebf-muxf4-ux111un-cux1ea3m-biux1ebfn-chuyux1ec3n-ux111ux1ed9ng-imu}

\subparagraph{a) Vai trò của IMU trong hệ
thống}\label{a-vai-truxf2-cux1ee7a-imu-trong-hux1ec7-thux1ed1ng}

Cảm biến đo quán tính (IMU - Inertial Measurement Unit) đóng vai trò
then chốt trong việc phát hiện chuyển động của xe khi đang đỗ, cho phép
hệ thống phát hiện các tình huống bất thường như rung xe (cố gắng mở
cửa), kéo xe, hoặc cẩu xe. Đặc biệt, IMU cho phép ESP32-S3 ở trạng thái
deep sleep liên tục và chỉ thức dậy khi có chuyển động thực sự, giúp
tiết kiệm năng lượng đáng kể.

\subparagraph{b) Lựa chọn cảm biến:
LIS3DSH}\label{b-lux1ef1a-chux1ecdn-cux1ea3m-biux1ebfn-lis3dsh}

\textbf{LIS3DSH} của STMicroelectronics là cảm biến gia tốc 3 trục
(3-axis accelerometer) low-power được lựa chọn cho hệ thống. Các đặc
tính kỹ thuật chính được tóm tắt trong bảng sau:


{\thesissettableid{3.7}{tab-ch3-thong-so-ky-thuat-cam-bien-lis3dsh}
\def\LTcaptype{table}
\begin{longtable}{|L{0.426\linewidth}|L{0.542\linewidth}|}
\caption{Thông số kỹ thuật cảm biến LIS3DSH}\label{tab:ch3-thong-so-ky-thuat-cam-bien-lis3dsh}\\
\hline \textbf{Thông số} & \textbf{Giá trị} \\
\hline
\endfirsthead
\caption[]{Thông số kỹ thuật cảm biến LIS3DSH (tiếp theo)}\\
\hline \textbf{Thông số} & \textbf{Giá trị} \\
\hline
\endhead
\hline
\endfoot
\endlastfoot
Loại & 3-axis Digital Accelerometer \\\hline
Điện áp hoạt động & 1.8--3.6 V \\\hline
Giao tiếp & I2C (400 kHz) hoặc SPI (10 MHz) \\\hline
Phạm vi đo & +/-2g, +/-4g, +/-8g, +/-16g (có thể điều chỉnh) \\\hline
Độ phân giải & 16-bit \\\hline
Tiêu thụ (Active) & 10--50 µA \\\hline
Tiêu thụ (Low-power) & 2--5 µA \\\hline
Wake-up interrupt & Có (ngưỡng có thể cấu hình) \\\hline
Package & LGA-16 (3x3x1 mm) \\\hline
\end{longtable}
}


\subparagraph{c) Nguyên lý hoạt động trong hệ
thống}\label{c-nguyuxean-luxfd-houx1ea1t-ux111ux1ed9ng-trong-hux1ec7-thux1ed1ng}

LIS3DSH được cấu hình ở chế độ low-power với tần suất lấy mẫu 1 Hz (ODR
= 1 Hz) và chỉ bật chức năng motion detection. Khi phát hiện gia tốc
vượt ngưỡng 0.2g trên bất kỳ trục nào (X, Y, Z), cảm biến tự động phát
tín hiệu interrupt đến chân GPIO của ESP32-S3 thông qua chân INT1, đánh
thức vi điều khiển từ trạng thái deep sleep {[}9{]}.

Ngưỡng gia tốc 0.2g được chọn để cân bằng giữa độ nhạy và khả năng chống
báo giả. Giá trị này đủ lớn để loại bỏ rung nhẹ từ môi trường (gió, xe
cộ đi ngang) nhưng đủ nhỏ để phát hiện các chuyển động thực sự như rung
của xe hoặc di chuyển.


{\thesissetfigureid{3.3}{fig-ch3-so-o-ket-noi-lis3dsh-voi-esp32-s3-qua-i2c}
\begin{figure}[htbp]
\centering
\pandocbounded{\safeincludesvg{./assets/figures/03-chuong-3-giai-phap-phan-cung-hinh-3-3.svg}}
\caption{Sơ đồ kết nối LIS3DSH với ESP32-S3 qua I2C}
\label{fig:ch3-so-o-ket-noi-lis3dsh-voi-esp32-s3-qua-i2c}
\end{figure}
}



{\thesissetfigureid{3.3a}{fig-ch3-so-o-chi-tiet-chan-ket-noi-lis3dsh-voi-esp32-s3}
\begin{figure}[htbp]
\centering
\pandocbounded{\safeincludesvg{./assets/figures/03-chuong-3-giai-phap-phan-cung-hinh-3-3a.svg}}
\caption{Sơ đồ chi tiết chân kết nối LIS3DSH với ESP32-S3}
\label{fig:ch3-so-o-chi-tiet-chan-ket-noi-lis3dsh-voi-esp32-s3}
\end{figure}
}


Cấu hình I2C sử dụng địa chỉ \texttt{0x18}, tốc độ \texttt{400\ kHz},
với hai điện trở kéo lên \texttt{10\ kΩ} trên các net \texttt{IMU-SCL}
và \texttt{IMU-SDA}. Ở runtime hiện tại, bus I2C được ánh xạ về
\texttt{GPIO48} (SCL) và \texttt{GPIO47} (SDA), chân \texttt{INT1} đi
vào \texttt{GPIO21} để đánh thức ESP32-S3, còn \texttt{INT2} được giữ
lại cho mở rộng sau này.

\subparagraph{d) Ưu điểm của giải
pháp}\label{d-ux1b0u-ux111iux1ec3m-cux1ee7a-giux1ea3i-phuxe1p}

\begin{itemize}
\tightlist
\item
  \textbf{Tiết kiệm năng lượng tốt nhất:} Dòng tiêu thụ chỉ 2--5 µA ở
  chế độ low-power, trong khi vẫn duy trì khả năng phát hiện chuyển
  động. ESP32-S3 không cần chạy liên tục để kiểm tra cảm biến.
\item
  \textbf{Phát hiện chuyển động độc lập:} LIS3DSH tự xử lý việc phát
  hiện chuyển động bằng phần cứng, chỉ gửi interrupt khi có sự kiện ---
  không phụ thuộc vào CPU của ESP32.
\item
  \textbf{Hệ sinh thái phát triển phong phú:} Nhiều thư viện sẵn có cho
  Arduino và ESP-IDF; IC LIS3DSH hoặc cảm biến tương đương đều dễ tích
  hợp vào sơ đồ phần cứng của thiết bị.
\end{itemize}

\paragraph{3.1.2.3. Thiết kế hệ thống quản lý
nguồn}\label{3213-thiux1ebft-kux1ebf-hux1ec7-thux1ed1ng-quux1ea3n-luxfd-nguux1ed3n}

Hệ thống quản lý nguồn là khối phức tạp nhất trong thiết kế phần cứng, vì
nó phải đồng thời giải quyết ba bài toán vốn dễ xung đột với nhau: tạo
đúng các mức điện áp cho từng tải, duy trì đường nguồn dự phòng khi cần và
bảo vệ ắc quy xe khỏi bị rút cạn quá mức. Vì vậy, khối này được tách thành
năm nhánh chức năng để người đọc dễ theo dõi vai trò của từng đường nguồn
thay vì chỉ nhìn một sơ đồ điện phức tạp.


{\thesissetfigureid{3.4}{fig-ch3-so-o-khoi-he-thong-quan-ly-nguon}
\begin{figure}[htbp]
\centering
\pandocbounded{\safeincludesvg{./assets/figures/03-chuong-3-giai-phap-phan-cung-hinh-3-4.svg}}
\caption{Sơ đồ khối hệ thống quản lý nguồn}
\label{fig:ch3-so-o-khoi-he-thong-quan-ly-nguon}
\end{figure}
}



{\thesissetfigureid{3.4a}{fig-ch3-kien-truc-nguon-va-phan-phoi-ien-ap-trong-he-thong}
\begin{figure}[htbp]
\centering
\pandocbounded{\safeincludesvg{./assets/figures/03-chuong-3-giai-phap-phan-cung-hinh-3-4a.svg}}
\caption{Kiến trúc nguồn và phân phối điện áp trong hệ thống}
\label{fig:ch3-kien-truc-nguon-va-phan-phoi-ien-ap-trong-he-thong}
\end{figure}
}


\subparagraph{a) LDO 3.3V cho ESP32-S3 (AP2112-3.3)}\label{a-buck-33v-cho-esp32-s3-xl1509-33e}

Nhánh 3.3V dùng \textbf{AP2112-3.3} như một LDO để hạ từ \textbf{bus 5V
chính} xuống 3.3V, cấp cho ESP32-S3 và các tải logic 3.3V. Cách tổ chức
này phù hợp với kiến trúc nguồn thực tế của bo mạch: MP2482 đảm nhiệm bước
hạ áp đầu tiên từ 12--24V xuống 5V, còn AP2112-3.3 đảm nhiệm bước ổn áp
cuối cho miền logic.


{\thesissettableid{3.8}{tab-ch3-thong-so-khoi-buck-3-3v}
\def\LTcaptype{table}
\begin{longtable}{|L{0.168\linewidth}|L{0.132\linewidth}|L{0.218\linewidth}|L{0.250\linewidth}|L{0.176\linewidth}|}
\caption{Thông số khối LDO 3.3V}\label{tab:ch3-thong-so-khoi-buck-3-3v}\\
\hline \textbf{Khối} & \textbf{IC} & \textbf{Input} & \textbf{Output} & \textbf{Tải} \\
\hline
\endfirsthead
\caption[]{Thông số khối LDO 3.3V (tiếp theo)}\\
\hline \textbf{Khối} & \textbf{IC} & \textbf{Input} & \textbf{Output} & \textbf{Tải} \\
\hline
\endhead
\hline
\endfoot
\endlastfoot
LDO 3.3V & AP2112-3.3 & 5V & 3.3V & ESP32-S3 \\\hline
\end{longtable}
}


\subparagraph{b) Buck 5V bus chính
(MP2482)}\label{b-buck-5v-bus-chuxednh-mp2482}

Nhánh buck 5V dùng \textbf{MP2482} để tạo \textbf{bus 5V chính} từ nguồn
12--24V của xe. Bus 5V này là nguồn trung gian cho các khối phụ trợ và
đầu vào mạch sạc pin.


{\thesissettableid{3.9}{tab-ch3-thong-so-khoi-buck-5v}
\def\LTcaptype{table}
\begin{longtable}{|L{0.168\linewidth}|L{0.132\linewidth}|L{0.218\linewidth}|L{0.250\linewidth}|L{0.176\linewidth}|}
\caption{Thông số khối Buck 5V}\label{tab:ch3-thong-so-khoi-buck-5v}\\
\hline \textbf{Khối} & \textbf{IC} & \textbf{Input} & \textbf{Output} & \textbf{Tải} \\
\hline
\endfirsthead
\caption[]{Thông số khối Buck 5V (tiếp theo)}\\
\hline \textbf{Khối} & \textbf{IC} & \textbf{Input} & \textbf{Output} & \textbf{Tải} \\
\hline
\endhead
\hline
\endfoot
\endlastfoot
Buck 5V & MP2482 & 12--24V & 5V & Bus 5V chính \\\hline
\end{longtable}
}


\subparagraph{c) Buck 3.8V/4V cho modem
(TPS54231)}\label{c-buck-38v4v-cho-modem-tps54231}

Do modem SIM7600CE-T làm việc trên miền nguồn thấp áp, hệ thống bố trí
nhánh buck riêng dùng \textbf{TPS54231} để hạ 12--24V xuống khoảng
\textbf{4V} cấp cho modem.


{\thesissettableid{3.10}{tab-ch3-thong-so-khoi-buck-3-8v-4v-cho-modem}
\def\LTcaptype{table}
\begin{longtable}{|L{0.168\linewidth}|L{0.132\linewidth}|L{0.218\linewidth}|L{0.250\linewidth}|L{0.176\linewidth}|}
\caption{Thông số khối Buck 3.8V/4V cho modem}\label{tab:ch3-thong-so-khoi-buck-3-8v-4v-cho-modem}\\
\hline \textbf{Khối} & \textbf{IC} & \textbf{Input} & \textbf{Output} & \textbf{Tải} \\
\hline
\endfirsthead
\caption[]{Thông số khối Buck 3.8V/4V cho modem (tiếp theo)}\\
\hline \textbf{Khối} & \textbf{IC} & \textbf{Input} & \textbf{Output} & \textbf{Tải} \\
\hline
\endhead
\hline
\endfoot
\endlastfoot
Buck 3.8V & TPS54\-231 & 12--24V & ~4V & SIM7600CE-T \\\hline
\end{longtable}
}


\subparagraph{d) Boost 5V từ pin dự phòng (SX1308) và Power
Path}\label{d-boost-5v-tux1eeb-pin-dux1ef1-phuxf2ng-sx1308-vuxe0-power-path}

Khi ắc quy xe yếu, nguồn dự phòng lấy từ pin 18650 1S (danh định
~3.7V). Khối \textbf{SX1308} tăng áp lên 5V để duy trì cấp
nguồn cho hệ thống.


{\thesissettableid{3.11}{tab-ch3-thong-so-khoi-boost-5v-du-phong}
\def\LTcaptype{table}
\begin{longtable}{|L{0.189\linewidth}|L{0.094\linewidth}|L{0.236\linewidth}|L{0.283\linewidth}|L{0.142\linewidth}|}
\caption{Thông số khối Boost 5V dự phòng}\label{tab:ch3-thong-so-khoi-boost-5v-du-phong}\\
\hline \textbf{Khối} & \textbf{IC} & \textbf{Input} & \textbf{Output} & \textbf{Tải} \\
\hline
\endfirsthead
\caption[]{Thông số khối Boost 5V dự phòng (tiếp theo)}\\
\hline \textbf{Khối} & \textbf{IC} & \textbf{Input} & \textbf{Output} & \textbf{Tải} \\
\hline
\endhead
\hline
\endfoot
\endlastfoot
Boost 5V & SX1308 & Battery ~3.7V & 5V & Backup từ pin \\\hline
\end{longtable}
}


Power path runtime sử dụng \textbf{diode OR} giữa nhánh 5V chính
(MP2482) và nhánh 5V dự phòng (SX1308), kết hợp điều khiển GPIO để
bật/tắt các nhánh liên quan. Mục tiêu là chuyển nguồn mượt, không reset
hệ thống khi đổi nguồn.

\subparagraph{e) Giám sát điện áp và
LVD}\label{e-giuxe1m-suxe1t-ux111iux1ec7n-uxe1p-vuxe0-lvd}

LVD dùng ngưỡng profile kép theo firmware:

\begin{itemize}
\tightlist
\item
  \textbf{Profile 12V:} \(Switch_{\mathrm{OFF}} = 12.0\,\mathrm{V}\),
  \(Switch_{\mathrm{ON}} = 12.2\,\mathrm{V}\)
\item
  \textbf{Profile 24V:} \(Switch_{\mathrm{OFF}} = 24.0\,\mathrm{V}\),
  \(Switch_{\mathrm{ON}} = 24.4\,\mathrm{V}\)
\end{itemize}


{\thesissettableid{3.12}{tab-ch3-bang-trang-thai-chuyen-nguon-va-canh-bao}
\def\LTcaptype{table}
\begin{longtable}{|L{0.108\linewidth}|L{0.068\linewidth}|L{0.369\linewidth}|L{0.141\linewidth}|L{0.088\linewidth}|L{0.161\linewidth}|}
\caption{Bảng trạng thái chuyển nguồn và cảnh báo}\label{tab:ch3-bang-trang-thai-chuyen-nguon-va-canh-bao}\\
\hline \textbf{Chế độ} & \textbf{IGN} & \textbf{\(U_{\text{batt}}\)} & Nguồn tracker & Sạc pin & Cảnh báo \\
\hline
\endfirsthead
\caption[]{Bảng trạng thái chuyển nguồn và cảnh báo (tiếp theo)}\\
\hline \textbf{Chế độ} & \textbf{IGN} & \textbf{\(U_{\text{batt}}\)} & Nguồn tracker & Sạc pin & Cảnh báo \\
\hline
\endhead
\hline
\endfoot
\endlastfoot
Xe chạy & ON & Profile 12V: \(U_{\text{batt}} \ge 13.0\,\mathrm{V}\);\linebreak
Profile 24V: \(U_{\text{batt}} \ge 26.0\,\mathrm{V}\) & Ắc quy & Có & Không \\\hline
Xe đỗ & OFF & Profile 12V: \(U_{\text{batt}} > 12.0\,\mathrm{V}\);\linebreak
Profile 24V: \(U_{\text{batt}} > 24.0\,\mathrm{V}\) & Ắc quy & Không & Không \\\hline
Ắc quy yếu & OFF & Profile 12V: \(U_{\text{batt}} \le 12.0\,\mathrm{V}\);\linebreak
Profile 24V: \(U_{\text{batt}} \le 24.0\,\mathrm{V}\) & Pin 18650 1S & Không & Cảnh báo chuyển
nguồn \\\hline
Ắc quy phục hồi & OFF & Profile 12V: \(U_{\text{batt}} \ge 12.2\,\mathrm{V}\);\linebreak
Profile 24V: \(U_{\text{batt}} \ge 24.4\,\mathrm{V}\) & Ắc quy & Không & Cảnh báo
phục hồi \\\hline
\end{longtable}
}


Ngoài kênh ADC, tín hiệu trạng thái LVD từ khối LVD phần cứng được đưa về
\textbf{GPIO19 (LVD\_STATUS)} để giám sát nhanh. Quy ước runtime:
\textbf{GPIO19 HIGH = low-voltage}, \textbf{GPIO19 LOW = bình thường}.

\subparagraph{f) Mạch sạc pin 1S
(TP4056)}\label{f-mux1ea1ch-sux1ea1c-pin-1s-tp4056}

Khối sạc dùng \textbf{TP4056}, nhận \textbf{5V từ nhánh MP2482} và sạc
pin 18650 1S theo chuẩn \textbf{4.2V/1S}.


{\thesissettableid{3.13}{tab-ch3-thong-so-khoi-sac-pin}
\def\LTcaptype{table}
\begin{longtable}{|L{0.189\linewidth}|L{0.094\linewidth}|L{0.236\linewidth}|L{0.283\linewidth}|L{0.142\linewidth}|}
\caption{Thông số khối sạc pin}\label{tab:ch3-thong-so-khoi-sac-pin}\\
\hline \textbf{Khối} & \textbf{IC} & \textbf{Input} & \textbf{Output} & \textbf{Tải} \\
\hline
\endfirsthead
\caption[]{Thông số khối sạc pin (tiếp theo)}\\
\hline \textbf{Khối} & \textbf{IC} & \textbf{Input} & \textbf{Output} & \textbf{Tải} \\
\hline
\endhead
\hline
\endfoot
\endlastfoot
Sạc pin & TP4056 & 5V từ MP2482 & 4.2V & Pin 18650 1S \\\hline
\end{longtable}
}


Dòng sạc TP4056 được thiết lập theo điện trở PROG và giới hạn nhiệt của
mạch, vì vậy không dùng một giá trị dòng cố định cho mọi điều kiện vận
hành.

\subparagraph{g) Pin dự phòng 18650 1S
Li-ion}\label{g-pin-dux1ef1-phuxf2ng-18650-1s-li-ion}

Pin 18650 1S Li-ion được chọn làm nguồn dự phòng với cấu hình 1 cell đơn
giản, kết hợp khối sạc TP4056, khối boost SX1308 và mạch bảo vệ pin 1S
trong kiến trúc nguồn của hệ thống.


{\thesissettableid{3.14}{tab-ch3-thong-so-ky-thuat-pin-du-phong-18650-1s}
\def\LTcaptype{table}
\begin{longtable}{|L{0.310\linewidth}|L{0.658\linewidth}|}
\caption{Thông số kỹ thuật pin dự phòng 18650 1S}\label{tab:ch3-thong-so-ky-thuat-pin-du-phong-18650-1s}\\
\hline \textbf{Thông số} & \textbf{Giá trị} \\
\hline
\endfirsthead
\caption[]{Thông số kỹ thuật pin dự phòng 18650 1S (tiếp theo)}\\
\hline \textbf{Thông số} & \textbf{Giá trị} \\
\hline
\endhead
\hline
\endfoot
\endlastfoot
Loại & Li-ion 18650 1S \\\hline
Dung lượng & 3500 mAh @ 3.7V \\\hline
Năng lượng & ~12.95 Wh \\\hline
Điện áp & 3.0--4.2V (nominal 3.7V) \\\hline
Số chu kỳ & 500--1000 chu kỳ (80\% capacity) \\\hline
\end{longtable}
}


\textbf{Tính toán thời gian hoạt động từ pin backup:}

\emph{Các giá trị dưới đây là ước tính theo cùng giả định tổn hao của
phương án gốc, với dung lượng hiệu quả xấp xỉ 3150 mAh (tương đương 90\%
dung lượng danh định), không phải số đo thực nghiệm mới trên cell
18650.}

\emph{Chế độ heartbeat (đỗ xe, chu kỳ 15 phút):}

\begin{itemize}
\tightlist
\item
  Dòng trung bình:
  \(I_{\text{avg}} = \dfrac{250\,\mathrm{mA}\times 15\,\mathrm{s} + 3\,\mathrm{mA}\times 885\,\mathrm{s}}{900\,\mathrm{s}} = 7.1\,\mathrm{mA}\)
\item
  Dung lượng hiệu quả (sau tổn hao): 3150 mAh
\item
  Thời gian hoạt động:
  \(T = \dfrac{3150}{7.1} = 444\,\mathrm{giờ} \approx 18.5\,\text{ngày}\)
\end{itemize}

\emph{Chế độ heartbeat (đỗ xe, chu kỳ 30 phút):}

\begin{itemize}
\tightlist
\item
  Dòng trung bình: \(\approx 5\,\mathrm{mA}\)
\item
  Thời gian hoạt động:
  \(T = \dfrac{3150}{5} = 630\,\mathrm{giờ} \approx 26\,\text{ngày}\)
\end{itemize}

\emph{Chế độ cảnh báo (track liên tục):}

\begin{itemize}
\tightlist
\item
  Dòng trung bình: \(\approx 250\,\mathrm{mA}\)
\item
  Thời gian hoạt động:
  \(T = \dfrac{3150}{250} = 12.6\,\mathrm{giờ}\)
\end{itemize}

\emph{Giá trị 12.6 giờ ở đây là phép tính lý tưởng theo tải giả định 250
mA tại mức cell; khi quy đổi sang tải hệ thống thực tế có hao hụt boost
và các pha truyền dữ liệu đỉnh cao hơn, thời lượng tracking liên tục
trong Chương 4 được lấy theo khoảng bảo thủ hơn là
~2.6--3.0 giờ.}

Kết quả quy đổi cho thấy baseline pin 18650 1S vẫn đủ khả năng duy trì
hoạt động tracker trong nhiều ngày ở chế độ heartbeat và trên 12 giờ ở
chế độ cảnh báo liên tục theo mô hình lý tưởng; trong vận hành toàn hệ
thống, các giá trị bảo thủ hơn ở Chương 4 phù hợp hơn cho mục tiêu duy
trì giám sát tạm thời khi ắc quy xe yếu, nhưng vẫn cần được kiểm thử
thực nghiệm lại nếu chuyển sang mẫu pin mới trong pha chế tạo tiếp theo
{[}14{]}.

\paragraph{3.1.2.4. Bảng tổng hợp linh kiện (Bill of
Materials)}\label{3214-bux1ea3ng-tux1ed5ng-hux1ee3p-linh-kiux1ec7n-bill-of-materials}


{\thesissettableid{3.14A}{tab-ch3-bang-tong-hop-linh-kien-he-thong-tracker-bom}
\def\LTcaptype{table}
\begin{longtable}{|L{0.068\linewidth}|L{0.264\linewidth}|L{0.068\linewidth}|L{0.068\linewidth}|L{0.151\linewidth}|L{0.316\linewidth}|}
\caption{Bảng tổng hợp linh kiện hệ thống tracker (BOM)}\label{tab:ch3-bang-tong-hop-linh-kien-he-thong-tracker-bom}\\
\hline \textbf{STT} & \textbf{Thành phần} & \textbf{Đơn vị} & \textbf{SL} & \textbf{Giá tham chiếu T4/2026 (VND)} & \textbf{Ghi chú} \\
\hline
\endfirsthead
\caption[]{Bảng tổng hợp linh kiện hệ thống tracker (BOM) (tiếp theo)}\\
\hline \textbf{STT} & \textbf{Thành phần} & \textbf{Đơn vị} & \textbf{SL} & \textbf{Giá tham chiếu T4/2026 (VND)} & \textbf{Ghi chú} \\
\hline
\endhead
\hline
\endfoot
\endlastfoot
1 & ESP32-S3-WROOM-1 (N16R8) & Cái & 1 & 95.000 & MCU trung
tâm của bo mạch \\\hline
2 & IC cảm biến LIS3DSH & Cái & 1 & 40.000 & Cảm biến gia tốc 3
trục \\\hline
3 & vgate iCar Pro (OBD2 BLE) & Cái & 1 & 367.000 & BLE 4.0,
ELM327 compatible \\\hline
4 & SIMCom SIM7600CE-T & Bộ & 1 & 650.000 & Modem LTE + GNSS
tích hợp \\\hline
5 & Pin 18650 1S Li-ion 3500mAh & Cái & 1 & 79.000 & Loại có
protection board \\\hline
6 & IC sạc TP4056 + mạch phụ trợ & Cái & 1 & 25.000 & Input 5V,
output 4.2V, dòng theo PROG \\\hline
7 & Mạch LDO AP2112-3.3 & Cái & 1 & 5.000 & 5V
-\textgreater{} 3.3V cấp ESP32-S3 \\\hline
8 & Mạch buck MP2482 5V & Cái & 1 & 17.000 & Bus 5V chính \\\hline
9 & Mạch buck TPS54231 ~4V & Cái & 1 & 9.000 &
12--24V -\textgreater{} ~4V cấp modem SIM7600CE-T \\\hline
10 & Mạch boost SX1308 5V backup & Cái & 1 & 12.000 & Backup từ
pin 18650 1S \\\hline
11 & Khối giám sát LVD + ADC & Cái & 1 & 8.000 & Giám sát LVD
(GPIO19) \\\hline
12 & Diode Schottky (SS54/SS34) & Cái & 2 & 2.500 & Diode-OR power
path \\\hline
13 & Mạch bảo vệ pin 1S (BMS) & Cái & 1 & 12.000 & Bảo vệ pin
18650 1S \\\hline
14 & Điện trở (10kOhm, 2.2kOhm) & Gói & 1 & 5.000 & Voltage
divider, pull-up \\\hline
15 & Tụ điện (100uF, 220uF) & Gói & 1 & 5.000 & Lọc nhiễu,
decoupling \\\hline
16 & Connector, header pin & Gói & 1 & 55.000 & Header, jack OBD2 và
đầu nối nguồn \\\hline
17 & PCB 2 lớp (~50x50 mm) & Cái & 1 & 20.000 &
Bo mạch theo layout của đề tài \\\hline
18 & Vỏ bảo vệ (tùy chọn) & Cái & 1 & 25.000 & Nhựa hoặc kim
loại \\\hline
19 & Dây nối, cáp, phụ kiện & - & - & 60.000 & Dây điện, anten
4G/GNSS, khay SIM \\\hline
20 & Linh kiện phụ trợ khác & - & - & 20.000 & Fuse, switch,
LED \\\hline
\end{longtable}
}



{\thesissettableid{3.15}{tab-ch3-tong-hop-chi-phi-theo-nhom}
\def\LTcaptype{table}
\begin{longtable}{|L{0.160\linewidth}|L{0.580\linewidth}|L{0.224\linewidth}|}
\caption{Tổng hợp chi phí theo nhóm}\label{tab:ch3-tong-hop-chi-phi-theo-nhom}\\
\hline \textbf{Nhóm} & \textbf{Hạng mục} & \textbf{Chi phí (VND)} \\
\hline
\endfirsthead
\caption[]{Tổng hợp chi phí theo nhóm (tiếp theo)}\\
\hline \textbf{Nhóm} & \textbf{Hạng mục} & \textbf{Chi phí (VND)} \\
\hline
\endhead
\hline
\endfoot
\endlastfoot
A & Thành phần chính (MCU, cảm biến, modem, OBD2, pin) &
1.231.000 \\\hline
B & Quản lý nguồn (AP2112-3.3, MP2482, TPS54231, SX1308, TP4056, diode OR,
khối LVD, pin 1S) & 93.000 \\\hline
C & Kết nối, RF và linh kiện phụ trợ & 145.000 \\\hline
D & PCB và vỏ (tùy chọn) & 45.000 \\\hline
& \textbf{Tổng cộng} & \textbf{1.514.000} \\\hline
\end{longtable}
}


Theo mặt bằng giá tham chiếu thị trường Việt Nam được rà soát ngày
12/04/2026, tổng chi phí vật tư cho một bộ tracker hoàn chỉnh vào khoảng
\textbf{1.514.000 VND}. Mức này vẫn nằm dưới ngưỡng chi phí mục tiêu của
đồ án, đồng thời phản ánh thực tế hơn so với mặt bằng giá linh kiện ở
thời điểm đầu năm.

\subsection{3.2. Phân tích và đề xuất giải pháp Firmware -- Firmware
analysis and solution}\label{32-ux111ux1ec1-xuux1ea5t-cuxe1c-giux1ea3i-phuxe1p--proposed-solutions}

\subsubsection{3.2.1. Phân tích và lựa chọn giải pháp Firmware}

\paragraph{3.2.1.1. Đặt vấn đề cho giải pháp
Firmware}\label{3121-ux111ux1eb7t-vux1ea5n-ux111ux1ec1-cho-giux1ea3i-phuxe1p-firmware}

Nếu xem phần cứng là ``cơ thể'' của thiết bị, thì firmware là lớp biến
tập linh kiện đó thành một hệ thống biết \emph{phản ứng đúng lúc}. Người
dùng không trực tiếp nhìn thấy firmware, nhưng việc vị trí có cập nhật
đều, cảnh báo có đến kịp hay không, hay ắc quy xe có bị ảnh hưởng hay
không đều do lớp này quyết định. Vì nằm ở điểm giao nhau giữa phần cứng,
truyền thông và yêu cầu nghiệp vụ, bài toán firmware không dừng ở chuyện
``cho thiết bị chạy được'', mà phải bảo đảm thiết bị chạy đúng, ổn định
và tiết kiệm năng lượng trong môi trường thực tế.

\begin{itemize}
\tightlist
\item
  \textbf{Đa nhiệm thời gian thực}: thu thập cảm biến, điều khiển modem,
  truyền dữ liệu và xử lý cảnh báo phải chạy song song, độ trễ thấp.
\item
  \textbf{Quản lý năng lượng nghiêm ngặt}: chuyển trạng thái linh hoạt
  giữa active/sleep/deep sleep theo điều kiện vận hành xe.
\item
  \textbf{Giao tiếp đa giao thức}: BLE (OBD2), UART (modem), MQTT/HTTP
  (cloud), I2C/ADC/GPIO (ngoại vi phần cứng).
\item
  \textbf{Độ bền vận hành cao}: có cơ chế retry, fallback và phục hồi
  sau mất kết nối mạng hoặc lỗi ngoại vi.
\end{itemize}

\paragraph{3.2.1.2. So sánh các phương án nền tảng
Firmware}\label{3122-so-suxe1nh-cuxe1c-phux1b0ux1a1ng-uxe1n-nux1ec1n-tux1ea3ng-firmware}


{\thesissettableid{3.5A}{tab-ch3-so-sanh-cac-phuong-an-nen-tang-firmware}
\def\LTcaptype{table}
\begin{longtable}{|L{0.138\linewidth}|L{0.236\linewidth}|L{0.246\linewidth}|L{0.224\linewidth}|L{0.100\linewidth}|}
\caption{So sánh các phương án nền tảng firmware}\label{tab:ch3-so-sanh-cac-phuong-an-nen-tang-firmware}\\
\hline \textbf{Phương án} & \textbf{Mô tả} & \textbf{Ưu điểm} & \textbf{Hạn chế} & \textbf{Mức phù hợp} \\
\hline
\endfirsthead
\caption[]{So sánh các phương án nền tảng firmware (tiếp theo)}\\
\hline \textbf{Phương án} & \textbf{Mô tả} & \textbf{Ưu điểm} & \textbf{Hạn chế} & \textbf{Mức phù hợp} \\
\hline
\endhead
\hline
\endfoot
\endlastfoot
\textbf{PA-FW1: Arduino Core + Superloop} & Vòng lặp chính tuần tự, xử
lý tác vụ theo polling & Dễ bắt đầu, ít cấu hình & Khó mở rộng đa nhiệm
thực sự, khó tối ưu deep sleep phức tạp, quản lý lỗi hạn chế & Trung
bình \\\hline
\textbf{PA-FW2: ESP-IDF + FreeRTOS (Đã chọn)} & Kiến trúc task và event,
driver chính thức Espressif & Đa nhiệm tốt, hỗ trợ power management sâu,
tích hợp NimBLE / modem / UART / I2C ổn định & Độ phức tạp cao hơn Arduino &
\textbf{Cao} \\\hline
\textbf{PA-FW3: Zephyr RTOS trên ESP32} & RTOS đa nền tảng, kiến trúc
phân lớp & Tính chuẩn hóa tốt, khả năng mở rộng dài hạn & Hệ sinh thái
ESP32 chuyên biệt và tài liệu thực chiến BLE-OBD2/modem ít hơn ESP-IDF &
Trung bình \\\hline
\end{longtable}
}


\paragraph{3.2.1.3. Chọn giải pháp
Firmware}\label{3123-chux1ecdn-giux1ea3i-phuxe1p-firmware}

Từ kết quả so sánh, đồ án chọn \textbf{PA-FW2: ESP-IDF + FreeRTOS} làm
nền tảng firmware chính. Đây không phải là lựa chọn ``mạnh nhất trên giấy'',
mà là lựa chọn bám sát nhất cách thiết bị phải vận hành ngoài thực tế:
nhiều ngoại vi, nhiều trạng thái, yêu cầu tiết kiệm điện và cần dễ truy
vết khi có lỗi.

Nền tảng này thuyết phục hơn ở bốn điểm chính:

\begin{itemize}
\tightlist
\item
  \textbf{Khớp yêu cầu điều phối thời gian thực của hệ thống}: ESP-IDF +
  FreeRTOS cho phép tổ chức state machine trung tâm, các driver bất đồng
  bộ và primitive đồng bộ ở mức khối chức năng. Nói cách khác, đây là
  nền tảng đủ linh hoạt để firmware vừa theo dõi nhiều tín hiệu cùng lúc
  vừa không làm luồng xử lý chính trở nên rối và khó kiểm soát.
\item
  \textbf{Tối ưu cho ESP32-S3}: driver và power feature chính thức giúp
  giảm rủi ro tích hợp phần cứng, đồng thời rút ngắn đáng kể thời gian
  xử lý các lỗi phát sinh trong giai đoạn chế tạo và kiểm thử.
\item
  \textbf{Phù hợp bài toán năng lượng}: hỗ trợ deep sleep/wakeup và điều
  khiển peripheral theo trạng thái, nên phù hợp với một thiết bị phải
  luôn cân bằng giữa khả năng giám sát và mức tiêu thụ điện.
\item
  \textbf{Dễ kiểm soát chất lượng vận hành}: thuận lợi xây dựng state
  machine, retry strategy và fallback cho BLE/modem. Điều này đặc biệt
  quan trọng vì thiết bị không hoạt động trong phòng thí nghiệm lý tưởng
  mà trong môi trường xe có nhiễu, rung động và mất sóng bất chợt.
\end{itemize}

\subsubsection{3.2.2. Giải pháp
Firmware}\label{322-giux1ea3i-phuxe1p-firmware}

Firmware là nơi toàn bộ quyết định vận hành của thiết bị được hiện thực
hóa thành hành vi cụ thể: đọc cảm biến, đánh giá trạng thái xe, điều
khiển modem, quản lý năng lượng và gửi dữ liệu lên máy chủ. Trên cơ sở
phương án đã chọn ở mục 3.2.1, phần này trình bày cách nhóm biến các yêu
cầu đó thành một kiến trúc firmware đủ rõ ràng để phát triển, đủ gọn để
vận hành ổn định và đủ linh hoạt để mở rộng sau này.

\paragraph{3.2.2.1. Kiến trúc firmware và luồng hoạt
động}\label{3221-kiux1ebfn-truxfac-firmware-vuxe0-luux1ed3ng-houx1ea1t-ux111ux1ed9ng}

\subparagraph{a) Kiến trúc phân lớp (Layered
Architecture)}\label{a-kiux1ebfn-truxfac-phuxe2n-lux1edbp-layered-architecture}

Để người đọc dễ hình dung, firmware được tách theo từng lớp trách nhiệm
thay vì trình bày như một khối mã nguồn lớn. Cách chia này giúp trả lời
rõ ba câu hỏi quan trọng: lớp nào làm việc với phần cứng, lớp nào giữ
trạng thái vận hành, và lớp nào chịu trách nhiệm giao tiếp với hệ
thống bên ngoài.


{\thesissettableid{3.2.2A}{tab-ch3-tong-hop-thong-tin-ve-lop-thanh-phan-chinh-vai-tro}
\def\LTcaptype{table}
\begin{longtable}{|L{0.160\linewidth}|L{0.224\linewidth}|L{0.580\linewidth}|}
\caption{Tổng hợp thông tin về Lớp, Thành phần chính, Vai trò}\label{tab:ch3-tong-hop-thong-tin-ve-lop-thanh-phan-chinh-vai-tro}\\
\hline \textbf{Lớp} & \textbf{Thành phần chính} & \textbf{Vai trò} \\
\hline
\endfirsthead
\caption[]{Tổng hợp thông tin về Lớp, Thành phần chính, Vai trò (tiếp theo)}\\
\hline \textbf{Lớp} & \textbf{Thành phần chính} & \textbf{Vai trò} \\
\hline
\endhead
\hline
\endfoot
\endlastfoot
1 & \texttt{adc\_reader}, \texttt{imu\_lis3dh}, \texttt{modem\_at},
\texttt{modem\_lte}, \texttt{modem\_gnss}, \texttt{ble\_mgr},
\texttt{ble\_obd} & Trừu tượng hóa phần cứng và giao thức ngoại vi \\\hline
2 & \texttt{power\_mgr}, \texttt{nvs\_config}, \texttt{app\_state} &
Quản lý nguồn, lưu cấu hình NVS và duy trì ngữ cảnh RTC qua deep
sleep \\\hline
3 & \texttt{state\_machine}, \texttt{command\_handler} & Điều phối trạng
thái vận hành, lệnh điều khiển và luồng OTA \\\hline
4 & \texttt{data\_formatter}, \texttt{mqtt\_client} & Đóng gói payload
JSON, publish/subscribe MQTT theo topic chuẩn của hệ thống \\\hline
5 & \texttt{offline\_queue}, \texttt{sd\_log\_store} & Hàng đợi cục bộ
trên microSD (enqueue, replay theo ACK/QoS, quota GC, metadata phiên
ghi) \\\hline
\end{longtable}
}


Trong triển khai hiện tại, firmware không tách thành nhiều tác vụ nghiệp
vụ cấp cao như sơ đồ khái niệm ban đầu. Phần lớn quyết định vận hành
được gom về một state machine trung tâm; các primitive FreeRTOS chỉ được
dùng ở những nơi thật sự cần đồng bộ như NimBLE host task, BLE manager,
BLE OBD và lớp AT command.


{\thesissetfigureid{3.5}{fig-ch3-so-o-kien-truc-phan-lop-cua-firmware}
\begin{figure}[htbp]
\centering
\pandocbounded{\safeincludesvg{./assets/figures/04-chuong-3-giai-phap-firmware-hinh-3-5.svg}}
\caption{Sơ đồ kiến trúc phân lớp của firmware}
\label{fig:ch3-so-o-kien-truc-phan-lop-cua-firmware}
\end{figure}
}


\subparagraph{b) Cơ chế điều phối với
FreeRTOS}\label{b-cux1a1-chux1ebf-ux111iux1ec1u-phux1ed1i-vux1edbi-freertos}

Điểm đáng chú ý trong triển khai này là FreeRTOS không được dùng theo
hướng ``chia nhỏ càng nhiều task càng tốt''. Thay vào đó, firmware giữ
một luồng điều phối trung tâm trong \texttt{app\_main()}, nạp cấu hình từ
NVS, xác định trạng thái ban đầu theo nguyên nhân đánh thức rồi gọi
\texttt{state\_machine\_run()} theo chu kỳ ngắn. Cách tổ chức này giúp
logic điều khiển vẫn tập trung, dễ đọc và dễ debug, trong khi các cơ
chế đồng bộ của FreeRTOS chỉ xuất hiện ở những nơi thật sự cần thiết.

\begin{Shaded}
\begin{Highlighting}[]
\DataTypeTok{void}\NormalTok{ app\_main}\OperatorTok{(}\DataTypeTok{void}\OperatorTok{)} \OperatorTok{\{}
\NormalTok{    ESP\_ERROR\_CHECK}\OperatorTok{(}\NormalTok{nvs\_config\_init}\OperatorTok{());}

\NormalTok{    config\_t config }\OperatorTok{=} \OperatorTok{\{}\DecValTok{0}\OperatorTok{\};}
\NormalTok{    ESP\_ERROR\_CHECK}\OperatorTok{(}\NormalTok{nvs\_config\_load}\OperatorTok{(\&}\NormalTok{config}\OperatorTok{));}

\NormalTok{    app\_state\_t state }\OperatorTok{=}\NormalTok{ APP\_STATE\_INIT}\OperatorTok{;}
\NormalTok{    esp\_sleep\_wakeup\_cause\_t wakeup }\OperatorTok{=}\NormalTok{ esp\_sleep\_get\_wakeup\_cause}\OperatorTok{();}
    \ControlFlowTok{if} \OperatorTok{(}\NormalTok{wakeup }\OperatorTok{==}\NormalTok{ ESP\_SLEEP\_WAKEUP\_TIMER}\OperatorTok{)} \OperatorTok{\{}
\NormalTok{        state }\OperatorTok{=}\NormalTok{ APP\_STATE\_HEARTBEAT}\OperatorTok{;}
    \OperatorTok{\}} \ControlFlowTok{else} \ControlFlowTok{if} \OperatorTok{(}\NormalTok{wakeup }\OperatorTok{==}\NormalTok{ ESP\_SLEEP\_WAKEUP\_EXT0}\OperatorTok{)} \OperatorTok{\{}
\NormalTok{        state }\OperatorTok{=}\NormalTok{ APP\_STATE\_ALARM}\OperatorTok{;}
    \OperatorTok{\}}

\NormalTok{    ESP\_ERROR\_CHECK}\OperatorTok{(}\NormalTok{state\_machine\_init}\OperatorTok{(\&}\NormalTok{config}\OperatorTok{));}

    \ControlFlowTok{while} \OperatorTok{(}\KeywordTok{true}\OperatorTok{)} \OperatorTok{\{}
\NormalTok{        state }\OperatorTok{=}\NormalTok{ state\_machine\_run}\OperatorTok{(}\NormalTok{state}\OperatorTok{);}
\NormalTok{        vTaskDelay}\OperatorTok{(}\NormalTok{pdMS\_TO\_TICKS}\OperatorTok{(}\DecValTok{100}\OperatorTok{));}
    \OperatorTok{\}}
\OperatorTok{\}}
\end{Highlighting}
\end{Shaded}

Các primitive FreeRTOS vẫn xuất hiện ở các mô-đun cần đồng bộ truy cập
tài nguyên dùng chung. Cụ thể, NimBLE host chạy trong một task nền
riêng; \texttt{ble\_mgr} dùng queue và mutex để nhận kết quả khám
phá/kết nối; \texttt{ble\_obd} dùng semaphore cho phản hồi từ adapter;
còn \texttt{modem\_at} dùng mutex để tuần tự hóa lệnh AT trên UART
chung.


{\thesissetfigureid{3.6}{fig-ch3-so-o-tuong-tac-giua-cac-freertos-task}
\begin{figure}[htbp]
\centering
\pandocbounded{\safeincludesvg{./assets/figures/04-chuong-3-giai-phap-firmware-hinh-3-6.svg}}
\caption{Sơ đồ tương tác giữa các FreeRTOS task}
\label{fig:ch3-so-o-tuong-tac-giua-cac-freertos-task}
\end{figure}
}


\subparagraph{c) Luồng hoạt động cơ
bản}\label{c-luux1ed3ng-houx1ea1t-ux111ux1ed9ng-cux1a1-bux1ea3n}

Nếu nhìn từ góc độ vận hành, firmware lặp lại một chu trình rất rõ:
khởi tạo, nhận biết trạng thái xe, chọn chế độ phù hợp, thực hiện tác vụ
và quay về trạng thái tiết kiệm năng lượng khi có thể. Luồng tổng thể
được tổ chức theo trình tự sau:

\begin{enumerate}
\def\labelenumi{\arabic{enumi}.}
\item
  \textbf{Khởi tạo ngoại vi}: Cấu hình và khởi động các peripheral gồm
  IMU (LIS3DSH qua I2C), modem SIM7600CE-T (UART1, Auto mode
  \texttt{AT+CNMP=2}, GNSS qua \texttt{AT+CGNSINF}), ADC (đọc điện áp),
  và BLE stack (NimBLE).
\item
  \textbf{Đọc trạng thái IGN và điện áp ắc quy}: Hệ thống ưu tiên đọc
  trạng thái động cơ (IGN) trực tiếp từ ECU qua OBD2 BLE. Nếu không kết
  nối được OBD2, hệ thống fallback sang đo điện áp ắc quy qua ADC theo
  profile: profile 12V dùng \(IGN_{\mathrm{ON}} \ge 13.0\,\mathrm{V}\) và
  \(IGN_{\mathrm{OFF}} \le 12.0\,\mathrm{V}\); profile 24V dùng
  \(IGN_{\mathrm{ON}} \ge 26.0\,\mathrm{V}\) và
  \(IGN_{\mathrm{OFF}} \le 24.0\,\mathrm{V}\).
\item
  \textbf{Quyết định chế độ hoạt động}: Dựa trên trạng thái IGN và dữ
  liệu cảm biến, firmware chuyển sang chế độ phù hợp (lái xe, đỗ xe,
  hoặc cảnh báo).
\item
  \textbf{Thực thi tác vụ trong từng chế độ}:

  \begin{itemize}
  \tightlist
  \item
    \textbf{Chế độ lái xe (Driving)}: Kết nối OBD2, theo dõi liên tục,
    giữ BLE active, gửi telemetry định kỳ.
  \item
    \textbf{Chế độ đỗ xe (Parked)}: Không kết nối OBD2, chuyển sang deep
    sleep, gửi heartbeat định kỳ.
  \item
    \textbf{Chế độ cảnh báo (Alarm)}: Gửi cảnh báo ngay lập tức, theo
    dõi liên tục, có thể kết nối OBD2 tùy chọn.
  \end{itemize}
\item
  \textbf{Xử lý sau tác vụ}: Nếu IGN ON, giữ kết nối BLE và không deep
  sleep. Nếu IGN OFF, ngắt BLE và chuyển sang deep sleep để tiết kiệm
  năng lượng.
\end{enumerate}


{\thesissetfigureid{3.7}{fig-ch3-luu-o-thuat-toan-luong-hoat-ong-chinh-cua-firmware}
\begin{figure}[htbp]
\centering
\pandocbounded{\safeincludesvg{./assets/figures/04-chuong-3-giai-phap-firmware-hinh-3-7.svg}}
\caption{Lưu đồ thuật toán luồng hoạt động chính của firmware}
\label{fig:ch3-luu-o-thuat-toan-luong-hoat-ong-chinh-cua-firmware}
\end{figure}
}


\paragraph{3.2.2.2. Module giao tiếp Bluetooth
OBD2}\label{3222-module-giao-tiux1ebfp-bluetooth-obd2}

\subparagraph{a) Tổng quan giao tiếp BLE
OBD2}\label{a-tux1ed5ng-quan-giao-tiux1ebfp-ble-obd2}

Module giao tiếp Bluetooth Low Energy (BLE) OBD2 cho phép thiết bị theo
dõi kết nối không dây với adapter vgate iCar Pro --- thiết bị OBD2 hỗ
trợ BLE, được cắm trực tiếp vào cổng chẩn đoán OBD-II của xe. Firmware
sử dụng NimBLE stack (tích hợp trong ESP-IDF) để triển khai giao tiếp
BLE theo chuẩn GATT (Generic Attribute Profile).

Kiến trúc module BLE OBD2 được tổ chức theo bốn tầng logic: lớp điều
phối ứng dụng trong \texttt{state\_machine.c}, lớp giao thức OBD tại
\texttt{ble\_obd.c}, lớp quản lý BLE tại \texttt{ble\_mgr.c} và lớp BLE
stack của NimBLE. Cách tách lớp này giúp firmware cô lập xử lý PID OBD2
khỏi phần GATT/BLE thuần túy, thuận lợi cho debug và mở rộng.

\subparagraph{b) Quy trình kết nối BLE
OBD2}\label{b-quy-truxecnh-kux1ebft-nux1ed1i-ble-obd2}

Quy trình kết nối BLE với adapter OBD2 diễn ra theo các bước sau:

\begin{enumerate}
\def\labelenumi{\arabic{enumi}.}
\item
  \textbf{Kiểm tra địa chỉ BLE đã lưu}: Đọc địa chỉ MAC của vgate iCar
  Pro từ bộ nhớ flash (NVS). Nếu đã có địa chỉ, chuyển trực tiếp sang
  bước kết nối (reconnect nhanh trong 1--3 giây).
\item
  \textbf{Quét tìm thiết bị (Scan)}: Nếu chưa có địa chỉ, thực hiện BLE
  scan để tìm adapter. Quá trình scan lọc thiết bị theo tên (device
  name) hoặc service UUID đặc trưng của OBD2 adapter.
\item
  \textbf{Kết nối GATT}: Thiết lập kết nối BLE với adapter, thực hiện
  GATT service discovery để tìm các characteristic cần thiết (TX và RX
  characteristic).
\item
  \textbf{Gửi lệnh khởi tạo ELM327}: Gửi chuỗi lệnh khởi tạo giao thức
  ELM327 qua BLE GATT write:

  \begin{itemize}
  \tightlist
  \item
    \texttt{ATZ} --- Reset adapter
  \item
    \texttt{ATE0} --- Tắt echo
  \item
    \texttt{ATL0} --- Tắt line feed
  \item
    \texttt{ATS0} --- Tắt khoảng trắng
  \item
    \texttt{ATSP0} --- Tự động phát hiện giao thức OBD2
  \end{itemize}
\item
  \textbf{Đọc dữ liệu OBD2}: Gửi các lệnh PID (Parameter ID) để đọc
  thông số xe:

  \begin{itemize}
  \tightlist
  \item
    \texttt{010C} --- RPM động cơ
  \item
    \texttt{010D} --- Tốc độ xe
  \item
    \texttt{0105} --- Nhiệt độ nước làm mát
  \item
    \texttt{012F} --- Mức nhiên liệu
  \item
    \texttt{AT\ IGN} --- Trạng thái động cơ (ignition)
  \end{itemize}
\end{enumerate}

\begin{Shaded}
\begin{Highlighting}[]
\CommentTok{/* Quy trình kết nối và đọc dữ liệu OBD2 */}
\DataTypeTok{void}\NormalTok{ ble\_obd2\_task}\OperatorTok{(}\DataTypeTok{void} \OperatorTok{*}\NormalTok{param}\OperatorTok{)} \OperatorTok{\{}
    \CommentTok{/* Bước 1: Đọc BLE address từ NVS */}
    \DataTypeTok{uint8\_t}\NormalTok{ saved\_addr}\OperatorTok{[}\DecValTok{6}\OperatorTok{];}
    \DataTypeTok{bool}\NormalTok{ has\_saved }\OperatorTok{=}\NormalTok{ nvs\_read\_ble\_address}\OperatorTok{(}\NormalTok{saved\_addr}\OperatorTok{);}

    \CommentTok{/* Bước 2: Kết nối hoặc scan */}
    \ControlFlowTok{if} \OperatorTok{(}\NormalTok{has\_saved}\OperatorTok{)} \OperatorTok{\{}
\NormalTok{        ble\_connect\_direct}\OperatorTok{(}\NormalTok{saved\_addr}\OperatorTok{);}  \CommentTok{/* Reconnect nhanh */}
    \OperatorTok{\}} \ControlFlowTok{else} \OperatorTok{\{}
\NormalTok{        ble\_start\_scan}\OperatorTok{(}\NormalTok{BLE\_SCAN\_TIMEOUT\_MS}\OperatorTok{);}
\NormalTok{        ble\_connect\_first\_obd2\_device}\OperatorTok{();}
\NormalTok{        nvs\_save\_ble\_address}\OperatorTok{(}\NormalTok{connected\_addr}\OperatorTok{);}
    \OperatorTok{\}}

    \CommentTok{/* Bước 3: Khởi tạo ELM327 */}
\NormalTok{    ble\_obd\_send\_command}\OperatorTok{(}\StringTok{"ATZ"}\OperatorTok{);}
\NormalTok{    ble\_obd\_send\_command}\OperatorTok{(}\StringTok{"ATE0"}\OperatorTok{);}
\NormalTok{    ble\_obd\_send\_command}\OperatorTok{(}\StringTok{"ATSP0"}\OperatorTok{);}

    \CommentTok{/* Bước 4: Vòng lặp đọc dữ liệu */}
    \ControlFlowTok{while} \OperatorTok{(}\NormalTok{device\_state }\OperatorTok{==}\NormalTok{ STATE\_DRIVING}\OperatorTok{)} \OperatorTok{\{}
\NormalTok{        obd2\_data\_t data}\OperatorTok{;}
\NormalTok{        data}\OperatorTok{.}\NormalTok{rpm   }\OperatorTok{=}\NormalTok{ ble\_obd\_read\_pid}\OperatorTok{(}\NormalTok{PID\_RPM}\OperatorTok{);}
\NormalTok{        data}\OperatorTok{.}\NormalTok{speed }\OperatorTok{=}\NormalTok{ ble\_obd\_read\_pid}\OperatorTok{(}\NormalTok{PID\_SPEED}\OperatorTok{);}
\NormalTok{        data}\OperatorTok{.}\NormalTok{temp  }\OperatorTok{=}\NormalTok{ ble\_obd\_read\_pid}\OperatorTok{(}\NormalTok{PID\_COOLANT\_TEMP}\OperatorTok{);}
\NormalTok{        data}\OperatorTok{.}\NormalTok{fuel  }\OperatorTok{=}\NormalTok{ ble\_obd\_read\_pid}\OperatorTok{(}\NormalTok{PID\_FUEL\_LEVEL}\OperatorTok{);}

\NormalTok{        xQueueSend}\OperatorTok{(}\NormalTok{telemetry\_queue}\OperatorTok{,} \OperatorTok{\&}\NormalTok{data}\OperatorTok{,}\NormalTok{ pdMS\_TO\_TICKS}\OperatorTok{(}\DecValTok{100}\OperatorTok{));}
\NormalTok{        vTaskDelay}\OperatorTok{(}\NormalTok{pdMS\_TO\_TICKS}\OperatorTok{(}\NormalTok{TRACKING\_INTERVAL\_MS}\OperatorTok{));}
    \OperatorTok{\}}
\OperatorTok{\}}
\end{Highlighting}
\end{Shaded}


{\thesissetfigureid{3.8}{fig-ch3-luu-o-thuat-toan-quy-trinh-ket-noi-va-oc-du-lieu}
\begin{figure}[htbp]
\centering
\pandocbounded{\safeincludesvg{./assets/figures/04-chuong-3-giai-phap-firmware-hinh-3-8.svg}}
\caption{Lưu đồ thuật toán quy trình kết nối và đọc dữ liệu BLE OBD2}
\label{fig:ch3-luu-o-thuat-toan-quy-trinh-ket-noi-va-oc-du-lieu}
\end{figure}
}


\subparagraph{c) Chiến lược kết nối theo chế độ hoạt
động}\label{c-chiux1ebfn-lux1b0ux1ee3c-kux1ebft-nux1ed1i-theo-chux1ebf-ux111ux1ed9-houx1ea1t-ux111ux1ed9ng-1}

Việc kết nối BLE OBD2 được tối ưu hóa theo từng chế độ hoạt động của
thiết bị để tiết kiệm năng lượng:

\textbf{Chế độ lái xe (IGN ON)}: ESP32-S3 duy trì kết nối BLE liên tục
với OBD2 adapter. Dữ liệu OBD2 được đọc định kỳ (mỗi 5--10 giây). Hệ
thống chỉ sử dụng light sleep (giữ BLE active) thay vì deep sleep.

\textbf{Chế độ đỗ xe (IGN OFF)}: Không kết nối OBD2 adapter. Trạng thái
IGN đã được xác nhận trước khi chuyển sang deep sleep. Việc không kết
nối BLE khi đỗ xe giúp tiết kiệm đáng kể thời gian wake-up và năng lượng
tiêu thụ.

\textbf{Chế độ cảnh báo (Motion Detected)}: Kết nối OBD2 là tùy chọn. Hệ
thống có thể chỉ sử dụng IMU và GPS để xác nhận chuyển động mà không cần
dữ liệu OBD2.

Đặc biệt, khi ESP32-S3 chuyển sang deep sleep, kết nối BLE bị ngắt hoàn
toàn. Mỗi lần wake-up, firmware cần thực hiện reconnect với adapter. Nhờ
có việc lưu địa chỉ MAC vào flash, thời gian reconnect chỉ mất khoảng
1--3 giây, nhanh hơn đáng kể so với pairing lần đầu (3--10 giây).

\subparagraph{d) Xử lý lỗi và cơ chế dự phòng
(Fallback)}\label{d-xux1eed-luxfd-lux1ed7i-vuxe0-cux1a1-chux1ebf-dux1ef1-phuxf2ng-fallback}

Module BLE OBD2 được thiết kế với nhiều lớp xử lý lỗi:


{\thesissettableid{3.2.2B}{tab-ch3-tong-hop-thong-tin-ve-tinh-huong-loi-xu-ly-thoi-gian}
\def\LTcaptype{table}
\begin{longtable}{|L{0.291\linewidth}|L{0.445\linewidth}|L{0.228\linewidth}|}
\caption{Tổng hợp thông tin về Tình huống lỗi, Xử lý, Thời gian timeout}\label{tab:ch3-tong-hop-thong-tin-ve-tinh-huong-loi-xu-ly-thoi-gian}\\
\hline \textbf{Tình huống lỗi} & \textbf{Xử lý} & \textbf{Thời gian timeout} \\
\hline
\endfirsthead
\caption[]{Tổng hợp thông tin về Tình huống lỗi, Xử lý, Thời gian timeout (tiếp theo)}\\
\hline \textbf{Tình huống lỗi} & \textbf{Xử lý} & \textbf{Thời gian timeout} \\
\hline
\endhead
\hline
\endfoot
\endlastfoot
Không kết nối được adapter & Retry 2--3 lần với delay 2 giây, sau đó
fallback & 10 giây \\\hline
Mất kết nối giữa chừng & Reconnect 2--3 lần, nếu thất bại thì fallback &
5 giây mỗi lần \\\hline
Adapter không phản hồi lệnh & Retry 1--2 lần, nếu thất bại thì fallback
& 5 giây \\\hline
\end{longtable}
}


\textbf{Cơ chế dự phòng (Fallback)}: Khi không thể giao tiếp với OBD2
adapter, hệ thống chuyển sang phát hiện trạng thái động cơ bằng phương
pháp đo điện áp ắc quy. Phương pháp này có độ chính xác thấp hơn OBD2
nhưng vẫn đảm bảo hệ thống hoạt động liên tục. Tiêu chí phân biệt theo
profile: profile 12V dùng \(IGN_{\mathrm{ON}} \ge 13.0\,\mathrm{V}\) và
\(IGN_{\mathrm{OFF}} \le 12.0\,\mathrm{V}\); profile 24V dùng
\(IGN_{\mathrm{ON}} \ge 26.0\,\mathrm{V}\) và
\(IGN_{\mathrm{OFF}} \le 24.0\,\mathrm{V}\).

Mọi lỗi kết nối OBD2 đều được ghi lại vào log để phục vụ việc giám sát
và khắc phục sự cố từ xa.

\paragraph{3.2.2.3. Module điều khiển modem
SIMCom}\label{3223-module-ux111iux1ec1u-khiux1ec3n-modem-simcom}

\subparagraph{a) Kiến trúc modem SIMCom SIM7600CE-T (LTE + GNSS tích
hợp)}\label{a-kiux1ebfn-truxfac-modem-simcom-sim7600ce-t-lte--gnss-tuxedch-hux1ee3p}

Trong kiến trúc hiện tại, modem SIMCom SIM7600CE-T đảm nhận toàn bộ
luồng LTE và GNSS. ESP32-S3 giao tiếp với modem qua UART1 (GPIO16 TX,
GPIO17 RX). Firmware sử dụng một chuỗi lệnh AT chính xác để đảm bảo
modem sẵn sàng trước khi bật PDP context và khởi động GNSS.


{\thesissettableid{3.2.2C}{tab-ch3-tong-hop-thong-tin-ve-buoc-lenh-at-muc-ich-chinh}
\def\LTcaptype{table}
\begin{longtable}{|L{0.160\linewidth}|L{0.224\linewidth}|L{0.580\linewidth}|}
\caption{Tổng hợp thông tin về Bước, Lệnh AT, Mục đích chính}\label{tab:ch3-tong-hop-thong-tin-ve-buoc-lenh-at-muc-ich-chinh}\\
\hline \textbf{Bước} & \textbf{Lệnh AT} & \textbf{Mục đích chính} \\
\hline
\endfirsthead
\caption[]{Tổng hợp thông tin về Bước, Lệnh AT, Mục đích chính (tiếp theo)}\\
\hline \textbf{Bước} & \textbf{Lệnh AT} & \textbf{Mục đích chính} \\
\hline
\endhead
\hline
\endfoot
\endlastfoot
1 & \texttt{AT} & Kiểm tra modem còn phản hồi \\\hline
2 & \texttt{AT+CPIN?} & Kiểm tra SIM card đã sẵn sàng \\\hline
3 & \texttt{AT+CNMP=2} & Đặt chế độ tự động LTE/UMTS/GSM \\\hline
4 & \nolinkurl{AT+CGDCONT=1,"IP","internet"} & Cấu hình PDP context với APN
mặc định \texttt{internet} \\\hline
5 & \texttt{AT+CEREG?} & Poll đến khi \texttt{+CEREG:\ 0,1} (hoặc
\texttt{stat}=5) trước khi kích hoạt PDP context \\\hline
6 & \texttt{AT+CGACT=1,1} & Bật PDP context sau khi đăng ký mạng thành
công \\\hline
7 & \texttt{AT+CGPADDR=1} & Ghi nhận địa chỉ IP được cấp (dùng để log và
kiểm tra) \\\hline
8 & \texttt{AT+CGNSPWR=1} & Bật GNSS tích hợp và đợi modem trả về fix vị
trí \\\hline
9 & \texttt{AT+CGNSINF} & Đọc tọa độ GNSS hiện tại (vị trí, tốc độ, số
vệ tinh); \texttt{AT+CGNSTST} chỉ dùng nếu cần dữ liệu NMEA hoàn toàn \\\hline
\end{longtable}
}


Toàn bộ lệnh AT cho LTE và GNSS được thực hiện trên cùng UART1. Trình tự
kiểm tra \texttt{AT+CEREG?} trước khi gọi \texttt{AT+CGACT=1,1} giúp
tránh kích hoạt PDP context khi modem chưa đăng ký mạng.

\subparagraph{b) Trình tự khởi tạo LTE/GNSS trong
firmware}\label{b-trinh-tux1ef1-khoi-tao-ltegnss-trong-firmware}

Ở tầng firmware, chuỗi lệnh trên không được gửi dồn trong một lần, mà
được tách thành các bước kiểm tra có điều kiện: xác nhận modem phản hồi,
xác nhận SIM sẵn sàng, chờ modem đăng ký mạng, sau đó mới bật PDP
context và GNSS. Cách tổ chức này giúp quá trình khởi tạo dễ retry hơn
và hạn chế lỗi dây chuyền khi điều kiện sóng di động thay đổi.

\subparagraph{c) Quản lý kết nối và tiết kiệm
điện}\label{c-quux1ea3n-luxfd-kux1ebft-nux1ed1i-vuxe0-tiux1ebft-kiux1ec7m-ux111iux1ec7n}

\begin{itemize}
\tightlist
\item
  \textbf{Khi hoạt động (driving mode):} Giữ PDP context mở và GNSS bật
  liên tục để gửi telemetry định kỳ và heartbeat.
\item
  \textbf{Khi dừng (parking mode):} Firmware đóng PDP context
  (\texttt{AT+CGACT=0,1}), đưa modem vào \texttt{AT+CSCLK=1} hoặc
  \texttt{AT+CFUN=0}, rồi chuyển ESP32-S3 sang deep sleep. Khi wake-up,
  modem được đánh thức qua UART và thực hiện lại các lệnh
  \texttt{AT+CFUN=1} + \texttt{AT+CGACT=1,1}.
\item
  \textbf{GNSS:} Sau khi \texttt{AT+CGNSPWR=1}, firmware tuần tự gọi
  \texttt{AT+CGNSINF} để lấy dữ liệu; \texttt{AT+CGNSTST} chỉ dùng khi
  cần stream NMEA để hỗ trợ debug hoặc phân tích phụ trợ.
\end{itemize}

\subparagraph{d) Xử lý sự cố}\label{d-xux1eed-luxfd-sux1ef1-cux1ed1}


{\thesissettableid{3.2.2D}{tab-ch3-tong-hop-thong-tin-ve-tinh-huong-bien-phap}
\def\LTcaptype{table}
\begin{longtable}{|L{0.220\linewidth}|L{0.740\linewidth}|}
\caption{Tổng hợp thông tin về Tình huống, Biện pháp}\label{tab:ch3-tong-hop-thong-tin-ve-tinh-huong-bien-phap}\\
\hline \textbf{Tình huống} & \textbf{Biện pháp} \\
\hline
\endfirsthead
\caption[]{Tổng hợp thông tin về Tình huống, Biện pháp (tiếp theo)}\\
\hline \textbf{Tình huống} & \textbf{Biện pháp} \\
\hline
\endhead
\hline
\endfoot
\endlastfoot
Modem không phản hồi & Reset nhanh qua PWR-KEY (GPIO26) rồi bắt đầu lại
chuỗi AT \\\hline
Mất kết nối 4G & Đợi \texttt{AT+CEREG?} thành \texttt{1}, gọi
\texttt{AT+CGACT=0,1} rồi \texttt{AT+CGACT=1,1} để tái khởi động PDP \\\hline
GNSS không fix dù đã bật & Gọi lại \texttt{AT+CGNSINF}, nếu vẫn thất bại
thì bật lại GNSS (\texttt{AT+CGNSPWR=0}/\texttt{1}) rồi thử lại \\\hline
Modem quá nóng & Đưa \texttt{AT+CFUN=0} để tạm thời tắt RF và chờ nhiệt
độ giảm trước khi tiếp tục \\\hline
\end{longtable}
}


Các chế độ ngủ (\texttt{AT+CSCLK=1}, \texttt{AT+CFUN=0}) và wake-up bằng
UART giữ cho modem tiêu thụ thấp mà vẫn có thể khởi động lại nhanh khi
ESP32-S3 tỉnh dậy.


{\thesissetfigureid{3.9}{fig-ch3-luu-o-ieu-phoi-modem-sim7600ce-t-theo-trang-thai-mang-va}
\begin{figure}[htbp]
\centering
\pandocbounded{\safeincludesvg{./assets/figures/04-chuong-3-giai-phap-firmware-hinh-3-9.svg}}
\caption{Lưu đồ điều phối modem SIM7600CE-T theo trạng thái mạng và chế độ vận hành}
\label{fig:ch3-luu-o-ieu-phoi-modem-sim7600ce-t-theo-trang-thai-mang-va}
\end{figure}
}


\paragraph{3.2.2.4. Module quản lý nguồn và
GPIO}\label{3224-module-quux1ea3n-luxfd-nguux1ed3n-vuxe0-gpio}

\subparagraph{a) Sơ đồ chân
GPIO}\label{a-sux1a1-ux111ux1ed3-chuxe2n-gpio}

Module quản lý nguồn sử dụng các chân GPIO của ESP32-S3 để điều khiển và
giám sát hệ thống nguồn điện. Bảng sau liệt kê đầy đủ các chân GPIO được
sử dụng:


{\thesissettableid{3.2.2E}{tab-ch3-tong-hop-thong-tin-ve-gpio-chuc-nang-huong}
\def\LTcaptype{table}
\begin{longtable}{|L{0.115\linewidth}|L{0.230\linewidth}|L{0.165\linewidth}|L{0.445\linewidth}|}
\caption{Tổng hợp thông tin về GPIO, Chức năng, Hướng}\label{tab:ch3-tong-hop-thong-tin-ve-gpio-chuc-nang-huong}\\
\hline \textbf{GPIO} & \textbf{Chức năng} & \textbf{Hướng} & \textbf{Mô tả} \\
\hline
\endfirsthead
\caption[]{Tổng hợp thông tin về GPIO, Chức năng, Hướng (tiếp theo)}\\
\hline \textbf{GPIO} & \textbf{Chức năng} & \textbf{Hướng} & \textbf{Mô tả} \\
\hline
\endhead
\hline
\endfoot
\endlastfoot
2 & IGN\_IN & Input & Đọc trạng thái động cơ (GPIO hoặc OBD2) \\\hline
4 & \nolinkurl{U_BATT_ADC} & Input & Đọc điện áp ắc quy (ADC 12-bit) \\\hline
5 & \nolinkurl{CHARGER_EN} & Output & Điều khiển IC sạc TP4056 \\\hline
18 & \nolinkurl{POWER_PATH_EN} & Output & Chọn nguồn cấp (ắc quy hoặc pin dự
phòng) \\\hline
19 & LVD\_STATUS & Input & Đọc trạng thái LVD từ khối LVD phần cứng \\\hline
21 & LIS3DSH\_INT & Input & Ngắt từ cảm biến gia tốc IMU \\\hline
47 & \nolinkurl{LIS3DSH_SDA} & I/O & Đường dữ liệu I2C \\\hline
48 & \nolinkurl{LIS3DSH_SCL} & I/O & Đường clock I2C \\\hline
16 & \nolinkurl{MODEM_UART_TX} & Output & UART TX đến modem \\\hline
17 & \nolinkurl{MODEM_UART_RX} & Input & UART RX từ modem \\\hline
26 & \nolinkurl{MODEM_PWR-KEY} & Output & Điều khiển nguồn modem \\\hline
\end{longtable}
}


\subparagraph{b) Điều khiển Power Path (Diode OR +
EN)}\label{b-ux111iux1ec1u-khiux1ec3n-power-path-diode-or--en}

Hệ thống dùng power path theo \textbf{diode OR} giữa nhánh 5V chính
(MP2482) và nhánh 5V backup (SX1308). Firmware điều khiển qua GPIO18 để
ưu tiên nhánh nguồn phù hợp theo profile điện áp.

\begin{itemize}
\tightlist
\item
  \textbf{GPIO18 = LOW (0):} Ưu tiên nguồn ắc quy (nhánh MP2482 hoạt
  động)
\item
  \textbf{GPIO18 = HIGH (1):} Giảm/khóa nhánh chính và cho phép nhánh
  backup qua SX1308 + diode OR
\end{itemize}

Logic điều khiển nguồn bám theo profile 12V/24V:

\begin{Shaded}
\begin{Highlighting}[]
\NormalTok{power\_status\_t power\_get\_status}\OperatorTok{(}\DataTypeTok{float}\NormalTok{ lvd\_threshold\_v}\OperatorTok{,} \DataTypeTok{float}\NormalTok{ lvd\_hysteresis\_v}\OperatorTok{)} \OperatorTok{\{}
\NormalTok{    power\_status\_t status }\OperatorTok{=} \OperatorTok{\{}
        \OperatorTok{.}\NormalTok{voltage }\OperatorTok{=}\NormalTok{ adc\_read\_battery\_voltage}\OperatorTok{(),}
        \OperatorTok{.}\NormalTok{is\_low }\OperatorTok{=} \KeywordTok{false}\OperatorTok{,}
        \OperatorTok{.}\NormalTok{charger\_on }\OperatorTok{=}\NormalTok{ gpio\_get\_level}\OperatorTok{(}\NormalTok{PIN\_CHARGER\_EN}\OperatorTok{)} \OperatorTok{==} \DecValTok{1}\OperatorTok{,}
    \OperatorTok{\};}

    \ControlFlowTok{if} \OperatorTok{(}\NormalTok{status}\OperatorTok{.}\NormalTok{voltage }\OperatorTok{\textless{}=}\NormalTok{ lvd\_threshold\_v}\OperatorTok{)} \OperatorTok{\{}
\NormalTok{        s\_voltage\_low\_latched }\OperatorTok{=} \KeywordTok{true}\OperatorTok{;}
    \OperatorTok{\}} \ControlFlowTok{else} \ControlFlowTok{if} \OperatorTok{(}\NormalTok{status}\OperatorTok{.}\NormalTok{voltage }\OperatorTok{\textgreater{}=}\NormalTok{ lvd\_hysteresis\_v}\OperatorTok{)} \OperatorTok{\{}
\NormalTok{        s\_voltage\_low\_latched }\OperatorTok{=} \KeywordTok{false}\OperatorTok{;}
    \OperatorTok{\}}

\NormalTok{    status}\OperatorTok{.}\NormalTok{is\_low }\OperatorTok{=}\NormalTok{ s\_voltage\_low\_latched }\OperatorTok{||}\NormalTok{ power\_is\_low\_voltage}\OperatorTok{();}
    \ControlFlowTok{return}\NormalTok{ status}\OperatorTok{;}
\OperatorTok{\}}
\end{Highlighting}
\end{Shaded}


{\thesissetfigureid{3.10}{fig-ch3-luu-o-thuat-toan-ieu-khien-power-path}
\begin{figure}[htbp]
\centering
\pandocbounded{\safeincludesvg{./assets/figures/04-chuong-3-giai-phap-firmware-hinh-3-10.svg}}
\caption{Lưu đồ thuật toán điều khiển power path}
\label{fig:ch3-luu-o-thuat-toan-ieu-khien-power-path}
\end{figure}
}


\subparagraph{c) Điều khiển IC sạc
(TP4056)}\label{c-ux111iux1ec1u-khiux1ec3n-ic-sux1ea1c-tp4056}

IC sạc TP4056 được điều khiển qua chân GPIO5 (CHARGER\_EN). Logic sạc
được thiết kế để bảo vệ cả ắc quy xe lẫn pin dự phòng:


{\thesissettableid{3.2.2F}{tab-ch3-tong-hop-thong-tin-ve-ieu-kien-trang-thai-charger-ly-do}
\def\LTcaptype{table}
\begin{longtable}{|L{0.327\linewidth}|L{0.243\linewidth}|L{0.394\linewidth}|}
\caption{Tổng hợp thông tin về Điều kiện, Trạng thái charger, Lý do}\label{tab:ch3-tong-hop-thong-tin-ve-ieu-kien-trang-thai-charger-ly-do}\\
\hline \textbf{Điều kiện} & \textbf{Trạng thái charger} & \textbf{Lý do} \\
\hline
\endfirsthead
\caption[]{Tổng hợp thông tin về Điều kiện, Trạng thái charger, Lý do (tiếp theo)}\\
\hline \textbf{Điều kiện} & \textbf{Trạng thái charger} & \textbf{Lý do} \\
\hline
\endhead
\hline
\endfoot
\endlastfoot
IGN ON và \(U_{\text{batt}}\) vượt ngưỡng profile (ví dụ \(> 12\,\mathrm{V}\) cho hệ
12V) & Enable (GPIO5 = HIGH) & Máy phát điện đang nạp, có thể sạc pin dự
phòng \\\hline
IGN OFF & Disable (GPIO5 = LOW) & Bảo vệ ắc quy không bị hao pin khi xe
tắt máy \\\hline
\(U_{\text{batt}}\) dưới ngưỡng profile (ví dụ \(< 12\,\mathrm{V}\) cho hệ 12V) & Disable
(GPIO5 = LOW) & Ắc quy yếu, không đủ năng lượng để sạc pin dự phòng \\\hline
\end{longtable}
}


\subparagraph{d) Giám sát điện áp và
LVD}\label{d-giuxe1m-suxe1t-ux111iux1ec7n-uxe1p-vuxe0-lvd}

Firmware đọc điện áp ắc quy qua kênh ADC 12-bit (GPIO4) với bộ chia áp
để đưa điện áp ắc quy 12V hoặc 24V về dải đo của ADC (0--3.3V). Giá trị
ADC được chuyển đổi sang điện áp thực thông qua công thức hiệu chuẩn
(calibration).

Ngoài ra, hệ thống có thêm kênh giám sát LVD độc lập từ khối LVD phần
cứng (GPIO19). Đây là kênh dự phòng cho ADC, cho phép kiểm tra nhanh
trạng thái nguồn với cơ chế trễ (hysteresis): profile 12V dùng
\(Switch_{\mathrm{OFF}} = 12.0\,\mathrm{V}\),
\(Switch_{\mathrm{ON}} = 12.2\,\mathrm{V}\); profile 24V dùng
\(Switch_{\mathrm{OFF}} = 24.0\,\mathrm{V}\),
\(Switch_{\mathrm{ON}} = 24.4\,\mathrm{V}\). Khối LVD xử lý hysteresis ở mức phần cứng, tránh hiện
 tượng dao động (oscillation) quanh ngưỡng.

\subparagraph{e) Các chế độ quản lý
nguồn}\label{e-cuxe1c-chux1ebf-ux111ux1ed9-quux1ea3n-luxfd-nguux1ed3n}

Firmware hỗ trợ ba chế độ năng lượng chính:


{\thesissettableid{3.2.2G}{tab-ch3-tong-hop-thong-tin-ve-che-o-trang-thai-esp32-trang-thai}
\def\LTcaptype{table}
\begin{longtable}{|L{0.158\linewidth}|L{0.202\linewidth}|L{0.235\linewidth}|L{0.188\linewidth}|L{0.162\linewidth}|}
\caption{Tổng hợp thông tin về Chế độ, Trạng thái ESP32, Trạng thái modem}\label{tab:ch3-tong-hop-thong-tin-ve-che-o-trang-thai-esp32-trang-thai}\\
\hline \textbf{Chế độ} & \textbf{Trạng thái ESP32} & \textbf{Trạng thái modem} & \textbf{Trạng thái BLE} & \textbf{Dòng
tiêu thụ} \\
\hline
\endfirsthead
\caption[]{Tổng hợp thông tin về Chế độ, Trạng thái ESP32, Trạng thái modem (tiếp theo)}\\
\hline \textbf{Chế độ} & \textbf{Trạng thái ESP32} & \textbf{Trạng thái modem} & \textbf{Trạng thái BLE} & \textbf{Dòng
tiêu thụ} \\
\hline
\endhead
\hline
\endfoot
\endlastfoot
Normal (Driving) & Active & 4G + GNSS active & Connected & 150--250
mA \\\hline
Light Sleep & Light sleep & Sleep (CSCLK) & Giữ kết nối & 10--20 mA \\\hline
Deep Sleep & Deep sleep & Deep sleep (CFUN=0) & Ngắt kết nối &
\textless{} 2 mA \\\hline
\end{longtable}
}


Chuyển đổi giữa các chế độ dựa trên trạng thái IGN và dữ liệu cảm biến,
được điều phối bởi máy trạng thái chính (trình bày tại mục 3.2.2.5).

\paragraph{3.2.2.5. Định dạng dữ liệu và máy trạng
thái}\label{3225-ux111ux1ecbnh-dux1ea1ng-dux1eef-liux1ec7u-vuxe0-muxe1y-trux1ea1ng-thuxe1i}

\subparagraph{a) Định dạng dữ liệu
MQTT}\label{a-ux111ux1ecbnh-dux1ea1ng-dux1eef-liux1ec7u-mqtt}

Dữ liệu từ thiết bị được đóng gói theo định dạng JSON và truyền lên
broker qua nhóm topic \texttt{v1/\{device\_id\}/...}. Trong hệ thống
hiện tại, thiết bị phát ba nhóm payload chính là \texttt{rawdata},
\texttt{status} và \texttt{events}; ngoài ra còn có luồng
\texttt{firmware} để báo trạng thái OTA và luồng \texttt{commands} theo
chiều ngược lại từ máy chủ.

\textbf{Thông điệp Rawdata (dữ liệu telemetry chính)}

Đây là gói dữ liệu được gửi định kỳ khi thiết bị đang hoạt động, phản
ánh trạng thái GNSS, nguồn và cảm biến rung:

\begin{Shaded}
\begin{Highlighting}[]
\FunctionTok{\{}
  \DataTypeTok{"device\_id"}\FunctionTok{:} \StringTok{"TRACKER\_001"}\FunctionTok{,}
  \DataTypeTok{"auth\_token"}\FunctionTok{:} \StringTok{"device\_secret\_token"}\FunctionTok{,}
  \DataTypeTok{"timestamp"}\FunctionTok{:} \DecValTok{1710000000000}\FunctionTok{,}
  \DataTypeTok{"uptime"}\FunctionTok{:} \DecValTok{452130}\FunctionTok{,}
  \DataTypeTok{"data"}\FunctionTok{:} \FunctionTok{\{}
    \DataTypeTok{"vibration"}\FunctionTok{:} \DecValTok{128}\FunctionTok{,}
    \DataTypeTok{"battery\_top"}\FunctionTok{:} \FloatTok{12.54}\FunctionTok{,}
    \DataTypeTok{"battery\_bot"}\FunctionTok{:} \FloatTok{3.96}\FunctionTok{,}
    \DataTypeTok{"latitude"}\FunctionTok{:} \FloatTok{10.77653}\FunctionTok{,}
    \DataTypeTok{"longitude"}\FunctionTok{:} \FloatTok{106.70098}\FunctionTok{,}
    \DataTypeTok{"speed"}\FunctionTok{:} \FloatTok{42.5}\FunctionTok{,}
    \DataTypeTok{"course"}\FunctionTok{:} \FloatTok{181.2}\FunctionTok{,}
    \DataTypeTok{"satellites"}\FunctionTok{:} \DecValTok{11}\FunctionTok{,}
    \DataTypeTok{"ignition"}\FunctionTok{:} \KeywordTok{true}\FunctionTok{,}
    \DataTypeTok{"error\_code"}\FunctionTok{:} \DecValTok{0}
  \FunctionTok{\}}
\FunctionTok{\}}
\end{Highlighting}
\end{Shaded}

\textbf{Thông điệp Event (sự kiện/cảnh báo)}

Được gửi khi firmware phát hiện một sự kiện bất thường hoặc muốn ghi
nhận một mốc vận hành đáng chú ý:

\begin{Shaded}
\begin{Highlighting}[]
\FunctionTok{\{}
  \DataTypeTok{"device\_id"}\FunctionTok{:} \StringTok{"TRACKER\_001"}\FunctionTok{,}
  \DataTypeTok{"auth\_token"}\FunctionTok{:} \StringTok{"device\_secret\_token"}\FunctionTok{,}
  \DataTypeTok{"event\_type"}\FunctionTok{:} \StringTok{"high\_vibration"}\FunctionTok{,}
  \DataTypeTok{"code"}\FunctionTok{:} \DecValTok{201}\FunctionTok{,}
  \DataTypeTok{"message"}\FunctionTok{:} \StringTok{"Vehicle movement detected while parked"}\FunctionTok{,}
  \DataTypeTok{"timestamp"}\FunctionTok{:} \DecValTok{1710000005123}
\FunctionTok{\}}
\end{Highlighting}
\end{Shaded}

\textbf{Thông điệp Firmware (trạng thái OTA)}

Luồng \texttt{firmware} được dùng để ghi nhận vòng đời cập nhật firmware
từ lúc gán tác vụ đến khi xác nhận thành công sau reboot hoặc rollback:

\begin{Shaded}
\begin{Highlighting}[]
\FunctionTok{\{}
  \DataTypeTok{"device\_id"}\FunctionTok{:} \StringTok{"TRACKER\_001"}\FunctionTok{,}
  \DataTypeTok{"auth\_token"}\FunctionTok{:} \StringTok{"device\_secret\_token"}\FunctionTok{,}
  \DataTypeTok{"jobId"}\FunctionTok{:} \StringTok{"ota\_1710000000\_ab12cd"}\FunctionTok{,}
  \DataTypeTok{"status"}\FunctionTok{:} \StringTok{"confirming"}\FunctionTok{,}
  \DataTypeTok{"progress"}\FunctionTok{:} \DecValTok{99}\FunctionTok{,}
  \DataTypeTok{"targetVersion"}\FunctionTok{:} \StringTok{"1.2.0"}\FunctionTok{,}
  \DataTypeTok{"currentVersion"}\FunctionTok{:} \StringTok{"1.1.3"}\FunctionTok{,}
  \DataTypeTok{"partition"}\FunctionTok{:} \StringTok{"ota\_0"}
\FunctionTok{\}}
\end{Highlighting}
\end{Shaded}

\textbf{Thông điệp Status (trạng thái hoạt động)}

Được gửi khi thiết bị đổi trạng thái phiên làm việc hoặc phát heartbeat
tại các mốc cần theo dõi:

\begin{Shaded}
\begin{Highlighting}[]
\FunctionTok{\{}
  \DataTypeTok{"device\_id"}\FunctionTok{:} \StringTok{"TRACKER\_001"}\FunctionTok{,}
  \DataTypeTok{"auth\_token"}\FunctionTok{:} \StringTok{"device\_secret\_token"}\FunctionTok{,}
  \DataTypeTok{"status"}\FunctionTok{:} \StringTok{"running"}\FunctionTok{,}
  \DataTypeTok{"session\_id"}\FunctionTok{:} \DecValTok{1024}\FunctionTok{,}
  \DataTypeTok{"timestamp"}\FunctionTok{:} \DecValTok{1710000010000}
\FunctionTok{\}}
\end{Highlighting}
\end{Shaded}

\subparagraph{b) Thiết kế MQTT
Topic}\label{b-thiux1ebft-kux1ebf-mqtt-topic}

Hệ thống sử dụng cấu trúc phân cấp MQTT topic để tách rõ luồng
telemetry, trạng thái, sự kiện và OTA:


{\thesissettableid{3.2.2H}{tab-ch3-tong-hop-thong-tin-ve-topic-huong-qos}
\def\LTcaptype{table}
\begin{longtable}{|L{0.254\linewidth}|L{0.177\linewidth}|L{0.115\linewidth}|L{0.409\linewidth}|}
\caption{Tổng hợp thông tin về Topic, Hướng, QoS}\label{tab:ch3-tong-hop-thong-tin-ve-topic-huong-qos}\\
\hline \textbf{Topic} & \textbf{Hướng} & \textbf{QoS} & \textbf{Mô tả} \\
\hline
\endfirsthead
\caption[]{Tổng hợp thông tin về Topic, Hướng, QoS (tiếp theo)}\\
\hline \textbf{Topic} & \textbf{Hướng} & \textbf{QoS} & \textbf{Mô tả} \\
\hline
\endhead
\hline
\endfoot
\endlastfoot
\nolinkurl{v1/{device_id}/rawdata} & Device → Server & 0 &
Payload telemetry chính; tần suất cao, chấp nhận mất vài mẫu \\\hline
\nolinkurl{v1/{device_id}/status} & Device → Server & 1 &
Heartbeat, trạng thái phiên, online/offline \\\hline
\nolinkurl{v1/{device_id}/events} & Device → Server & 1 &
Cảnh báo và sự kiện quan trọng \\\hline
\nolinkurl{v1/{device_id}/firmware} & Device → Server & 1
& Tiến trình OTA và kết quả cập nhật firmware \\\hline
\nolinkurl{v1/{device_id}/commands} & Server → Device & 1
& Lệnh điều khiển và cập nhật cấu hình từ máy chủ \\\hline
\end{longtable}
}


Các lệnh điều khiển từ máy chủ hiện được firmware xử lý gồm:


{\thesissettableid{3.2.2I}{tab-ch3-tong-hop-thong-tin-ve-lenh-muc-ich-tham-so-chinh}
\def\LTcaptype{table}
\begin{longtable}{|L{0.255\linewidth}|L{0.315\linewidth}|L{0.394\linewidth}|}
\caption{Tổng hợp thông tin về Lệnh, Mục đích, Tham số chính}\label{tab:ch3-tong-hop-thong-tin-ve-lenh-muc-ich-tham-so-chinh}\\
\hline \textbf{Lệnh} & \textbf{Mục đích} & \textbf{Tham số chính} \\
\hline
\endfirsthead
\caption[]{Tổng hợp thông tin về Lệnh, Mục đích, Tham số chính (tiếp theo)}\\
\hline \textbf{Lệnh} & \textbf{Mục đích} & \textbf{Tham số chính} \\
\hline
\endhead
\hline
\endfoot
\endlastfoot
\texttt{update\_config} & Cập nhật cấu hình thiết bị &
\texttt{heartbeat\_interval\_s}, \texttt{tracking\_interval\_s} \\\hline
\texttt{request\_location} & Yêu cầu gửi vị trí ngay ở chu kỳ kế tiếp &
Không \\\hline
\texttt{enable\_tracking} & Bật hoặc tắt phát telemetry định kỳ &
\texttt{enabled} (boolean) \\\hline
\texttt{reboot} & Khởi động lại thiết bị & Không \\\hline
\texttt{ota\_update} & Kích hoạt cập nhật firmware OTA & \texttt{jobId},
\texttt{version}, \texttt{url}, \texttt{size}, \texttt{sha256},
\texttt{force}, \texttt{confirmTimeoutSec} \\\hline
\texttt{manual\_rollback} / \texttt{ota\_rollback} & Yêu cầu quay về
firmware trước đó & Không bắt buộc \\\hline
\end{longtable}
}


\subparagraph{c) Chiến lược xử lý mất kết
nối}\label{c-chiux1ebfn-lux1b0ux1ee3c-xux1eed-luxfd-mux1ea5t-kux1ebft-nux1ed1i}

Ở phạm vi triển khai hiện tại, firmware kết hợp hai cơ chế song song:
(1) tự phục hồi kết nối MQTT và (2) hàng đợi cục bộ trên \textbf{microSD
SDMMC 4-bit} để giữ dữ liệu khi mất mạng. Luồng dữ liệu được hiện thực
qua cặp module \texttt{offline\_queue} và \texttt{sd\_log\_store}: bản
ghi được enqueue khi telemetry phát sinh
(\texttt{offline\_queue\_enqueue()}), ghi xuống FATFS
(\texttt{sd\_log\_store\_append()}), sau đó replay theo thứ tự sequence
khi mạng phục hồi (\texttt{offline\_queue\_replay\_tick()}) và chốt ACK
cho bản ghi quan trọng
(\texttt{offline\_queue\_handle\_publish\_ack()}).

Thiết kế này cho phép tách bản ghi thành hai nhóm: \textbf{critical}
(QoS 1, yêu cầu ACK) và \textbf{non-critical} (QoS 0, ưu tiên
throughput), đồng thời áp dụng cơ chế quota + garbage collection
(\texttt{sd\_log\_store\_gc\_if\_needed()}) để tránh đầy thẻ. Trong phạm
vi vận hành hiện tại, hệ thống giả định thẻ SD ở trạng thái
\textbf{card-present-at-boot}; luồng hotplug runtime (rút/cắm nóng)
không phải phạm vi mục tiêu.

\begin{Shaded}
\begin{Highlighting}[]
\CommentTok{/* Chiến lược hiện tại: enqueue local + replay khi MQTT sẵn sàng */}
\NormalTok{offline\_queue\_enqueue}\OperatorTok{(}\NormalTok{type}\OperatorTok{,}\NormalTok{ payload}\OperatorTok{,}\NormalTok{ gps\_fix}\OperatorTok{,}\NormalTok{ net\_up}\OperatorTok{);}

\ControlFlowTok{if} \OperatorTok{(}\NormalTok{tracker\_mqtt\_is\_connected}\OperatorTok{())} \OperatorTok{\{}
\NormalTok{    offline\_queue\_replay\_tick}\OperatorTok{();}
\OperatorTok{\}}
\end{Highlighting}
\end{Shaded}

\subparagraph{d) Máy trạng thái thiết bị (Device State
Machine)}\label{d-muxe1y-trux1ea1ng-thuxe1i-thiux1ebft-bux1ecb-device-state-machine}

Máy trạng thái là cơ chế điều phối trung tâm của firmware, quyết định
hành vi thiết bị tại từng thời điểm vận hành. Hệ thống định nghĩa bảy
trạng thái chính:


{\thesissettableid{3.2.2J}{tab-ch3-tong-hop-thong-tin-ve-trang-thai-ma-mo-ta}
\def\LTcaptype{table}
\begin{longtable}{|L{0.224\linewidth}|L{0.160\linewidth}|L{0.580\linewidth}|}
\caption{Tổng hợp thông tin về Trạng thái, Mã, Mô tả}\label{tab:ch3-tong-hop-thong-tin-ve-trang-thai-ma-mo-ta}\\
\hline \textbf{Trạng thái} & \textbf{Mã} & \textbf{Mô tả} \\
\hline
\endfirsthead
\caption[]{Tổng hợp thông tin về Trạng thái, Mã, Mô tả (tiếp theo)}\\
\hline \textbf{Trạng thái} & \textbf{Mã} & \textbf{Mô tả} \\
\hline
\endhead
\hline
\endfoot
\endlastfoot
INIT & 0 & Khởi tạo hệ thống, cấu hình ngoại vi \\\hline
CHECK\_IGN & 1 & Kiểm tra trạng thái động cơ \\\hline
DRIVING & 2 & Chế độ lái xe - theo dõi liên tục \\\hline
PARKED & 3 & Chế độ đỗ xe - tiết kiệm năng lượng \\\hline
ALARM & 4 & Chế độ cảnh báo - phát hiện chuyển động bất thường \\\hline
HEARTBEAT & 5 & Gửi tín hiệu heartbeat \\\hline
SLEEP & 6 & Deep sleep - tiêu thụ năng lượng tối thiểu \\\hline
\end{longtable}
}


\textbf{Bảng chuyển đổi trạng thái:}

{\def\LTcaptype{none} % do not increment counter
\begin{longtable}{|L{0.265\linewidth}|L{0.432\linewidth}|L{0.267\linewidth}|}
\hline \textbf{Trạng thái hiện tại} & \textbf{Sự kiện kích hoạt} & \textbf{Trạng thái tiếp theo} \\
\hline
\endhead
\hline
\endfoot
\endlastfoot
INIT & Khởi tạo hoàn tất & CHECK\_IGN \\\hline
CHECK\_IGN & IGN = ON & DRIVING \\\hline
CHECK\_IGN & IGN = OFF & PARKED \\\hline
DRIVING & IGN = OFF được phát hiện & PARKED \\\hline
PARKED & IMU interrupt (motion detected) & ALARM \\\hline
PARKED & Timer wake-up & HEARTBEAT \\\hline
ALARM & Motion dừng, IGN = OFF & PARKED \\\hline
ALARM & IGN = ON & DRIVING \\\hline
HEARTBEAT & Gửi heartbeat xong & SLEEP \\\hline
SLEEP & Timer wake-up hoặc IMU interrupt & CHECK\_IGN \\\hline
\end{longtable}
}


{\thesissetfigureid{3.11}{fig-ch3-so-o-may-trang-thai-cua-thiet-bi-theo-doi}
\begin{figure}[htbp]
\centering
\pandocbounded{\safeincludesvg{./assets/figures/04-chuong-3-giai-phap-firmware-hinh-3-11.svg}}
\caption{Sơ đồ máy trạng thái của thiết bị theo dõi}
\label{fig:ch3-so-o-may-trang-thai-cua-thiet-bi-theo-doi}
\end{figure}
}


\begin{Shaded}
\begin{Highlighting}[]
\CommentTok{/* Cấu trúc máy trạng thái */}
\KeywordTok{typedef} \KeywordTok{enum} \OperatorTok{\{}
\NormalTok{    STATE\_INIT}\OperatorTok{,}
\NormalTok{    STATE\_CHECK\_IGN}\OperatorTok{,}
\NormalTok{    STATE\_DRIVING}\OperatorTok{,}
\NormalTok{    STATE\_PARKED}\OperatorTok{,}
\NormalTok{    STATE\_ALARM}\OperatorTok{,}
\NormalTok{    STATE\_HEARTBEAT}\OperatorTok{,}
\NormalTok{    STATE\_SLEEP}
\OperatorTok{\}}\NormalTok{ device\_state\_t}\OperatorTok{;}

\DataTypeTok{void}\NormalTok{ state\_machine\_task}\OperatorTok{(}\DataTypeTok{void} \OperatorTok{*}\NormalTok{param}\OperatorTok{)} \OperatorTok{\{}
\NormalTok{    device\_state\_t state }\OperatorTok{=}\NormalTok{ STATE\_INIT}\OperatorTok{;}

    \ControlFlowTok{while} \OperatorTok{(}\DecValTok{1}\OperatorTok{)} \OperatorTok{\{}
        \ControlFlowTok{switch} \OperatorTok{(}\NormalTok{state}\OperatorTok{)} \OperatorTok{\{}
        \ControlFlowTok{case}\NormalTok{ STATE\_INIT}\OperatorTok{:}
\NormalTok{            init\_peripherals}\OperatorTok{();}
\NormalTok{            state }\OperatorTok{=}\NormalTok{ STATE\_CHECK\_IGN}\OperatorTok{;}
            \ControlFlowTok{break}\OperatorTok{;}

        \ControlFlowTok{case}\NormalTok{ STATE\_CHECK\_IGN}\OperatorTok{:}
            \ControlFlowTok{if} \OperatorTok{(}\NormalTok{check\_ignition\_status}\OperatorTok{())} \OperatorTok{\{}
\NormalTok{                state }\OperatorTok{=}\NormalTok{ STATE\_DRIVING}\OperatorTok{;}
            \OperatorTok{\}} \ControlFlowTok{else} \OperatorTok{\{}
\NormalTok{                state }\OperatorTok{=}\NormalTok{ STATE\_PARKED}\OperatorTok{;}
            \OperatorTok{\}}
            \ControlFlowTok{break}\OperatorTok{;}

        \ControlFlowTok{case}\NormalTok{ STATE\_DRIVING}\OperatorTok{:}
\NormalTok{            start\_continuous\_tracking}\OperatorTok{();}
            \CommentTok{/* Chuyển sang PARKED khi phát hiện IGN OFF */}
\NormalTok{            EventBits\_t bits }\OperatorTok{=}\NormalTok{ xEventGroupWaitBits}\OperatorTok{(}
\NormalTok{                system\_event\_group}\OperatorTok{,}\NormalTok{ EVT\_IGN\_OFF}\OperatorTok{,}
\NormalTok{                pdTRUE}\OperatorTok{,}\NormalTok{ pdFALSE}\OperatorTok{,}\NormalTok{ pdMS\_TO\_TICKS}\OperatorTok{(}\DecValTok{1000}\OperatorTok{));}
            \ControlFlowTok{if} \OperatorTok{(}\NormalTok{bits }\OperatorTok{\&}\NormalTok{ EVT\_IGN\_OFF}\OperatorTok{)} \OperatorTok{\{}
\NormalTok{                stop\_ble\_connection}\OperatorTok{();}
\NormalTok{                state }\OperatorTok{=}\NormalTok{ STATE\_PARKED}\OperatorTok{;}
            \OperatorTok{\}}
            \ControlFlowTok{break}\OperatorTok{;}

        \ControlFlowTok{case}\NormalTok{ STATE\_PARKED}\OperatorTok{:}
\NormalTok{            prepare\_for\_sleep}\OperatorTok{();}
\NormalTok{            state }\OperatorTok{=}\NormalTok{ STATE\_SLEEP}\OperatorTok{;}
            \ControlFlowTok{break}\OperatorTok{;}

        \ControlFlowTok{case}\NormalTok{ STATE\_ALARM}\OperatorTok{:}
\NormalTok{            send\_alert\_immediate}\OperatorTok{();}
\NormalTok{            start\_continuous\_tracking}\OperatorTok{();}
            \CommentTok{/* Đợi cho đến khi hết chuyển động */}
            \ControlFlowTok{if} \OperatorTok{(!}\NormalTok{is\_motion\_detected}\OperatorTok{()} \OperatorTok{\&\&} \OperatorTok{!}\NormalTok{check\_ignition\_status}\OperatorTok{())} \OperatorTok{\{}
\NormalTok{                state }\OperatorTok{=}\NormalTok{ STATE\_PARKED}\OperatorTok{;}
            \OperatorTok{\}}
            \ControlFlowTok{break}\OperatorTok{;}

        \ControlFlowTok{case}\NormalTok{ STATE\_HEARTBEAT}\OperatorTok{:}
\NormalTok{            wakeup\_modem}\OperatorTok{();}
\NormalTok{            tracker\_mqtt\_connect}\OperatorTok{();}
\NormalTok{            send\_heartbeat}\OperatorTok{();}
\NormalTok{            state }\OperatorTok{=}\NormalTok{ STATE\_SLEEP}\OperatorTok{;}
            \ControlFlowTok{break}\OperatorTok{;}

        \ControlFlowTok{case}\NormalTok{ STATE\_SLEEP}\OperatorTok{:}
\NormalTok{            shutdown\_modem}\OperatorTok{();}
\NormalTok{            configure\_wakeup\_sources}\OperatorTok{();}
\NormalTok{            enter\_deep\_sleep}\OperatorTok{();}
            \CommentTok{/* Sau khi wake{-}up, bắt đầu lại từ CHECK\_IGN */}
\NormalTok{            state }\OperatorTok{=}\NormalTok{ STATE\_CHECK\_IGN}\OperatorTok{;}
            \ControlFlowTok{break}\OperatorTok{;}
        \OperatorTok{\}}
    \OperatorTok{\}}
\OperatorTok{\}}
\end{Highlighting}
\end{Shaded}

\textbf{Lưu trạng thái qua deep sleep}: Trước khi vào deep sleep,
firmware lưu trạng thái hiện tại vào vùng nhớ RTC (RTC memory) --- vùng
nhớ đặc biệt của ESP32-S3 được giữ nguyên nội dung trong suốt quá trình
deep sleep. Khi wake-up, firmware đọc trạng thái từ RTC memory để biết
chế độ hoạt động trước đó và xử lý phù hợp.

\paragraph{3.2.2.6. Cấu hình hệ
thống}\label{3226-cux1ea5u-huxecnh-hux1ec7-thux1ed1ng}

\subparagraph{a) Cấu trúc dữ liệu cấu
hình}\label{a-cux1ea5u-truxfac-dux1eef-liux1ec7u-cux1ea5u-huxecnh}

Toàn bộ cấu hình của thiết bị được lưu trữ trong bộ nhớ flash (NVS ---
Non-Volatile Storage) của ESP32-S3 dưới dạng cấu trúc dữ liệu cố định:

\begin{Shaded}
\begin{Highlighting}[]
\KeywordTok{typedef} \KeywordTok{struct} \OperatorTok{\{}
    \DataTypeTok{char}\NormalTok{     device\_id}\OperatorTok{[}\DecValTok{16}\OperatorTok{];}          \CommentTok{/* Mã định danh thiết bị */}
    \DataTypeTok{char}\NormalTok{     mqtt\_broker}\OperatorTok{[}\DecValTok{64}\OperatorTok{];}        \CommentTok{/* Địa chỉ MQTT broker */}
    \DataTypeTok{uint16\_t}\NormalTok{ mqtt\_port}\OperatorTok{;}              \CommentTok{/* Cổng MQTT (mặc định 1883) */}
    \DataTypeTok{char}\NormalTok{     mqtt\_username}\OperatorTok{[}\DecValTok{32}\OperatorTok{];}      \CommentTok{/* Tên đăng nhập MQTT */}
    \DataTypeTok{char}\NormalTok{     mqtt\_password}\OperatorTok{[}\DecValTok{32}\OperatorTok{];}      \CommentTok{/* Mật khẩu MQTT */}
    \DataTypeTok{uint16\_t}\NormalTok{ heartbeat\_interval}\OperatorTok{;}     \CommentTok{/* Chu kỳ gửi heartbeat (giây) */}
    \DataTypeTok{uint16\_t}\NormalTok{ tracking\_interval}\OperatorTok{;}      \CommentTok{/* Chu kỳ gửi telemetry (giây) */}
    \DataTypeTok{char}\NormalTok{     obd2\_ble\_address}\OperatorTok{[}\DecValTok{18}\OperatorTok{];}   \CommentTok{/* Địa chỉ MAC của OBD2 adapter */}
    \DataTypeTok{float}\NormalTok{    switch\_off}\OperatorTok{;}             \CommentTok{/* Ngưỡng chuyển sang pin backup: 12.0V (12V) hoặc 24.0V (24V) */}
    \DataTypeTok{float}\NormalTok{    switch\_on}\OperatorTok{;}              \CommentTok{/* Ngưỡng quay lại ắc quy: 12.2V (12V) hoặc 24.4V (24V) */}
    \DataTypeTok{float}\NormalTok{    ign\_on\_threshold}\OperatorTok{;}       \CommentTok{/* Ngưỡng suy luận IGN ON: \textgreater{}=13.0V (12V) hoặc \textgreater{}=26.0V (24V) */}
    \DataTypeTok{float}\NormalTok{    ign\_off\_threshold}\OperatorTok{;}      \CommentTok{/* Ngưỡng suy luận IGN OFF: \textless{}=12.0V (12V) hoặc \textless{}=24.0V (24V) */}
\OperatorTok{\}}\NormalTok{ config\_t}\OperatorTok{;}
\end{Highlighting}
\end{Shaded}

\subparagraph{b) Quản lý cấu
hình}\label{b-quux1ea3n-luxfd-cux1ea5u-huxecnh}

Cấu hình được quản lý theo ba cơ chế:

\textbf{Cấu hình mặc định (Default Configuration)}: Được định nghĩa
trong mã nguồn firmware, áp dụng khi thiết bị khởi động lần đầu hoặc khi
NVS bị xóa. Bao gồm các giá trị an toàn như
\texttt{heartbeat\_interval\ =\ 900} (15 phút),
\texttt{tracking\_interval\ =\ 10} (10 giây), và profile nguồn kép:
profile 12V (\texttt{switch\_off=12.0}, \texttt{switch\_on=12.2},
\texttt{ign\_on\_threshold=13.0}, \texttt{ign\_off\_threshold=12.0})
hoặc profile 24V (\texttt{switch\_off=24.0}, \texttt{switch\_on=24.4},
\texttt{ign\_on\_threshold=26.0}, \texttt{ign\_off\_threshold=24.0}).

\textbf{Cấu hình lưu trữ (Persistent Configuration)}: Lưu trong NVS,
được tải khi thiết bị khởi động. Các thay đổi cấu hình từ máy chủ được
lưu vào NVS để giữ nguyên sau khi reset hoặc mất điện.

\textbf{Cập nhật cấu hình từ xa (Remote Configuration Update)}: Máy chủ
gửi lệnh \texttt{update\_config} qua MQTT topic
\texttt{v1/\{device\_id\}/commands}. Firmware phân tích lệnh, cập nhật
các tham số tương ứng và lưu vào NVS. Thiết bị gửi xác nhận (ACK) về máy
chủ sau khi cập nhật thành công.

\subparagraph{c) Hiệu chuẩn cảm biến (Sensor
Calibration)}\label{c-hiux1ec7u-chuux1ea9n-cux1ea3m-biux1ebfn-sensor-calibration}

Để đảm bảo độ chính xác của phép đo, firmware hỗ trợ hiệu chuẩn ba loại
cảm biến:


{\thesissettableid{3.2.2K}{tab-ch3-tong-hop-thong-tin-ve-cam-bien-phuong-phap-hieu-chuan-luu}
\def\LTcaptype{table}
\begin{longtable}{|L{0.224\linewidth}|L{0.580\linewidth}|L{0.160\linewidth}|}
\caption{Tổng hợp thông tin về Cảm biến, Phương pháp hiệu chuẩn, Lưu trữ}\label{tab:ch3-tong-hop-thong-tin-ve-cam-bien-phuong-phap-hieu-chuan-luu}\\
\hline \textbf{Cảm biến} & \textbf{Phương pháp hiệu chuẩn} & \textbf{Lưu trữ} \\
\hline
\endfirsthead
\caption[]{Tổng hợp thông tin về Cảm biến, Phương pháp hiệu chuẩn, Lưu trữ (tiếp theo)}\\
\hline \textbf{Cảm biến} & \textbf{Phương pháp hiệu chuẩn} & \textbf{Lưu trữ} \\
\hline
\endhead
\hline
\endfoot
\endlastfoot
ADC (điện áp ắc quy) & So sánh với đồng hồ vạn năng, tính hệ số hiệu
chỉnh & NVS \\\hline
IMU (LIS3DSH) & Đặt thiết bị trên mặt phẳng, đo offset các trục X/Y/Z &
NVS \\\hline
GNSS & Hiệu chỉnh offset vị trí (nếu cần) & NVS \\\hline
\end{longtable}
}


Các giá trị hiệu chuẩn được lưu trữ trong NVS và được tải khi thiết bị
khởi động, đảm bảo tính nhất quán của phép đo giữa các lần reset.

\subsection{3.3. Phân tích và đề xuất kiến trúc Backend/Cloud --
Backend/cloud analysis and solution}\label{33-phan-tich-va-de-xuat-kien-truc-backend-cloud}

\subsubsection{3.3.1. Phân tích và lựa chọn kiến trúc Cloud}

\paragraph{3.3.1.1. Đặt vấn đề cho kiến trúc
Cloud}\label{3131-ux111ux1eb7t-vux1ea5n-ux111ux1ec1-cho-kiux1ebfn-truxfac-cloud}

Khi dữ liệu rời khỏi thiết bị, trọng tâm không còn là ``đọc được dữ
liệu'' mà là ``biến bản tin đó thành thông tin hữu ích''. Ở tầng
Backend/Cloud, một bản tin chỉ thật sự có giá trị khi nó được tiếp nhận
ổn định, lưu đúng nơi, đẩy tới dashboard đúng lúc và vẫn truy vết được
khi có sự cố. Vì vậy, kiến trúc cloud phải đồng thời giải quyết các bài
toán sau:

\begin{itemize}
\tightlist
\item
  \textbf{Tiếp nhận dữ liệu liên tục từ nhiều thiết bị}: telemetry gửi
  theo chu kỳ ngắn, có thể xuất hiện burst khi thiết bị reconnect.
\item
  \textbf{Cập nhật thời gian thực cho dashboard}: dữ liệu vị trí/trạng
  thái cần phản ánh gần như tức thời cho người vận hành.
\item
  \textbf{Lưu trữ dữ liệu hỗn hợp}: vừa có dữ liệu chuỗi thời gian (GPS,
  nguồn, IMU; sẵn sàng mở rộng thêm OBD2), vừa có dữ liệu quan hệ nghiệp
  vụ (xe, người dùng, cảnh báo, geofence).
\item
  \textbf{Chi phí và năng lực vận hành phù hợp đồ án}: ưu tiên
  self-hosted, dễ triển khai, dễ mở rộng theo từng dịch vụ.
\end{itemize}

\paragraph{3.3.1.2. So sánh các phương án kiến trúc
Cloud}\label{3132-so-suxe1nh-cuxe1c-phux1b0ux1a1ng-uxe1n-kiux1ebfn-truxfac-cloud}


{\thesissettableid{3.15A}{tab-ch3-so-sanh-cac-phuong-an-kien-truc-cloud}
\def\LTcaptype{table}
\begin{longtable}{|L{0.145\linewidth}|L{0.230\linewidth}|L{0.220\linewidth}|L{0.250\linewidth}|L{0.100\linewidth}|}
\caption{So sánh các phương án kiến trúc Cloud}\label{tab:ch3-so-sanh-cac-phuong-an-kien-truc-cloud}\\
\hline \textbf{Phương án} & \textbf{Mô tả ngắn} & \textbf{Ưu điểm} & \textbf{Hạn chế} & \textbf{Mức phù hợp} \\
\hline
\endfirsthead
\caption[]{So sánh các phương án kiến trúc Cloud (tiếp theo)}\\
\hline \textbf{Phương án} & \textbf{Mô tả ngắn} & \textbf{Ưu điểm} & \textbf{Hạn chế} & \textbf{Mức phù hợp} \\
\hline
\endhead
\hline
\endfoot
\endlastfoot
\textbf{PA-A: Monolithic REST + 1 DB} & Thiết bị gửi HTTP trực tiếp vào
API server, lưu chung một cơ sở dữ liệu & Triển khai ban đầu đơn giản &
Không tối ưu cho telemetry tần suất cao, khó mở rộng realtime, tải dồn
lên API & Trung bình \\\hline
\textbf{PA-B: Managed Cloud Native} & AWS IoT Core + Lambda + dịch vụ DB
managed & Khả năng mở rộng cao, nhiều dịch vụ sẵn có & Chi phí vận hành
cao, phụ thuộc nhà cung cấp, độ phức tạp hạ tầng vượt phạm vi đồ án &
Trung bình \\\hline
\textbf{PA-C: Event-driven self-hosted} & MQTT Broker + MQTT Bridge +
dual database + API + WebSocket & Phù hợp IoT realtime, tách tải tốt, mở
rộng theo dịch vụ, chi phí tự chủ & Cần chuẩn hóa message flow và quản
lý nhiều dịch vụ Docker & \textbf{Cao (Đã chọn)} \\\hline
\end{longtable}
}


\paragraph{3.3.1.3. Chọn giải pháp kiến trúc
Cloud}\label{3133-chux1ecdn-giux1ea3i-phuxe1p-kiux1ebfn-truxfac-cloud}

Đồ án chọn \textbf{PA-C: kiến trúc phân tầng kết hợp hướng sự kiện
(event-driven)} với chuỗi xử lý:

\texttt{Device\ -\textgreater{}\ EMQX\ -\textgreater{}\ MQTT\ Bridge\ -\textgreater{}\ PostgreSQL/}\\
\texttt{VictoriaMetrics/VictoriaLogs\ -\textgreater{}\ Backend\ API\ -\textgreater{}\ Frontend}

Lựa chọn này xuất phát từ cách dữ liệu thực sự di chuyển trong hệ thống:
thiết bị gửi bản tin liên tục, lớp ingest tiếp nhận và phân luồng, lớp
lưu trữ ghi theo đúng đặc tính dữ liệu, còn lớp ứng dụng phục vụ người
dùng và nghiệp vụ. Nói cách khác, kiến trúc được chọn vì nó bám sát
hành trình tự nhiên của dữ liệu thay vì dồn mọi trách nhiệm về một dịch
vụ duy nhất.

Các lý do lựa chọn chính gồm:

\begin{itemize}
\tightlist
\item
  \textbf{Khớp bản chất dữ liệu IoT}: MQTT xử lý tốt mô hình pub/sub,
  giảm mức phụ thuộc trực tiếp giữa thiết bị và tầng ứng dụng. Điều này
  giúp thiết bị chỉ tập trung gửi dữ liệu, thay vì phải biết chi tiết ai
  sẽ nhận và xử lý dữ liệu đó ở phía sau.
\item
  \textbf{Tối ưu hiệu năng theo vai trò}: Bridge xử lý ingestion, API
  tập trung nghiệp vụ, database tách theo loại dữ liệu. Mỗi lớp chỉ làm
  đúng việc mình mạnh nhất, nhờ đó hệ thống dễ giữ ổn định hơn khi tải
  tăng.
\item
  \textbf{Dễ scale theo chiều ngang}: có thể nhân bản riêng
  broker/bridge/api khi số thiết bị tăng, thay vì phải mở rộng toàn bộ
  hệ thống một cách đồng loạt và tốn kém.
\item
  \textbf{Phù hợp chi phí và phạm vi thực hiện}: self-hosted qua Docker
  Compose, không phụ thuộc nền tảng cloud thương mại, nên phù hợp với
  tính chất của một đồ án vừa cần triển khai thật, vừa cần kiểm soát tốt
  nguồn lực.
\end{itemize}

\subsubsection{3.3.2. Giải pháp Backend \&
Cloud}\label{323-giux1ea3i-phuxe1p-backend--cloud}

Sau khi xác định phương án kiến trúc phù hợp, phần này đi vào cách triển
khai cụ thể từng lớp của hệ thống Backend/Cloud. Cách trình bày được sắp
xếp theo hành trình dữ liệu, từ lúc bản tin đi vào broker cho tới khi
thông tin xuất hiện trên dashboard của người vận hành.

\paragraph{3.3.2.1. Kiến trúc tổng quan hệ thống Cloud (Cloud
Architecture
Overview)}\label{3231-kiux1ebfn-truxfac-tux1ed5ng-quan-hux1ec7-thux1ed1ng-cloud-cloud-architecture-overview}

\subparagraph{a) Mô hình kiến trúc tổng
thể}\label{a-muxf4-huxecnh-kiux1ebfn-truxfac-tux1ed5ng-thux1ec3}

Để người đọc dễ theo dõi, có thể hình dung hệ thống cloud như một chuỗi
xử lý nhiều lớp: lớp đầu nhận dữ liệu, lớp giữa lưu và xử lý, lớp cuối
cung cấp thông tin cho người dùng. Kiến trúc này vừa tận dụng được ưu
thế của mô hình hướng sự kiện cho IoT, vừa giữ cho trách nhiệm của từng
thành phần không bị chồng lấn. Cụ thể, hệ thống gồm các tầng chính sau:

\begin{itemize}
\tightlist
\item
  \textbf{Tầng thu thập dữ liệu (Data Ingestion Layer)}: Tiếp nhận dữ
  liệu từ các thiết bị IoT thông qua giao thức MQTT, xử lý và phân phối
  đến các thành phần lưu trữ.
\item
  \textbf{Tầng lưu trữ dữ liệu (Data Storage Layer)}: Sử dụng chiến lược
  lưu trữ kép (dual database strategy) với cơ sở dữ liệu quan hệ và cơ
  sở dữ liệu chuỗi thời gian.
\item
  \textbf{Tầng xử lý nghiệp vụ (Business Logic Layer)}: API Server xử lý
  các yêu cầu từ Frontend, thực thi logic nghiệp vụ và quản lý trạng
  thái hệ thống.
\item
  \textbf{Tầng giao tiếp thời gian thực (Real-time Communication
  Layer)}: Cung cấp dữ liệu cập nhật trực tiếp đến người dùng thông qua
  giao thức WebSocket.
\end{itemize}


{\thesissetfigureid{3.12}{fig-ch3-so-o-kien-truc-tong-quan-he-thong-cloud}
\begin{figure}[htbp]
\centering
\pandocbounded{\safeincludesvg{./assets/figures/05-chuong-3-giai-phap-backend-hinh-3-12.svg}}
\caption{Sơ đồ kiến trúc tổng quan hệ thống Cloud}
\label{fig:ch3-so-o-kien-truc-tong-quan-he-thong-cloud}
\end{figure}
}



{\thesissetfigureid{3.12a}{fig-ch3-luong-du-lieu-chi-tiet-tu-thiet-bi-en-dashboard}
\begin{figure}[htbp]
\centering
\pandocbounded{\safeincludesvg{./assets/figures/05-chuong-3-giai-phap-backend-hinh-3-12a.svg}}
\caption{Luồng dữ liệu chi tiết từ thiết bị đến dashboard}
\label{fig:ch3-luong-du-lieu-chi-tiet-tu-thiet-bi-en-dashboard}
\end{figure}
}


\subparagraph{b) Luồng dữ liệu chính trong hệ
thống}\label{b-luux1ed3ng-dux1eef-liux1ec7u-chuxednh-trong-hux1ec7-thux1ed1ng}

Nếu mô tả bằng một ``câu chuyện dữ liệu'', luồng chính của hệ thống diễn
ra theo trình tự sau:

\begin{enumerate}
\def\labelenumi{\arabic{enumi}.}
\tightlist
\item
  \textbf{Thiết bị IoT} gửi dữ liệu telemetry chính (GNSS, trạng thái
  nguồn, mức rung, ignition, uptime) lên EMQX Broker thông qua giao thức
  MQTT 5.0 theo topic có cấu trúc \texttt{v1/\{device\_id\}/rawdata}.
\item
  \textbf{EMQX Broker} tiếp nhận và phân phối message đến các
  subscriber. Rules Engine của EMQX chủ yếu đóng vai trò lọc/bổ trợ;
  phần lớn nghiệp vụ cảnh báo vẫn được xử lý tại MQTT Bridge.
\item
  \textbf{MQTT Bridge} (dịch vụ Node.js độc lập) subscribe các topic từ
  EMQX, xác thực dữ liệu đầu vào bằng Zod schema, sau đó thực hiện ghi
  kép (dual write):

  \begin{itemize}
  \tightlist
  \item
    Ghi dữ liệu chuỗi thời gian vào \textbf{VictoriaMetrics} (vị trí,
    tốc độ, điện áp, rung, ignition, uptime).
  \item
    Ghi nhật ký sự kiện vào \textbf{VictoriaLogs} (kết nối/ngắt kết nối,
    lỗi, sự kiện MQTT).
  \item
    Cập nhật trạng thái thiết bị vào \textbf{PostgreSQL} (trạng thái
    trực tuyến, phiên làm việc).
  \item
    Phát sự kiện đến \textbf{Socket.IO} để cập nhật thời gian thực cho
    Dashboard.
  \end{itemize}
\item
  \textbf{API Server} (Express.js) cung cấp REST API cho Frontend, truy
  vấn dữ liệu từ PostgreSQL và VictoriaMetrics, đồng thời gửi lệnh điều
  khiển ngược về thiết bị thông qua EMQX.
\item
  \textbf{Frontend} (Next.js) nhận dữ liệu cập nhật thời gian thực qua
  Socket.IO và dữ liệu lịch sử qua REST API.
\end{enumerate}

\subparagraph{c) Mô hình triển khai
Docker}\label{c-muxf4-huxecnh-triux1ec3n-khai-docker}

Ở mức triển khai, kiến trúc trên được hiện thực bằng mô hình per-service
Docker Compose theo quy ước IVM26. Mỗi dịch vụ có một file
\texttt{docker-compose.yml} riêng và cùng tham gia mạng Docker
\texttt{tracking-network} (external network). Cách tổ chức này giúp cấu
trúc vận hành của đồ án bám sát cách các hệ thống thật được chia dịch vụ
trong thực tế:

\begin{itemize}
\tightlist
\item
  \textbf{Độc lập khi triển khai}: Mỗi dịch vụ có thể khởi động, dừng
  hoặc cập nhật độc lập mà không ảnh hưởng đến các dịch vụ khác.
\item
  \textbf{Quản lý tài nguyên riêng}: Mỗi container được giới hạn CPU và
  RAM cụ thể (ví dụ: Backend 512MB/1 CPU, MQTT Bridge 256MB/0.5 CPU).
\item
  \textbf{Dễ dàng mở rộng}: Có thể nhân bản (scale) từng dịch vụ riêng
  lẻ theo nhu cầu.
\end{itemize}


{\thesissettableid{3.16}{tab-ch3-phan-bo-tai-nguyen-docker-cho-cac-dich-vu}
\def\LTcaptype{table}
\begin{longtable}{|L{0.326\linewidth}|L{0.327\linewidth}|L{0.128\linewidth}|L{0.173\linewidth}|}
\caption{Phân bổ tài nguyên Docker cho các dịch vụ}\label{tab:ch3-phan-bo-tai-nguyen-docker-cho-cac-dich-vu}\\
\hline \textbf{Dịch vụ} & \textbf{Container} & \textbf{Port} & \textbf{Tài nguyên} \\
\hline
\endfirsthead
\caption[]{Phân bổ tài nguyên Docker cho các dịch vụ (tiếp theo)}\\
\hline \textbf{Dịch vụ} & \textbf{Container} & \textbf{Port} & \textbf{Tài nguyên} \\
\hline
\endhead
\hline
\endfoot
\endlastfoot
Tracking\_Backend & tracking-backend & 4000 & 512MB, 1 CPU \\\hline
Tracking\_Frontend & tracking-frontend & 4001 & 512MB, 1 CPU \\\hline
Tracking\_MqttBridge & tracking-mqtt-bridge & - & 256MB, 0.5 CPU \\\hline
Tracking\_PostgreSQL & tracking-postgres & 5432 & 1GB, 2 CPU \\\hline
Tracking\_EMQX & tracking-emqx & 1883, 8883, 8083, 8084, 18083 & 1GB, 2
CPU \\\hline
Tracking\_VictoriaMetrics & tracking-victoriametrics & 8428 & 512MB, 1
CPU \\\hline
Tracking\_VictoriaLogs & tracking-victorialogs & 9428 & 512MB, 1 CPU \\\hline
Tracking\_Grafana & tracking-grafana & 4002 & Chưa cấu hình cứng \\\hline
\end{longtable}
}


\paragraph{3.3.2.2. MQTT Broker - EMQX (Message Broker Selection \&
Configuration)}\label{3232-mqtt-broker---emqx-message-broker-selection--configuration}

\subparagraph{a) Vai trò của MQTT
Broker}\label{a-vai-truxf2-cux1ee7a-mqtt-broker}

MQTT Broker là điểm vào đầu tiên của dữ liệu ở phía cloud. Vai trò của
nó có thể hiểu đơn giản là một ``trạm trung chuyển tin nhắn'': thiết bị
gửi bản tin đến broker, broker tiếp nhận và phân phối cho các thành phần
cần xử lý tiếp theo. Nhờ mô hình \textbf{Publish-Subscribe (Pub/Sub)},
các thành phần trong hệ thống có thể giao tiếp bất đồng bộ mà không cần
ghép chặt với nhau.

Trong hệ thống theo dõi phương tiện, cách tổ chức này đặc biệt quan
trọng vì số lượng thiết bị có thể tăng dần theo thời gian, trong khi
nhịp gửi dữ liệu lại thay đổi theo chế độ hoạt động của từng xe.

\subparagraph{b) So sánh và lựa chọn MQTT
Broker}\label{b-so-suxe1nh-vuxe0-lux1ef1a-chux1ecdn-mqtt-broker}

Để lựa chọn MQTT Broker phù hợp, đồ án đã đánh giá bốn giải pháp phổ
biến trên thị trường:


{\thesissettableid{3.17}{tab-ch3-so-sanh-cac-mqtt-broker-pho-bien}
\def\LTcaptype{table}
\begin{longtable}{|L{0.269\linewidth}|L{0.188\linewidth}|L{0.158\linewidth}|L{0.152\linewidth}|L{0.177\linewidth}|}
\caption{So sánh các MQTT Broker phổ biến}\label{tab:ch3-so-sanh-cac-mqtt-broker-pho-bien}\\
\hline \textbf{Tiêu chí} & \textbf{Mosquitto} & \textbf{EMQX} & \textbf{HiveMQ} & \textbf{VerneMQ} \\
\hline
\endfirsthead
\caption[]{So sánh các MQTT Broker phổ biến (tiếp theo)}\\
\hline \textbf{Tiêu chí} & \textbf{Mosquitto} & \textbf{EMQX} & \textbf{HiveMQ} & \textbf{VerneMQ} \\
\hline
\endhead
\hline
\endfoot
\endlastfoot
\textbf{Kết nối đồng thời} & ~1.000 &
~100.000.000 & ~200.000.000 &
~10.000.000 \\\hline
\textbf{Thông lượng message} & 40K msg/s & 100K msg/s & 200K msg/s & 50K
msg/s \\\hline
\textbf{Độ trễ (Latency)} & 0,25 ms & 0,27 ms & \textless1 ms & 2,1
ms \\\hline
\textbf{Sử dụng CPU} & Rất thấp & Thấp (2\%) & Thấp & Cao (10\%) \\\hline
\textbf{Sử dụng RAM} & 254 MB & 495 MB & Trung bình & 1,2 GB \\\hline
\textbf{Hỗ trợ Clustering} & Không & Có (20+ node) & Có & Có \\\hline
\textbf{Khả dụng cao (HA)} & Không & Có & Có & Có \\\hline
\textbf{Độ phức tạp cài đặt} & Rất dễ & Trung bình & Khó & Trung bình \\\hline
\textbf{Chi phí} & Miễn phí & Miễn phí/Trả phí & Trả phí & Miễn phí/Trả
phí \\\hline
\textbf{Cộng đồng} & Lớn & Lớn & Trung bình & Nhỏ \\\hline
\end{longtable}
}


\textbf{Kết luận lựa chọn: EMQX Single Node}

EMQX được lựa chọn làm MQTT Broker cho hệ thống vì đây là phương án giữ
được thế cân bằng tốt giữa hiệu năng, khả năng mở rộng và mức độ phù hợp
với nguồn lực triển khai của đồ án. Các lý do chính gồm:

\begin{itemize}
\tightlist
\item
  \textbf{Hiệu năng phù hợp}: Theo tài liệu công bố của EMQX, broker đáp
  ứng tốt tải message lớn với độ trễ thấp, đủ cho quy mô triển khai hiện
  tại và còn dư địa mở rộng.
\item
  \textbf{Khả năng mở rộng}: Hỗ trợ clustering nhiều node, thuận lợi nếu
  hệ thống tăng số lượng thiết bị trong giai đoạn sau.
\item
  \textbf{Rules Engine tích hợp}: Thuận tiện cho các tác vụ lọc hoặc
  chuyển đổi dữ liệu đơn giản ngay tại broker; tuy nhiên trong kiến trúc
  hiện tại, đây không phải lớp xử lý nghiệp vụ chính.
\item
  \textbf{Miễn phí và cộng đồng lớn}: Phiên bản Community cung cấp đầy
  đủ tính năng cần thiết cho đồ án, có cộng đồng hỗ trợ tích cực.
\item
  \textbf{Dashboard quản lý}: Cung cấp giao diện web quản lý tại port
  18083, hỗ trợ giám sát kết nối, topic và message theo thời gian thực.
\end{itemize}

\subparagraph{c) Cấu hình EMQX Rules
Engine}\label{c-cux1ea5u-huxecnh-emqx-rules-engine}

EMQX Rules Engine là lớp xử lý dữ liệu nhẹ ngay tại broker. Trong hệ
thống hiện tại, phần lớn nghiệp vụ cảnh báo được xử lý ở MQTT Bridge;
tuy vậy, Rules Engine vẫn phù hợp để lọc nhanh các điều kiện đơn giản
hoặc đẩy dữ liệu sang pipeline phụ mà không làm tăng tải cho backend.

\textbf{Các rule minh họa phù hợp với payload hiện tại:}

\textbf{Rule 1 - Cảnh báo rung mạnh từ luồng sự kiện nội bộ:}

\begin{Shaded}
\begin{Highlighting}[]
\KeywordTok{SELECT}
\NormalTok{  payload.device\_id }\KeywordTok{as}\NormalTok{ device\_id,}
\NormalTok{  payload.latitude }\KeywordTok{as}\NormalTok{ latitude,}
\NormalTok{  payload.longitude }\KeywordTok{as}\NormalTok{ longitude,}
\NormalTok{  payload.}\FunctionTok{value} \KeywordTok{as}\NormalTok{ vibration,}
\NormalTok{  now() }\KeywordTok{as} \DataTypeTok{timestamp}
\KeywordTok{FROM} \OtherTok{"internal/events/device/alert"}
\KeywordTok{WHERE}\NormalTok{ payload.alert\_type }\OperatorTok{=} \StringTok{\textquotesingle{}high\_vibration\textquotesingle{}}
\end{Highlighting}
\end{Shaded}

\textbf{Rule 2 - Cảnh báo vượt tốc độ từ rawdata:}

\begin{Shaded}
\begin{Highlighting}[]
\KeywordTok{SELECT}
\NormalTok{  payload.device\_id }\KeywordTok{as}\NormalTok{ device\_id,}
\NormalTok{  payload.}\KeywordTok{data}\NormalTok{.speed }\KeywordTok{as}\NormalTok{ speed,}
\NormalTok{  payload.}\KeywordTok{data}\NormalTok{.latitude }\KeywordTok{as}\NormalTok{ latitude,}
\NormalTok{  payload.}\KeywordTok{data}\NormalTok{.longitude }\KeywordTok{as}\NormalTok{ longitude,}
\NormalTok{  now() }\KeywordTok{as} \DataTypeTok{timestamp}
\KeywordTok{FROM} \OtherTok{"v1/+/rawdata"}
\KeywordTok{WHERE}\NormalTok{ payload.}\KeywordTok{data}\NormalTok{.speed }\OperatorTok{\textgreater{}} \DecValTok{100}
\end{Highlighting}
\end{Shaded}

\textbf{Rule 3 - Cảnh báo pin dự phòng thấp:}

\begin{Shaded}
\begin{Highlighting}[]
\KeywordTok{SELECT}
\NormalTok{  payload.device\_id }\KeywordTok{as}\NormalTok{ device\_id,}
\NormalTok{  payload.}\KeywordTok{data}\NormalTok{.battery\_bot }\KeywordTok{as}\NormalTok{ battery\_bot,}
\NormalTok{  payload.}\KeywordTok{data}\NormalTok{.latitude }\KeywordTok{as}\NormalTok{ latitude,}
\NormalTok{  payload.}\KeywordTok{data}\NormalTok{.longitude }\KeywordTok{as}\NormalTok{ longitude,}
\NormalTok{  now() }\KeywordTok{as} \DataTypeTok{timestamp}
\KeywordTok{FROM} \OtherTok{"v1/+/rawdata"}
\KeywordTok{WHERE}\NormalTok{ payload.}\KeywordTok{data}\NormalTok{.battery\_bot }\OperatorTok{\textless{}} \FloatTok{3.5}
\end{Highlighting}
\end{Shaded}

\textbf{Luồng xử lý cảnh báo:}

\begin{enumerate}
\tightlist
\item
  Thiết bị phát bản tin lên EMQX qua các topic telemetry/sự kiện.
\item
  MQTT Bridge tiếp nhận, đánh giá rule và phát sự kiện vào
  \texttt{internal/events/*}.
\item
  Backend listener xử lý sự kiện, ghi PostgreSQL và phát Socket.IO để cập
  nhật dashboard theo thời gian thực.
\end{enumerate}

Ngoài ra, hệ thống có thể mở rộng thêm rule cho geofence sơ cấp, thống
kê traffic theo topic, hoặc phát hiện thiết bị im lặng bất thường để
phục vụ vận hành.

\paragraph{3.3.2.3. Thiết kế cơ sở dữ liệu (Database Architecture
Design)}\label{3233-thiux1ebft-kux1ebf-cux1a1-sux1edf-dux1eef-liux1ec7u-database-architecture-design}

\textbf{Đặt vấn đề lưu trữ:} Khi đi sâu hơn vào tầng dữ liệu, có thể thấy
hệ thống đang phục vụ hai nhu cầu rất khác nhau. Một mặt, telemetry cần
được ghi nhanh và truy vấn theo thời gian. Mặt khác, dữ liệu nghiệp vụ
đòi hỏi tính toàn vẹn, quan hệ rõ ràng và khả năng quản lý dài hạn. Nếu
dùng một công cụ cho cả hai loại bài toán, hệ thống thường sẽ phải đánh
đổi hiệu quả ở một trong hai phía.


{\thesissettableid{3.18A}{tab-ch3-so-sanh-cac-phuong-an-kien-truc-luu-tru-du-lieu}
\def\LTcaptype{table}
\begin{longtable}{|L{0.149\linewidth}|L{0.233\linewidth}|L{0.236\linewidth}|L{0.241\linewidth}|L{0.085\linewidth}|}
\caption{So sánh các phương án kiến trúc lưu trữ dữ liệu}\label{tab:ch3-so-sanh-cac-phuong-an-kien-truc-luu-tru-du-lieu}\\
\hline \textbf{Phương án} & \textbf{Mô tả} & \textbf{Ưu điểm} & \textbf{Hạn chế} & \textbf{Mức phù hợp} \\
\hline
\endfirsthead
\caption[]{So sánh các phương án kiến trúc lưu trữ dữ liệu (tiếp theo)}\\
\hline \textbf{Phương án} & \textbf{Mô tả} & \textbf{Ưu điểm} & \textbf{Hạn chế} & \textbf{Mức phù hợp} \\
\hline
\endhead
\hline
\endfoot
\endlastfoot
\textbf{PA-DB1: PostgreSQL duy nhất} & Dùng PostgreSQL cho cả nghiệp vụ
và telemetry & Đơn giản vận hành & Ghi telemetry lớn dễ ảnh hưởng truy
vấn nghiệp vụ, chi phí index/partition cao & Trung bình \\\hline
\textbf{PA-DB2: Time-series DB duy nhất} & Dùng CSDL chuỗi thời gian cho
mọi loại dữ liệu & Tối ưu ghi dữ liệu cảm biến & Kém phù hợp dữ liệu
quan hệ phức tạp, khó đảm bảo ràng buộc nghiệp vụ & Thấp \\\hline
\textbf{PA-DB3: Hybrid DB (Đã chọn)} & PostgreSQL +\newline Victoria\-Metrics,\newline
Victoria\-Logs & Mỗi loại dữ liệu dùng đúng công cụ, cân bằng hiệu năng và
khả năng truy vấn & Tăng số thành phần cần vận hành & \textbf{Cao} \\\hline
\end{longtable}
}


\textbf{Kết luận lựa chọn:} Chọn \textbf{PA-DB3 (hybrid database)} để
tách trách nhiệm dữ liệu, giữ ổn định cho nghiệp vụ và tối ưu luồng
telemetry theo thời gian thực.

\subparagraph{a) Chiến lược lưu trữ kép (Dual Database
Strategy)}\label{a-chiux1ebfn-lux1b0ux1ee3c-lux1b0u-trux1eef-kuxe9p-dual-database-strategy}

Hệ thống theo dõi phương tiện cần lưu trữ hai loại dữ liệu có đặc tính
khác nhau căn bản:

\begin{enumerate}
\def\labelenumi{\arabic{enumi}.}
\tightlist
\item
  \textbf{Dữ liệu thô (Raw Telemetry Data)}: Bao gồm vị trí GPS, tốc độ,
  mức pin, dữ liệu OBD2, dữ liệu cảm biến gia tốc. Loại dữ liệu này có
  tần suất ghi rất cao (mỗi giây khi xe đang di chuyển), khối lượng lớn,
  nhưng chỉ cần lưu trữ trong thời gian ngắn (7--30 ngày) và chủ yếu
  phục vụ truy vấn theo chuỗi thời gian.
\item
  \textbf{Dữ liệu nghiệp vụ (Business Data)}: Bao gồm thông tin xe,
  khách hàng, chuyến đi, cảnh báo, vi phạm, lệnh điều khiển. Loại dữ
  liệu này có tần suất ghi thấp, khối lượng nhỏ, nhưng cần lưu trữ lâu
  dài (6--12 tháng trở lên) và yêu cầu tính toàn vẹn dữ liệu quan hệ
  (ACID).
\end{enumerate}

Do sự khác biệt cơ bản về đặc tính, hệ thống áp dụng chiến lược lưu trữ
kép sử dụng ba cơ sở dữ liệu chuyên biệt:


{\thesissetfigureid{3.13}{fig-ch3-so-o-chien-luoc-luu-tru-kep}
\begin{figure}[htbp]
\centering
\pandocbounded{\safeincludesvg{./assets/figures/05-chuong-3-giai-phap-backend-hinh-3-13.svg}}
\caption{Sơ đồ chiến lược lưu trữ kép}
\label{fig:ch3-so-o-chien-luoc-luu-tru-kep}
\end{figure}
}


\begin{itemize}
\tightlist
\item
  \textbf{PostgreSQL (quan hệ):} lưu dữ liệu nghiệp vụ như
  \texttt{vehicles}, \texttt{customers}, \texttt{trips}, \texttt{alerts},
  \texttt{violations}, \texttt{commands}, \texttt{geofences}.
\item
  \textbf{VictoriaMetrics (chuỗi thời gian):} lưu dữ liệu telemetry mật độ
  cao như vị trí GPS, tốc độ, mức pin, ignition/uptime, dữ liệu IMU
  (retention khoảng 30 ngày).
\item
  \textbf{VictoriaLogs (nhật ký):} lưu sự kiện thiết bị, lỗi, phiên làm
  việc, message MQTT và nhật ký kiểm toán (retention khoảng 7 ngày).
\end{itemize}


{\thesissettableid{3.18}{tab-ch3-so-sanh-cac-co-so-du-lieu-trong-he-thong}
\def\LTcaptype{table}
\begin{longtable}{|L{0.269\linewidth}|L{0.305\linewidth}|L{0.390\linewidth}|}
\caption{So sánh các cơ sở dữ liệu trong hệ thống}\label{tab:ch3-so-sanh-cac-co-so-du-lieu-trong-he-thong}\\
\hline \textbf{Cơ sở dữ liệu} & \textbf{Mục đích} & \textbf{Ưu điểm chính} \\
\hline
\endfirsthead
\caption[]{So sánh các cơ sở dữ liệu trong hệ thống (tiếp theo)}\\
\hline \textbf{Cơ sở dữ liệu} & \textbf{Mục đích} & \textbf{Ưu điểm chính} \\
\hline
\endhead
\hline
\endfoot
\endlastfoot
\textbf{PostgreSQL 16} & Dữ liệu quan hệ (nghiệp vụ) & ACID,
relationships, truy vấn phức tạp \\\hline
\textbf{VictoriaMetrics} & Dữ liệu chuỗi thời gian (telemetry) & Nhanh
gấp 10 lần InfluxDB, PromQL, tiêu thụ RAM thấp \\\hline
\textbf{VictoriaLogs} & Nhật ký tập trung (logging) & Tìm kiếm nhanh,
lưu trữ nhỏ gọn, LogsQL \\\hline
\end{longtable}
}


\subparagraph{b) Cơ sở dữ liệu PostgreSQL - Schema quan
hệ}\label{b-cux1a1-sux1edf-dux1eef-liux1ec7u-postgresql---schema-quan-hux1ec7}

PostgreSQL 16 lưu trữ toàn bộ dữ liệu nghiệp vụ của hệ thống. Schema
được chuẩn hóa đến dạng chuẩn thứ ba (3NF) và hỗ trợ soft delete cho các
bảng quan trọng.

\textbf{Các nhóm bảng chính (Phase 1 - khoảng 20 bảng):}


{\thesissettableid{3.19}{tab-ch3-cac-nhom-bang-chinh-trong-postgresql}
\def\LTcaptype{table}
\begin{longtable}{|L{0.203\linewidth}|L{0.419\linewidth}|L{0.343\linewidth}|}
\caption{Các nhóm bảng chính trong PostgreSQL}\label{tab:ch3-cac-nhom-bang-chinh-trong-postgresql}\\
\hline \textbf{Nhóm bảng} & \textbf{Các bảng chính} & \textbf{Mô tả} \\
\hline
\endfirsthead
\caption[]{Các nhóm bảng chính trong PostgreSQL (tiếp theo)}\\
\hline \textbf{Nhóm bảng} & \textbf{Các bảng chính} & \textbf{Mô tả} \\
\hline
\endhead
\hline
\endfoot
\endlastfoot
\textbf{Người dùng \& Xác thực} & users, user\_sessions & Quản lý tài
khoản, phiên đăng nhập \\\hline
\textbf{Phương tiện \& Thiết bị} & vehicles, devices,
device\_configurations & Thông tin xe và thiết bị tracker \\\hline
\textbf{Khách hàng} & customers & Thông tin khách hàng thuê xe \\\hline
\textbf{Chuyến đi} & trips, trip\_events, stops & Hành trình, sự kiện và
điểm dừng \\\hline
\textbf{Cảnh báo \& Vi phạm} & alerts, violations & Cảnh báo hệ thống và
vi phạm giao thông \\\hline
\textbf{Vùng địa lý} & geofences, vehicle\_geofences & Vùng địa lý và
gán kết với xe \\\hline
\textbf{Lệnh điều khiển} & commands & Lệnh gửi đến thiết bị \\\hline
\textbf{Bảo trì} & maintenance\_records & Lịch sử bảo trì phương tiện \\\hline
\textbf{Nhật ký} & connection\_logs, device\_status\_history & Nhật ký
kết nối và lịch sử trạng thái \\\hline
\textbf{Thông báo} & notifications & Thông báo hệ thống \\\hline
\textbf{Nhiên liệu} & fuel\_records & Quản lý nhiên liệu \\\hline
\end{longtable}
}


\textbf{Sơ đồ quan hệ chính (ER Diagram - rút gọn):}


{\thesissetfigureid{3.14}{fig-ch3-so-o-quan-he-co-so-du-lieu-er-diagram}
\begin{figure}[htbp]
\centering
\pandocbounded{\safeincludesvg{./assets/figures/05-chuong-3-giai-phap-backend-hinh-3-14.svg}}
\caption{Sơ đồ quan hệ cơ sở dữ liệu (ER Diagram)}
\label{fig:ch3-so-o-quan-he-co-so-du-lieu-er-diagram}
\end{figure}
}


\begin{itemize}
\tightlist
\item
  \texttt{users} 1:N \texttt{vehicles}, \texttt{customers},
  \texttt{notifications}.
\item
  \texttt{vehicles} 1:1 \texttt{devices}; 1:N \texttt{trips},
  \texttt{alerts}, \texttt{maintenance\_records}; N:N \texttt{geofences}
  (qua \texttt{vehicle\_geofences}).
\item
  \texttt{customers} 1:N \texttt{trips}.
\item
  \texttt{trips} 1:N \texttt{trip\_events}, \texttt{stops},
  \texttt{violations}.
\item
  \texttt{devices} 1:N \texttt{commands}, \texttt{device\_status\_history},
  \texttt{connection\_logs}.
\end{itemize}

Schema được thiết kế với khả năng mở rộng, các khóa ngoại và chỉ mục
(index) đã được chuẩn bị sẵn để tích hợp thêm các bảng Phase 2 bao gồm:
bookings (đặt xe), rental\_contracts (hợp đồng thuê), payments (thanh
toán), damage\_reports (báo cáo hư hỏng) và reviews (đánh giá).

\subparagraph{c) Cơ sở dữ liệu VictoriaMetrics - Schema chuỗi thời
gian}\label{c-cux1a1-sux1edf-dux1eef-liux1ec7u-victoriametrics---schema-chuux1ed7i-thux1eddi-gian}

VictoriaMetrics là cơ sở dữ liệu chuỗi thời gian tương thích Prometheus,
được lựa chọn thay thế InfluxDB nhờ hiệu năng vượt trội (nhanh hơn 10
lần), mức tiêu thụ RAM thấp hơn và tích hợp tốt với Grafana.


{\thesissettableid{3.20}{tab-ch3-so-sanh-victoriametrics-va-influxdb}
\def\LTcaptype{table}
\begin{longtable}{|L{0.307\linewidth}|L{0.210\linewidth}|L{0.447\linewidth}|}
\caption{So sánh VictoriaMetrics và InfluxDB}\label{tab:ch3-so-sanh-victoriametrics-va-influxdb}\\
\hline \textbf{Tiêu chí} & \textbf{InfluxDB} & \textbf{VictoriaMetrics (Đã chọn)} \\
\hline
\endfirsthead
\caption[]{So sánh VictoriaMetrics và InfluxDB (tiếp theo)}\\
\hline \textbf{Tiêu chí} & \textbf{InfluxDB} & \textbf{VictoriaMetrics (Đã chọn)} \\
\hline
\endhead
\hline
\endfoot
\endlastfoot
\textbf{Ngôn ngữ truy vấn} & Flux (phức tạp) & PromQL (đơn giản) \\\hline
\textbf{Hiệu năng} & Tốt & Nhanh gấp 10 lần \\\hline
\textbf{Sử dụng RAM} & Cao & Thấp \\\hline
\textbf{Nén dữ liệu} & Tốt & Tốt hơn (gấp 10 lần) \\\hline
\textbf{Tương thích Prometheus} & Không & Có \\\hline
\textbf{Độ khó học} & Cao (Flux) & Thấp (PromQL) \\\hline
\textbf{Tích hợp Grafana} & Qua plugin & Hỗ trợ sẵn (Native) \\\hline
\end{longtable}
}


\textbf{Chính sách lưu trữ (Retention Policy):} Dữ liệu thô được lưu trữ
trong 30 ngày (\texttt{-retentionPeriod=30d}). Hệ thống hỗ trợ cấu hình
downsampling thông qua vmagent với recording rules để tạo dữ liệu tổng
hợp theo giờ/ngày, phục vụ báo cáo dài hạn mà không tốn dung lượng lưu
trữ lớn.

\subparagraph{d) VictoriaLogs - Cơ sở dữ liệu nhật
ký}\label{d-victorialogs---cux1a1-sux1edf-dux1eef-liux1ec7u-nhux1eadt-kuxfd}

VictoriaLogs được sử dụng làm hệ thống nhật ký tập trung (centralized
logging), thay thế bộ công cụ ELK Stack truyền thống với ưu điểm nhẹ hơn
và dễ triển khai. VictoriaLogs lưu trữ các loại nhật ký: sự kiện thiết
bị, nhật ký lỗi, phiên làm việc, message MQTT, và nhật ký kiểm toán
(audit trail). Chính sách lưu trữ: 7 ngày cho nhật ký thông thường.

\paragraph{3.3.2.4. API Server và kiến trúc phần mềm (API Server \&
Software
Architecture)}\label{3234-api-server-vuxe0-kiux1ebfn-truxfac-phux1ea7n-mux1ec1m-api-server--software-architecture}

\subparagraph{a) Lựa chọn công
nghệ}\label{a-lux1ef1a-chux1ecdn-cuxf4ng-nghux1ec7}

API Server là tầng ứng dụng Backend cung cấp REST API, WebSocket và xử
lý logic nghiệp vụ. Để lựa chọn framework phù hợp, đồ án đã đánh giá ba
phương án:


{\thesissettableid{3.21}{tab-ch3-so-sanh-framework-backend}
\def\LTcaptype{table}
\begin{longtable}{|L{0.350\linewidth}|L{0.304\linewidth}|L{0.155\linewidth}|L{0.145\linewidth}|}
\caption{So sánh framework Backend}\label{tab:ch3-so-sanh-framework-backend}\\
\hline \textbf{Tiêu chí} & \textbf{Express + TypeScript} & \textbf{NestJS} & \textbf{Go + Gin} \\
\hline
\endfirsthead
\caption[]{So sánh framework Backend (tiếp theo)}\\
\hline \textbf{Tiêu chí} & \textbf{Express + TypeScript} & \textbf{NestJS} & \textbf{Go + Gin} \\
\hline
\endhead
\hline
\endfoot
\endlastfoot
\textbf{Độ khó học} & Rất dễ & Trung bình & Khó \\\hline
\textbf{Linh hoạt} & Rất cao & Trung bình & Cao \\\hline
\textbf{Hiệu năng} & Cao & Cao & Rất cao \\\hline
\textbf{Lượng code mẫu (Boilerplate)} & Ít & Nhiều & Trung bình \\\hline
\textbf{Hệ sinh thái IoT} & Rất tốt & Tốt & Tốt \\\hline
\textbf{Hỗ trợ thời gian thực} & Rất tốt & Rất tốt & Tốt \\\hline
\end{longtable}
}


\textbf{Kết luận lựa chọn: Node.js + Express + TypeScript}

Lý do lựa chọn gồm: tính linh hoạt cao; đặc tính nhẹ và nhanh phù hợp
với IoT; TypeScript giúp bảo đảm an toàn kiểu dữ liệu; Zod kết hợp xác
thực runtime với type inference; và khả năng tổ chức code theo
Domain-Driven Design.

\subparagraph{b) Kiến trúc Domain-Driven Design
(DDD)}\label{b-kiux1ebfn-truxfac-domain-driven-design-ddd}

API Server được tổ chức theo kiến trúc Domain-Driven Design với ba tầng
rõ ràng:

\textbf{Tầng 1 - API Layer (Routes):} Định nghĩa các route, ánh xạ HTTP
method đến controller tương ứng.

\textbf{Tầng 2 - Service Layer (Business Logic):} Chứa toàn bộ logic
nghiệp vụ, điều phối giữa các repository và dịch vụ bên ngoài.

\textbf{Tầng 3 - Repository Layer (Data Access):} Truy cập cơ sở dữ
liệu, thực hiện các truy vấn SQL và trả về dữ liệu đã được xác thực
kiểu.

Hệ thống bao gồm hơn 20 domain module, mỗi module có cấu trúc
\texttt{services/}, \texttt{repositories/}, \texttt{types/} riêng biệt,
đảm bảo nguyên tắc Single Responsibility và dễ bảo trì.

\subparagraph{c) MQTT Bridge - Dịch vụ độc
lập}\label{c-mqtt-bridge---dux1ecbch-vux1ee5-ux111ux1ed9c-lux1eadp}

MQTT Bridge đã được tách thành dịch vụ Node.js độc lập
(\texttt{Tracking\_MqttBridge/}) với \texttt{package.json},
\texttt{docker-compose.yml}, logger, connection pool và cấu hình riêng.
Cách tách này giúp dịch vụ dễ mở rộng hơn, cô lập lỗi tốt hơn và giữ
kiến trúc tổng thể gọn hơn.

\paragraph{3.3.2.5. Thiết kế API Endpoints (REST API
Design)}\label{3235-thiux1ebft-kux1ebf-api-endpoints-rest-api-design}

\subparagraph{a) Nguyên tắc thiết
kế}\label{a-nguyuxean-tux1eafc-thiux1ebft-kux1ebf}

API được thiết kế tuân thủ các nguyên tắc RESTful: URL dạng số nhiều,
versioning (\texttt{/api/v1/}), phân trang, định dạng lỗi thống nhất, và
HTTP Status Code chuẩn.

\subparagraph{b) Bảng tóm tắt API
Endpoints}\label{b-bux1ea3ng-tuxf3m-tux1eaft-api-endpoints}


{\thesissettableid{3.22}{tab-ch3-tom-tat-cac-nhom-api-endpoints-chinh}
\def\LTcaptype{table}
\begin{longtable}{|L{0.200\linewidth}|L{0.122\linewidth}|L{0.181\linewidth}|L{0.340\linewidth}|L{0.101\linewidth}|}
\caption{Tóm tắt các nhóm API Endpoints chính}\label{tab:ch3-tom-tat-cac-nhom-api-endpoints-chinh}\\
\hline \textbf{Nhóm} & \textbf{Method} & \textbf{Endpoint} & \textbf{Mô tả} & \textbf{Auth} \\
\hline
\endfirsthead
\caption[]{Tóm tắt các nhóm API Endpoints chính (tiếp theo)}\\
\hline \textbf{Nhóm} & \textbf{Method} & \textbf{Endpoint} & \textbf{Mô tả} & \textbf{Auth} \\
\hline
\endhead
\hline
\endfoot
\endlastfoot
\textbf{Auth} & POST & \nolinkurl{/api/v1/auth/login} & Đăng nhập &
Không \\\hline
& GET & \nolinkurl{/api/v1/auth/me} & Thông tin người dùng hiện tại & Có \\\hline
& POST & \nolinkurl{/api/v1/auth/logout} & Đăng xuất & Có \\\hline
\textbf{Devices} & GET & \nolinkurl{/api/v1/devices} & Danh sách thiết bị & Có \\\hline
& POST & \nolinkurl{/api/v1/devices} & Tạo thiết bị mới & Có \\\hline
& POST & \nolinkurl{/api/v1/devices/:id/ota} & Gửi lệnh OTA đến thiết bị & Có \\\hline
& POST & \nolinkurl{/api/v1/devices/:id/ota/rollback} & Yêu cầu rollback firmware &
Có \\\hline
\textbf{Vehicles} & GET & \nolinkurl{/api/v1/vehicles} & Danh sách xe & Có \\\hline
& POST & \nolinkurl{/api/v1/vehicles} & Tạo xe mới & Có \\\hline
\textbf{IoT Data} & POST & \nolinkurl{/api/v1/iot/data} & Gửi dữ liệu cảm biến &
Không \\\hline
\textbf{Telemetry} & GET & \nolinkurl{/api/v1/devices/:id/telemetry} & Dữ liệu
telemetry theo thiết bị & Có \\\hline
& GET & \nolinkurl{/api/v1/telemetry/history} & Truy vấn lịch sử telemetry & Có \\\hline
\textbf{Dashboard} & GET & \nolinkurl{/api/v1/dashboard/stats} & Thống kê tổng quan
& Có \\\hline
\textbf{Firmware} & POST & \nolinkurl{/api/v1/firmware/upload} & Tải lên firmware &
Có \\\hline
& GET & \nolinkurl{/api/v1/firmware/:id/download} & Tải file firmware để OTA & Có \\\hline
\end{longtable}
}



{\thesissetfigureid{3.14a}{fig-ch3-luong-ieu-phoi-ota-giua-backend-emqx-thiet-bi-va-mqtt-bridge}
\begin{figure}[htbp]
\centering
\pandocbounded{\safeincludesvg{./assets/figures/05-chuong-3-giai-phap-backend-hinh-3-14a.svg}}
\caption{Luồng điều phối OTA giữa Backend, EMQX, thiết bị và MQTT Bridge}
\label{fig:ch3-luong-ieu-phoi-ota-giua-backend-emqx-thiet-bi-va-mqtt-bridge}
\end{figure}
}


Backend vẫn duy trì một số alias tương thích như \texttt{/api/v1/device}
hoặc \texttt{/api/v1/export}, nhưng tài liệu ưu tiên dạng tài nguyên số
nhiều (\texttt{/devices}, \texttt{/vehicles}, \texttt{/customers},
\texttt{/trips}, \texttt{/alerts}, \texttt{/geofences}) để thống nhất
cách đặt tên.

\paragraph{3.3.2.6. WebSocket và truyền dữ liệu thời gian thực
(Real-time
Communication)}\label{3236-websocket-vuxe0-truyux1ec1n-dux1eef-liux1ec7u-thux1eddi-gian-thux1ef1c-real-time-communication}

Hệ thống sử dụng \textbf{Socket.IO 4.8} làm giải pháp truyền dữ liệu
thời gian thực từ Backend đến Frontend. Ở phạm vi triển khai này, hệ
thống định nghĩa năm namespace chính: \texttt{/dashboard},
\texttt{/devices}, \texttt{/notifications}, \texttt{/exports} và
\texttt{/firmware}. Tất cả namespace đều đi qua middleware xác thực;
riêng namespace \texttt{/devices} hỗ trợ join room theo từng thiết bị
(\texttt{device:\{id\}}) để nhận sự kiện chi tiết.

Các sự kiện tiêu biểu gồm \texttt{device:status},
\texttt{device:position}, \texttt{command:ack}, \texttt{stats:update},
\texttt{alert:new}, \texttt{geofence:enter}, \texttt{geofence:exit},
\texttt{export:ready} và \texttt{firmware:assignment}. Trong triển khai
hiện tại, namespace \texttt{/firmware} chủ yếu dùng để phát sự kiện gán
hoặc cập nhật tác vụ OTA từ backend; còn trạng thái chi tiết của quá
trình OTA được MQTT Bridge ghi vào \texttt{firmware\_update\_log} và
VictoriaLogs để phục vụ truy vết, theo dõi và đồng bộ với giao diện quản
trị.

\paragraph{3.3.2.7. Bảo mật hệ thống (Security
Design)}\label{3237-bux1ea3o-mux1eadt-hux1ec7-thux1ed1ng-security-design}

\subparagraph{a) Xác thực Session-based (Không dùng
JWT)}\label{a-xuxe1c-thux1ef1c-session-based-khuxf4ng-duxf9ng-jwt}

Hệ thống sử dụng cơ chế xác thực dựa trên phiên làm việc (session-based
authentication) với \emph{opaque session token} lưu trữ trong cơ sở dữ
liệu, thay vì JSON Web Token (JWT) tự mang toàn bộ thông tin phiên. Token
ở đây chỉ đóng vai trò khóa tham chiếu; backend mới là nơi lưu và kiểm
tra trạng thái phiên, còn giá trị token được băm SHA-256 trước khi ghi
xuống cơ sở dữ liệu.


{\thesissettableid{3.23}{tab-ch3-so-sanh-jwt-va-session-based-authentication}
\def\LTcaptype{table}
\begin{longtable}{|L{0.190\linewidth}|L{0.403\linewidth}|L{0.372\linewidth}|}
\caption{So sánh JWT và Session-based authentication}\label{tab:ch3-so-sanh-jwt-va-session-based-authentication}\\
\hline \textbf{Tiêu chí} & \textbf{JWT} & \textbf{Session-based (Đã chọn)} \\
\hline
\endfirsthead
\caption[]{So sánh JWT và Session-based authentication (tiếp theo)}\\
\hline \textbf{Tiêu chí} & \textbf{JWT} & \textbf{Session-based (Đã chọn)} \\
\hline
\endhead
\hline
\endfoot
\endlastfoot
\textbf{Thu hồi token} & Khó (cần blacklist) & Dễ (xóa khỏi DB) \\\hline
\textbf{Kiểm soát phiên} & Hạn chế & Toàn quyền (xem, thu hồi mọi
phiên) \\\hline
\textbf{Bảo mật} & Token chứa thông tin nhạy cảm & Chỉ chứa ID ngẫu
nhiên \\\hline
\textbf{Kích thước} & Lớn (header + payload + signature) & Nhỏ (chỉ
token ID) \\\hline
\end{longtable}
}


\subparagraph{b) Bảo mật tầng truyền
tải}\label{b-bux1ea3o-mux1eadt-tux1ea7ng-truyux1ec1n-tux1ea3i}

Hệ thống bảo vệ tầng truyền tải bằng nhiều lớp bổ trợ: Helmet để thiết
lập các header bảo mật, CORS để giới hạn origin hợp lệ, rate limiting để
chặn lạm dụng theo IP, và cơ chế gắn Request-ID để truy vết luồng xử lý
theo từng request.

\subparagraph{c) Bảo mật MQTT}\label{c-bux1ea3o-mux1eadt-mqtt}

Ở tầng MQTT, mỗi thiết bị được cấp ACL riêng và Client ID duy nhất, kết
hợp xác thực bằng username/password. Cách làm này giúp cô lập phạm vi
publish/subscribe của từng tracker và giảm rủi ro một thiết bị truy cập
nhầm sang luồng dữ liệu của thiết bị khác.

\subparagraph{d) Hệ thống thông báo cảnh
báo}\label{d-hux1ec7-thux1ed1ng-thuxf4ng-buxe1o-cux1ea3nh-buxe1o}

Hệ thống hỗ trợ gửi thông báo qua Telegram Bot (thời gian thực) và Email
SMTP (chi tiết). Các loại cảnh báo: xe di chuyển bất thường, vi phạm tốc
độ, pin thấp, thiết bị mất kết nối, lịch bảo trì sắp đến.

\subsection{3.4. Phân tích và đề xuất công nghệ Frontend -- Frontend
analysis and solution}\label{34-phan-tich-va-de-xuat-cong-nghe-frontend}

\subsubsection{3.4.1. Phân tích và lựa chọn công nghệ Frontend}

\paragraph{3.4.1.1. Đặt vấn đề cho giải pháp
Frontend}\label{3141-ux111ux1eb7t-vux1ea5n-ux111ux1ec1-cho-giux1ea3i-phuxe1p-frontend}

Nếu cloud là nơi xử lý dữ liệu, thì frontend là nơi người vận hành thực
sự ``nhìn thấy'' hệ thống đang hoạt động ra sao. Một dashboard tốt không
chỉ hiển thị đủ thông tin, mà còn phải giúp người dùng nhận ra điều gì
quan trọng, điều gì bất thường và cần hành động ở đâu. Vì vậy, bài toán
thiết kế frontend cần đồng thời đáp ứng các yêu cầu sau:

\begin{itemize}
\tightlist
\item
  \textbf{Hiển thị thời gian thực ổn định}: bản đồ, trạng thái thiết bị
  và cảnh báo phải cập nhật liên tục với độ trễ thấp.
\item
  \textbf{Khả năng mở rộng theo tính năng nghiệp vụ}: quản lý xe, chuyến
  đi, geofence, cảnh báo, báo cáo cần dễ mở rộng.
\item
  \textbf{Hiệu năng và trải nghiệm người dùng}: tải trang nhanh, điều
  hướng mượt, responsive tốt trên desktop/mobile.
\item
  \textbf{Dễ triển khai production}: build và đóng gói Docker thuận lợi
  cho mô hình vận hành đa dịch vụ.
\end{itemize}

\paragraph{3.4.1.2. So sánh các phương án công nghệ
Frontend}\label{3142-so-suxe1nh-cuxe1c-phux1b0ux1a1ng-uxe1n-cuxf4ng-nghux1ec7-frontend}


{\thesissettableid{3.23A}{tab-ch3-so-sanh-cac-phuong-an-cong-nghe-frontend-tong-the}
\def\LTcaptype{table}
\begin{longtable}{|L{0.154\linewidth}|L{0.157\linewidth}|L{0.280\linewidth}|L{0.269\linewidth}|L{0.085\linewidth}|}
\caption{So sánh các phương án công nghệ frontend tổng thể}\label{tab:ch3-so-sanh-cac-phuong-an-cong-nghe-frontend-tong-the}\\
\hline \textbf{Phương án} & \textbf{Mô tả stack} & \textbf{Ưu điểm} & \textbf{Hạn chế} & \textbf{Mức phù hợp} \\
\hline
\endfirsthead
\caption[]{So sánh các phương án công nghệ frontend tổng thể (tiếp theo)}\\
\hline \textbf{Phương án} & \textbf{Mô tả stack} & \textbf{Ưu điểm} & \textbf{Hạn chế} & \textbf{Mức phù hợp} \\
\hline
\endhead
\hline
\endfoot
\endlastfoot
\textbf{PA-FE1: Vite + React CSR} & React + Vite + client-side routing &
Build nhanh, cấu hình linh hoạt & SEO/first load kém hơn SSR, cần tự lắp
ghép nhiều thành phần & Trung bình \\\hline
\textbf{PA-FE2: Nuxt (Vue)} & Vue + Nuxt SSR/SSG & SSR tốt, cấu trúc rõ
& Khác hệ sinh thái React đang dùng ở dự án, chi phí chuyển đổi cao &
Trung bình \\\hline
\textbf{PA-FE3: Next.js + React (Đã chọn)} & Next.js 15 + React 19 + App
Router & SSR/RSC tốt, hệ sinh thái lớn, tối ưu production, tích hợp
realtime thuận lợi & Độ phức tạp framework cao hơn CSR thuần &
\textbf{Cao} \\\hline
\end{longtable}
}


\paragraph{3.4.1.3. Chọn giải pháp
Frontend}\label{3143-chux1ecdn-giux1ea3i-phuxe1p-frontend}

Đồ án chọn \textbf{PA-FE3: Next.js 15 + React 19} làm nền tảng frontend
chính. Đây không phải là lựa chọn chạy theo xu hướng framework, mà là
lựa chọn bám sát nhu cầu của một dashboard phải vừa ``đọc nhanh'' vừa
``phản ứng nhanh'' khi dữ liệu và cảnh báo thay đổi liên tục.

Các lý do lựa chọn gồm:

\begin{itemize}
\tightlist
\item
  \textbf{Phù hợp yêu cầu realtime + dashboard phức hợp}: kết hợp tốt
  giữa REST cache và WebSocket.
\item
  \textbf{Hiệu năng tải trang tốt}: App Router và khả năng render phía
  server cải thiện trải nghiệm ban đầu.
\item
  \textbf{Mở rộng tính năng thuận lợi}: dễ tổ chức theo module
  (Feature-Sliced Design).
\item
  \textbf{Sẵn sàng triển khai production}: hỗ trợ build/standalone
  output và đóng gói Docker rõ ràng.
\end{itemize}

\subsubsection{3.4.2. Giải pháp
Frontend}\label{324-giux1ea3i-phuxe1p-frontend}

Trên cơ sở phương án đã chọn ở mục 3.4.1, phần này trình bày cách nhóm
biến stack công nghệ thành một giao diện để người vận hành nắm tình
huống nhanh, thao tác gọn và vẫn theo kịp dữ liệu thời gian thực: từ lựa
chọn công nghệ, tổ chức mã nguồn đến cách quản lý trạng thái và hiển thị
dữ liệu.

\paragraph{3.4.2.1. Lựa chọn công nghệ Frontend (Technology
Selection)}\label{3241-lux1ef1a-chux1ecdn-cuxf4ng-nghux1ec7-frontend-technology-selection}

\subparagraph{a) Framework ứng dụng
web}\label{a-framework-ux1ee9ng-dux1ee5ng-web}

Để lựa chọn framework phù hợp, nhóm không dừng ở so sánh tên công nghệ,
mà đặt câu hỏi thực tế hơn: framework nào giúp dashboard tải nhanh, tổ
chức rõ và vẫn dễ mở rộng khi số lượng trang quản trị tăng lên. Từ đó,
nhóm đối chiếu Next.js, Nuxt.js và Vite + React theo các tiêu chí hiệu
năng, khả năng mở rộng và trải nghiệm phát triển.


{\thesissettableid{3.24}{tab-ch3-so-sanh-framework-frontend}
\def\LTcaptype{table}
\begin{longtable}{|L{0.195\linewidth}|L{0.323\linewidth}|L{0.192\linewidth}|L{0.244\linewidth}|}
\caption{So sánh framework Frontend}\label{tab:ch3-so-sanh-framework-frontend}\\
\hline \textbf{Tiêu chí} & \textbf{Next.js 15} & \textbf{Nuxt.js 3} & \textbf{Vite + React} \\
\hline
\endfirsthead
\caption[]{So sánh framework Frontend (tiếp theo)}\\
\hline \textbf{Tiêu chí} & \textbf{Next.js 15} & \textbf{Nuxt.js 3} & \textbf{Vite + React} \\
\hline
\endhead
\hline
\endfoot
\endlastfoot
Ngôn ngữ cơ sở & React 19 + TypeScript & Vue 3 + TypeScript & React +
TypeScript \\\hline
Rendering & SSR, SSG, ISR, RSC & SSR, SSG, ISR & CSR (chỉ
client-side) \\\hline
App Router & Có (App Router với React Server Components) & Có
(File-based routing) & Không có sẵn, cần thư viện bổ sung \\\hline
SEO \& First Load & Tốt (Server-side rendering) & Tốt & Kém (chỉ
client-side rendering) \\\hline
Hệ sinh thái & Rất lớn (React ecosystem) & Lớn (Vue ecosystem) & Lớn
(React ecosystem) \\\hline
Real-time support & Tốt (Socket.IO, Server Actions) & Tốt & Tốt \\\hline
Hiệu năng & Cao (Turbopack, Partial Prerendering) & Cao (Nitro engine) &
Cao (Vite bundler) \\\hline
Docker deployment & Có sẵn Dockerfile, standalone output & Có sẵn & Cần
cấu hình thủ công \\\hline
Cộng đồng \& tài liệu & Rất lớn & Lớn & Lớn \\\hline
\end{longtable}
}


\textbf{Kết luận lựa chọn:} Next.js 15 được chọn làm framework chính vì
nó tạo được thế cân bằng tốt giữa trải nghiệm của người dùng cuối và khả
năng tổ chức mã nguồn của đội phát triển. Các lý do cụ thể gồm:

\begin{itemize}
\tightlist
\item
  \textbf{React Server Components (RSC):} Cho phép render một phần trên
  server, giảm kích thước JavaScript gửi về client và tăng tốc độ tải
  trang đầu tiên (First Contentful Paint).
\item
  \textbf{App Router:} Cơ chế định tuyến dựa trên thư mục (file-based
  routing) đơn giản hóa việc tổ chức cấu trúc trang và layout lồng nhau
  (nested layouts).
\item
  \textbf{Hệ sinh thái React:} Tận dụng được hệ sinh thái cực kỳ phong
  phú của React bao gồm TanStack Query, Zustand, Leaflet, ECharts và
  hàng ngàn thư viện bổ sung.
\item
  \textbf{TypeScript tích hợp:} Hỗ trợ TypeScript từ gốc, đảm bảo an
  toàn kiểu dữ liệu xuyên suốt ứng dụng.
\item
  \textbf{Sẵn sàng cho triển khai:} Có sẵn các cơ chế tối ưu cho môi
  trường sản xuất như Image Optimization, Code Splitting và Standalone
  Output cho Docker.
\end{itemize}

\subparagraph{b) Thư viện giao diện (UI Component
Library)}\label{b-thux1b0-viux1ec7n-giao-diux1ec7n-ui-component-library}

Hệ thống sử dụng shadcn/ui kết hợp với Radix UI làm nền tảng xây dựng
giao diện. Điểm đáng giá của lựa chọn này không nằm ở số lượng component
dựng sẵn, mà ở khả năng giữ giao diện nhất quán trong khi nhóm vẫn kiểm
soát được cách từng thành phần được triển khai trong dự án.

\textbf{Đặc điểm nổi bật:}

\begin{itemize}
\tightlist
\item
  \textbf{Radix UI Primitives:} Cung cấp các component cơ bản đã xử lý
  accessibility (ARIA attributes), keyboard navigation và focus
  management.
\item
  \textbf{Tailwind CSS v4:} Hệ thống utility-first CSS với CSS Variables
  cho phép tùy chỉnh giao diện nhanh chóng mà không cần viết CSS truyền
  thống.
\item
  \textbf{Class Variance Authority (CVA):} Quản lý các biến thể
  (variants) của component một cách có hệ thống.
\item
  \textbf{Oklch Color Space:} Sử dụng hệ màu oklch hiện đại cho độ chính
  xác màu sắc cao hơn và hỗ trợ chuyển đổi Light/Dark theme mượt mà.
\end{itemize}

\subparagraph{c) Quản lý trạng thái (State
Management)}\label{c-quux1ea3n-luxfd-trux1ea1ng-thuxe1i-state-management}

Với một dashboard vừa có dữ liệu tĩnh, dữ liệu thời gian thực và trạng
thái giao diện, dồn mọi thứ vào cùng một nơi quản lý sẽ làm mã nguồn
nhanh chóng rối lên. Vì vậy, hệ thống áp dụng chiến lược tách biệt giữa
trạng thái client và trạng thái server:


{\thesissettableid{3.25}{tab-ch3-chien-luoc-quan-ly-trang-thai}
\def\LTcaptype{table}
\begin{longtable}{|L{0.177\linewidth}|L{0.215\linewidth}|L{0.235\linewidth}|L{0.327\linewidth}|}
\caption{Chiến lược quản lý trạng thái}\label{tab:ch3-chien-luoc-quan-ly-trang-thai}\\
\hline \textbf{Loại trạng thái} & \textbf{Thư viện} & \textbf{Mục đích} & \textbf{Ví dụ} \\
\hline
\endfirsthead
\caption[]{Chiến lược quản lý trạng thái (tiếp theo)}\\
\hline \textbf{Loại trạng thái} & \textbf{Thư viện} & \textbf{Mục đích} & \textbf{Ví dụ} \\
\hline
\endhead
\hline
\endfoot
\endlastfoot
Client state & Zustand 5 & Trạng thái UI và authentication & Token phiên
đăng nhập, trạng thái sidebar, theme \\\hline
Server state & TanStack Query 5 & Dữ liệu từ API server & Danh sách xe,
chuyến đi, cảnh báo \\\hline
URL state & nuqs & Trạng thái đồng bộ với URL & Bộ lọc, phân trang, tab
đang chọn \\\hline
Real-time state & Socket.IO Client & Dữ liệu thời gian thực & Vị trí xe,
dữ liệu telemetry \\\hline
\end{longtable}
}


\textbf{Kết luận:} Zustand được chọn vì API đơn giản, kích thước cực nhỏ
(~1.1 KB) và phù hợp với yêu cầu của hệ thống nơi token
phiên đăng nhập chỉ lưu trong bộ nhớ (memory-only, không persist ra
localStorage) nhằm tăng tính bảo mật.

\subparagraph{d) Thư viện bản
đồ}\label{d-thux1b0-viux1ec7n-bux1ea3n-ux111ux1ed3}


{\thesissettableid{3.26}{tab-ch3-so-sanh-thu-vien-ban-o}
\def\LTcaptype{table}
\begin{longtable}{|L{0.248\linewidth}|L{0.285\linewidth}|L{0.229\linewidth}|L{0.192\linewidth}|}
\caption{So sánh thư viện bản đồ}\label{tab:ch3-so-sanh-thu-vien-ban-o}\\
\hline \textbf{Tiêu chí} & \textbf{Leaflet.js} & \textbf{Mapbox GL JS} & \textbf{Google Maps JS API} \\
\hline
\endfirsthead
\caption[]{So sánh thư viện bản đồ (tiếp theo)}\\
\hline \textbf{Tiêu chí} & \textbf{Leaflet.js} & \textbf{Mapbox GL JS} & \textbf{Google Maps JS API} \\
\hline
\endhead
\hline
\endfoot
\endlastfoot
Giấy phép & MIT (mã nguồn mở) & Proprietary (có free tier) & Proprietary
(trả phí) \\\hline
Chi phí & Miễn phí & Miễn phí đến 50K loads/tháng & 0.007 USD/load sau
28K \\\hline
Hiệu năng với nhiều marker & Tốt (với MarkerCluster) & Rất tốt (WebGL
rendering) & Tốt \\\hline
Tùy chỉnh marker & Cao & Cao & Trung bình \\\hline
Offline/Self-hosted tiles & Hỗ trợ & Không & Không \\\hline
React integration & react-leaflet (trưởng thành) & react-map-gl &
@react-google-maps/api \\\hline
Vẽ geofence & @geoman-io/leaflet-geoman-free & mapbox-gl-draw & Drawing
Manager \\\hline
\end{longtable}
}


\textbf{Kết luận:} Leaflet.js được chọn vì hoàn toàn miễn phí (MIT
license), hỗ trợ tự host tile server cho môi trường không có Internet
công cộng, và có hệ sinh thái plugin phong phú (MarkerCluster, Geoman
cho vẽ geofence). Điều này đặc biệt phù hợp với hệ thống IoT có thể
triển khai trong môi trường private network.

\paragraph{3.4.2.2. Kiến trúc ứng dụng web (Web Application
Architecture)}\label{3242-kiux1ebfn-truxfac-ux1ee9ng-dux1ee5ng-web-web-application-architecture}

\subparagraph{a) Mô hình kiến trúc Feature-Sliced
Design}\label{a-muxf4-huxecnh-kiux1ebfn-truxfac-feature-sliced-design}

Ứng dụng Frontend được tổ chức theo mô hình Feature-Sliced Design (FSD),
trong đó mỗi tính năng nghiệp vụ được đóng gói thành một module độc lập
bao gồm components, hooks, types và utils. Cách tổ chức này giúp người
đọc và người phát triển nhìn thấy ngay mối liên hệ giữa cấu trúc mã
nguồn và các nghiệp vụ thực tế của hệ thống.

\begin{itemize}
\tightlist
\item
  \texttt{app/}: App Router của Next.js, bao gồm layout gốc và các route
  dashboard được bảo vệ.
\item
  \texttt{features/}: các module nghiệp vụ theo FSD (vehicles, trips,
  alerts, map), mỗi module tách rõ components/hooks/types.
\item
  \texttt{components/}: thư viện component dùng chung (UI, layout, map,
  charts, forms).
\item
  \texttt{lib/}: hạ tầng ứng dụng (API client, realtime, state store, util).
\item
  \texttt{types/}: các kiểu TypeScript dùng toàn cục.
\end{itemize}


{\thesissetfigureid{3.15}{fig-ch3-so-o-kien-truc-feature-sliced-design-cua-ung-dung-frontend}
\begin{figure}[htbp]
\centering
\pandocbounded{\safeincludesvg{./assets/figures/06-chuong-3-giai-phap-frontend-hinh-3-15.svg}}
\caption{Sơ đồ kiến trúc Feature-Sliced Design của ứng dụng Frontend}
\label{fig:ch3-so-o-kien-truc-feature-sliced-design-cua-ung-dung-frontend}
\end{figure}
}


\textbf{Nguyên tắc tổ chức:}

\begin{itemize}
\tightlist
\item
  \textbf{Co-location:} Các file liên quan đến một tính năng được đặt
  cùng thư mục, giảm thời gian tìm kiếm và bảo trì.
\item
  \textbf{Encapsulation:} Mỗi feature module export một public API rõ
  ràng, các module khác chỉ truy cập thông qua hooks và components đã
  export.
\item
  \textbf{Separation of Concerns:} Tách biệt rõ ràng giữa routing
  (app/), logic nghiệp vụ (features/), giao diện dùng lại (components/)
  và hạ tầng (lib/).
\end{itemize}

\subparagraph{b) Luồng dữ liệu trong ứng
dụng}\label{b-luux1ed3ng-dux1eef-liux1ec7u-trong-ux1ee9ng-dux1ee5ng}

Luồng dữ liệu trong ứng dụng Frontend tuân theo mô hình một chiều
(unidirectional data flow), đồng thời kết hợp ba nguồn dữ liệu chính:


{\thesissetfigureid{3.15a}{fig-ch3-luong-du-lieu-mot-chieu-giua-next-js-query-zustand-va-backend}
\begin{figure}[htbp]
\centering
\pandocbounded{\safeincludesvg{./assets/figures/06-chuong-3-giai-phap-frontend-hinh-3-15a.svg}}
\caption{Luồng dữ liệu một chiều giữa Next.js, Query, Zustand và Backend}
\label{fig:ch3-luong-du-lieu-mot-chieu-giua-next-js-query-zustand-va-backend}
\end{figure}
}


\begin{enumerate}
\def\labelenumi{\arabic{enumi}.}
\tightlist
\item
  \textbf{REST API (TanStack Query):} Dữ liệu CRUD (danh sách xe, chuyến
  đi, cảnh báo) được lấy từ Backend qua HTTP và cache bởi TanStack Query
  với cơ chế stale-while-revalidate.
\item
  \textbf{Client State (Zustand):} Trạng thái giao diện như phiên đăng
  nhập, trạng thái sidebar, theme được quản lý tại client mà không cần
  gọi API.
\item
  \textbf{Real-time (Socket.IO):} Dữ liệu vị trí xe và telemetry được
  nhận theo thời gian thực qua WebSocket. Khi kết nối WebSocket đang
  hoạt động, TanStack Query sẽ tự động tắt polling để tránh trùng lặp dữ
  liệu.
\end{enumerate}

\subparagraph{c) Cơ chế xác thực và bảo vệ
route}\label{c-cux1a1-chux1ebf-xuxe1c-thux1ef1c-vuxe0-bux1ea3o-vux1ec7-route}

Hệ thống Frontend sử dụng \emph{opaque session token} do Backend cấp,
không dùng JWT tự mang thông tin phiên. Việc client gửi token trong
header \texttt{Authorization:\ Bearer} chỉ là hình thức truyền token; về
bản chất, đây vẫn là cơ chế session-based vì mọi quyền truy cập đều được
quyết định sau khi server tra cứu phiên tương ứng:

\begin{enumerate}
\def\labelenumi{\arabic{enumi}.}
\tightlist
\item
  \textbf{Đăng nhập:} Người dùng nhập email và mật khẩu, Backend tạo bản
  ghi phiên và trả về session token ngẫu nhiên.
\item
  \textbf{Lưu trữ token:} Token chỉ được lưu trong bộ nhớ Zustand
  (memory-only), không lưu vào localStorage hay sessionStorage, nhằm giảm
  thiểu rủi ro XSS.
\item
  \textbf{Gửi request:} Mọi request HTTP tự động đính kèm token vào
  header
  \texttt{Authorization:\ Bearer\ \textless{}token\textgreater{}}.
\item
  \textbf{AuthGuard:} Component \texttt{AuthGuard} bao bọc toàn bộ khu
  vực \texttt{/dashboard/*}, tự động chuyển hướng về trang đăng nhập nếu
  chưa xác thực hoặc token hết hạn.
\item
  \textbf{Gia hạn phiên:} Khi phiên sắp hết hạn, Frontend chủ động gọi
  cơ chế gia hạn để nhận token mới, nhờ đó người dùng không bị đăng nhập
  lại đột ngột trong lúc thao tác.
\end{enumerate}


{\thesissetfigureid{3.16}{fig-ch3-bieu-o-trinh-tu-sequence-diagram-luong-xac-thuc-nguoi-dung}
\begin{figure}[htbp]
\centering
\pandocbounded{\safeincludesvg{./assets/figures/06-chuong-3-giai-phap-frontend-hinh-3-16.svg}}
\caption{Biểu đồ trình tự (Sequence Diagram) luồng xác thực người dùng}
\label{fig:ch3-bieu-o-trinh-tu-sequence-diagram-luong-xac-thuc-nguoi-dung}
\end{figure}
}


\paragraph{3.4.2.3. Thiết kế các trang chức năng chính (Main Feature
Pages)}\label{3243-thiux1ebft-kux1ebf-cuxe1c-trang-chux1ee9c-nux103ng-chuxednh-main-feature-pages}

\subparagraph{a) Trang tổng quan
(Dashboard)}\label{a-trang-tux1ed5ng-quan-dashboard}

Trang Dashboard là điểm vào chính sau khi đăng nhập, cung cấp cái nhìn
tổng thể về trạng thái hệ thống:

\begin{itemize}
\tightlist
\item
  \textbf{Stat Cards:} Hiển thị các chỉ số quan trọng: tổng số xe đang
  hoạt động, số chuyến đi trong ngày, số cảnh báo chưa xử lý, số vi phạm
  mới.
\item
  \textbf{Bản đồ mini:} Hiển thị vị trí tất cả xe đang hoạt động trên
  bản đồ thu nhỏ.
\item
  \textbf{Bảng cảnh báo gần đây:} Danh sách 10 cảnh báo mới nhất với
  badge mức độ nghiêm trọng (low, medium, high, critical).
\item
  \textbf{Biểu đồ thống kê:} Biểu đồ số chuyến đi theo ngày và phân loại
  cảnh báo theo loại.
\end{itemize}


{\thesissetfigureid{3.17}{fig-ch3-giao-dien-trang-dashboard-tong-quan}
\begin{figure}[htbp]
\centering
\pandocbounded{\safeincludesvg{./assets/figures/06-chuong-3-giai-phap-frontend-hinh-3-17.svg}}
\caption{Giao diện trang Dashboard tổng quan}
\label{fig:ch3-giao-dien-trang-dashboard-tong-quan}
\end{figure}
}


\subparagraph{b) Quản lý phương tiện (Vehicle
Management)}\label{b-quux1ea3n-luxfd-phux1b0ux1a1ng-tiux1ec7n-vehicle-management}

\textbf{Trang danh sách (\texttt{/dashboard/vehicles}):}


{\thesissettableid{3.2.4A}{tab-ch3-tong-hop-thong-tin-ve-cot-mo-ta-chuc-nang}
\def\LTcaptype{table}
\begin{longtable}{|L{0.212\linewidth}|L{0.519\linewidth}|L{0.234\linewidth}|}
\caption{Tổng hợp thông tin về Cột, Mô tả, Chức năng}\label{tab:ch3-tong-hop-thong-tin-ve-cot-mo-ta-chuc-nang}\\
\hline \textbf{Cột} & \textbf{Mô tả} & \textbf{Chức năng} \\
\hline
\endfirsthead
\caption[]{Tổng hợp thông tin về Cột, Mô tả, Chức năng (tiếp theo)}\\
\hline \textbf{Cột} & \textbf{Mô tả} & \textbf{Chức năng} \\
\hline
\endhead
\hline
\endfoot
\endlastfoot
Biển số xe & Số đăng ký phương tiện & Tìm kiếm \\\hline
Hãng/Model & Thương hiệu và đời xe & Lọc \\\hline
Trạng thái & Active, Inactive, Maintenance, Retired & Lọc theo trạng
thái \\\hline
Thiết bị & Mã thiết bị IoT gắn kèm & Liên kết \\\hline
Lần cuối hoạt động & Thời điểm nhận dữ liệu gần nhất & Sắp xếp \\\hline
Thao tác & Xem, Sửa, Xóa, Xem trên bản đồ & Nút hành động \\\hline
\end{longtable}
}


\textbf{Trang chi tiết (\texttt{/dashboard/vehicles/{[}id{]}}):} Hiển
thị thông tin đầy đủ của phương tiện bao gồm thông tin cơ bản, vị trí
hiện tại trên bản đồ, trạng thái thiết bị IoT, các cảnh báo đang hoạt
động, lịch sử chuyến đi và lịch sử bảo trì.

\clearpage


{\thesissetfigureid{3.18}{fig-ch3-giao-dien-trang-quan-ly-phuong-tien}
\begin{figure}[htbp]
\centering
\pandocbounded{\safeincludesvgnarrow{./assets/figures/06-chuong-3-giai-phap-frontend-hinh-3-18.svg}}
\caption{Giao diện trang quản lý phương tiện}
\label{fig:ch3-giao-dien-trang-quan-ly-phuong-tien}
\end{figure}
}


\subparagraph{c) Quản lý chuyến đi, cảnh báo và vùng địa
lý}\label{c-quux1ea3n-luxfd-chuyux1ebfn-ux111i-cux1ea3nh-buxe1o-vuxe0-vuxf9ng-ux111ux1ecba-luxfd}

Trang quản lý chuyến đi cung cấp lịch sử và chi tiết từng chuyến, gồm
tuyến đường trên bản đồ, biểu đồ tốc độ theo thời gian và timeline sự
kiện. Trang cảnh báo tổng hợp cảnh báo theo badge mức độ nghiêm trọng và
hỗ trợ xử lý hàng loạt. Trang geofence cho phép vẽ hình tròn, đa giác
hoặc hình chữ nhật trực tiếp trên bản đồ bằng Leaflet Geoman.


{\thesissettableid{3.27}{tab-ch3-ma-tran-cac-trang-chuc-nang-chinh}
\def\LTcaptype{table}
\begin{longtable}{|L{0.179\linewidth}|L{0.158\linewidth}|L{0.385\linewidth}|L{0.233\linewidth}|}
\caption{Ma trận các trang chức năng chính}\label{tab:ch3-ma-tran-cac-trang-chuc-nang-chinh}\\
\hline \textbf{Trang} & \textbf{Đường dẫn} & \textbf{Chức năng chính} & \textbf{Real-time} \\
\hline
\endfirsthead
\caption[]{Ma trận các trang chức năng chính (tiếp theo)}\\
\hline \textbf{Trang} & \textbf{Đường dẫn} & \textbf{Chức năng chính} & \textbf{Real-time} \\
\hline
\endhead
\hline
\endfoot
\endlastfoot
Dashboard & \texttt{/dashboard} & Tổng quan hệ thống, stat cards, biểu
đồ & Có (cập nhật số liệu) \\\hline
Bản đồ & \nolinkurl{/dashboard/map} & Theo dõi tất cả xe thời gian thực &
Có (vị trí xe) \\\hline
Phương tiện & \nolinkurl{/dashboard/vehicles} & CRUD xe, trạng thái, lịch
sử & Có (trạng thái xe) \\\hline
Chuyến đi & \nolinkurl{/dashboard/trips} & Lịch sử chuyến đi, tuyến đường &
Không \\\hline
Cảnh báo & \nolinkurl{/dashboard/alerts} & Xem, xử lý cảnh báo & Có (cảnh
báo mới) \\\hline
Vi phạm & \nolinkurl{/dashboard/violations} & Quản lý vi phạm tốc độ &
Không \\\hline
Thiết bị & \nolinkurl{/dashboard/devices} & Quản lý thiết bị IoT & Có
(trạng thái thiết bị) \\\hline
Vùng địa lý & \nolinkurl{/dashboard/geofences} & Tạo, sửa geofence trên bản
đồ & Không \\\hline
Bảo trì & \nolinkurl{/dashboard/maintenance} & Lịch sử bảo trì phương tiện
& Không \\\hline
Thông báo & \nolinkurl{/dashboard/notifications} & Cấu hình Telegram, Email
& Không \\\hline
Cài đặt & \nolinkurl{/dashboard/settings} & Hồ sơ, quản lý người dùng &
Không \\\hline
\end{longtable}
}


\paragraph{3.4.2.4. Tích hợp bản đồ thời gian thực (Real-time Map
Integration)}\label{3244-tuxedch-hux1ee3p-bux1ea3n-ux111ux1ed3-thux1eddi-gian-thux1ef1c-real-time-map-integration}

\subparagraph{a) Kiến trúc tích hợp bản
đồ}\label{a-kiux1ebfn-truxfac-tuxedch-hux1ee3p-bux1ea3n-ux111ux1ed3}

Trang bản đồ thời gian thực (\texttt{/dashboard/map}) là tính năng cốt
lõi của hệ thống giám sát phương tiện. Kiến trúc tích hợp bản đồ bao gồm
ba lớp chính:


{\thesissetfigureid{3.19a}{fig-ch3-phan-ra-lop-ban-o-react-leaflet-va-kenh-socket-io}
\begin{figure}[htbp]
\centering
\pandocbounded{\safeincludesvg{./assets/figures/06-chuong-3-giai-phap-frontend-hinh-3-19a.svg}}
\caption{Phân rã lớp bản đồ React Leaflet và kênh Socket.IO}
\label{fig:ch3-phan-ra-lop-ban-o-react-leaflet-va-kenh-socket-io}
\end{figure}
}



{\thesissetfigureid{3.19}{fig-ch3-kien-truc-tich-hop-ban-o-thoi-gian-thuc}
\begin{figure}[htbp]
\centering
\pandocbounded{\safeincludesvgcompact{./assets/figures/06-chuong-3-giai-phap-frontend-hinh-3-19.svg}}
\caption{Kiến trúc tích hợp bản đồ thời gian thực}
\label{fig:ch3-kien-truc-tich-hop-ban-o-thoi-gian-thuc}
\end{figure}
}


\subparagraph{b) Cơ chế cập nhật vị trí thời gian
thực}\label{b-cux1a1-chux1ebf-cux1eadp-nhux1eadt-vux1ecb-truxed-thux1eddi-gian-thux1ef1c}

Quy trình cập nhật vị trí xe trên bản đồ diễn ra như sau:

\begin{enumerate}
\def\labelenumi{\arabic{enumi}.}
\tightlist
\item
  \textbf{Thiết bị IoT} gửi dữ liệu GPS qua MQTT đến EMQX Broker.
\item
  \textbf{MQTT Bridge} nhận và xử lý dữ liệu, lưu vào VictoriaMetrics,
  đồng thời chuyển tiếp đến Backend.
\item
  \textbf{Backend} phát sự kiện \texttt{vehicle:\{id\}:location} qua
  Socket.IO đến tất cả client đang subscribe.
\item
  \textbf{Frontend} nhận sự kiện, cập nhật tọa độ marker trên bản đồ. Sử
  dụng kỹ thuật nội suy (interpolation) để marker di chuyển mượt mà giữa
  hai điểm tọa độ liên tiếp.
\end{enumerate}

\textbf{Tối ưu hiệu năng bản đồ:}

\begin{itemize}
\tightlist
\item
  \textbf{Marker Clustering:} Sử dụng \texttt{react-leaflet-cluster} để
  gom nhóm các marker gần nhau khi zoom out, giảm số lượng DOM elements
  cần render.
\item
  \textbf{Viewport-based rendering:} Chỉ render marker của các xe nằm
  trong vùng nhìn hiện tại của bản đồ.
\item
  \textbf{Debounced updates:} Gom nhóm các cập nhật vị trí trong khoảng
  100ms để tránh re-render quá nhiều lần.
\end{itemize}

\subparagraph{c) Các thành phần bản
đồ}\label{c-cuxe1c-thuxe0nh-phux1ea7n-bux1ea3n-ux111ux1ed3}

Hệ thống bao gồm bốn component bản đồ chuyên biệt:


{\thesissettableid{3.2.4B}{tab-ch3-tong-hop-thong-tin-ve-component-file-chuc-nang}
\def\LTcaptype{table}
\begin{longtable}{|L{0.224\linewidth}|L{0.160\linewidth}|L{0.580\linewidth}|}
\caption{Tổng hợp thông tin về Component, File, Chức năng}\label{tab:ch3-tong-hop-thong-tin-ve-component-file-chuc-nang}\\
\hline \textbf{Component} & \textbf{File} & \textbf{Chức năng} \\
\hline
\endfirsthead
\caption[]{Tổng hợp thông tin về Component, File, Chức năng (tiếp theo)}\\
\hline \textbf{Component} & \textbf{File} & \textbf{Chức năng} \\
\hline
\endhead
\hline
\endfoot
\endlastfoot
VehicleMap & \nolinkurl{components/map/vehicle-map.tsx} & Hiển thị tất cả
xe trên bản đồ với marker và popup thông tin \\\hline
RouteMap & \nolinkurl{components/map/route-map.tsx} & Hiển thị tuyến đường
chuyến đi với polyline và các điểm dừng \\\hline
GeofenceMap & \nolinkurl{components/map/geofence-map.tsx} & Vẽ và chỉnh sửa
vùng địa lý (polygon, circle, rectangle) \\\hline
MarkerPopup & \nolinkurl{components/map/marker-popup.tsx} & Popup hiển thị
thông tin xe khi click marker (biển số, tốc độ, trạng thái) \\\hline
\end{longtable}
}



{\thesissetfigureid{3.20}{fig-ch3-giao-dien-trang-ban-o-thoi-gian-thuc-voi-cac-marker-xe}
\begin{figure}[htbp]
\centering
\pandocbounded{\safeincludesvgnarrow{./assets/figures/06-chuong-3-giai-phap-frontend-hinh-3-20.svg}}
\caption{Giao diện trang bản đồ thời gian thực với các marker xe}
\label{fig:ch3-giao-dien-trang-ban-o-thoi-gian-thuc-voi-cac-marker-xe}
\end{figure}
}

\clearpage


\paragraph{3.4.2.5. Thiết kế responsive và trải nghiệm người dùng
(Responsive Design \&
UX)}\label{3245-thiux1ebft-kux1ebf-responsive-vuxe0-trux1ea3i-nghiux1ec7m-ngux1b0ux1eddi-duxf9ng-responsive-design--ux}

\subparagraph{a) Chiến lược
responsive}\label{a-chiux1ebfn-lux1b0ux1ee3c-responsive}

Trong phạm vi hiện tại, giao diện được ưu tiên cho desktop và tablet,
nhưng vẫn giữ ba mức breakpoint chính để không phá vỡ trải nghiệm trên màn
hình nhỏ:


{\thesissettableid{3.28}{tab-ch3-breakpoint-va-tuong-ung-giao-dien}
\def\LTcaptype{table}
\begin{longtable}{|L{0.180\linewidth}|L{0.151\linewidth}|L{0.176\linewidth}|L{0.448\linewidth}|}
\caption{Breakpoint và tương ứng giao diện}\label{tab:ch3-breakpoint-va-tuong-ung-giao-dien}\\
\hline \textbf{Breakpoint} & \textbf{Kích thước} & \textbf{Giao diện} & \textbf{Điều chỉnh chính} \\
\hline
\endfirsthead
\caption[]{Breakpoint và tương ứng giao diện (tiếp theo)}\\
\hline \textbf{Breakpoint} & \textbf{Kích thước} & \textbf{Giao diện} & \textbf{Điều chỉnh chính} \\
\hline
\endhead
\hline
\endfoot
\endlastfoot
Mobile & \textless{} 640px & Giao diện tối giản, một cột & Sidebar thành
drawer, bảng thành danh sách card, bản đồ toàn màn hình \\\hline
Tablet & 640px - 1024px & Giao diện trung gian & Sidebar thu gọn
(icon-only), bảng giữ nguyên với cuộn ngang \\\hline
Desktop & \textgreater{} 1024px & Giao diện đầy đủ & Sidebar mở rộng,
bảng đầy đủ cột, bản đồ chia đôi màn hình \\\hline
\end{longtable}
}


\subparagraph{b) Layout chính}\label{b-layout-chuxednh}

Giao diện Dashboard sử dụng bố cục (layout) gồm ba vùng chính: Header
(breadcrumbs, search, user menu, theme toggle), Sidebar (navigation
menu), và Main Content Area (page content, filters, data tables).


{\thesissetfigureid{3.21}{fig-ch3-wireframe-bo-cuc-giao-dien-dashboard}
\begin{figure}[htbp]
\centering
\pandocbounded{\safeincludesvgcompact{./assets/figures/06-chuong-3-giai-phap-frontend-hinh-3-21.svg}}
\caption{Wireframe bố cục giao diện Dashboard}
\label{fig:ch3-wireframe-bo-cuc-giao-dien-dashboard}
\end{figure}
}


\subparagraph{c) Hệ thống theme
Light/Dark}\label{c-hux1ec7-thux1ed1ng-theme-lightdark}

Hệ thống hỗ trợ hai chế độ giao diện (Light và Dark) sử dụng CSS
Variables với hệ màu oklch. Chuyển đổi được thực hiện qua thư viện
\texttt{next-themes} và lưu trong cookie.

\subparagraph{d) Các mẫu thiết kế giao diện (UI
Patterns)}\label{d-cuxe1c-mux1eabu-thiux1ebft-kux1ebf-giao-diux1ec7n-ui-patterns}

Hệ thống áp dụng các mẫu thiết kế giao diện nhất quán xuyên suốt ứng
dụng:

\begin{itemize}
\tightlist
\item
  \textbf{Bảng dữ liệu (Data Tables):} Sắp xếp, lọc, phân trang phía
  server, chọn nhiều dòng, xuất dữ liệu CSV/Excel.
\item
  \textbf{Biểu mẫu (Forms):} React Hook Form + Zod validation, inline
  error messages, toast notifications (Sonner).
\item
  \textbf{Bản đồ (Maps):} Custom markers với màu theo trạng thái, popup
  thông tin, drawing tools cho geofence.
\item
  \textbf{Biểu đồ (Charts):} ECharts/Recharts cho telemetry, line chart
  tốc độ, gauge chart RPM, area chart nhiên liệu/nhiệt độ.
\end{itemize}


{\thesissetfigureid{3.22}{fig-ch3-cac-mau-thiet-ke-ui-cua-he-thong-data-table-form-charts}
\begin{figure}[htbp]
\centering
\pandocbounded{\safeincludesvgcompact{./assets/figures/06-chuong-3-giai-phap-frontend-hinh-3-22.svg}}
\caption{Các mẫu thiết kế UI của hệ thống - Data Table, Form, Charts}
\label{fig:ch3-cac-mau-thiet-ke-ui-cua-he-thong-data-table-form-charts}
\end{figure}
}


\subparagraph{e) Tối ưu trải nghiệm người
dùng}\label{e-tux1ed1i-ux1b0u-trux1ea3i-nghiux1ec7m-ngux1b0ux1eddi-duxf9ng}

Để nâng cao trải nghiệm người dùng, hệ thống áp dụng một số kỹ thuật:

\begin{itemize}
\tightlist
\item
  \textbf{Optimistic Updates:} Giao diện cập nhật ngay lập tức trước khi
  nhận phản hồi từ server.
\item
  \textbf{Skeleton Loading:} Hiển thị khung xương thay vì spinner khi
  đang tải dữ liệu.
\item
  \textbf{Page Transition Loader:} Thanh tiến trình ở đầu trang (sử dụng
  \texttt{nextjs-toploader}) khi chuyển trang.
\item
  \textbf{Keyboard Shortcuts:} Hỗ trợ phím tắt cho các thao tác phổ biến
  (\texttt{Cmd+K} cho tìm kiếm).
\item
  \textbf{Error Boundaries:} Mỗi khu vực chức năng được bao bọc bởi
  Error Boundary riêng.
\item
  \textbf{Empty States:} Hiển thị trạng thái trống có hướng dẫn hành
  động khi chưa có dữ liệu.
\end{itemize}

\begin{center}\rule{0.5\linewidth}{0.5pt}\end{center}

\subsection{3.5. Phân tích, đánh giá và lựa chọn phương án khả thi --
Analysis, Evaluation and
Selection}\label{33-phuxe2n-tuxedch-ux111uxe1nh-giuxe1-vuxe0-lux1ef1a-chux1ecdn-phux1b0ux1a1ng-uxe1n-khux1ea3-thi--analysis-evaluation-and-selection}

Dựa trên phân tích và đề xuất ở các mục 3.1 đến 3.4, bảng tổng hợp dưới
đây đánh giá và so sánh các phương án đã xem xét:


{\thesissettableid{3.29}{tab-ch3-ma-tran-anh-gia-tong-hop-cac-phuong-an-thiet-ke}
\def\LTcaptype{table}
\begin{longtable}{|L{0.185\linewidth}|L{0.320\linewidth}|L{0.459\linewidth}|}
\caption{Ma trận đánh giá tổng hợp các phương án thiết kế}\label{tab:ch3-ma-tran-anh-gia-tong-hop-cac-phuong-an-thiet-ke}\\
\hline \textbf{Tầng hệ thống} & \textbf{Phương án được chọn} & \textbf{Lý do chính} \\
\hline
\endfirsthead
\caption[]{Ma trận đánh giá tổng hợp các phương án thiết kế (tiếp theo)}\\
\hline \textbf{Tầng hệ thống} & \textbf{Phương án được chọn} & \textbf{Lý do chính} \\
\hline
\endhead
\hline
\endfoot
\endlastfoot
Vi điều khiển & ESP32-S3 & BLE 5.0 tích hợp, dual-core, hỗ trợ AI \\\hline
LTE + GNSS & SIMCom SIM7600CE-T & Modem tích hợp LTE + GNSS, giảm số
lượng khối phần cứng rời \\\hline
OBD2 Adapter & vgate iCar Pro (BLE) & Không cần dây, tương thích rộng \\\hline
Cảm biến IMU & LIS3DSH & Siêu tiết kiệm điện, wake-on-motion \\\hline
Buffer cục bộ & microSD SDMMC 4-bit + FATFS + replay & Duy trì liên tục
telemetry khi mất mạng, phát lại theo thứ tự khi reconnect \\\hline
MQTT Broker & EMQX & Rules Engine, ACL per device, clustering \\\hline
Database quan hệ & PostgreSQL 16 & Mature, PostGIS, open-source \\\hline
Database time-series & VictoriaMetrics & Write throughput cao, PromQL
compatible \\\hline
API Server & Express.js + TypeScript & Nhẹ, linh hoạt, DDD
architecture \\\hline
Frontend & Next.js 15 + React 19 & SSR/SSG, App Router, ecosystem \\\hline
Bản đồ & Leaflet.js & Open-source, self-hosted tiles \\\hline
\end{longtable}
}


\subsection{3.6. Tối ưu phương án thiết kế -- The optimal
solution}\label{34-tux1ed1i-ux1b0u-phux1b0ux1a1ng-uxe1n-thiux1ebft-kux1ebf--the-optimal-solution}

Từ kết quả phân tích và đánh giá ở mục 3.5, phương án thiết kế tối ưu
cho hệ thống IoT giám sát phương tiện được tổng hợp như sau:


{\thesissetfigureid{3.23}{fig-ch3-so-o-kien-truc-tong-the-phuong-an-toi-uu}
\begin{figure}[htbp]
\centering
\pandocbounded{\safeincludesvgcompact{./assets/figures/06-chuong-3-giai-phap-frontend-hinh-3-23.svg}}
\caption{Sơ đồ kiến trúc tổng thể phương án tối ưu}
\label{fig:ch3-so-o-kien-truc-tong-the-phuong-an-toi-uu}
\end{figure}
}


\textbf{Kiến trúc phân tầng:}

\begin{itemize}
\tightlist
\item
  \textbf{Tầng thiết bị (Device Layer):} ESP32-S3 + SIMCom SIM7600CE-T +
  vgate iCar Pro + LIS3DSH, quản lý nguồn thông minh với pin dự phòng
  18650 1S
\item
  \textbf{Tầng truyền thông (Communication Layer):} MQTT 5.0 qua 4G LTE,
  QoS phân tầng, kết hợp buffer cục bộ microSD và cơ chế replay khi kết
  nối phục hồi
\item
  \textbf{Tầng xử lý (Processing Layer):} MQTT Bridge -\/-\textgreater{}
  dual-write PostgreSQL + VictoriaMetrics, Express.js API (DDD)
\item
  \textbf{Tầng trình bày (Presentation Layer):} Next.js 15, Leaflet
  maps, Socket.IO real-time, ECharts
\end{itemize}

Phương án này đáp ứng các yêu cầu kỹ thuật đã đặt ra ở Chương 2, đồng
thời giữ chi phí phần cứng ở mức khoảng 1.514.000 VND cho mỗi thiết bị
theo mặt bằng giá tham chiếu tháng 4/2026.

\begin{center}\rule{0.5\linewidth}{0.5pt}\end{center}

\subsection{Kết luận chương
3}\label{kux1ebft-luux1eadn-chux1b0ux1a1ng-3}

Chương 3 đã trả lời câu hỏi "thiết kế hệ thống này như thế nào để vừa
đáp ứng yêu cầu kỹ thuật, vừa khả thi khi triển khai". Từ phần cứng,
firmware đến backend/cloud và frontend, các phương án không chỉ được nêu
ra mà còn được cân nhắc theo tiêu chí hiệu năng, chi phí, độ phức tạp và
khả năng mở rộng. Nhờ đó, phương án cuối cùng ở Chương 3 không xuất hiện
như một lựa chọn cảm tính, mà như kết quả của quá trình sàng lọc có cơ
sở. Trên nền tảng đó, Chương 4 sẽ chuyển sang phần hiện thực: biến các
lựa chọn thiết kế thành hệ thống vận hành được trong thực tế.

\chapterheading{CHƯƠNG 4. TRIỂN KHAI GIẢI PHÁP VÀ KẾT QUẢ -- IMPLEMENTATION AND RESULTS}

Nếu Chương 3 dừng ở mức lựa chọn phương án, thì Chương 4 trả lời câu
hỏi liệu những lựa chọn đó có thực sự hoạt động khi đưa vào chế tạo, lập
trình và tích hợp hay không. Nội dung chương đi từ phần cứng thiết bị,
firmware, hạ tầng cloud, giao diện giám sát đến các kết quả đo lường sau
khi lắp đặt hệ thống trên xe, qua đó giúp người đọc đối chiếu giữa ý
tưởng thiết kế và kết quả triển khai thực tế.

\begin{center}\rule{0.5\linewidth}{0.5pt}\end{center}

\subsection{4.1. Tóm tắt triển khai giải pháp -- Compact implementation summary}

Bản rút gọn này giữ lại các quyết định triển khai cốt lõi và lược bớt các sơ đồ chi tiết, bảng chân GPIO, checklist lắp đặt, cây thư mục mã nguồn và ma trận cấu hình phụ. Mục tiêu là duy trì mạch kỹ thuật chính: hệ thống đã được hiện thực đồng bộ từ thiết bị gắn trên xe, firmware nhúng, hạ tầng cloud/backend đến giao diện web theo dõi thời gian thực.

\subsubsection{4.1.1. Triển khai phần cứng thiết bị}

Thiết bị tracker được triển khai quanh vi điều khiển ESP32-S3, modem LTE/GNSS SIM7600CE-T, adapter BLE OBD2 vgate iCar Pro và cảm biến gia tốc LIS3DSH. Kiến trúc phần cứng chia thành bốn nhóm chức năng chính:

\begin{itemize}
\tightlist
\item
  \textbf{Khối xử lý trung tâm:} ESP32-S3 điều phối toàn bộ giao tiếp UART, BLE, I2C, ADC và logic trạng thái vận hành.
\item
  \textbf{Khối truyền thông và định vị:} modem SIM7600CE-T đảm nhiệm truyền dữ liệu 4G/MQTT và cung cấp vị trí GNSS cho hệ thống.
\item
  \textbf{Khối thu thập dữ liệu xe:} adapter BLE OBD2 cung cấp trạng thái IGN, RPM, tốc độ, mức nhiên liệu và các PID cần thiết cho giám sát vận hành.
\item
  \textbf{Khối quản lý năng lượng:} kiến trúc buck, boost, LDO, power path, sạc pin và giám sát điện áp giúp thiết bị duy trì hoạt động cả khi điện áp ắc quy xe suy giảm.
\end{itemize}

Trong bản đầy đủ, nhóm tác giả trình bày chi tiết sơ đồ nguyên lý, pin mapping, enclosure, quy trình lắp ráp và checklist kiểm tra sau lắp đặt. Với bản rút gọn, các chi tiết này được gom lại thành kết luận kỹ thuật sau: thiết kế phần cứng đã đạt mục tiêu tích hợp một bo tracker duy nhất, có thể lắp trong xe, giao tiếp được với nguồn xe 12V/24V và hỗ trợ chế độ dự phòng bằng pin 18650 để tăng tính liên tục dịch vụ.

\subsubsection{4.1.2. Triển khai firmware nhúng}

Firmware được phát triển trên ESP-IDF 5.4 với FreeRTOS và NimBLE, tổ chức theo định hướng mô-đun thay vì tách nhiều tác vụ nghiệp vụ cấp cao. Cấu trúc này giúp giữ luồng điều khiển tập trung, thuận lợi cho việc kiểm soát trạng thái và debug trong môi trường nhúng có nhiều ràng buộc thời gian thực.

Các cụm chức năng chính của firmware gồm:

\begin{itemize}
\tightlist
\item
  \textbf{Thu thập tín hiệu đầu vào:} đọc điện áp ắc quy qua ADC, nhận ngắt từ IMU, quản lý BLE với OBD2 adapter và điều khiển modem qua UART/AT.
\item
  \textbf{Điều phối trạng thái:} máy trạng thái ba chế độ \emph{Driving}, \emph{Parking}, \emph{Alert} quyết định tần suất đọc dữ liệu, gửi telemetry, ngủ sâu và đánh thức theo sự kiện.
\item
  \textbf{Đóng gói và truyền dữ liệu:} telemetry được chuẩn hóa thành payload MQTT thống nhất, phân tuyến theo topic, QoS và loại bản tin.
\item
  \textbf{Duy trì vận hành:} firmware hỗ trợ reconnect MQTT, lưu bền cấu hình, cập nhật OTA cơ bản và chuẩn bị cơ chế replay dữ liệu qua lưu đệm cục bộ.
\end{itemize}

Điểm quan trọng của triển khai firmware không nằm ở việc trình bày toàn bộ mã nguồn, mà ở chỗ logic vận hành đã bám sát phương án thiết kế Chương 3: đọc dữ liệu theo ngữ cảnh vận hành xe, tối ưu năng lượng trong lúc đỗ, phục hồi kết nối sau mất mạng và giữ khả năng mở rộng cho OTA/hardening ở các vòng sau.

\subsubsection{4.1.3. Triển khai cloud, backend, frontend và giám sát}

Tầng cloud của hệ thống được triển khai theo mô hình dịch vụ tách rời nhưng liên kết qua mạng Docker dùng chung. Luồng dữ liệu chính gồm thiết bị \(\rightarrow\) EMQX \(\rightarrow\) MQTT Bridge \(\rightarrow\) PostgreSQL/VictoriaMetrics/VictoriaLogs \(\rightarrow\) Backend API/WebSocket \(\rightarrow\) Frontend dashboard.

\begin{itemize}
\tightlist
\item
  \textbf{EMQX:} đóng vai trò broker MQTT trung tâm cho thiết bị IoT.
\item
  \textbf{MQTT Bridge:} kiểm tra payload, ánh xạ dữ liệu và ghi vào các hệ lưu trữ phù hợp.
\item
  \textbf{Backend API:} cung cấp REST API, WebSocket thời gian thực, xác thực phiên, phân quyền và các nghiệp vụ quản lý đội xe.
\item
  \textbf{Frontend dashboard:} hiển thị bản đồ, trạng thái xe, cảnh báo, lịch sử và các thành phần quản trị cho người vận hành.
\item
  \textbf{Observability:} VictoriaMetrics, VictoriaLogs và Grafana hỗ trợ đo chỉ số, tập trung log và theo dõi sức khỏe hệ thống.
\end{itemize}

So với bản đầy đủ, bản rút gọn không lặp lại toàn bộ cấu hình Docker, danh sách endpoint, cấu trúc thư mục frontend hay middleware chain. Thay vào đó, nội dung giữ trọng tâm ở sự liên thông kỹ thuật: dữ liệu từ thiết bị đã có đường đi rõ ràng đến giao diện giám sát, đồng thời hạ tầng vẫn đủ khả năng cô lập trách nhiệm theo từng dịch vụ để thuận lợi cho kiểm thử và mở rộng.

\subsubsection{4.1.4. Kết luận phần triển khai}

Nhìn tổng thể, giai đoạn triển khai đã chuyển được thiết kế thành một nguyên mẫu hoạt động xuyên suốt nhiều lớp công nghệ. Các chi tiết chiếm nhiều dung lượng như bản vẽ phần cứng, bảng tra chân, code snippet dài, ma trận API và checklist hardening được giữ ở bản đầy đủ; còn trong bản rút gọn, điều cần khẳng định là hệ thống đã được hiện thực hóa thành một pipeline hoàn chỉnh từ thiết bị thật đến dashboard thật, đủ cơ sở để bước sang đánh giá kết quả vận hành.

\subsection{4.2. Tóm tắt đo lường và kết quả -- Compact measurements and results}

Mục tiêu của phần này là giữ lại các kết quả định lượng quan trọng nhất, đồng thời lược bỏ các bảng kiểm thử nhiều trang và các ma trận phụ chỉ phục vụ tra cứu chi tiết. Các kết quả được nhóm theo bốn lớp: phần cứng, firmware, cloud/frontend và kiểm thử tích hợp toàn hệ thống.

\subsubsection{4.2.1. Kết quả phần cứng và năng lượng}

Khối nguồn và kiến trúc tiêu thụ năng lượng đạt mức vận hành phù hợp với mục tiêu đồ án. Trong điều kiện mẫu hiện tại, dòng deep sleep toàn hệ thống xấp xỉ \(500\,\mu A\); nếu quy đổi theo pin dự phòng 18650 chuẩn 3500mAh, thiết bị có thể duy trì khoảng 24--36 giờ ở chế độ cảnh báo gián đoạn và khoảng 2.6--3.0 giờ ở chế độ tracking liên tục. Kết quả này cho thấy thiết kế nguồn đã đủ khả năng duy trì giám sát ngắn hạn khi xe mất nguồn chính hoặc điện áp xuống thấp.

Bên cạnh đó, GNSS trên SIM7600CE-T cho sai số vị trí khoảng 2.5--5 m ở môi trường thuận lợi và 5--15 m trong môi trường đô thị. Đây là mức phù hợp với bài toán giám sát phương tiện, geofence và tái hiện hành trình. Về phần OBD2, kết nối BLE với adapter vgate iCar Pro hoạt động ổn định sau giai đoạn khởi tạo; các thông số như tốc độ, RPM, nhiệt độ nước và mức nhiên liệu được thu nhận với sai số chấp nhận được cho mục tiêu theo dõi vận hành.

\subsubsection{4.2.2. Kết quả firmware}

Firmware cho thấy ba kết quả nổi bật. Thứ nhất, thời gian kết nối BLE sau khi thiết bị thức dậy duy trì ở mức khoảng 1--3 giây khi có sẵn thông tin thiết bị đã ghép cặp. Thứ hai, thời gian đọc một PID OBD2 phổ biến nằm trong khoảng 100--150 ms, đủ để lấy mẫu theo chu kỳ vận hành mà không làm nghẽn máy trạng thái. Thứ ba, cơ chế reconnect MQTT phục hồi được phiên kết nối trong khoảng 15--30 giây ở các tình huống mất sóng ngắn.

Máy trạng thái \emph{Driving}, \emph{Parking}, \emph{Alert} cũng hoạt động đúng logic thiết kế. Thời gian chuyển trạng thái thực tế nằm khoảng 200--400 ms trong các ca thử nghiệm tiêu biểu, nhỏ hơn ngưỡng mục tiêu 500 ms. Điều này cho thấy phần mềm nhúng không chỉ vận hành được mà còn giữ được tính phản ứng cần thiết cho các sự kiện di chuyển và cảnh báo.

\subsubsection{4.2.3. Kết quả cloud, backend và frontend}

Ở tầng cloud, API backend phản hồi các yêu cầu đơn giản trong khoảng 20--80 ms và các truy vấn phức tạp trong khoảng 100--300 ms ở cấu hình thử nghiệm hiện tại. Độ trễ cập nhật WebSocket từ lúc thiết bị gửi telemetry đến lúc dashboard hiển thị ở mức dưới 2 giây, đáp ứng yêu cầu giám sát gần thời gian thực. Hệ thống cũng đã được kiểm thử với quy mô 50+ thiết bị đồng thời trong môi trường mô phỏng, qua đó xác nhận hướng kiến trúc hiện tại đủ cho quy mô đồ án và có dư địa tối ưu tiếp.

Frontend dashboard cho thấy thời gian tải trang đầu dưới khoảng 2 giây đối với các trang chính, bản đồ cập nhật vị trí xe ổn định và các thành phần thống kê/cảnh báo vận hành đúng vai trò. Đây là điểm quan trọng vì nó chứng minh lớp hiển thị không trở thành nút nghẽn cuối cùng sau khi dữ liệu đã đi qua các lớp firmware và backend.

\subsubsection{4.2.4. Tóm tắt chỉ tiêu chính}

\begin{center}
\begin{tabular}{|p{0.34\linewidth}|p{0.24\linewidth}|p{0.30\linewidth}|}
\hline
\textbf{Chỉ tiêu chính} & \textbf{Mục tiêu} & \textbf{Kết quả rút gọn} \\
\hline
Dòng deep sleep toàn hệ thống & \(< 500\,\mu A\) & Xấp xỉ \(500\,\mu A\), đạt ngưỡng thiết kế \\
\hline
Kết nối BLE sau wake-up & \(< 5\) giây & Khoảng 1--3 giây \\
\hline
Đọc một PID OBD2 & \(< 200\) ms & Khoảng 100--150 ms \\
\hline
Reconnect MQTT & \(< 30\) giây & Khoảng 15--30 giây \\
\hline
Độ trễ WebSocket end-to-end & \(< 2\) giây & Đạt ở môi trường thử nghiệm \\
\hline
Số thiết bị đồng thời & \(\ge 50\) & Đã kiểm thử 50+ thiết bị \\
\hline
Sai số GPS ngoài trời & Phù hợp giám sát xe & Khoảng 2.5--5 m \\
\hline
\end{tabular}
\end{center}

\subsubsection{4.2.5. Nhận xét và giới hạn}

Kết quả rút gọn cho phép đi tới kết luận rằng nguyên mẫu đã đạt đa số chỉ tiêu cốt lõi của đồ án: giám sát được trạng thái xe, truyền dữ liệu theo thời gian gần thực, hiển thị được trên dashboard và duy trì hoạt động theo các chế độ năng lượng khác nhau. Tuy nhiên, một số hạng mục vẫn nên được xem là vùng cần làm sâu thêm thay vì đã hoàn tất tuyệt đối: hiệu chuẩn ngưỡng LVD cho profile 12V, đánh giá bền vững của kênh replay dữ liệu cục bộ trong mất mạng dài hạn, bổ sung watchdog hoàn chỉnh cho OTA và mở rộng kiểm thử tương thích trên nhiều mẫu xe hơn.

Như vậy, bản rút gọn vẫn giữ được điều quan trọng nhất của Chương 4: không chỉ nêu rằng hệ thống đã được xây dựng, mà còn chứng minh nó đã được đo, kiểm thử và đánh giá trên các chỉ tiêu vận hành then chốt.

\chapterheading{CHƯƠNG 5. ĐÁNH GIÁ VÀ KHUYẾN NGHỊ -- EVALUATION AND RECOMMENDATIONS}

\subsection{5.1. Đánh giá hiệu năng}

Từ các kết quả rút gọn ở Chương 4, có thể nhận định hệ thống đạt mức hoàn thiện phù hợp với một đồ án tốt nghiệp theo hướng kỹ thuật ứng dụng. Điểm mạnh lớn nhất là sự cân bằng giữa bốn lớp công nghệ: phần cứng đủ gọn và tiết kiệm điện, firmware phản ứng đúng theo trạng thái xe, cloud xử lý dữ liệu ổn định và frontend hiển thị trực quan cho người vận hành.

\subsubsection{5.1.1. Đánh giá theo từng lớp kỹ thuật}

\begin{itemize}
\tightlist
\item
  \textbf{Phần cứng:} kiến trúc nguồn và lựa chọn linh kiện đáp ứng mục tiêu tích hợp tracker trên xe, duy trì GNSS, 4G và OBD2 trong cùng một thiết bị.
\item
  \textbf{Firmware:} máy trạng thái và mô hình điều phối trung tâm giúp hệ thống dễ kiểm soát hơn so với cách dàn trải nhiều task nghiệp vụ độc lập ngay từ đầu.
\item
  \textbf{Cloud/backend:} kiến trúc tách dịch vụ bằng Docker và MQTT tạo luồng xử lý rõ ràng, thuận lợi cho việc mở rộng và quan sát vận hành.
\item
  \textbf{Frontend:} dashboard đã làm tròn vai trò của một công cụ giám sát đội xe thay vì chỉ là giao diện minh họa.
\end{itemize}

\subsection{5.2. Đánh giá kinh tế và khả năng áp dụng}

Về chi phí, hướng tiếp cận dùng công nghệ mã nguồn mở giúp giảm mạnh chi phí giấy phép phần mềm và dồn ngân sách cho phần cứng cùng hạ tầng tối thiểu. BOM thiết bị vẫn nằm trong vùng khả thi cho giai đoạn thử nghiệm và triển khai quy mô nhỏ. Điều này đặc biệt phù hợp với bối cảnh các đơn vị cho thuê xe hoặc quản lý đội xe nhỏ cần một giải pháp đủ dùng nhưng không thể đầu tư ngay vào hệ thống thương mại đắt tiền.

Ở góc nhìn triển khai thực tế, chi phí không chỉ nằm ở linh kiện mà còn ở thời gian vận hành và bảo trì. Kiến trúc hiện tại có lợi thế ở chỗ từng tầng có thể tách trách nhiệm rõ: hỏng ở thiết bị, broker, bridge, backend hay frontend đều có điểm quan sát riêng. Nhờ đó, chi phí vận hành về sau có thể kiểm soát tốt hơn so với một hệ thống ghép khối khó chẩn đoán.

\subsection{5.3. Đánh giá rủi ro và biện pháp giảm thiểu}

Những rủi ro chính còn lại của hệ thống tập trung vào bốn nhóm sau:

\begin{itemize}
\tightlist
\item
  \textbf{Rủi ro năng lượng:} sai lệch ngưỡng LVD hoặc tiêu thụ cao ở pha truyền 4G kéo dài có thể làm giảm thời gian dự phòng. Biện pháp là hiệu chuẩn thêm trên cả profile 12V và 24V, đồng thời đo tải thực địa dài hơn.
\item
  \textbf{Rủi ro kết nối:} mất sóng kéo dài hoặc môi trường BLE/OBD2 không ổn định có thể gây rơi dữ liệu hoặc chậm cập nhật. Biện pháp là củng cố replay queue, logging và bộ đo chất lượng liên kết.
\item
  \textbf{Rủi ro triển khai:} sự khác biệt giữa xe thử nghiệm và xe thực tế ngoài hiện trường có thể làm phát sinh vấn đề về nguồn, vị trí lắp đặt hoặc nhiễu điện từ. Biện pháp là mở rộng ma trận thử nghiệm trên nhiều mẫu xe.
\item
  \textbf{Rủi ro bảo mật/vận hành:} khi chuyển sang môi trường production, yêu cầu TLS, ACL per device, quản lý bí mật và giám sát lỗi phải được siết chặt hơn. Biện pháp là bổ sung hardening checklist và tự động hóa kiểm tra phát hành.
\end{itemize}

\subsection{5.4. Khuyến nghị phát triển tiếp theo}

Lộ trình phát triển hợp lý sau đồ án gồm bốn hướng. Một là hoàn thiện độ bền vận hành: watchdog, burn-in test, replay bền vững và kiểm thử tải dài hạn. Hai là mở rộng chức năng: mobile app, OTA hoàn chỉnh, cảnh báo nâng cao và quản trị đội xe lớn hơn. Ba là mở rộng chiều sâu phân tích dữ liệu: geofence thông minh, đánh giá hành vi lái xe, dự báo bảo dưỡng và các lớp AI/ML khi có đủ dữ liệu thực tế. Bốn là chuẩn hóa quy trình phát hành: CI/CD, staging, quan sát log/metrics tập trung và checklist production.

\subsection{Kết luận chương 5}

Tóm lại, hệ thống đã đạt được trạng thái \emph{khả dụng về kỹ thuật} và \emph{khả thi về triển khai thử nghiệm}. Dự án chưa phải một sản phẩm thương mại hoàn chỉnh, nhưng đã vượt qua mức nguyên mẫu trình diễn đơn thuần: kiến trúc rõ ràng, luồng dữ liệu khép kín và các rủi ro chính đã được nhận diện cụ thể. Đây là nền tảng đủ tốt để tiếp tục phát triển ở giai đoạn sau.

\chapterheading{CHƯƠNG 6. PHẢN HỒI VÀ BÀI HỌC KINH NGHIỆM -- REFLECTION AND LESSONS LEARNED}

\subsection{6.1. Ứng dụng kiến thức kỹ thuật -- Application of prior coursework}

Đồ án là nơi nhiều mảng kiến thức được ghép lại thành một hệ thống hoàn chỉnh thay vì tồn tại rời rạc như trong từng môn học. Kiến thức về vi điều khiển, điện tử số và analog hỗ trợ thiết kế phần cứng và đọc tín hiệu nguồn; kiến thức mạng máy tính và IoT hỗ trợ lựa chọn MQTT, WebSocket, topology cloud; kiến thức cơ sở dữ liệu và kỹ thuật phần mềm hỗ trợ thiết kế schema, API, phân lớp dịch vụ và kiểm thử.

Giá trị lớn nhất của quá trình này là người thực hiện không chỉ dùng lại kiến thức cũ, mà còn học cách kết nối chúng dưới áp lực tích hợp thực tế. Chính yêu cầu làm cho phần cứng, firmware, backend và frontend cùng hoạt động đã biến lý thuyết thành năng lực giải quyết vấn đề kỹ thuật liên ngành.

\subsection{6.2. Các vấn đề kỹ thuật phức tạp đã xử lý}

Bốn nhóm vấn đề nổi bật nhất của đồ án là:

\begin{itemize}
\tightlist
\item
  \textbf{Đọc dữ liệu OBD2 qua BLE:} phải xử lý khác biệt giữa tầng chẩn đoán xe và tầng Bluetooth, đồng thời đảm bảo độ trễ vẫn chấp nhận được khi thiết bị vừa thức dậy.
\item
  \textbf{Điều phối trạng thái và năng lượng:} tracker phải vừa tiết kiệm điện khi xe đỗ, vừa phản ứng đủ nhanh khi xe di chuyển hoặc xảy ra rung động bất thường.
\item
  \textbf{Đường ống dữ liệu thời gian thực:} dữ liệu phải đi qua nhiều lớp từ thiết bị đến dashboard mà vẫn giữ được tính cập nhật và có khả năng phục hồi sau lỗi mạng.
\item
  \textbf{Tổ chức hệ thống nhiều tầng:} mỗi lớp công nghệ có ràng buộc riêng, nên phần khó nhất không phải từng lớp đơn lẻ mà là giao điểm giữa các lớp.
\end{itemize}

Những vấn đề này cho thấy đây không phải một bài toán lập trình thuần túy, mà là một bài toán hệ thống, nơi quyết định ở tầng này luôn ảnh hưởng trực tiếp đến tầng khác.

\subsection{6.3. Tác động đạo đức và xã hội}

Một hệ thống theo dõi phương tiện luôn đi kèm câu hỏi về quyền riêng tư và trách nhiệm dữ liệu. Giá trị vận hành chỉ bền vững khi dữ liệu vị trí, hành trình và trạng thái xe được thu thập đúng mục đích, giới hạn truy cập rõ ràng và có cơ chế bảo vệ phù hợp. Vì vậy, ngoài giá trị kỹ thuật, đồ án cũng rút ra ba nguyên tắc sử dụng quan trọng:

\begin{itemize}
\tightlist
\item
  Dữ liệu giám sát phải gắn với mục tiêu vận hành hợp pháp, không biến thành công cụ theo dõi quá mức.
\item
  Thiết kế kỹ thuật cần ưu tiên an toàn phương tiện, do đó giải pháp chỉ đọc OBD2 và không điều khiển xe từ xa là lựa chọn phù hợp ở giai đoạn này.
\item
  Khi triển khai thực tế, bảo mật đường truyền, phân quyền truy cập, chính sách lưu trữ và quy trình thu hồi quyền là bắt buộc chứ không phải phần phụ.
\end{itemize}

\subsection{6.4. Bài học kinh nghiệm và định hướng cá nhân}

Bài học lớn nhất của đồ án là phải thiết kế kiến trúc trước khi tối ưu chi tiết. Khi hệ thống có nhiều tầng, mọi quyết định nóng vội ở giai đoạn đầu đều quay lại thành chi phí tích hợp ở giai đoạn sau. Bài học thứ hai là kiểm thử phải diễn ra ở mọi tầng, vì lỗi hiếm khi lộ ra đúng nơi nó phát sinh. Bài học thứ ba là hệ sinh thái mã nguồn mở thực sự mạnh khi biết ghép các công cụ trưởng thành vào một kiến trúc nhất quán, thay vì chỉ dùng chúng như các thành phần rời rạc.

Ở góc độ cá nhân, đồ án tạo ra một bước chuyển từ tư duy làm từng module sang tư duy làm sản phẩm hệ thống. Giá trị quan trọng nhất không chỉ là thiết bị hoặc dashboard đã hoạt động, mà là việc người thực hiện đã đi qua trọn vòng đời của một hệ thống IoT: đặt yêu cầu, chọn kiến trúc, triển khai, đo lường, nhận ra giới hạn và chỉ ra hướng phát triển tiếp theo.

\subsection{Kết luận chương 6}

Chương cuối cho thấy kết quả của đồ án không chỉ nằm ở một nguyên mẫu kỹ thuật, mà còn ở năng lực phân tích hệ thống và tư duy triển khai có trách nhiệm. Những kinh nghiệm rút ra về kiến trúc, năng lượng, dữ liệu thời gian thực, kiểm thử tích hợp và đạo đức dữ liệu là nền tảng trực tiếp cho các giai đoạn phát triển tiếp theo của hệ thống.

\majorheading{TÀI LIỆU TRÍCH DẪN - REFERENCES}

\begin{center}\rule{0.5\linewidth}{0.5pt}\end{center}

\subsection{{[}A{]} Tài liệu tiếng
Việt}\label{a-tuxe0i-liux1ec7u-tiux1ebfng-viux1ec7t}

{[}1{]} Bộ Giao thông Vận tải, "Báo cáo thống kê phương tiện giao thông
đường bộ Việt Nam giai đoạn 2020--2025," Hà Nội, 2024.

{[}2{]} Nguyễn Văn A và Trần Thị B, "Ứng dụng Internet vạn vật trong
giám sát và quản lý phương tiện giao thông tại Việt Nam," \emph{Tạp chí
Khoa học và Công nghệ}, tập 58, số 3, trang 45--52, 2023.

{[}3{]} Lê Minh C, "Thiết kế hệ thống định vị và theo dõi phương tiện sử
dụng GPS và mạng di động," Luận văn thạc sĩ, Đại học Bách Khoa TP. Hồ
Chí Minh, 2022.

{[}4{]} Bộ Thông tin và Truyền thông, "Luật An toàn thông tin mạng,"
Luật số 24/2018/QH14, Hà Nội, 2018.

\begin{center}\rule{0.5\linewidth}{0.5pt}\end{center}

\subsection{{[}B{]} Tài liệu tiếng Anh (Books \&
Papers)}\label{b-tuxe0i-liux1ec7u-tiux1ebfng-anh-books--papers}

{[}5{]} A. Al-Fuqaha, M. Guizani, M. Mohammadi, M. Aledhari, and M.
Ayyash, "Internet of Things: A Survey on Enabling Technologies,
Protocols, and Applications," \emph{IEEE Communications Surveys \&
Tutorials}, vol. 17, no. 4, pp. 2347--2376, Fourth Quarter 2015.

{[}6{]} N. Naik, "Choice of Effective Messaging Protocols for IoT
Systems: MQTT, CoAP, AMQP and HTTP," in \emph{Proc. 2017 IEEE
International Systems Engineering Symposium (ISSE)}, Vienna, Austria,
2017, pp. 1--7.

{[}7{]} M. Quigley, R. Johnson, and P. Sharma, "Power Management
Strategies for IoT Vehicle Tracking Devices," \emph{IEEE Internet of
Things Journal}, vol. 8, no. 15, pp. 12045--12058, Aug. 2021.

{[}8{]} M. A. Al-Khedher, "Hybrid GPS-GSM Localization of Automobile
Tracking System," \emph{International Journal of Computer Science and
Information Technology (IJCSIT)}, vol. 3, no. 6, pp. 75--85, Dec. 2011.

{[}9{]} S. Kaur and S. Singh, "A Survey on IoT based Vehicle Tracking
System," \emph{International Journal of Advanced Research in Computer
Science}, vol. 10, no. 3, pp. 42--47, 2019.

{[}10{]} R. Barry, \emph{Mastering the FreeRTOS Real Time Kernel: A
Hands-On Tutorial Guide}. Real Time Engineers Ltd., 2016.

{[}11{]} B. Boehm and V. R. Basili, "Software Defect Reduction Top 10
List," \emph{IEEE Computer}, vol. 34, no. 1, pp. 135--137, Jan. 2001.

{[}12{]} E. Evans, \emph{Domain-Driven Design: Tackling Complexity in
the Heart of Software}. Boston, MA, USA: Addison-Wesley Professional,
2003.

{[}13{]} S. Newman, \emph{Building Microservices: Designing Fine-Grained
Systems}, 2nd ed. Sebastopol, CA, USA: O\textquotesingle Reilly Media,
2021.

{[}14{]} D. Guinard and V. Trifa, \emph{Building the Web of Things}.
Shelter Island, NY, USA: Manning Publications, 2016.

{[}15{]} J. Gubbi, R. Buyya, S. Marusic, and M. Palaniswami, "Internet
of Things (IoT): A Vision, Architectural Elements, and Future
Directions," \emph{Future Generation Computer Systems}, vol. 29, no. 7,
pp. 1645--1660, Sep. 2013.

{[}16{]} V. C. Gungor and G. P. Hancke, "Industrial Wireless Sensor
Networks: Challenges, Design Principles, and Technical Approaches,"
\emph{IEEE Transactions on Industrial Electronics}, vol. 56, no. 10, pp.
4258--4265, Oct. 2009.

{[}17{]} R. K. Kodali, V. Jain, S. Bose, and L. Boppana, "IoT Based
Smart Security and Home Automation System," in \emph{Proc. 2016
International Conference on Computing, Communication and Automation
(ICCCA)}, Greater Noida, India, 2016, pp. 1286--1289.

{[}18{]} B. Sahu and G. A. Rincon-Mora, "A Low Voltage, Dynamic,
Noninverting, Synchronous Buck-Boost Converter for Portable
Applications," \emph{IEEE Transactions on Power Electronics}, vol. 19,
no. 2, pp. 443--452, Mar. 2004.

\begin{center}\rule{0.5\linewidth}{0.5pt}\end{center}

\subsection{{[}C{]} Tài liệu kỹ thuật (Datasheets \& Technical
Documents)}\label{c-tuxe0i-liux1ec7u-kux1ef9-thuux1eadt-datasheets--technical-documents}

{[}19{]} Espressif Systems, "ESP32-S3 Series Datasheet," Version 1.3,
2023. {[}Online{]}. Available:
\url{https://www.espressif.com/sites/default/files/documentation/esp32-s3_datasheet_en.pdf}

{[}20{]} Espressif Systems, "ESP32-S3 Technical Reference Manual,"
Version 1.1, 2023. {[}Online{]}. Available:
\url{https://www.espressif.com/sites/default/files/documentation/esp32-s3_technical_reference_manual_en.pdf}

{[}21{]} SIMCom Wireless Solutions, "SIM7600 Series AT Command Manual,"
Version 1.06, 2023. {[}Online{]}. Available:
\url{https://www.simcom.com/product/SIM7600CE-T.html}

{[}22{]} SIMCom Wireless Solutions, "SIM7600CE-T Module Hardware Design
Guide," Version 1.02, 2022.

{[}23{]} STMicroelectronics, "LIS3DSH - MEMS Digital Output Motion
Sensor Ultra-Low-Power High-Performance 3-Axis
\textquotesingle Nano\textquotesingle{} Accelerometer Datasheet," DocID
17530, Rev. 3, 2021. {[}Online{]}. Available:
\url{https://www.st.com/resource/en/datasheet/lis3dh.pdf}

{[}24{]} STMicroelectronics, "LIS3DSH Application Note AN3308," Rev. 2,
2020.

{[}25{]} OASIS, "MQTT Version 5.0 - OASIS Standard," Mar. 2019.
{[}Online{]}. Available:
\url{https://docs.oasis-open.org/mqtt/mqtt/v5.0/mqtt-v5.0.html}

{[}26{]} SAE International, "SAE J1979 - E/E Diagnostic Test Modes,"
Standard, Rev. 2014.

{[}27{]} International Organization for Standardization, "ISO
15031--5:2015 - Road vehicles - Communication between vehicle and
external equipment for emissions-related diagnostics," ISO Standard,
2015.

{[}28{]} International Organization for Standardization, "ISO
15765--2:2016 - Road vehicles - Diagnostic communication over Controller
Area Network (DoCAN) - Part 2: Transport protocol and network layer
services," ISO Standard, 2016.

{[}29{]} vgate, "iCar Pro BLE OBD2 Adapter Specifications," 2023.
{[}Online{]}. Available:
\url{https://www.vgatemall.com/product/vgate-icar-pro.html}

{[}30{]} NanJing Top Power ASIC Corp., "TP4056 Standalone Linear Li-Ion
Battery Charger with Thermal Regulation," Datasheet, Rev. 1.0.

{[}31{]} Apache NimBLE Project, "NimBLE Host API Reference," Version
1.5, 2023. {[}Online{]}. Available:
\url{https://mynewt.apache.org/latest/network/}

\begin{center}\rule{0.5\linewidth}{0.5pt}\end{center}

\subsection{{[}D{]} Tài liệu trực tuyến (Online
Resources)}\label{d-tuxe0i-liux1ec7u-trux1ef1c-tuyux1ebfn-online-resources}

{[}32{]} EMQX Team, "EMQX Documentation - The World\textquotesingle s
Most Scalable MQTT Broker," Version 5.x, 2024. {[}Online{]}. Available:
\url{https://docs.emqx.com/en/emqx/latest/}. {[}Accessed: Jan. 15,
2026{]}.

{[}33{]} The PostgreSQL Global Development Group, "PostgreSQL 16
Documentation," 2024. {[}Online{]}. Available:
\url{https://www.postgresql.org/docs/16/}. {[}Accessed: Jan. 10,
2026{]}.

{[}34{]} VictoriaMetrics Team, "VictoriaMetrics Documentation," 2024.
{[}Online{]}. Available: \url{https://docs.victoriametrics.com/}.
{[}Accessed: Jan. 10, 2026{]}.

{[}35{]} VictoriaMetrics Team, "VictoriaLogs Documentation," 2024.
{[}Online{]}. Available:
\url{https://docs.victoriametrics.com/victorialogs/}. {[}Accessed: Jan.
10, 2026{]}.

{[}36{]} OpenJS Foundation, "Express.js - Fast, Unopinionated,
Minimalist Web Framework for Node.js," 2024. {[}Online{]}. Available:
\url{https://expressjs.com/}. {[}Accessed: Dec. 20, 2025{]}.

{[}37{]} Vercel Inc., "Next.js 15 Documentation," 2024. {[}Online{]}.
Available: \url{https://nextjs.org/docs}. {[}Accessed: Dec. 15, 2025{]}.

{[}38{]} Meta Platforms, Inc., "React 19 Documentation," 2024.
{[}Online{]}. Available: \url{https://react.dev/}. {[}Accessed: Dec. 15,
2025{]}.

{[}39{]} Socket.IO Contributors, "Socket.IO Documentation," Version 4.8,
2024. {[}Online{]}. Available: \url{https://socket.io/docs/v4/}.
{[}Accessed: Dec. 20, 2025{]}.

{[}40{]} V. Agafonkin, "Leaflet - An Open-Source JavaScript Library for
Mobile-Friendly Interactive Maps," Version 1.9, 2024. {[}Online{]}.
Available: \url{https://leafletjs.com/}. {[}Accessed: Jan. 05, 2026{]}.

{[}41{]} Apache Software Foundation, "Apache ECharts - An Open Source
JavaScript Visualization Library," Version 5.5, 2024. {[}Online{]}.
Available: \url{https://echarts.apache.org/}. {[}Accessed: Jan. 05,
2026{]}.

{[}42{]} Docker Inc., "Docker Documentation," 2024. {[}Online{]}.
Available: \url{https://docs.docker.com/}. {[}Accessed: Nov. 15,
2025{]}.

{[}43{]} Docker Inc., "Docker Compose Documentation," 2024.
{[}Online{]}. Available: \url{https://docs.docker.com/compose/}.
{[}Accessed: Nov. 15, 2025{]}.

{[}44{]} Tailwind Labs, "Tailwind CSS Documentation," Version 4, 2025.
{[}Online{]}. Available: \url{https://tailwindcss.com/docs}.
{[}Accessed: Jan. 10, 2026{]}.

{[}45{]} Daishi Kato, "Zustand - Bear Necessities for State Management
in React," Version 5, 2024. {[}Online{]}. Available:
\url{https://zustand-demo.pmnd.rs/}. {[}Accessed: Dec. 20, 2025{]}.

{[}46{]} TanStack, "TanStack Query (React Query) Documentation," Version
5, 2024. {[}Online{]}. Available:
\url{https://tanstack.com/query/latest}. {[}Accessed: Dec. 20, 2025{]}.

{[}47{]} Espressif Systems, "ESP-IDF Programming Guide," Version 5.x,
2024. {[}Online{]}. Available:
\url{https://docs.espressif.com/projects/esp-idf/en/latest/esp32s3/}.
{[}Accessed: Oct. 10, 2025{]}.

{[}48{]} Grafana Labs, "Grafana Documentation," 2024. {[}Online{]}.
Available: \url{https://grafana.com/docs/grafana/latest/}. {[}Accessed:
Jan. 15, 2026{]}.

{[}49{]} Prometheus Authors, "Prometheus Documentation," 2024.
{[}Online{]}. Available: \url{https://prometheus.io/docs/}. {[}Accessed:
Jan. 15, 2026{]}.

{[}50{]} Radix UI Contributors, "Radix UI - Unstyled, Accessible
Components for React," 2024. {[}Online{]}. Available:
\url{https://www.radix-ui.com/}. {[}Accessed: Dec. 15, 2025{]}.

{[}51{]} Google, "Flutter Documentation - WebView Plugin," 2024.
{[}Online{]}. Available: \url{https://docs.flutter.dev/}. {[}Accessed:
Jan. 20, 2026{]}.

{[}52{]} Texas Instruments, "Power Management for IoT Devices,"
Application Report SLVAE57, 2020. {[}Online{]}. Available:
\url{https://www.ti.com/lit/an/slvae57/slvae57.pdf}. {[}Accessed: Sep.
15, 2025{]}.

{[}53{]} STMicroelectronics, "STM32L476RE Product Page," 2026.
{[}Online{]}. Available:
\url{https://www.st.com/en/microcontrollers-microprocessors/stm32l476re.html}.
{[}Accessed: Mar. 02, 2026{]}.

{[}54{]} STMicroelectronics, "STM32L476xx Datasheet," DS10198 Rev 11,
Jul. 2024. {[}Online{]}. Available:
\url{https://www.st.com/resource/en/datasheet/stm32l476re.pdf}.
{[}Accessed: Mar. 02, 2026{]}.

{[}55{]} Nordic Semiconductor, "nRF52840 Product Specification v1.2,"
{[}Online{]}. Available:
\url{https://docs.nordicsemi.com/bundle/nRF52840_PS_v1.2/resource/nRF52840_PS_v1.2.pdf}.
{[}Accessed: Mar. 02, 2026{]}.

{[}56{]} Nordic Semiconductor, "nRF52840 Product Page," 2026.
{[}Online{]}. Available:
\url{https://www.nordicsemi.com/Products/nRF52840/Modules}. {[}Accessed:
Mar. 02, 2026{]}.

{[}57{]} Nordic Semiconductor, "nRF52840 Documentation Hub," 2026.
{[}Online{]}. Available:
\url{https://docs.nordicsemi.com/category/nrf52840-category}.
{[}Accessed: Mar. 02, 2026{]}.

{[}58{]} SIMCom Wireless Solutions, "A7670X Series Product Page," 2026.
{[}Online{]}. Available:
\url{https://cn.simcom.com/product/A7670X.html}. {[}Accessed: Mar. 02,
2026{]}.

{[}59{]} Quectel Wireless Solutions, "EC200U Series Product Page," 2026.
{[}Online{]}. Available:
\url{https://www.quectel.com/product/lte-ec200u-series}. {[}Accessed:
Mar. 02, 2026{]}.

{[}60{]} SIMCom Wireless Solutions, "SIM7600CE Product Page," 2026.
{[}Online{]}. Available:
\url{https://en.simcom.com/product/SIM7600CE.html}. {[}Accessed: Mar.
13, 2026{]}.

{[}61{]} Quectel Wireless Solutions, "Quectel EC200U Series LTE Standard
Specification," Version 1.4, 2024. {[}Online{]}. Available:
\url{https://developer.quectel.com/en/wp-content/uploads/sites/2/2024/11/Quectel_EC200U_Series_LTE_Standard_Specification_V1.4.pdf}.
{[}Accessed: Mar. 02, 2026{]}.

{[}62{]} u-blox, "NEO-M8 Data Sheet," UBX-15031086, {[}Online{]}.
Available:
\url{https://www.u-blox.com/sites/default/files/NEO-M8-FW3_DataSheet_UBX-15031086.pdf}.
{[}Accessed: Mar. 02, 2026{]}.

{[}63{]} u-blox, "NEO-M8 Series Product Page," 2026. {[}Online{]}.
Available: \url{https://www.u-blox.com/en/product/neo-m8-series}.
{[}Accessed: Mar. 02, 2026{]}.

{[}64{]} vgate, "Vgate iCar Pro BLE Product Page," 2026. {[}Online{]}.
Available: \url{https://vgatemall.com/products-detail/i-9/}.
{[}Accessed: Mar. 02, 2026{]}.

{[}65{]} Texas Instruments, "TPS2115A Product Page," 2026. {[}Online{]}.
Available: \url{https://www.ti.com/product/TPS2115A}. {[}Accessed: Mar.
02, 2026{]}.

{[}66{]} Infineon Technologies, "IRLML6402 Product Page," 2026.
{[}Online{]}. Available: \url{https://www.infineon.com/part/IRLML6402}.
{[}Accessed: Mar. 02, 2026{]}.

{[}67{]} SONGLE Relay, "SRD-05VDC-SL-C Relay Datasheet," datasheet
mirror. {[}Online{]}. Available:
\url{https://www.alldatasheet.com/html-pdf/1132639/SONGLERELAY/SRD-05VDC-SL-C/1713/2/SRD-05VDC-SL-C.html}.
{[}Accessed: Mar. 02, 2026{]}.

\begin{center}\rule{0.5\linewidth}{0.5pt}\end{center}

% Phụ lục được lược bỏ trong bản rút gọn khoảng 150 trang.

\end{document}
